% Merged from ten separately delivered manuscripts on quaternionic analysis:
% one expository survey and nine independent attempts at the same follow-up
% problem.  See MERGE_NOTES.md for the provenance and the notation
% reconciliation.  Every symbol below is fixed by that reconciliation; the
% source manuscripts disagreed on most of them.
\documentclass[11pt,a4paper]{article}
\usepackage[T1]{fontenc}
\usepackage{lmodern}
\usepackage[margin=26mm]{geometry}
\usepackage{amsmath,amssymb,amsthm,mathtools}
\usepackage{booktabs,array,longtable}
\usepackage{microtype}
\usepackage{enumitem}
\usepackage{xcolor}
\usepackage{fancyhdr}
\usepackage{hyperref}
\usepackage{bookmark}

\definecolor{Ink}{HTML}{233240}
\hypersetup{colorlinks=true,linkcolor=Ink,citecolor=Ink,urlcolor=Ink,
 pdftitle={Global Fueter Primitives: affine monodromy, spherical residues, and
 the energy of a singular sphere},
 pdfauthor={Research notes prepared for Vladimir Reshetnikov},
 pdfsubject={One document merged from ten manuscripts; slice regularity,
 Cauchy-Fueter analysis, and the global inverse Fueter problem}}

\pagestyle{fancy}
\fancyhf{}
\fancyhead[L]{\small Global Fueter primitives}
\fancyhead[R]{\small\leftmark}
\fancyfoot[C]{\thepage}
\renewcommand{\sectionmark}[1]{\markboth{#1}{}}
\setlength{\headheight}{14pt}
\setlength{\emergencystretch}{2em}

\numberwithin{equation}{section}
\newtheorem{theorem}{Theorem}[section]
\newtheorem{proposition}[theorem]{Proposition}
\newtheorem{lemma}[theorem]{Lemma}
\newtheorem{corollary}[theorem]{Corollary}
\theoremstyle{definition}
\newtheorem{definition}[theorem]{Definition}
\newtheorem{conjecture}[theorem]{Conjecture}
\newtheorem{example}[theorem]{Example}
\theoremstyle{remark}
\newtheorem{remark}[theorem]{Remark}
\theoremstyle{plain}

% ---------------------------------------------------------------- algebra ---
\newcommand{\HH}{\mathbb H}
\newcommand{\HC}{\mathbb H_{\mathbb C}}
\newcommand{\RR}{\mathbb R}
\newcommand{\CC}{\mathbb C}
\newcommand{\ZZ}{\mathbb Z}
\newcommand{\SSS}{\mathbb S}
\newcommand{\Cplus}{\CC^{+}}
\DeclareMathOperator{\Ind}{Ind}
\DeclareMathOperator{\Rea}{Re}
\DeclareMathOperator{\Ima}{Im}
\newcommand{\ReH}{\Rea_{\HH}}
\newcommand{\ImH}{\Ima_{\HH}}
\newcommand{\ReC}{\Rea_{\CC}}

% The Cauchy-Fueter operator and the Fueter map.  The Fueter map is written
% out as \lap\I(\cdot) everywhere and is NEVER abbreviated to a single letter:
% the sources used T for this map, for the holomorphic invariant, and for the
% normalized inverse, three incompatible meanings in one commuting diagram.
\newcommand{\Dop}{\mathcal D}
\newcommand{\lap}{\Delta_4}
\newcommand{\I}{\mathcal I}

% ------------------------------------------------------------ the spaces ---
\newcommand{\SR}{\mathcal{SR}}
\newcommand{\AM}{\mathcal{AM}}
\newcommand{\Om}{\Omega}
\newcommand{\Stem}{\mathcal S}
\newcommand{\Aff}{\mathrm{Aff}_{\HH}}
\newcommand{\Omom}{\mathcal O_{\mathrm{mom}}}

% ---------------------------------------------- the two forms and periods ---
% omega is the slope form, theta the constant form.  theta is computed from the
% TARGET, with no first primitive chosen; that is the whole point of the
% theory.  eta_C below is the primitive-DEPENDENT form of the original local
% construction, kept only to display the identity eta_C - d(xC) = theta.
\newcommand{\om}{\omega}
\newcommand{\thet}{\theta}
\newcommand{\etaC}{\eta_{C}}
\newcommand{\Per}{\mathrm{Per}}
\newcommand{\pot}{C}          % slope potential  C = B/y
\newcommand{\potV}{V}         % constant potential V = A - xB/y

% ------------------------------------------------- singular-sphere fixtures --
% Centre p = u + iota v with v>0; offsets xi = x-u, eta = y-v.  The letters
% s and t are NOT used for either: three manuscripts gave them three
% mutually exclusive meanings.
\newcommand{\Sp}{S_{p}}
\newcommand{\Rp}{\mathcal R_{p}}      % residue norm
\newcommand{\Hg}{\mathcal H_{g}}      % holomorphic invariant
\newcommand{\nn}{\mathfrak n}         % pair norm on H^2

% ------------------------------------------------------------ the two norms --
% THE MOST DANGEROUS COLLISION IN THE SOURCES.  M_g meant the spherical
% supremum in three manuscripts and the axial pair norm in another, and the
% two obey laws with genuinely different constants.  They get different
% symbols here and are never conflated.
\newcommand{\axn}[1]{\lVert #1\rVert_{\mathrm{ax}}}   % axial pair norm
\newcommand{\Msph}{M^{\mathrm{sph}}}                  % spherical supremum
\newcommand{\Max}{M^{\mathrm{ax}}}                    % axial-pair maximum

% ------------------------------------------------------------------ energy --
\newcommand{\Eg}{\mathcal E_{g}}
\newcommand{\Jg}{J_{g}}
\newcommand{\Eren}{E^{\mathrm{ren}}}
\newcommand{\Tube}{T}

\newcommand{\abs}[1]{\lvert #1\rvert}
\newcommand{\norm}[1]{\lVert #1\rVert}
\newcommand{\dd}{\mathrm d}
\newcommand{\Res}{\operatorname{Res}}
\newcommand{\Resaff}{\operatorname{Res}^{\mathrm{aff}}}

\title{\textbf{Global Fueter Primitives}\\[4pt]
\large Affine monodromy, spherical residues, and the energy of a singular
sphere\\[2pt]
\normalsize with a self-contained account of slice-regular and
Cauchy--Fueter analysis}
\author{Research notes prepared for Vladimir Reshetnikov}
\date{September 2026}

\begin{document}
\maketitle
\setcounter{tocdepth}{2}
\tableofcontents
\clearpage

%% ==================== part0_1.tex ====================

% =====================================================================
%  Part 0 and Part I of the merged document.
%  Writer key: a.   All labels defined here begin with "a:".
%  Lead source: manuscript 03 (the expository original) for Part I;
%  Part 0 draws its status ledger from 01, 04, 06, 08 and its notation
%  from the merge contract.
% =====================================================================

\section{Part 0. The question, the status ledger, and the conventions}
\label{a:sec:front}

\subsection{The single question}
\label{a:sec:question}

Two notions of regularity coexist over the quaternions, and a second-order
differential operator --- the Euclidean Laplacian of $\RR^4$ --- maps one into
the other. Fueter's mapping theorem says that if $f$ is slice regular then
$\lap f$ is Cauchy--Fueter regular. The whole of this document is organized
around the inverse of that map, and around one question about it.

\medskip
\noindent\textbf{The problem.}\ \ Let $U$ be a connected open subset of the
upper half-plane $\Cplus=\{x+\iota y: y>0\}$, let
\[
 \Om_U=\{x+Iy:\ x+\iota y\in U,\ I\in\SSS\}
\]
be its circularization, and let $g$ be a single-valued axially monogenic
function on $\Om_U$, that is, a smooth $g(x+Iy)=P(x,y)+IQ(x,y)$ with
$\Dop g=0$. \emph{When does there exist a single-valued slice-regular $f$ on
the same domain $\Om_U$ with $\lap f=g$, and when there does not, what
measures the failure?} The same question is asked for a connected
conjugation-invariant planar set $D\subseteq\CC$ that meets $\RR$, whose
circularization $\Om_D$ is a genuine slice domain rather than a product
domain.

\medskip
The emphasis on \emph{the same domain} is the entire content. Locally the
answer is classical and affirmative: on a simply connected base one integrates
two closed one-forms in succession and obtains a primitive, unique up to
$q\mapsto qa+b$ (Theorem~\ref{a:thm:local-inverse} below, which is the last
theorem of Part~I and the seam at which the rest of the document begins).
What the local construction leaves open is exactly what happens when a local
primitive is continued around a hole of $U$: it can return changed by an
affine function, and the Laplacian cannot see that change. Local surjectivity
is therefore compatible with a nonzero global cokernel, and an integral
formula --- however explicit --- does not by itself remove monodromy.

The answers, proved in Parts~II--IX, are these. Two closed quaternion-valued
one-forms $\om_g$ and $\thet_g$, both computed directly from the target $g$
with no primitive chosen, have periods $(a_\gamma,b_\gamma)\in\HH^2$ that
vanish on every loop precisely when a single-valued primitive exists; they are
precisely the coefficients of the affine monodromy
$f\mapsto f+qa_\gamma+b_\gamma$. On a base with a finite winding basis of size
$m$ the cokernel of $\lap$ has real dimension exactly $8m$, realized by
explicit logarithmic modes. All of this data is carried by a single
holomorphic $\HC$-valued function $\Hg$. At a missing nonreal sphere the same
eight real parameters reappear as an affine spherical residue pair, as the
density of a distributional surface source, as two conserved three-dimensional
fluxes, and as the exact coefficient of a logarithmic divergence of the
$L^2$ energy.

\subsection{What is background and what is proposed}
\label{a:sec:status}

This document is assembled from ten separately written manuscripts: one
expository synthesis of the two regularity theories, and nine independent
attacks on the problem just stated. Part~I is the expository trunk; Parts~II
through~IX are the extension. One status paragraph is stated here and is not
repeated anywhere else.

Established background, claimed by nobody here as new: the quaternionic
algebra and the stem formalism of Ghiloni--Perotti~\cite{GP}; slice regularity and the
classical function theory built on it; Cauchy--Fueter regularity, the Green
identity and the Cauchy--Fueter formula; Fueter's mapping theorem; the affine
kernel $\Aff$; the axial Vekua system; the local inverse Fueter theory of
Colombo--Sabadini--Sommen~\cite{CSS} and Colombo--Pe\~na Pe\~na--Sabadini--Sommen~\cite{CPSS},
including the arctangent primitives and the sphere-integrated Cauchy kernels;
Perotti's surjectivity theorem under the hypothesis that every component of
the planar generating set is simply connected; the classical scalar Runge~\cite{Runge} and
Mittag--Leffler theorems; Weyl's lemma; and the standard logarithmic-cutoff
capacity argument for $L^2$ removability. Every one of these is used, and
several are reproved in Part~I to keep the document self-contained; none is
advertised as a contribution.

Proposed as contributions: the two simultaneously defined, target-computable
closed forms and the resulting same-domain solvability criterion with no
topological hypothesis; its identification with an additive affine monodromy
$H_1(U;\ZZ)\to\HH\oplus\HH$; the reflection-sensitive split range theorem and
the $8m$ cokernel with the correct index; the holomorphic encoding $\Hg$ with
its exact range, its exact kernel, and the reality constraint on its Laurent
coefficients; the singularity classification along a nonreal sphere under the
weakest hypothesis (local integrability) together with the growth filtration
above it and the exact leading constants in both norms; the distributional
surface source and the two conserved fluxes; the exact energy law with its
sharp growth-free lower bound; and the Mittag--Leffler realization theorem.

\begin{center}
\small
\renewcommand{\arraystretch}{1.18}
\begin{tabular}{@{}p{0.31\textwidth}p{0.62\textwidth}@{}}
\toprule
Established, used here & Proposed, proved here\\
\midrule
Stems, slice regularity, the classical slice apparatus (Part~I)
 & Two target-computable closed forms; same-domain criterion; affine
   monodromy (Part~II)\\
Cauchy--Fueter theory: Green, Cauchy--Pompeiu, $o(r^{-3})$ removability,
 flux residues (Part~I)
 & The holomorphic invariant $\Hg$, its exact range $\Omom$, its kernel,
   and the residue reality constraint (Part~III)\\
Fueter's mapping theorem; $\ker(\lap|_{\SR})=\Aff$ (Part~I)
 & Logarithmic modes, split exact sequence, $\dim_\RR\operatorname{coker}
   \lap=8m$ (Parts~IV--V)\\
Local inverse on a simply connected base (Part~I)
 & $L^1$ and growth classifications at a nonreal sphere with exact
   constants (Part~VII)\\
Arctangent and spherical Cauchy kernels; Runge; Mittag--Leffler; Weyl;
 the $L^2$ cutoff method
 & Surface source, conserved fluxes, exact energy law, realization
   theorem (Parts~VIII--IX)\\
\bottomrule
\end{tabular}
\end{center}

\medskip
The novelty claim is qualified and is not certified. The literature audit
behind it was carried out on 21 September 2026 and is targeted, not
exhaustive; it is described once, in Part~XI. Every mathematical claim below
rests on its proof and not on that audit. One editorial position is taken
once and held throughout: this document states a necessary and sufficient
period criterion, exhibits explicit obstructed axially monogenic fields on
$\HH\setminus\SSS$, and observes that an \emph{unqualified} same-domain
surjectivity assertion is therefore false as stated; it declines to adjudicate
whether any particular published theorem, read with its own domain, locality
and branch hypotheses, is incorrect. The topologically restricted
surjectivity theorems cited above are correct as hypothesized, and the
contribution here is the removal of the hypothesis, not the correction of the
theorem.

\subsection{Conventions, and the collisions they resolve}
\label{a:sec:notation}

The ten sources disagreed on nearly every symbol, and several of the
disagreements are dangerous rather than cosmetic: the same letter was used for
two quantities obeying laws with genuinely different constants. The
conventions below are fixed for the whole document and are never varied.

\paragraph{Quaternions.}
$q=x_0+\mathbf i x_1+\mathbf j x_2+\mathbf k x_3$ with
$\mathbf i^2=\mathbf j^2=\mathbf k^2=\mathbf i\mathbf j\mathbf k=-1$;
$\bar q$ is quaternionic conjugation, which reverses products;
$\abs q^2=q\bar q$; $\Rea q=\tfrac12(q+\bar q)$ and $\Ima q=q-\Rea q$;
$\ReH$ and $\ImH$ denote the same operations when the ambient algebra could
be misread. The sphere of imaginary units and its planes are
\[
 \SSS=\{I\in\HH: I^2=-1\},\qquad \CC_I=\RR+I\RR .
\]
Every nonreal quaternion is uniquely $x+Iy$ with $y>0$ and $I\in\SSS$, and its
conjugacy class is the two-sphere $[q]=x+y\SSS$.

\paragraph{The auxiliary complex unit.}
$\HC=\HH\oplus\iota\HH$ with $\iota$ \emph{central}, $\iota^2=-1$; the
complex conjugation of this algebra is $\kappa(A+\iota B)=A-\iota B$ and is
\emph{not} quaternionic conjugation; $\ReC(A+\iota B)=A$ and
$\Ima_{\CC}(A+\iota B)=B$ take values in $\HH$, and
$\norm{A+\iota B}_{\HC}^2=\abs A^2+\abs B^2$.
\textbf{Warning.} $\iota$ is never an element of $\SSS$. Confusing the two is
the single most common way to destroy a statement in this subject: the kernel
of $\lap$ on slice-regular functions consists of $qa+b$ with $a,b\in\HH$ and
would be twice as large if $a,b$ were allowed in $\HC$
(Theorem~\ref{a:thm:fueter-map} and Remark~\ref{a:rem:complex-affine}).

\paragraph{Sides.}
All slice-regular functions are \emph{left} slice regular and all
quaternionic vector spaces are \emph{right} vector spaces; every quaternionic
coefficient multiplies on the right, in power series, in period pairs, in
Laurent coefficients and in every integral formula. A coefficient moved to
the left changes the problem.

\paragraph{The operator and the map.}
\[
 \Dop=\partial_{x_0}+\mathbf i\partial_{x_1}+\mathbf j\partial_{x_2}
       +\mathbf k\partial_{x_3},
 \qquad
 \overline{\Dop}=\partial_{x_0}-\mathbf i\partial_{x_1}
       -\mathbf j\partial_{x_2}-\mathbf k\partial_{x_3},
 \qquad
 \overline{\Dop}\,\Dop=\Dop\,\overline{\Dop}=\lap=\sum_{\mu=0}^{3}
       \partial_{x_\mu}^2 .
\]
There is \emph{no} factor $\tfrac12$ anywhere in $\Dop$ or in $\lap$;
consequently $\Dop q=-2$, $\lap q^2=-4$, and every constant in the energy law
of Part~VIII is normalized against these. The Fueter map on stems is written
out as $F\mapsto\lap\I(F)$ throughout and is \textbf{never} abbreviated by a
single letter. Three of the sources called it $T$, one called $T$ the
holomorphic invariant, and one called $T$ the normalized inverse --- three
incompatible arrows of the same commuting diagram. Likewise the letter $J$ is
not used: it named the splitting section in two sources and the invariant in
four. The section is $S$, the projection is $\Pi=\mathrm{id}-S\Per$, and the
invariant is $\Hg$.

\paragraph{Bases and coordinates.}
$U$ always denotes a connected open subset of $\Cplus$, and $D$ a connected
conjugation-invariant open subset of $\CC$ meeting $\RR$; $\Om_U$ and $\Om_D$
are their circularizations. The planar coordinate is $z=x+\iota y$ with
$y>0$, and the quaternionic point is $q=x+Iy$. \emph{The axial radius is $y$};
the letters $r$ and $R$ are reserved for radii of balls and tubes. Five of the
sources wrote $r$ for the axial radius and five wrote $y$; the second choice
is adopted.

\paragraph{The singular sphere.}
The centre is $p=u+\iota v$ with $v>0$; the sphere is $\Sp=u+v\SSS$; the
offsets are $\xi=x-u$ and $\eta=y-v$; and
$\rho=\abs{z-p}=\operatorname{dist}(q,\Sp)$.
\textbf{Warning.} The letters $s$ and $t$ are not used for any of these.
Three sources gave the pair $(s,t)$ three mutually exclusive meanings --- the
offsets in one order, the offsets in the opposite order, and the centre --- in
formulas that this document must place side by side.

\paragraph{The table.}
Everything else is collected in Table~\ref{a:tab:notation}, which is the
reference for the rest of the document.

{\small
\renewcommand{\arraystretch}{1.25}
\begin{longtable}{@{}p{0.30\textwidth}p{0.64\textwidth}@{}}
\caption{Notation fixed for the whole document.}
\label{a:tab:notation}\\
\toprule
Symbol & Meaning\\
\midrule
\endfirsthead
\toprule
Symbol & Meaning\\
\midrule
\endhead
$\Stem(U)=\mathcal O(U,\HC)$ & holomorphic stems $F=A+\iota B$ on $U$, with
 $A_x=B_y$, $A_y=-B_x$\\
$\I(F)(x+Iy)=A+IB$ & the induced slice function; $\SR$ and $\AM$ are the
 spaces of slice-regular and axially monogenic functions\\
$g=P+IQ$ & the axially monogenic target; $P,Q:U\to\HH$\\
$f=\I(F)=A+IB$ & a slice-regular primitive, $\lap f=g$\\
$\pot=B/y$, \ $\potV=A-xB/y$ & the slope potential and the constant
 potential; $f(q)=q\pot+\potV$ (Part~II)\\
$\om_g=\tfrac12(-P\dd x+Q\dd y)$ & the first (slope) closed form\\
$\thet_g=\tfrac12\bigl((xP+yQ)\dd x+(yP-xQ)\dd y\bigr)$ & the second
 (constant) closed form, computed from $g$ alone, with \emph{no} primitive
 chosen\\
$\etaC=(\pot+\tfrac{yQ}2)\dd x+\tfrac{yP}2\dd y$ & the
 primitive-\emph{dependent} second form of the classical local construction
 \eqref{a:eq:inverse-second}. \textbf{Warning:} $\etaC\ne\thet_g$; they
 differ by $\dd(x\pot)$, and one source called $\etaC$ what four others called
 $\thet_g$\\
$a_\gamma=\oint_\gamma\om_g$, \ $b_\gamma=\oint_\gamma\thet_g$ & the slope
 period and the constant period; $\Per(g)$ is the period map\\
$\Aff=\{q\mapsto qa+b:\ a,b\in\HH\}$ & the affine kernel; $a,b$ are
 \emph{quaternions}\\
$S$, \ $\Pi=\mathrm{id}-S\Per$ & the splitting section and the projection
 onto the range\\
$\Hg=-\tfrac12(P+yP_y)-\tfrac{\iota y}{2}P_x$ & the holomorphic invariant;
 $\Hg=F''$ for every local primitive stem\\
$\Omom(U)$ & the moment-real class: holomorphic $h:U\to\HC$ with
 $\oint_\gamma h\dd z\in\HH$ and $-\oint_\gamma zh\dd z\in\HH$ for every loop\\
$G_{p;a,b}=\lap\I\bigl(\tfrac{za+b}{2\pi\iota}\log(z-p)\bigr)$ & the
 logarithmic mode with period pair $\Ind(\gamma,p)\,(a,b)$\\
$\Ind(\gamma,p)$ & winding number\\
$\axn{g}(z)=(\abs P^2+\abs Q^2)^{1/2}$ & the \emph{axial pair norm}\\
$\sup_{I\in\SSS}\abs{P+IQ}$ & the \emph{spherical supremum} at one base point;
 $\axn g\le\sup_I\abs{P+IQ}\le\sqrt2\,\axn g$\\
$\Max(\rho)$, \ $\Msph(\rho)$ & the maxima of the two norms over
 $\abs{z-p}=\rho$. \textbf{Warning, the worst collision in the sources:} one
 letter $M_g$ denoted the first in one manuscript and the second in three
 others, and the two obey leading-size laws with \emph{different constants}
 (Part~VII). They are never conflated here\\
$\Rp(a,b)^2=\abs{ua+b}^2+v^2\abs a^2=\norm{pa+b}_{\HC}^2$ & the residue norm;
 five sources gave it five names\\
$\nn(\alpha,\beta)=\max_{I\in\SSS}\abs{\beta+I\alpha}
 =(\abs\alpha^2+\abs\beta^2+2\abs{\ImH(\beta\bar\alpha)})^{1/2}$ & the pair
 norm on $\HH^2$ governing the spherical supremum\\
$\Eg(\varepsilon,R)=\int_{\varepsilon<\rho<R}\abs g^2\dd V$ & the truncated
 \emph{energy}. \textbf{Warning:} ``energy'' means the squared $L^2$ norm,
 never a Dirichlet integral\\
$\Jg(\rho)$, \ $\Eren$, \ $\Tube_R(\Sp)$ & the shell density, the
 renormalized constant, the tube of radius $R$ about $\Sp$\\
$E(q)=\bar q/(2\pi^2\abs q^4)$ & the Cauchy--Fueter kernel. The letter $E$
 also began ``energy'' in the sources; here the energy always carries a
 subscript or the decoration $\Eg,\Eren$\\
$m$ & the rank of a winding basis of the upper meridian
 $U=D\cap\{y>0\}$, equivalently the number of conjugation-exchanged pairs of
 complementary components of $D$. These two counts agree; the equivalence is
 proved in Part~V and is assumed nowhere before it\\
$\Res_c f$ & the slice Laurent residue at a \emph{real} isolated singularity
 (\S\ref{a:sec:slice-laurent})\\
$\Res^{F}_a f$ & the Cauchy--Fueter flux residue at an \emph{isolated point}
 (\S\ref{a:sec:fueter-residue})\\
$\Resaff_{\Sp}(g)=(a,b)$ & the \emph{affine spherical residue} at a nonreal
 sphere (Part~VII). These three residues are different objects and are never
 identified\\
$\mathcal P_k(q)=-\lap(q^{k+2})/\bigl(2(k+1)(k+2)\bigr)$ & the entire axially
 monogenic polynomials of Part~IX, normalized by
 $\mathcal P_k(x)=x^k$ on $\RR$. The calligraphic letter distinguishes them
 from the axial component $P$\\
\bottomrule
\end{longtable}}

\medskip
Two further conventions are stated once. Integration of a quaternion-valued
object always means integration of its four real components with the displayed
order of quaternionic factors retained. And a period pair $(a,b)\in\HH^2$ is
an \emph{ordered pair of quaternions}, carrying eight real parameters, not a
single quaternion and not a complex number.

% =====================================================================

\section{Part I. Two regularity theories and the bridge between them}
\label{a:sec:two-theories}

This part supplies, once, the background that Parts~II--XI silently assume,
and ends with the local theorem that they globalize. It follows the
expository source throughout; the nine extension manuscripts reprove none of
it, yet each of them uses the stem formalism, the affine kernel, real-centred
Taylor expansion (for smoothness across the real axis), Laurent theory,
harmonic Liouville and the Green identity. Everything here is established
mathematics; the proofs are included so that the later parts depend on no
unproved quaternionic machinery. The standard scalar theorems of one-variable
complex analysis --- Cauchy, the maximum principle, Laurent, Montel, Hurwitz,
the normalized Riemann mapping theorem --- are prerequisites, as are the
divergence theorem and the existence of smooth mollifiers.

\subsection{Why there are two theories}
\label{a:sec:why-two}

Over $\CC$, differentiability, local power series, the Cauchy--Riemann
equation and several geometric properties fit into one theory. Over $\HH$ the
literal translations do not fit together, and the reason is a rigidity
theorem.

\begin{proposition}[Rigidity of a direction-independent quotient]
\label{a:prop:rigidity}
Let $\Om\subseteq\HH$ be connected and open and let $f\in C^2(\Om,\HH)$.
If for every $q\in\Om$ the limit
\[
 a(q)=\lim_{h\to0}h^{-1}\bigl(f(q+h)-f(q)\bigr)
\]
exists independently of the direction of $h$, then $f(q)=qa+b$ for fixed
$a,b\in\HH$.
\end{proposition}

\begin{proof}
Put $e_0=1$, $e_1=\mathbf i$, $e_2=\mathbf j$, $e_3=\mathbf k$. Increments
$h=te_\mu$ give $\partial_\mu f=e_\mu a$ and $a=\partial_0f$. Equality of the
mixed second derivatives with $\partial_0$ gives
$\partial_\mu a=e_\mu\partial_0a$, and then equality of the derivatives with
indices $1,2$ gives $\mathbf i\mathbf j\,\partial_0a
=\mathbf j\mathbf i\,\partial_0a$. Since
$\mathbf i\mathbf j-\mathbf j\mathbf i=2\mathbf k$ is invertible,
$\partial_0a=0$; hence all first derivatives of $a$ vanish, connectedness
makes $a$ constant, and all first derivatives of $f(q)-qa$ vanish.
\end{proof}

The quotient with $h^{-1}$ on the right gives $f(q)=aq+b$ instead. Either
way the class excludes $q^2$, so one does not force the complex difference
quotient to do everything. \emph{Slice regularity} keeps power series in one
quaternionic variable; \emph{Cauchy--Fueter regularity} keeps a first-order
elliptic system and a boundary integral in four real dimensions. The relation
between them is a differential transform, not an identification of function
spaces, and that transform is the subject of \S\ref{a:sec:bridge}.

\subsection{Slice regularity, representation, and stems}
\label{a:sec:slice}

Throughout this subsection $D\subseteq\CC$ is connected, open, invariant under
complex conjugation and meets $\RR$; $\Om_D$ is its circularization. It is an
axially symmetric slice domain: it contains a whole conjugacy sphere whenever
it contains one of its points, and each section $\Om_D\cap\CC_I$ is connected.

\begin{definition}
\label{a:def:slice-regular}
A $C^1$ function $f:\Om_D\to\HH$ is \emph{slice regular} if for each
$I\in\SSS$ its restriction $f_I$ satisfies
\begin{equation}
 \bar\partial_If_I=\tfrac12(\partial_x+I\partial_y)f_I=0 .
 \label{a:eq:CRslice}
\end{equation}
Its slice (Cullen) derivative is
$\partial_cf=\tfrac12(\partial_x-I\partial_y)f=\partial_xf$.
\end{definition}

\begin{lemma}[Splitting]
\label{a:lem:splitting}
Fix orthogonal $I,J\in\SSS$. Every quaternion-valued function on $\CC_I$ is
uniquely $F+GJ$ with $F,G$ taking values in $\CC_I$. It satisfies
\eqref{a:eq:CRslice} if and only if $F$ and $G$ are holomorphic, and
$\abs{F+GJ}^2=\abs F^2+\abs G^2$.
\end{lemma}

\begin{proof}
The real basis $1,I,J,IJ$ gives the decomposition; substituting it into
\eqref{a:eq:CRslice} the two $\CC_I$ components vanish independently and give
ordinary Cauchy--Riemann equations. The norm identity is orthogonality of the
four basis elements.
\end{proof}

\begin{theorem}[Identity principle]
\label{a:thm:identity}
If two slice-regular functions on $\Om_D$ agree on a subset of one
$\CC_I$-section having an accumulation point in that section, they agree on
$\Om_D$.
\end{theorem}

\begin{proof}
Apply scalar uniqueness to both splitting components of the difference: it
vanishes on the whole $I$-section, hence on a real interval of $D$. Its
restriction to any other section has those same real zeros, and the splitting
lemma with scalar uniqueness finishes.
\end{proof}

The phrase ``one section'' is essential: an arbitrary accumulation of zeros in
$\HH$ does not suffice, since a whole two-sphere can be a zero set
(Theorem~\ref{a:thm:zeros}).

\begin{theorem}[Representation formula and stems]
\label{a:thm:representation}
For slice-regular $f$ on $\Om_D$ and all $I,K\in\SSS$,
\begin{equation}
 f(x+Ky)=\frac{1-KI}{2}f(x+Iy)+\frac{1+KI}{2}f(x-Iy).
 \label{a:eq:representation}
\end{equation}
Equivalently there are unique $\HH$-valued $A(x,y),B(x,y)$ with
\begin{equation}
 f(x+Ky)=A(x,y)+KB(x,y),
 \label{a:eq:AB}
\end{equation}
$A$ even and $B$ odd in $y$, satisfying
\begin{equation}
 A_x=B_y,\qquad A_y=-B_x .
 \label{a:eq:stemCR}
\end{equation}
\end{theorem}

\begin{proof}
On one section put $h_\pm(x,y)=f(x\pm Iy)$ and
$A=\tfrac12(h_++h_-)$, $B=-\tfrac I2(h_+-h_-)$; the parities are immediate and
\eqref{a:eq:CRslice} gives $(h_\pm)_y=\pm I(h_\pm)_x$, whence
\eqref{a:eq:stemCR}. Thus $A+KB$ is holomorphic on each $K$-section and equals
$f(x)$ at $y=0$, so scalar uniqueness identifies it with $f$. Expanding
$A+KB$ gives \eqref{a:eq:representation}; evaluating at $K$ and $-K$ gives
uniqueness.
\end{proof}

The function $F(x+\iota y)=A(x,y)+\iota B(x,y)$ is holomorphic with values in
$\HC$ and satisfies
\begin{equation}
 F(\bar z)=\kappa(F(z)),
 \label{a:eq:stem-symmetry}
\end{equation}
and is called a \emph{holomorphic stem}; conversely a stem induces a
slice-regular function by \eqref{a:eq:AB}, an operation written $f=\I(F)$.
Away from $\RR$ this is \eqref{a:eq:stemCR}; at a real point the complex
Taylor coefficients of $F$ lie in $\HH$ by \eqref{a:eq:stem-symmetry}, so the
induced series is real analytic --- which is also what justifies $C^1$
regularity at real points in the converse construction. A slice function is
\emph{intrinsic} (slice preserving) when $A,B$ are real valued; its values on
each $\CC_I$ lie in that plane, and real-coefficient power series are
intrinsic.

\begin{proposition}[The two derivatives at a nonreal point]
\label{a:prop:differential}
Let $q=x+Iy$ with $y>0$ and split $v\in\HH$ as $v_\parallel+v_\perp$ with
$v_\parallel\in\CC_I$ and $v_\perp\perp\CC_I$. Then
\begin{equation}
 \dd f_q(v)=v_\parallel\,\partial_cf(q)
  +v_\perp\,\partial_{\mathrm{sph}}f(q),
 \qquad
 \partial_{\mathrm{sph}}f(q)=\frac{B(x,y)}{y}
  =(2Iy)^{-1}\bigl(f(q)-f(\bar q)\bigr).
 \label{a:eq:differential}
\end{equation}
At a real point, $\dd f_x(v)=v\,\partial_cf(x)$.
\end{proposition}

\begin{proof}
Within $\CC_I$ the Cauchy--Riemann equation gives $\partial_yf=I\partial_xf$,
which is the parallel formula. In a perpendicular direction the real part and
the length of the imaginary part have zero first variation while $I$ varies by
$v_\perp/y$; differentiate \eqref{a:eq:AB}. At a real point use the
real-centred power series.
\end{proof}

For $f(q)=q^2$ the two right multipliers are $2q$ and $2x$; their norms differ,
so slice regularity does not imply four-dimensional conformality. The
spherical derivative $B/y$ is exactly the slope potential $\pot$ of
Part~II, a coincidence that is not accidental: the first obstruction of the
global theory is the obstruction to integrating $\pot$.

\subsection{Cauchy formula, Taylor series, Morera}
\label{a:sec:slice-cauchy}

For $q\notin[s]$ set
\begin{equation}
 \mathcal C(q,s)=-\bigl(q^2-2\Rea(s)q+\abs s^2\bigr)^{-1}(q-\bar s),
 \label{a:eq:kernel}
\end{equation}
the inverse being the ordinary quaternionic inverse of the displayed value.
For fixed $s$ this is slice regular in $q$ off $[s]$: it is induced by the
stem obtained by replacing $q$ with the central variable $z$. When
$q,s\in\CC_I$ commutativity gives $\mathcal C(q,s)=(s-q)^{-1}$, and for
$\abs q<\abs s$,
\begin{equation}
 \mathcal C(q,s)=\sum_{n\ge0}q^ns^{-n-1}.
 \label{a:eq:kernelseries}
\end{equation}
Indeed the absolutely convergent sum $R$ satisfies $Rs-qR=1$, and
$s^2-2\Rea(s)s+\abs s^2=0$ then gives
$0=R(s^2-2\Rea(s)s+\abs s^2)=(q^2-2\Rea(s)q+\abs s^2)R+q-\bar s$. This is not
a licence to commute $q$ with $s$.

\begin{theorem}[Slice Cauchy formula]
\label{a:thm:cauchy-slice}
Let $D$ be bounded with piecewise $C^1$ boundary and $f$ slice regular on an
axially symmetric slice domain containing $\overline{\Om_D}$. Give
$\partial D_I$ its positive complex orientation (with the usual negative
orientation on holes) and set $\dd s_I=-I\dd s$. Then for all $q\in\Om_D$,
\begin{equation}
 f(q)=\frac1{2\pi}\int_{\partial D_I}\mathcal C(q,s)\,\dd s_I\,f(s),
 \label{a:eq:cauchy-slice}
\end{equation}
independently of $I$. In particular
$\int_{\partial V_I}\dd s_I\,f(s)=0$ for any relatively compact planar region
$V_I$ in a section on which $f$ is regular up to the boundary.
\end{theorem}

\begin{proof}
On the $I$-section use Lemma~\ref{a:lem:splitting} and the ordinary Cauchy
formula on both components; the normalization $\dd s_I=-I\dd s$ places the
complex factor $1/I$ correctly. For arbitrary $q$, both sides are slice
functions of $q$ obeying \eqref{a:eq:representation} --- the integral does
because every kernel does and the factors $\dd s_I f(s)$ sit on the right ---
so equality on one section gives equality everywhere. Conjugation invariance
of $D$ keeps the kernel nonsingular for $s\in\partial D_I$. Cauchy's theorem
is the same argument applied to the two splitting components.
\end{proof}

\begin{theorem}[Taylor theorem at real centres]
\label{a:thm:taylor}
Let $c\in\RR$ and let $f$ be slice regular on $B(c,R)$. There are unique
$a_n\in\HH$ with
\begin{equation}
 f(q)=\sum_{n\ge0}(q-c)^na_n,\qquad a_n=\frac{\partial_c^nf(c)}{n!},
 \label{a:eq:taylor}
\end{equation}
converging absolutely and uniformly on smaller closed balls, and for
$0<r<R$,
\begin{equation}
 a_n=\frac1{2\pi}\int_{\abs{s-c}=r,\,s\in\CC_I}(s-c)^{-n-1}\dd s_I\,f(s),
 \qquad
 \abs{a_n}\le\frac{M_I(r)}{r^n},\quad
 M_I(r)=\max_{\abs{s-c}=r,\,s\in\CC_I}\abs{f(s)} .
 \label{a:eq:cauchy-estimate}
\end{equation}
\end{theorem}

\begin{proof}
Take ordinary Taylor expansions of both splitting components on one section;
their quaternionic combinations give $a_n$, and real differentiation at $c$
identifies them independently of the section. Scalar Cauchy gives the integral
formula, whose norm estimate is $(2\pi)^{-1}(2\pi r)r^{-n-1}M_I(r)$. For
$\abs{q-c}\le\varrho<r$ the resulting series is dominated by a geometric
series, defines a slice-regular function agreeing with $f$ on the chosen
section, and hence agrees everywhere by Theorem~\ref{a:thm:identity}.
\end{proof}

Consequently every slice-regular function is real analytic: the stem
description proves this off $\RR$ and \eqref{a:eq:taylor} proves it at $\RR$.
This is precisely the fact used in Part~V, where the two closed one-forms are
extended smoothly across the real axis. Ordinary powers centred at a nonreal
quaternion are \emph{not} asserted to give a Taylor theorem on a Euclidean
ball.

\begin{theorem}[Morera]
\label{a:thm:morera-slice}
A continuous $f:\Om_D\to\HH$ with $\int_{\partial T}\dd s_I\,f(s)=0$ for every
$I\in\SSS$ and every triangle $T$ with closure in the $I$-section is slice
regular.
\end{theorem}

\begin{proof}
Scalar Morera makes both splitting components of $f_I$ holomorphic; construct
$A,B$ from that section as in Theorem~\ref{a:thm:representation} to get a
real-analytic slice-regular extension. On every other section the extension
and $f$ are holomorphic and agree on a real interval, hence coincide.
\end{proof}

\subsection{Maximum modulus, Liouville, Schwarz}
\label{a:sec:max}

\begin{theorem}[Maximum modulus]
\label{a:thm:max}
If $f\in\SR(\Om_D)$ and $\abs f$ has a local maximum at an interior point,
then $f$ is constant.
\end{theorem}

\begin{proof}
Let the maximum be at $p\in\CC_I$ with $M=\abs{f(p)}$. If $M=0$ apply
Theorem~\ref{a:thm:identity}. Otherwise put $u=f(p)/M$, $g=f\bar u$, so $g$ is
regular, $g(p)=M$, $\abs g=\abs f$. Split $g_I=F+GJ$; near $p$,
$\abs F\le\abs g\le M$ and $F(p)=M$, so scalar maximum modulus gives
$F\equiv M$ locally and $\abs g^2=\abs F^2+\abs G^2\le M^2$ forces $G=0$.
Theorem~\ref{a:thm:identity} finishes.
\end{proof}

\begin{theorem}[Liouville and polynomial growth]
\label{a:thm:liouville-slice}
A bounded entire slice-regular function is constant. More generally if
$\abs{f(q)}\le C(1+\abs q)^\alpha$ on $\HH$ with $\alpha\ge0$, then $f$ is a
right-coefficient polynomial of degree at most $\lfloor\alpha\rfloor$.
\end{theorem}

\begin{proof}
Expand at $0$; \eqref{a:eq:cauchy-estimate} gives
$\abs{a_n}\le C(1+r)^\alpha r^{-n}$, and $r\to\infty$ kills every $n>\alpha$.
\end{proof}

\begin{theorem}[Higher-order Schwarz lemma]
\label{a:thm:schwarz}
Let $f:\mathbb B\to\mathbb B$ be slice regular on the unit ball
$\mathbb B=B(0,1)$, with the first $m\ge1$ Taylor coefficients vanishing.
Then
\begin{equation}
 \abs{f(q)}\le\abs q^m,\qquad \frac{\abs{\partial_c^mf(0)}}{m!}\le1,
 \label{a:eq:schwarz}
\end{equation}
with equality in the first inequality at a nonzero point, or in the second,
exactly when $f(q)=q^mu$ for a fixed $\abs u=1$.
\end{theorem}

\begin{proof}
Write $f(q)=q^mg(q)$ from the power series; for $\abs q=r<1$,
$\abs{g(q)}\le r^{-m}$, and maximum modulus on $B(0,r)$ with $r\uparrow1$
gives $\abs g\le1$. Either equality forces $\abs g$ to attain $1$ inside, so
Theorem~\ref{a:thm:max} makes $g$ a unit constant.
\end{proof}

These proofs illustrate a principle used repeatedly below: use scalar
holomorphic components for linear analytic questions, but preserve the
quaternionic norm whenever a sharp scalar bound is wanted.

\subsection{The regular product, symmetrization, and the geometry of zeros}
\label{a:sec:algebra}

Pointwise multiplication does not preserve slice regularity: the constant
$\mathbf i$ and the identity are regular, but on the $\mathbf j$-section
$(\partial_x+\mathbf j\partial_y)(\mathbf iq)=2\mathbf i\ne0$. The replacement
is multiplication before evaluating the stem.

\begin{definition}
\label{a:def:star}
For $f=\I(F)$ and $g=\I(G)$ put $f*g=\I(FG)$. Thus with $f=A+IB$ and
$g=C'+ID'$,
\begin{equation}
 (f*g)(x+Iy)=AC'-BD'+I(AD'+BC').
 \label{a:eq:starAB}
\end{equation}
\end{definition}

Stem symmetry is preserved because $\kappa$ is an algebra automorphism, and
holomorphy because multiplication is complex bilinear; so $*$ is associative
on $\SR$, and for power series
\begin{equation}
 \Bigl(\sum_{n\ge0}q^na_n\Bigr)*\Bigl(\sum_{n\ge0}q^nb_n\Bigr)
 =\sum_{n\ge0}q^n\sum_{k=0}^na_kb_{n-k}.
 \label{a:eq:convolution}
\end{equation}
Intrinsic functions have central stems: if $h$ is intrinsic then $h*f=f*h$ and
this equals the \emph{pointwise} product $h(q)f(q)$ --- but need not equal
$f(q)h(q)$.

\begin{proposition}[Evaluation of a regular product]
\label{a:prop:evaluation}
For $q\in\Om_D$,
\begin{equation}
 (f*g)(q)=
 \begin{cases}
  f(q)\,g\bigl(f(q)^{-1}qf(q)\bigr),& f(q)\ne0,\\
  0,& f(q)=0.
 \end{cases}
 \label{a:eq:evaluation}
\end{equation}
\end{proposition}

\begin{proof}
Write $q=x+Iy$, $u=f(q)=A+IB$; for $u\ne0$ put $I'=u^{-1}Iu$. Then
$u(C'+I'D')=uC'+IuD'=AC'-BD'+I(BC'+AD')$, which is \eqref{a:eq:starAB}; the
same expression vanishes when $u=0$.
\end{proof}

The argument of the second factor rotates within $[q]$: this is exactly why
roots of successive factors need not be roots of their product.

Extend quaternionic conjugation complex-linearly to stems by
$(a+\iota b)^c=\bar a+\iota\bar b$, an antiautomorphism, and set
$f^c=\I(F^c)$, $f^s=f*f^c=f^c*f$. The equality of the two products follows
from
\begin{align}
 \mathcal N_f(z)&:=F(z)F(z)^c=\abs A^2-\abs B^2+2\iota\Rea(A\bar B),
 \label{a:eq:normstem}\\
 f^s(x+Iy)&=\abs A^2-\abs B^2+2I\Rea(A\bar B).
 \label{a:eq:symmetrization}
\end{align}
The \emph{scalar norm stem} $\mathcal N_f$ is complex-scalar valued and
holomorphic, so $f^s$ is intrinsic; in a power series $f^c(q)=\sum q^n\bar a_n$
and $f^s(q)=\sum q^n\sum_k a_k\bar a_{n-k}$ with real coefficients. On the
real axis
\begin{equation}
 \mathcal N_f(x)=\abs{f(x)}^2\qquad(x\in D\cap\RR),
 \label{a:eq:normonreal}
\end{equation}
but for general $q$, $f^s(q)\ne\abs{f(q)}^2$: for $f(q)=q-\mathbf i$ one has
$f^s(q)=q^2+1$, and at $q=\mathbf j$ these are $0$ and $2$.

\begin{proposition}[Multiplicativity of the norm stem]
\label{a:prop:norm-multiplicative}
$\mathcal N_{f*g}=\mathcal N_f\mathcal N_g$ and $(f*g)^s=f^sg^s$; and
$f\not\equiv0$ implies $\mathcal N_f\not\equiv0$.
\end{proposition}

\begin{proof}
$\mathcal N_g$ is central, so
$(FG)(FG)^c=FGG^cF^c=\mathcal N_gFF^c=\mathcal N_g\mathcal N_f$. If
$\mathcal N_f\equiv0$ then \eqref{a:eq:normonreal} makes $f$ vanish on a real
interval and Theorem~\ref{a:thm:identity} gives $f\equiv0$.
\end{proof}

\begin{theorem}[Zero geometry and detection by symmetrization]
\label{a:thm:zeros}
On a nonreal sphere $S=x+y\SSS$ a slice function has no zero, exactly one
zero, or is identically zero there; and
\begin{equation}
 f^s\big|_S=0\iff f\text{ has a zero on }S,
 \label{a:eq:zero-detection}
\end{equation}
while for real $x$, $f^s(x)=0$ iff $f(x)=0$. For regular $f\not\equiv0$ the
zero-bearing conjugacy classes are locally finite in $\Om_D$; the zeros
consist of isolated real or nonreal points and isolated two-spheres.
\end{theorem}

\begin{proof}
On $S$ write $f=A+IB$. If $B=0$ the claim is immediate. If $B\ne0$ the only
possible zero is at $I=-AB^{-1}$, which is an imaginary unit exactly when
$\abs A=\abs B$ and $\Rea(A\bar B)=0$ --- its norm is $\abs A/\abs B$ and its
real part is $-\Rea(A\bar B)/\abs B^2$ --- and these are exactly the two real
conditions for \eqref{a:eq:symmetrization} to vanish. For local finiteness,
$\mathcal N_f$ is a nonzero scalar holomorphic function with discrete zeros;
each conjugate pair is one nonreal sphere and each real zero one real point,
and $q\mapsto\Rea q\pm\iota\abs{\Ima q}$ carries compacta to compacta.
\end{proof}

For example $q^2+1$ vanishes on all of $\SSS$. This is the structural fact
that forces the geometric unit of the later theory to be a conjugacy sphere
rather than a point.

\begin{proposition}[Division by a spherical factor]
\label{a:prop:spherical-division}
Let $p=x+Iy$ be nonreal and $\Delta_p(q)=q^2-2xq+x^2+y^2$. A regular $f$
vanishes identically on $[p]$ if and only if $f=\Delta_p*g=\Delta_p\,g$ for a
regular $g$ on $\Om_D$.
\end{proposition}

\begin{proof}
Vanishing on the sphere means $A=B=0$ there, i.e. the stem vanishes at
$z_0=x+\iota y$ and $\bar z_0$; divide the four complex components by
$(z-z_0)(z-\bar z_0)$, whose coefficients are real, so both apparent
singularities are removable and the quotient keeps the stem symmetry
\eqref{a:eq:stem-symmetry}.
\end{proof}

\subsection{Factorization, total multiplicity, the argument principle, Rouch\'e}
\label{a:sec:counting}

\begin{lemma}[Regular division on a ball]
\label{a:lem:division}
Let $f(q)=\sum_{n\ge0}q^na_n$ converge on $B(0,R)$ and $p\in B(0,R)$. There is
a unique regular $g$ on that ball with
\begin{equation}
 f=(q-p)*g+f(p),\qquad
 g(q)=\sum_{n\ge0}q^nb_n,\quad b_n=\sum_{k=n+1}^{\infty}p^{k-n-1}a_k .
 \label{a:eq:division}
\end{equation}
For a polynomial of degree $n\ge1$, $g$ has degree $n-1$.
\end{lemma}

\begin{proof}
For $\abs p<r<R$ Cauchy's estimate gives $\abs{a_k}\le M(r)r^{-k}$, so
$\abs{b_n}\le M(r)r^{-n-1}/(1-\abs p/r)$ and the series converges on
$\abs q<r$, hence on the ball. Directly, $b_{n-1}-pb_n=a_n$ for $n\ge1$ and
$-pb_0=a_0-f(p)$, which are the coefficients of \eqref{a:eq:division}. If two
quotients existed their difference $h$ would satisfy $(q-p)*h=0$, and
evaluation at real $x\ne p$ with Theorem~\ref{a:thm:identity} gives $h=0$.
\end{proof}

\begin{theorem}[Quaternionic fundamental theorem of algebra]
\label{a:thm:FTA}
Every nonconstant right-coefficient quaternionic polynomial has a zero in
$\HH$, and if $P(q)=q^na_n+\dots+a_0$ with $a_n\ne0$ then
\begin{equation}
 P=(q-p_1)*\cdots*(q-p_n)\,a_n
 \label{a:eq:factorization}
\end{equation}
for suitable $p_1,\dots,p_n\in\HH$.
\end{theorem}

\begin{proof}
The real-coefficient polynomial $P^s$ has degree $2n$ and leading coefficient
$\abs{a_n}^2$, so its scalar stem has a root $z$; a real root gives $P(z)=0$
and a nonreal root gives a zero on the corresponding sphere by
Theorem~\ref{a:thm:zeros}. Divide out one linear regular factor by
Lemma~\ref{a:lem:division} and repeat; leading coefficients identify the
remaining constant.
\end{proof}

This concerns the central-variable, right-coefficient polynomial algebra; it
is not a claim about expressions with $q$ interspersed among coefficients.

\begin{example}[Factors do not simply list roots]
\label{a:ex:factors}
$P=(q-\mathbf i)*(q-\mathbf j)=q^2-q(\mathbf i+\mathbf j)+\mathbf k$ satisfies
$P(\mathbf i)=0$ but $P(\mathbf j)=2\mathbf k\ne0$. Its symmetrization is
$(q^2+1)^2$, and on $\SSS$ its coefficient $B=-(\mathbf i+\mathbf j)$ is
nonzero, so Theorem~\ref{a:thm:zeros} allows only the one pointwise zero
$\mathbf i$. By contrast $(q-\mathbf i)*(q+\mathbf i)=q^2+1$ vanishes on an
entire sphere with the \emph{same} symmetrization.
\end{example}

For nonzero regular $f$ define the total multiplicity of a zero-bearing class
by
\begin{equation}
 \mu_f([p])=
 \begin{cases}
  \operatorname{ord}_{x+\iota y}\mathcal N_f,&p=x+Iy,\ y>0,\\[2pt]
  \tfrac12\operatorname{ord}_x\mathcal N_f,&p=x\in\RR .
 \end{cases}
 \label{a:eq:multiplicity}
\end{equation}
The two nonreal conjugate zeros have equal order, and the real order is even:
if $f=(q-x)^mu$ with $u(x)\ne0$ then
$\mathcal N_f(z)=(z-x)^{2m}\mathcal N_u(z)$ with
$\mathcal N_u(x)=\abs{u(x)}^2>0$. A simple isolated linear factor has total
multiplicity one and the spherical factor $\Delta_p$ has total multiplicity
two; this counts algebraic mass in a class, not the cardinality of the
pointwise zero set. (The letter $\mu$ is used here because $m$ is reserved
for the hole count of Table~\ref{a:tab:notation}.)

\begin{corollary}[Degree equals total multiplicity]
\label{a:cor:degreecount}
For a degree-$n$ polynomial, $\sum_{[p]}\mu_P([p])=n$.
\end{corollary}

\begin{proof}
By \eqref{a:eq:factorization} and Proposition~\ref{a:prop:norm-multiplicative},
$P^s(q)=\abs{a_n}^2\prod_{j=1}^n\Delta_{p_j}(q)$; each nonreal factor
contributes one to its class and each real factor is $(q-p_j)^2$, contributing
one under the real convention.
\end{proof}

Let $U'\Subset D$ be bounded, conjugation invariant, with finitely many
piecewise $C^1$ boundary curves (not necessarily connected), put
$\Gamma_{U'}=\{x+Iy:x+\iota y\in\partial U',\ I\in\SSS\}$, suppose $f$ has no
zero on $\Gamma_{U'}$, and let $\nu_f(U')$ be the sum of
\eqref{a:eq:multiplicity} over zero-bearing classes in $\Om_{U'}$.

\begin{theorem}[Argument principle]
\label{a:thm:argument}
Under these assumptions
\begin{equation}
 \nu_f(U')=\frac1{4\pi\iota}\int_{\partial U'}
   \frac{\mathcal N_f'(z)}{\mathcal N_f(z)}\dd z .
 \label{a:eq:argument}
\end{equation}
\end{theorem}

\begin{proof}
Theorem~\ref{a:thm:zeros} makes $\mathcal N_f$ nonzero on $\partial U'$. Near a
zero of order $\mu$ write $\mathcal N_f=(z-z_0)^\mu v$ with $v(z_0)\ne0$; the
ordinary argument principle counts complex multiplicities, which occur in
conjugate pairs except at real zeros, whose orders are even. Halving the
complex count is exactly \eqref{a:eq:multiplicity}.
\end{proof}

The integrand is scalar complex valued; a pointwise quaternionic expression
such as $f'(q)f(q)^{-1}$ cannot replace it in this proof.

\begin{theorem}[Rouch\'e on circular boundaries]
\label{a:thm:rouche}
If $f,g$ are regular on $\Om_D$ and $\abs{g(q)-f(q)}<\abs{f(q)}$ for all
$q\in\Gamma_{U'}$, then $\nu_g(U')=\nu_f(U')$.
\end{theorem}

\begin{proof}
Put $h_t=f+t(g-f)$ for $t\in[0,1]$; the strict inequality gives
$\abs{h_t}\ge\abs f-t\abs{g-f}>0$ on $\Gamma_{U'}$, so $\mathcal N_{h_t}$ has
no zero on $\partial U'$. The integral \eqref{a:eq:argument} depends
continuously on $t$ and is integer valued, hence constant.
\end{proof}

The hypothesis is imposed on \emph{all} quaternionic points of each boundary
sphere, not only on one complex contour; that is what makes the
zero-detection argument valid.

\begin{corollary}[Polynomial localization]
\label{a:cor:localization}
If $P(q)=q^n+\sum_{k<n}q^ka_k$ and $\sum_{k<n}R^k\abs{a_k}<R^n$, then all
zeros lie in $\abs q<R$ with total multiplicity $n$; in particular every zero
satisfies $\abs q\le1+\max_{k<n}\abs{a_k}$.
\end{corollary}

\begin{proof}
Apply Theorem~\ref{a:thm:rouche} to $q^n$ and $P$ on $\abs q=R$ together with
Corollary~\ref{a:cor:degreecount}; with $M=\max_{k<n}\abs{a_k}$ and $R>1+M$,
$\sum_{k<n}R^k\abs{a_k}\le M(R^n-1)/(R-1)<R^n$, and let $R\downarrow1+M$.
\end{proof}

\subsection{The regular reciprocal, minimum modulus, and open mappings}
\label{a:sec:reciprocal}

Let $Z(f^s)$ be the union of zero-bearing classes of a nonzero regular $f$.
On $\Om_D\setminus Z(f^s)$ put
\begin{equation}
 f^{-*}=(f^s)^{-1}f^c ,
 \label{a:eq:reciprocal}
\end{equation}
the first factor being intrinsic and regular; stem multiplication gives
$f*f^{-*}=f^{-*}*f=1$ there. For instance
$(q-p)^{-*}=\Delta_p(q)^{-1}(q-\bar p)$, whose excluded set is the whole
sphere $[p]$ even though $q-p$ has only the one pointwise zero $p$.

\begin{proposition}[The conjugation change of variable]
\label{a:prop:T}
On $\Om_D\setminus Z(f^s)$ set $T_f(q)=f^c(q)^{-1}qf^c(q)$. This is a
real-analytic diffeomorphism of that domain onto itself preserving each
conjugacy class, with inverse $T_{f^c}$, and
\begin{equation}
 f^{-*}(q)=f(T_f(q))^{-1},\qquad
 (f^{-*}*g)(q)=f(T_f(q))^{-1}g(T_f(q)).
 \label{a:eq:quotient-eval}
\end{equation}
\end{proposition}

\begin{proof}
Nonvanishing of $f^s(q)$ gives nonvanishing of $f^c(q)$ by
Theorem~\ref{a:thm:zeros}. Let $c=f^c(q)$, $r=T_f(q)$, $n=f^s(q)$;
Proposition~\ref{a:prop:evaluation} applied to $f^c*f$ gives $n=cf(r)$. Since
$n$ is intrinsic it commutes with $q$, so
$T_{f^c}(r)=f(r)^{-1}rf(r)=n^{-1}c(c^{-1}qc)c^{-1}n=q$. Also
$n^{-1}c=(cf(r))^{-1}c=f(r)^{-1}$, which is the first identity; for the
second, the conjugating value is $u=f^{-*}(q)=n^{-1}c$ and
$u^{-1}qu=c^{-1}qc=T_f(q)$.
\end{proof}

So the reciprocal is a pointwise inverse only \emph{after} a
sphere-preserving change of variable.

\begin{theorem}[Minimum modulus]
\label{a:thm:min}
If $f\in\SR(\Om_D)$ and $\abs f$ has a local minimum at $p$, then $f(p)=0$ or
$f$ is constant.
\end{theorem}

\begin{proof}
Assume $f(p)\ne0$. First, $[p]$ contains no zero: on a nonreal sphere
$\abs{A+IB}^2=\abs A^2+\abs B^2-2\Rea(A\bar BI)$ is an affine real function of
$I\in\SSS$, and a local minimum of an affine function on a sphere is global,
so the value at $p$ is no larger than anywhere on $[p]$. Next,
$D'=D\setminus Z(\mathcal N_f)$ is connected (perturb a polygonal path around
a locally finite set), conjugation invariant and still meets $\RR$, so
$f^{-*}$ is regular on $\Om_{D'}$. By Proposition~\ref{a:prop:T},
$\abs{f^{-*}}=1/\abs{f\circ T_f}$ has a local maximum, so
Theorem~\ref{a:thm:max} makes $f^{-*}$ a nonzero constant $c$; then
$f(q)c=1$, and Theorem~\ref{a:thm:identity} propagates this.
\end{proof}

Define the \emph{degenerate set}
$\Sigma_f=\bigcup\{x+y\SSS: y>0,\ f\text{ constant on }x+y\SSS\}$, the union
of spheres on which $B(x,y)=0$.

\begin{theorem}[Open mapping with spherical qualifications]
\label{a:thm:open}
Let $f\in\SR(\Om_D)$ be nonconstant.
\begin{enumerate}[label=\textup{(\roman*)}]
\item At every $p\notin\Sigma_f$ the image of every sufficiently small
Euclidean ball centred at $p$ contains a neighbourhood of $f(p)$;
consequently $f$ restricted to $\Om_D\setminus\overline{\Sigma_f}$ is open.
\item If $V\subseteq\Om_D$ is open and axially symmetric then $f(V)$ is open;
in particular $f(\Om_D)$ is open.
\end{enumerate}
\end{theorem}

\begin{proof}
For (i), $h=f-f(p)$ has an isolated zero at $p$ (a real zero is isolated; a
nonreal zero is the unique zero on a nondegenerate sphere, and zero-bearing
classes are locally finite). Take a closed ball whose boundary carries no
zero of $h$ and put $m_0=\min_{\partial B(p,r)}\abs h>0$. If
$\abs{w-f(p)}<m_0/2$ then $\abs{f-w}>m_0/2$ on the boundary while
$\abs{f(p)-w}<m_0/2$, so $\abs{f-w}$ attains an interior minimum, which by
Theorem~\ref{a:thm:min} is zero. For (ii), choose a relatively compact
axially symmetric neighbourhood of the whole class $[p]$ inside $V$ whose
boundary avoids the zeros of $f-f(p)$, and repeat.
\end{proof}

\begin{example}[Why the unqualified open mapping theorem fails]
\label{a:ex:q2}
The entire regular $f(q)=q^2$ is not open at $p=\mathbf i$. If $q^2=-1+t\mathbf j$
with $t\ne0$ and $q=x+v$, then $\Ima(q^2)=2xv=t\mathbf j$, so $x\ne0$ and $v$ is
parallel to $\mathbf j$; in particular $\abs{q-\mathbf i}\ge1$. Hence
$f(B(\mathbf i,1/2))$ contains $-1$ but no point $-1+t\mathbf j$, $t\ne0$. The
missing behaviour sits on a sphere that $f$ collapses to a point. The closure
in (i) is also needed: for $q^2$ the degenerate spheres form the nonzero
imaginary hyperplane, whose closure contains $0$, and removing only the
nonreal ones leaves $0$ while deleting all preimages of negative reals.
\end{example}

\subsection{Schwarz--Pick and an intrinsic Riemann mapping theorem}
\label{a:sec:geometry}

For $a\in\mathbb B$ define the regular M\"obius map
\begin{equation}
 \mathcal M_a(q)=(1-q\bar a)^{-*}*(q-a).
 \label{a:eq:Mobius}
\end{equation}
Its denominator has no pointwise zero in $\mathbb B$ because
$\abs{q\bar a}<1$, so Theorem~\ref{a:thm:zeros} makes the regular reciprocal
well defined on the whole ball; the two factors commute under $*$ because
their coefficients lie in one commutative plane.

\begin{lemma}[Ball and boundary estimates]
\label{a:lem:Mobius}
$\mathcal M_a$ maps $\mathbb B$ into $\mathbb B$ with the single zero $a$, and
for $\abs a<r<1$,
\begin{equation}
 \min_{\abs q=r}\abs{\mathcal M_a(q)}=\frac{r-\abs a}{1-r\abs a}.
 \label{a:eq:Mobius-min}
\end{equation}
\end{lemma}

\begin{proof}
With $h(q)=1-q\bar a$, \eqref{a:eq:quotient-eval} gives
$\mathcal M_a(q)=(1-\varpi\bar a)^{-1}(\varpi-a)$ with $\varpi=T_h(q)$, and the
elementary identity
\begin{equation}
 \abs{1-\varpi\bar a}^2-\abs{\varpi-a}^2=(1-\abs\varpi^2)(1-\abs a^2)
 \label{a:eq:ballidentity}
\end{equation}
bounds the norm by one. $T_h$ is a bijection and $h^c(a)=1-a^2$ commutes with
$a$, so $T_h(a)=a$. For $\abs\varpi=r$ and $t=\Rea(\varpi\bar a)$,
$\abs{\varpi-a}^2/\abs{1-\varpi\bar a}^2=(r^2+\abs a^2-2t)/(1+r^2\abs a^2-2t)$
is decreasing in $t$ (numerator minus denominator is $-(1-r^2)(1-\abs a^2)$),
so the minimum is at $t=r\abs a$.
\end{proof}

\begin{theorem}[Schwarz lemma with an arbitrary zero]
\label{a:thm:schwarz-zero}
If $F:\mathbb B\to\mathbb B$ is regular with $F(a)=0$, then
$\abs{F(q)}\le\abs{\mathcal M_a(q)}$ on $\mathbb B$, with equality at some
$q\ne a$ exactly when $F=\mathcal M_au$ for a constant $\abs u=1$.
\end{theorem}

\begin{proof}
Regular division gives $F=(q-a)*G$; put $H=(1-q\bar a)*G$, regular on
$\mathbb B$, so $F=\mathcal M_a*H$ and $H=\mathcal M_a^{-*}*F$ off the
zero-bearing sphere. For $r>\abs a$, quotient evaluation and
\eqref{a:eq:Mobius-min} give $\abs{H}\le(1-r\abs a)/(r-\abs a)$ on
$\abs q=r$; maximum modulus and $r\uparrow1$ give $\abs H\le1$, and
Proposition~\ref{a:prop:evaluation} finishes. Equality forces $\abs H=1$ at an
interior point.
\end{proof}

\begin{theorem}[Regular Schwarz--Pick inequality]
\label{a:thm:pick}
Let $f:\mathbb B\to\mathbb B$ be regular, $a\in\mathbb B$, $b=f(a)$, and
$\Phi_{f,b}=(1-f\bar b)^{-*}*(f-b)$. Then
$\abs{\Phi_{f,b}(q)}\le\abs{\mathcal M_a(q)}$ on $\mathbb B$, with equality at
some $q\ne a$ exactly when $\Phi_{f,b}=\mathcal M_au$, $\abs u=1$. If $a$ is
real,
\begin{equation}
 \frac{\abs{\partial_cf(a)}}{1-\abs{f(a)}^2}\le\frac1{1-a^2}.
 \label{a:eq:real-pick}
\end{equation}
\end{theorem}

\begin{proof}
$h=1-f\bar b$ is regular and zero free since $\abs{f(q)}\abs b<1$, so its
reciprocal is regular; quotient evaluation and \eqref{a:eq:ballidentity} give
$\abs{\Phi_{f,b}}<1$. Evaluating $h*\Phi_{f,b}=f-b$ at $a$ gives
$(1-\abs b^2)\Phi_{f,b}(a)=0$, because $h(a)=1-\abs b^2$ is a nonzero real
scalar and does not rotate the argument; now apply
Theorem~\ref{a:thm:schwarz-zero}. For real $a$, restrict
$h*\Phi_{f,b}=f-b$ to $\RR$ and differentiate: $\Phi_{f,b}(a)=0$ gives
$\partial_c\Phi_{f,b}(a)=\partial_cf(a)/(1-\abs b^2)$, while
$\mathcal M_a(x)=(x-a)/(1-ax)$ has $\mathcal M_a'(a)=1/(1-a^2)$.
\end{proof}

At nonreal base points the real differential has the two pieces of
\eqref{a:eq:differential}, and the simple differential inequality
\eqref{a:eq:real-pick} is deliberately \emph{not} asserted there.

\begin{theorem}[Intrinsic Riemann mapping theorem]
\label{a:thm:riemann}
Let $D\subsetneq\CC$ be simply connected, open, conjugation invariant, and let
$c\in D\cap\RR$. There is a unique intrinsic slice-regular bijection
$\varphi:\Om_D\to\mathbb B$ with $\varphi(c)=0$ and $\partial_c\varphi(c)>0$
real; its inverse is intrinsic and slice regular.
\end{theorem}

\begin{proof}
The normalized scalar Riemann mapping theorem gives a unique biholomorphism
$\phi:D\to\{\abs z<1\}$ with $\phi(c)=0$, $\phi'(c)>0$; since
$z\mapsto\overline{\phi(\bar z)}$ has the same normalization, uniqueness gives
$\phi(\bar z)=\overline{\phi(z)}$, and likewise for $\psi=\phi^{-1}$. Both
therefore have intrinsic slice extensions, which on every $\CC_I$ are the
scalar maps under $x+\iota y\leftrightarrow x+Iy$, so they are mutually
inverse section by section. For uniqueness, an intrinsic bijection with this
normalization restricts to a biholomorphism on each section --- a nonreal
point cannot map to a real value, since its whole sphere would then have the
same image --- so scalar uniqueness on one section and
\eqref{a:eq:representation} finish.
\end{proof}

This asserts a slice-analytic equivalence, not conformal uniformization in
$\RR^4$: for an intrinsic function induced by $\phi=u+\iota v$ the two stretch
factors in \eqref{a:eq:differential} are $\abs{\phi'(z)}$ and
$\abs{v(x,y)/y}$, and for $q^2$ at $1+\mathbf i$ they are $2\sqrt2$ and $2$.
Note also that ordinary composition of regular functions need not be regular:
$q\mapsto q\mathbf i$ and $q\mapsto q^2$ are regular but their composite is
$-\abs q^2$ on $\CC_{\mathbf j}$. Composition works when the inner function is
intrinsic.

\subsection{Laurent expansions, residues, and essential singularities}
\label{a:sec:slice-laurent}

The statements here concern isolated singularities at \emph{real} points,
which is what makes ordinary radial annuli and ordinary integer powers the
right objects. Nonreal singularities require an entire conjugacy sphere and
are the subject of Part~VII, whose residue is a different object.

\begin{theorem}[Laurent expansion on a spherical annulus]
\label{a:thm:laurent}
Let $f$ be slice regular on $A_{r,R}=\{r<\abs q<R\}$, $0\le r<R\le\infty$.
Then $f(q)=\sum_{n\in\ZZ}q^na_n$ with unique $a_n\in\HH$, converging
absolutely and uniformly on compact subannuli, and for $r<\varrho<R$,
\begin{equation}
 a_n=\frac1{2\pi}\int_{\abs s=\varrho,\,s\in\CC_I}s^{-n-1}\dd s_I\,f(s),
 \qquad \abs{a_n}\le M_I(\varrho)\varrho^{-n},
 \label{a:eq:laurentcoef}
\end{equation}
independently of $\varrho$ and $I$.
\end{theorem}

\begin{proof}
Apply the scalar Laurent theorem to the four complex coordinates of the stem
$F$; the symmetry \eqref{a:eq:stem-symmetry} and uniqueness of Laurent
coefficients make each coefficient $\kappa$-fixed, hence in $\HH$. Evaluating
the stem series gives the expansion, and the splitting lemma gives
\eqref{a:eq:laurentcoef}. Coefficient estimates on a smaller radius control
the negative tail and on a larger radius the positive tail, both
geometrically.
\end{proof}

Translation by a real centre $c$ replaces $q^n$ by $(q-c)^n$ verbatim.

\begin{corollary}[Classification and Riemann removability at a real point]
\label{a:cor:removable-slice}
For $f$ regular on $0<\abs q<R$ the singularity at $0$ is removable if all
negative coefficients vanish, a pole if finitely many but at least one are
nonzero, and essential otherwise. A bounded singularity is removable, and for
a pole of order $\mathfrak m$,
\begin{equation}
 \abs{f(q)}=\abs q^{-\mathfrak m}\bigl(\abs{a_{-\mathfrak m}}+o(1)\bigr)
 \qquad(q\to0).
 \label{a:eq:polegrowth}
\end{equation}
\end{corollary}

\begin{proof}
Under a local bound $M$, \eqref{a:eq:laurentcoef} gives
$\abs{a_{-\mathfrak m}}\le M\varrho^{\mathfrak m}\to0$. For a pole,
$q^{\mathfrak m}f(q)$ extends regularly with nonzero value
$a_{-\mathfrak m}$, and multiplicativity of the norm gives
\eqref{a:eq:polegrowth}.
\end{proof}

\begin{theorem}[Residue theorem for isolated real singularities]
\label{a:thm:residue-slice}
Let $f$ be regular near $\overline{\Om_{U'}}\setminus\{c_1,\dots,c_N\}$ with
real interior $c_j$ and $U'$ having piecewise smooth positively oriented
boundary, and let $\Res_{c_j}f$ be the coefficient of $(q-c_j)^{-1}$. Then
\begin{equation}
 \frac1{2\pi}\int_{\partial U'_I}\dd s_I\,f(s)=\sum_{j=1}^N\Res_{c_j}f .
 \label{a:eq:residue-slice}
\end{equation}
\end{theorem}

\begin{proof}
Remove small disks around the $c_j$ in the $I$-section and apply Cauchy's
theorem to both splitting components; on each small circle the Laurent series
converges uniformly, all powers integrate to zero except $(s-c_j)^{-1}$ whose
normalized integral is $1$, and the right coefficients stay on the right.
\end{proof}

\begin{theorem}[Casorati--Weierstrass at a real essential singularity]
\label{a:thm:casorati}
If $f$ has an essential singularity at $0\in\RR$ then $f(\{0<\abs q<r\})$ is
dense in $\HH$ for every $0<r<R$.
\end{theorem}

\begin{proof}
If a ball $B(a,\varepsilon)$ were omitted, $g=f-a$ would be zero free with
$\abs g\ge\varepsilon$, so
$\abs{g^{-*}(q)}=\abs{g(T_g(q))}^{-1}\le\varepsilon^{-1}$ because $T_g$
preserves each radius. By removability $h=g^{-*}$ extends regularly and is not
identically zero. If $h(0)\ne0$ then $h^{-*}$ is regular near $0$ and agrees
with $g$, so the singularity is removable; if $h(0)=0$ write $h=q^{\mathfrak m}u$
with $u(0)\ne0$, whence $g=h^{-*}=q^{-\mathfrak m}u^{-*}$ has a pole. Both
contradict essentiality.
\end{proof}

\subsection{Normal families and Hurwitz}
\label{a:sec:normal}

\begin{theorem}[Weierstrass convergence and Montel]
\label{a:thm:montel}
A locally uniform limit of slice-regular functions on $\Om_D$ is slice
regular, and every family bounded uniformly on compacta is normal.
\end{theorem}

\begin{proof}
Split each restriction to $\CC_I$ into two complex components: scalar
Weierstrass convergence makes the limit components holomorphic, and their
slice extension equals the limit by \eqref{a:eq:representation}. For
normality, scalar Montel on both components with successive extraction gives
a subsequence converging on compacta of that section; to propagate to a
compact $K\subset\Om_D$, use the compact two-point projection
$K_I=\{\Rea q\pm I\abs{\Ima q}:q\in K\}$, and note that the coefficients in
\eqref{a:eq:representation} have norm at most one, so the uniform difference
of two members on $K$ is at most twice their uniform difference on $K_I$.
\end{proof}

\begin{theorem}[Hurwitz]
\label{a:thm:hurwitz}
If $f_n\in\SR(\Om_D)$ have no zeros and converge locally uniformly to $f$,
then $f\equiv0$ or $f$ is zero free.
\end{theorem}

\begin{proof}
The representation formula recovers $A_n,B_n$ from values on a fixed section,
so \eqref{a:eq:normstem} gives $\mathcal N_{f_n}\to\mathcal N_f$ locally
uniformly on $D$, and each $\mathcal N_{f_n}$ is zero free by
Theorem~\ref{a:thm:zeros}. Scalar Hurwitz then gives the dichotomy, and
\eqref{a:eq:normonreal} with Theorem~\ref{a:thm:identity} handles the
degenerate case.
\end{proof}

Viewing $f_I$ as a holomorphic map into $\CC^2$ would not suffice: the maps
$z\mapsto(z,1/n)$ are nonvanishing but converge to $(z,0)$, which vanishes.
Zero-freeness on \emph{all} spheres, encoded by $\mathcal N_f$, is the extra
information.

\subsection{Cauchy--Fueter regularity and the four-dimensional Cauchy formula}
\label{a:sec:fueter-cauchy}

Now change the notion of regularity. Let $\Om\subset\HH\cong\RR^4$ be an
arbitrary connected open set; no axial symmetry is assumed. A $C^1$ function
$f:\Om\to\HH$ is \emph{left Cauchy--Fueter regular} (left monogenic) when
$\Dop f=0$; for a factor on the other side one uses the right-acting operator
$g\Dop=\partial_{x_0}g+(\partial_{x_1}g)\mathbf i+(\partial_{x_2}g)\mathbf j
+(\partial_{x_3}g)\mathbf k$. Anticommutation gives
$\overline{\Dop}\Dop=\Dop\overline{\Dop}=\lap$ as identities of
constant-coefficient operators. Note at once that $\Dop q=1-1-1-1=-2$: the
identity function, the basic example of slice analysis, is not Fueter regular.

\begin{lemma}[Quaternionic Green identity]
\label{a:lem:green}
Let $\Om$ be bounded with $C^1$ boundary, $n=n_0+\mathbf in_1+\mathbf jn_2
+\mathbf kn_3$ the outward unit normal, and $f,g$ of class $C^1$ near
$\overline\Om$. Then
\begin{equation}
 \int_{\partial\Om}g(y)n(y)f(y)\dd S_y
 =\int_\Om\bigl((g\Dop)(y)f(y)+g(y)(\Dop f)(y)\bigr)\dd V_y .
 \label{a:eq:green}
\end{equation}
\end{lemma}

\begin{proof}
Apply the divergence theorem componentwise to the four-vector with
quaternionic components $ge_\mu f$: its divergence is
$\sum_\mu(\partial_{x_\mu}g)e_\mu f+\sum_\mu ge_\mu\partial_{x_\mu}f
=(g\Dop)f+g(\Dop f)$, which fixes the order of the three factors on the
boundary.
\end{proof}

In particular a Fueter-regular function satisfies the Cauchy theorem
\begin{equation}
 \int_{\partial\Om}n(y)f(y)\dd S_y=0,
 \label{a:eq:fueter-cauchy-theorem}
\end{equation}
the integration cycle being three-dimensional. The fundamental kernel is
\begin{equation}
 E(q)=\frac{\bar q}{2\pi^2\abs q^4},\qquad q\ne0,
 \label{a:eq:fueter-kernel}
\end{equation}
left and right Fueter regular away from $0$: indeed $\Dop\bar q=4$ and
$\sum_\mu e_\mu\bar q\,\partial_{x_\mu}\abs q^{-4}=-4q\bar q\abs q^{-6}
=-4\abs q^{-4}$ cancels it, and the right calculation is identical with the
order reversed. On $\abs q=r$ the outward normal is $q/r$, so
\begin{equation}
 E(q)n(q)=n(q)E(q)=\frac1{2\pi^2r^3},
 \qquad \int_{\abs q=r}E(q)n(q)\dd S_q=1,
 \label{a:eq:kernel-flux}
\end{equation}
since the unit three-sphere has area $2\pi^2$.

\begin{theorem}[Cauchy--Pompeiu and Cauchy--Fueter]
\label{a:thm:fueter-cauchy}
Under the boundary hypotheses of Lemma~\ref{a:lem:green}, for $f\in C^1$ near
$\overline\Om$ and $x\in\Om$,
\begin{equation}
 f(x)=\int_{\partial\Om}E(y-x)n(y)f(y)\dd S_y
      -\int_\Om E(y-x)(\Dop f)(y)\dd V_y ,
 \label{a:eq:pompeiu}
\end{equation}
and consequently, if $\Dop f=0$,
\begin{equation}
 f(x)=\int_{\partial\Om}E(y-x)n(y)f(y)\dd S_y .
 \label{a:eq:fueter-cauchy}
\end{equation}
\end{theorem}

\begin{proof}
Apply \eqref{a:eq:green} with $g(y)=E(y-x)$ on
$\Om\setminus\overline{B(x,\varepsilon)}$, where $g\Dop=0$; the inner normal
is the negative of the small ball's outward normal. The inner boundary term
tends to $f(x)$ by \eqref{a:eq:kernel-flux} and continuity, and the volume
integral converges because $\abs{E(y-x)}=O(\abs{y-x}^{-3})$ is locally
integrable in four dimensions.
\end{proof}

Equation \eqref{a:eq:fueter-cauchy} represents functions already known to be
regular; it is not a claim that arbitrary boundary values admit a regular
extension.

\subsection{Fueter analogues of the basic analytic theorems}
\label{a:sec:fueter-consequences}

\begin{proposition}
\label{a:prop:fueter-analytic}
Every $C^1$ Fueter-regular function is real analytic, and each of its four
real components is harmonic.
\end{proposition}

\begin{proof}
Use \eqref{a:eq:fueter-cauchy} on a closed ball inside the domain; the kernel
is smooth in $x$ uniformly for $y$ on the boundary, so all derivatives pass
under the integral. For analyticity near the centre $a$, expand
\[
 \abs{y-x}^{-4}=\abs{y-a}^{-4}
 \Bigl(1+\frac{-2\langle y-a,x-a\rangle+\abs{x-a}^2}{\abs{y-a}^2}\Bigr)^{-2},
\]
whose binomial series converges absolutely and uniformly on the integration
sphere for $\abs{x-a}$ small; after multiplication by $\overline{y-x}$ it is a
convergent power series in the four coordinates of $x-a$. Finally
$\lap f=\overline{\Dop}(\Dop f)=0$.
\end{proof}

\begin{theorem}[Cauchy derivative estimates]
\label{a:thm:fueter-estimates}
If $f$ is Fueter regular near $\overline{B(a,R)}$ and
$M_R=\max_{\abs{q-a}=R}\abs{f(q)}$, then for every multi-index $\alpha$
\begin{equation}
 \abs{\partial^\alpha f(a)}\le C_\alpha\frac{M_R}{R^{\abs\alpha}},
 \qquad C_\alpha=2\pi^2\max_{\abs\omega=1}\abs{\partial^\alpha E(\omega)},
 \label{a:eq:fueter-estimate}
\end{equation}
with $C_0=1$ and $C_\alpha=5$ admissible for each first partial derivative.
\end{theorem}

\begin{proof}
Differentiate \eqref{a:eq:fueter-cauchy}; $E$ is homogeneous of degree $-3$,
so $\partial^\alpha E$ is homogeneous of degree $-3-\abs\alpha$, and the area
factor is $2\pi^2R^3$. Explicitly
$\partial_{x_\mu}E(q)=(2\pi^2)^{-1}(\bar e_\mu\abs q^{-4}
-4x_\mu\bar q\abs q^{-6})$, of norm at most $5/(2\pi^2\abs q^4)$. No
sharpness is asserted for these constants.
\end{proof}

\begin{corollary}[Liouville and polynomial growth]
\label{a:cor:fueter-liouville}
A bounded entire Fueter-regular function is constant; if
$\abs{f(q)}\le C(1+\abs q)^\beta$ with $\beta\ge0$ then $f$ is a polynomial in
$x_0,x_1,x_2,x_3$ of total degree at most $\lfloor\beta\rfloor$ satisfying
$\Dop f=0$.
\end{corollary}

\begin{proof}
Let $R\to\infty$ in \eqref{a:eq:fueter-estimate}: every derivative of order
$>\beta$ vanishes at every centre, and the real Taylor expansion terminates.
\end{proof}

The conclusion differs from Theorem~\ref{a:thm:liouville-slice}: here a
polynomial means a polynomial in the four real coordinates, not necessarily a
right-coefficient polynomial in $q$. The axially monogenic entire functions of
polynomial growth are classified in Part~IX by the polynomials
$\mathcal P_k$.

\begin{theorem}[Mean value and maximum modulus]
\label{a:thm:fueter-max}
If $f$ is Fueter regular near $\overline{B(a,R)}$ then
\begin{equation}
 f(a)=\frac1{2\pi^2R^3}\int_{\abs{y-a}=R}f(y)\dd S_y ,
 \label{a:eq:fueter-mean}
\end{equation}
and if $\abs f$ has a local maximum in a connected domain then $f$ is
constant.
\end{theorem}

\begin{proof}
\eqref{a:eq:fueter-mean} is \eqref{a:eq:fueter-cauchy} with
\eqref{a:eq:kernel-flux}. If $\abs f\le M=\abs{f(a)}$ near $a$ with $M>0$, the
real-linear functional $\ell(w)=\Rea(w\overline{f(a)})/M$ satisfies
$\ell(w)\le\abs w$, and on every small sphere the average of $\ell(f)$ is $M$
while $\ell(f)\le\abs f\le M$; continuity forces equality pointwise, and
equality in real Cauchy--Schwarz with $\abs f\le M$ gives $f\equiv f(a)$.
Real analyticity propagates.
\end{proof}

\begin{theorem}[Fueter--Morera]
\label{a:thm:fueter-morera}
A continuous $f:\Om\to\HH$ is Fueter regular if and only if
$\int_{\partial B}n(y)f(y)\dd S_y=0$ for every ball with closure in $\Om$.
\end{theorem}

\begin{proof}
Necessity is \eqref{a:eq:fueter-cauchy-theorem}. For sufficiency mollify:
by Fubini the flux of $f_\varepsilon$ through a ball is the average of fluxes
of $f$ through translated balls, hence zero, so the divergence theorem gives
$\int_B\Dop f_\varepsilon\dd V=0$ for every such ball and shrinking balls
gives $\Dop f_\varepsilon=0$. Apply \eqref{a:eq:fueter-cauchy} to
$f_\varepsilon$ on a fixed smaller ball and pass to the limit using uniform
boundary convergence; the limit integral is regular inside because
$\Dop_xE(y-x)=0$ there.
\end{proof}

\subsection{Removable singularities, the flux residue, and two counterexamples}
\label{a:sec:fueter-residue}

\begin{theorem}[Sharp growth threshold for removability at a point]
\label{a:thm:fueter-removable}
Suppose $f$ is Fueter regular on $B(a,R)\setminus\{a\}$ and
\begin{equation}
 \sup_{\abs{q-a}=r}\abs{f(q)}=o(r^{-3})\qquad(r\downarrow0).
 \label{a:eq:fueter-removable-growth}
\end{equation}
Then $f$ extends uniquely to a Fueter-regular function on $B(a,R)$; in
particular a bounded isolated singularity is removable. The exponent $3$
cannot be replaced by a statement allowing all $O(r^{-3})$ growth.
\end{theorem}

\begin{proof}
Apply \eqref{a:eq:fueter-cauchy} on the punctured ball with outer radius
$R_0<R$ and inner radius $\varepsilon$ at a fixed $x\ne a$. For small
$\varepsilon$ the kernel $E(y-x)$ is bounded on the inner sphere, so that
contribution is at most a constant times
$\varepsilon^3\sup_{\abs{y-a}=\varepsilon}\abs{f(y)}\to0$. Hence $f(x)$ is the
outer boundary integral alone, which is regular throughout $B(a,R_0)$.
Sharpness: $E(q-a)c$ with $c\ne0$ is a nonremovable example of exact order
$\abs{q-a}^{-3}$.
\end{proof}

The exponent $3$ here belongs to an \emph{isolated point} in $\RR^4$. Part~VII
proves a threshold with exponent $1$ at a singular two-sphere; the two
exponents describe singular sets of real codimension four and two and are
never to be confused.

\begin{definition}[Flux residue]
\label{a:def:flux-residue}
For a Fueter-regular function with an isolated singularity at $a$ put
\begin{equation}
 \Res^{F}_af=\int_{\abs{q-a}=r}n(q)f(q)\dd S_q ,
 \label{a:eq:fueter-residue}
\end{equation}
for any sufficiently small $r>0$, the normal pointing away from $a$. The
normalization keeps the $2\pi^2$ in the kernel, not in the residue.
\end{definition}

The value is independent of $r$ by the Green identity on an annulus.

\begin{theorem}[Fueter residue theorem]
\label{a:thm:fueter-residue}
Let $\Om$ be a bounded $C^1$ domain and $f$ Fueter regular near
$\overline\Om$ except at finitely many interior points $a_1,\dots,a_N$, where
it is regular in punctured neighbourhoods. Then
\begin{equation}
 \int_{\partial\Om}n(q)f(q)\dd S_q=\sum_{j=1}^N\Res^F_{a_j}f ,
 \label{a:eq:fueter-residue-theorem}
\end{equation}
and $\Res^F_a(E(q-a)c)=c$.
\end{theorem}

\begin{proof}
Remove disjoint small balls and apply
\eqref{a:eq:fueter-cauchy-theorem}; the inner normals are the negatives of
those in \eqref{a:eq:fueter-residue}. The last assertion follows from
$nE=1/(2\pi^2r^3)$ on the sphere.
\end{proof}

The next two examples are the reason every quaternionic statement in this
document has the shape it has. Both are needed later and neither is reproved
anywhere else.

\begin{example}[A zero flux residue without removability]
\label{a:ex:dxE}
$\partial_{x_0}E$ is Fueter regular away from $0$, is nonzero, is homogeneous
of degree $-4$, and is singular at $0$. Since $E$ is odd, $\partial_{x_0}E$ is
even, whereas $n$ is odd on a centred sphere; hence its flux, and therefore
its residue \eqref{a:eq:fueter-residue}, vanishes. So
\[
 \Res^F_0(\partial_{x_0}E)=0
 \quad\text{while the singularity at }0\text{ is not removable.}
\]
A single flux residue does not encode the whole principal part of a
singularity. This is the isolated-point ancestor of the fact, proved in
Part~VIII, that the total spherical flux $8\pi v(ua+b)$ misses the slope
residue $a$: it is exactly why the spherical residue of Part~VII must be a
\emph{pair} $(a,b)$, defined by two period integrals rather than by one flux,
and why a genuinely second conserved flux with a nonconstant regular weight is
needed to recover $a$.
\end{example}

\begin{example}[Vanishing on a plane: the identity and open mapping theorems fail]
\label{a:ex:fueter-plane}
The function
\begin{equation}
 L(q)=x_1-\mathbf ix_0
 \label{a:eq:fueter-plane}
\end{equation}
is Fueter regular, since $\Dop L=-\mathbf i+\mathbf i=0$. It vanishes on the
entire plane $\{x_0=x_1=0\}$ without being identically zero, and its image
lies inside the single slice $\CC_{\mathbf i}$ and hence has empty interior in
$\HH$. Therefore \emph{neither} the accumulation-point identity theorem
(contrast Theorem~\ref{a:thm:identity}) \emph{nor} the nonconstant open
mapping theorem (contrast Theorem~\ref{a:thm:open}) holds for arbitrary
Cauchy--Fueter-regular functions. The appropriate elementary uniqueness
statement is vanishing on a nonempty open set, which does force global
vanishing by Proposition~\ref{a:prop:fueter-analytic}.
\end{example}

\subsection{The bridge: axial identities, the Vekua system, and Fueter's
mapping theorem}
\label{a:sec:bridge}

The two theories are connected by a second-order operator. Everything in this
subsection is used by every later part; it is stated once here.

Away from the real axis write $q=x+Iy$ with $y=\abs{\Ima q}>0$ and consider an
\emph{axial} function $f(q)=A(x,y)+IB(x,y)$ with $\HH$-valued $A,B$.

\begin{lemma}[Axial differential identities and the Vekua system]
\label{a:lem:axial}
For smooth axial $A+IB$ with $y>0$,
\begin{align}
 \Dop(A+IB)&=\Bigl(A_x-B_y-\frac{2B}{y}\Bigr)+I\bigl(B_x+A_y\bigr),
 \label{a:eq:axial-D}\\
 \lap(A+IB)&=A_{xx}+A_{yy}+\frac{2A_y}{y}
  +I\Bigl(B_{xx}+B_{yy}+\frac{2B_y}{y}-\frac{2B}{y^2}\Bigr).
 \label{a:eq:axial-lap}
\end{align}
Consequently $g=P+IQ$ is axially monogenic if and only if the
\emph{Vekua system}
\begin{equation}
 P_x-Q_y=\frac{2Q}{y},\qquad Q_x+P_y=0
 \label{a:eq:vekua}
\end{equation}
holds on the base. For a slice-regular $f=\I(F)$ with $F=A+\iota B$,
\begin{equation}
 \Dop f=-\frac{2B}{y},\qquad
 \lap f=\frac{2A_y}{y}+I\Bigl(\frac{2B_y}{y}-\frac{2B}{y^2}\Bigr),
 \label{a:eq:slice-D-lap}
\end{equation}
that is, $g=\lap\I(F)=P+IQ$ with $P=2A_y/y$ and $Q=2B_y/y-2B/y^2$;
equivalently, in complexified form,
\begin{equation}
 Q-\iota P=\frac2yF'-\frac2{y^2}\Ima_{\CC}F,
 \qquad
 \axn g\le\frac2y\norm{F'}_{\HC}+\frac2{y^2}\norm{F}_{\HC}.
 \label{a:eq:Fu-first}
\end{equation}
\end{lemma}

\begin{proof}
Put $\mathbf v=\Ima q=\mathbf ix_1+\mathbf jx_2+\mathbf kx_3$, so
$I=\mathbf v/y$, and for $\ell=1,2,3$
\[
 \partial_{x_\ell}y=\frac{x_\ell}{y},\qquad
 \partial_{x_\ell}I=\frac{e_\ell}{y}-\frac{\mathbf vx_\ell}{y^3}.
\]
Hence $\sum_{\ell=1}^3e_\ell\partial_{x_\ell}I=-2/y$, the angular Laplacian
contributes $\Delta_{\RR^3}I=-2I/y^2$, and the mixed radial--angular term
$\sum_\ell(\partial_{x_\ell}I)(\partial_{x_\ell}y)$ vanishes. The product rule
together with the ordinary radial Laplacian in three dimensions gives
\eqref{a:eq:axial-D}--\eqref{a:eq:axial-lap}, with all quaternionic factors on
their displayed sides. An expression $C_1+IC_2$ vanishing for every $I\in\SSS$
has $C_1=C_2=0$, by comparing $I$ and $-I$; applied to \eqref{a:eq:axial-D}
this gives \eqref{a:eq:vekua}. For a slice-regular $f$ the stem equations
\eqref{a:eq:stemCR} give $A_x=B_y$, $A_y=-B_x$ and
$A_{xx}+A_{yy}=B_{xx}+B_{yy}=0$, which turn
\eqref{a:eq:axial-D}--\eqref{a:eq:axial-lap} into
\eqref{a:eq:slice-D-lap}. Finally $F'=A_x+\iota B_x=B_y-\iota A_y$, so
$\tfrac2yF'-\tfrac2{y^2}\Ima_\CC F
=\tfrac{2B_y}{y}-\tfrac{2B}{y^2}-\iota\tfrac{2A_y}{y}=Q-\iota P$, and the
estimate follows since $\norm{Q-\iota P}_{\HC}^2=\abs P^2+\abs Q^2=\axn g^2$.
\end{proof}

All ten source manuscripts contain this computation with identical signs and
coefficients; it is stated here once and used without further comment in
Parts~II--IX. Identity \eqref{a:eq:Fu-first} explains, incidentally, why off
the real axis the order of a singularity increases by \emph{one} under a map
defined by a four-dimensional operator of order two.

\begin{theorem}[Fueter mapping, biharmonicity, and the affine kernel]
\label{a:thm:fueter-map}
If $f\in\SR(\Om_D)$ then
\begin{equation}
 \Dop(\lap f)=0,\qquad \lap^2f=0 .
 \label{a:eq:fueter-map}
\end{equation}
On a connected slice domain the kernel of $f\mapsto\lap f$ is exactly
\[
 \Aff=\{q\mapsto qa+b:\ a,b\in\HH\},
\]
with $a$ and $b$ \emph{quaternions}. Moreover a function that is both slice
regular and left Fueter regular is constant.
\end{theorem}

\begin{proof}
Away from $\RR$ write $\lap f=P+IQ$ with $P=2A_y/y$ and $Q=2B_y/y-2B/y^2$ as
in \eqref{a:eq:slice-D-lap}. The stem equations give
\[
 P_x=\frac{2B_{yy}}{y}=Q_y+\frac{2Q}{y},
 \qquad
 Q_x+P_y=\frac{2(B_{yx}+A_{yy})}{y}-\frac{2(B_x+A_y)}{y^2}=0,
\]
which is exactly the Vekua system \eqref{a:eq:vekua}; by
\eqref{a:eq:axial-D} this says $\Dop(\lap f)=0$, and applying
$\overline{\Dop}$ gives $\lap^2f=0$. Slice-regular functions are smooth
across $\RR$ by Theorem~\ref{a:thm:taylor}, so continuity extends both
identities across it.

If $\lap f=0$ then comparing the values at $I$ and $-I$ gives $P=Q=0$, i.e.
$A_y=0$ and $yB_y=B$. Locally in the upper half-plane $B=yb(x)$, and the stem
equations then give $b'(x)=0$ and $A_x=b$, so $f(q)=qa+b_0$ on a nonempty open
set for fixed $a,b_0\in\HH$; Theorem~\ref{a:thm:identity} propagates this
through $\Om_D$. Conversely affine functions have zero Laplacian. Finally
$\Dop f=0$ in \eqref{a:eq:slice-D-lap} forces $B=0$, the stem equations force
$A$ to be locally constant, and Theorem~\ref{a:thm:identity} makes it
globally constant.
\end{proof}

\begin{remark}[Why the coefficients may not be complexified]
\label{a:rem:complex-affine}
The kernel would be twice as large, and the entire obstruction theory of
Parts~II--IX would disappear, if $a,b$ were allowed to lie in $\HC$. They may
not. A direct computation on the product-domain stem $\iota(za+b)$, whose
components are $A=-ya$ and $B=xa+b$, gives
\begin{equation}
 \lap\I\bigl(\iota(za+b)\bigr)=-\frac{2a}{y}-I\,\frac{2(xa+b)}{y^2}\ne0
 \qquad\text{unless }a=b=0 .
 \label{a:eq:complex-affine}
\end{equation}
The special case $a=0$, $b=1$ is the stem $F=\iota$, giving $f(q)=I$ and
$\lap f=-2I/y^2$. Note that such an $f$ exists only on a product domain: on a
slice domain meeting $\RR$ it is not even continuous. Part~III identifies the
right-hand side of \eqref{a:eq:complex-affine} as exactly the kernel of the
holomorphic invariant $\Hg$, and Part~X uses it as a separating example.
\end{remark}

The two Cauchy kernels are linked by the same operator. Writing
$\mathbf v=\Ima q$ and $x=x_0$, direct computation gives
\begin{equation}
 \lap1=\lap q=0,\qquad
 \lap q^2=-4,\qquad
 \lap q^3=-4(3x+\mathbf v),\qquad
 \lap(q^{-1})=-\frac{4\bar q}{\abs q^4}=-8\pi^2E(q)\quad(q\ne0).
 \label{a:eq:kernel-link}
\end{equation}
For the last identity write $q^{-1}=\bar q\abs q^{-2}$; in four dimensions
$\lap\abs q^{-2}=0$ away from $0$, and the cross term of the product rule is
$2\sum_\mu\bar e_\mu\partial_{x_\mu}\abs q^{-2}=-4\bar q\abs q^{-4}$. Thus the
slice kernel of homogeneity $-1$ becomes, under $\lap$, the Fueter kernel of
homogeneity $-3$. The third identity in \eqref{a:eq:kernel-link} is not
decoration: Part~VIII uses $-\tfrac14\lap\bigl((q-u)^3\bigr)=3(\Rea q-u)
+\Ima q$ as the unique weight making the second spherical flux conserved, and
Part~IX identifies it with the degree-one entire axially monogenic polynomial
$3\mathcal P_1(q-u)$. The normalization $\lap q^2=-4$ recorded here is what
fixes every numerical constant in those parts.

\subsection{The local inverse by two integrations, and the seam}
\label{a:sec:local-inverse}

A local inverse of $\lap$ should only be expected within the axial class. In
the following theorem the planar set lies strictly above the real axis, so its
circularization is a \emph{product domain} rather than a slice domain meeting
$\RR$; the stem definition is used there, and the global identity and maximum
principles proved above are not asserted for arbitrary product domains.

\begin{theorem}[Constructive axial inverse]
\label{a:thm:local-inverse}
Let $U\subset\Cplus$ be nonempty, open and \emph{simply connected}, and let
\[
 g(x+Iy)=P(x,y)+IQ(x,y)
\]
be a smooth axially monogenic function on $\Om_U$. Then there is a
slice-regular axial $f$ on the same set with $\lap f=g$, unique up to addition
of $qa+b$ with $a,b\in\HH$.
\end{theorem}

\begin{proof}
By Lemma~\ref{a:lem:axial} the target satisfies the Vekua system
\eqref{a:eq:vekua}. The quaternion-valued one-form
\begin{equation}
 \om_g=-\frac P2\dd x+\frac Q2\dd y
 \label{a:eq:inverse-first}
\end{equation}
is closed, by the second Vekua equation. Integrating its four real components
on the simply connected $U$ gives a potential $\pot$ with
$\pot_x=-P/2$ and $\pot_y=Q/2$; set $B=y\pot$. Next consider
\begin{equation}
 \etaC=\Bigl(\pot+\frac{yQ}{2}\Bigr)\dd x+\frac{yP}{2}\dd y .
 \label{a:eq:inverse-second}
\end{equation}
It is closed because, using the first Vekua equation,
\[
 \partial_y\Bigl(\pot+\frac{yQ}{2}\Bigr)
 =\frac Q2+\frac Q2+\frac{yQ_y}{2}
 =\frac{yP_x}{2}
 =\partial_x\Bigl(\frac{yP}{2}\Bigr).
\]
Choose a primitive $A$ of $\etaC$. Then
\[
 A_x=\pot+\frac{yQ}{2}=B_y,\qquad A_y=\frac{yP}{2}=-B_x ,
\]
so $F=A+\iota B$ is a holomorphic stem on $U$, and
\eqref{a:eq:slice-D-lap} gives
\[
 \lap\I(F)=\frac{2A_y}{y}+I\Bigl(\frac{2B_y}{y}-\frac{2B}{y^2}\Bigr)
 =P+IQ=g .
\]
Changing $\pot$ by a constant quaternion $a$ changes $B$ by $ya$ and changes
$A$ by $xa+b$, which is precisely the addition of $qa+b$; conversely the
kernel computation in Theorem~\ref{a:thm:fueter-map}, carried out locally on
$U$, shows that the difference of any two primitives has this form, and
connectedness makes the constants global.
\end{proof}

\begin{remark}[The seam]
\label{a:rem:seam}
This construction explains both the affine ambiguity and the topological
hypothesis, and it is exactly where the rest of the document begins. On a
base that is \emph{not} simply connected, the closed forms
\eqref{a:eq:inverse-first} and \eqref{a:eq:inverse-second} can have nonzero
periods, and the theorem claims no global primitive without period conditions;
nor does it claim a primitive for a general nonaxial Cauchy--Fueter-regular
function. Two features of the construction are what the extension replaces.

First, the two integrations here are \emph{sequential}: the second form
$\etaC$ of \eqref{a:eq:inverse-second} depends on the first primitive $\pot$,
so on a multiply connected base it is not even defined until the first
integration has succeeded. Part~II removes this by exhibiting a
primitive-\emph{independent} closed form $\thet_g$, computed from $P$ and $Q$
alone, related to $\etaC$ by $\etaC-\dd(x\pot)=\thet_g$; the two obstructions
then become simultaneous rather than sequential, and the resulting criterion
needs no simple connectivity, no finite topology and no boundary regularity.
The reader is warned that one of the ten source manuscripts used the letter
$\eta$ for $\thet_g$ and four used it for $\etaC$; in this document $\etaC$
means \eqref{a:eq:inverse-second} and nothing else.

Second, the normalization of the primitive. Fixing $\pot(z_*)=\potV(z_*)=0$
at a base point $z_*=x_*+\iota y_*$ selects a unique primitive, and Part~II
records the two equivalent descriptions of that choice: it is the unique
primitive vanishing on the \emph{entire} sphere $x_*+y_*\SSS$, equivalently
$F(z_*)=0$ in $\HC$ --- \emph{eight} real conditions, not the four real
conditions $f(x_*+I_*y_*)=0$ at a single quaternionic point.
\end{remark}

\subsection{What survives, what changes, and three cautions}
\label{a:sec:comparison}

The proofs above identify two different mechanisms. Slice analysis retains
one-complex-variable holomorphy in each section and coordinates the sections
through a stem and the representation formula. Fueter analysis retains a
first-order elliptic operator and its fundamental solution, but no longer
admits all ordinary powers of $q$.

\begin{center}
\small
\renewcommand{\arraystretch}{1.22}
\begin{tabular}{@{}p{0.17\textwidth}p{0.38\textwidth}p{0.38\textwidth}@{}}
\toprule
Classical feature & Slice-regular counterpart & Fueter-regular counterpart\\
\midrule
Cauchy formula & planar contour, regular resolvent kernel
 (\S\ref{a:sec:slice-cauchy}) & three-dimensional boundary, Dirac kernel
 (\S\ref{a:sec:fueter-cauchy})\\
Taylor theory & powers $(q-c)^na_n$ at real centres
 (\S\ref{a:sec:slice-cauchy}) & real-coordinate analyticity and derivative
 estimates (\S\ref{a:sec:fueter-consequences})\\
Harmonicity & biharmonic: $\lap^2f=0$ (\S\ref{a:sec:bridge}) & harmonic:
 $\lap f=0$ (\S\ref{a:sec:fueter-consequences})\\
Zeros & isolated points or isolated whole spheres; the scalar norm stem
 controls multiplicity (\S\S\ref{a:sec:algebra}--\ref{a:sec:counting}) &
 a nonzero function can vanish on a plane
 (Example~\ref{a:ex:fueter-plane})\\
Maximum, Liouville & both hold (\S\ref{a:sec:max}) & both hold
 (\S\ref{a:sec:fueter-consequences})\\
Open mapping & holds off the closure of the degenerate spheres and for
 symmetric open sets (\S\ref{a:sec:reciprocal}) & fails in general
 (Example~\ref{a:ex:fueter-plane})\\
Singularities & Laurent series and residues at real isolated singularities
 (\S\ref{a:sec:slice-laurent}) & flux residues and a sharp $o(r^{-3})$
 removability criterion (\S\ref{a:sec:fueter-residue})\\
Conformal mapping & intrinsic slice-wise Riemann mapping, not arbitrary
 four-dimensional conformality (\S\ref{a:sec:geometry}) & no corresponding
 claim follows from the Fueter equation alone\\
\bottomrule
\end{tabular}
\end{center}

\medskip
Three cautions recur, and they govern every statement in Parts~II--XI.
\emph{First}, a right coefficient cannot be moved to the left without changing
the problem; this is why the period pair $(a_\gamma,b_\gamma)$ acts by
$f\mapsto f+qa_\gamma+b_\gamma$ with the coefficients on the right, and why
the monodromy of Part~II is a translation in a fixed eight-real-dimensional
space rather than a multiplicative holonomy. \emph{Second}, the geometric unit
for many slice theorems is the whole conjugacy sphere and not an individual
point; this is why the singular set of Part~VII is a two-sphere of real
codimension two, with critical exponent $1$, and not an isolated point with
critical exponent $3$. \emph{Third}, a quaternionic theorem must specify which
regularity it uses --- the source $g$ is axially \emph{monogenic} and the
primitive $f$ is slice \emph{regular}, and by Theorem~\ref{a:thm:fueter-map}
nothing nonconstant is both.

%% ==================== part2_3.tex ====================

% ===========================================================================
%  PART II and PART III of the merged document.
%  Writer key: b.   All labels begin with "b:".
%  Notation follows the merge contract without exception:
%    omega_g = \om_g (slope form), theta_g = \thet_g (constant form, computed
%    from the TARGET), eta_C = \etaC (the primitive-DEPENDENT form of the old
%    local construction, kept only to display the identity linking the two);
%    potentials \pot = B/y and \potV = A - xB/y; centre p = u + \iota v;
%    invariant \Hg; affine kernel \Aff; moment-real class \Omom.
%  The Fueter map is written \lap\I(\cdot) and is never abbreviated.
% ===========================================================================

\section{Two closed forms and the complete global criterion}
\label{b:sec:forms}

\subsection{The setting, and what is recalled from Part~I}
\label{b:sub:setting}

Throughout this part $U$ is a connected open subset of the upper half-plane
$\Cplus=\{z=x+\iota y:y>0\}$, and
\[
 \Om_U=\{q=x+Iy:\ x+\iota y\in U,\ I\in\SSS\}
\]
is its circularization.  A holomorphic $F=A+\iota B:U\to\HC$ induces the
slice-regular function $\I(F)(x+Iy)=A(x,y)+IB(x,y)$; all quaternionic
coefficients multiply on the right, and $\iota$ is the central imaginary unit
of $\HC$, never an element of $\SSS$.  We write $\SR(\Om_U)$ for the induced
slice-regular functions and $\AM(\Om_U)$ for the right $\HH$-vector space of
smooth axially monogenic fields $g(x+Iy)=P(x,y)+IQ(x,y)$, that is, those
annihilated by $\Dop$.  Domains meeting the real axis are treated separately
in Part~V; nothing in the present part requires $U$ to be simply connected,
finitely connected, bounded, or to have regular boundary.

Two facts established in Part~I are used constantly.  The first is the Vekua
system: for $y>0$, $g=P+IQ$ is axially monogenic exactly when
\begin{equation}
 P_x-Q_y=\frac{2Q}{y},\qquad Q_x+P_y=0.
 \label{b:vekua}
\end{equation}
The second is the Fueter formula: for $f=\I(F)=A+IB$ with $F$ holomorphic,
\begin{equation}
 \lap f=\frac{2A_y}{y}+I\Bigl(\frac{2B_y}{y}-\frac{2B}{y^2}\Bigr),
 \label{b:fueter}
\end{equation}
and $\lap f$ is again axially monogenic.  In particular $\lap q^2=-4$: the
operator $\Dop=\partial_{x_0}+i\partial_{x_1}+j\partial_{x_2}+k\partial_{x_3}$
carries no factor $\tfrac12$, and neither does $\lap=\overline{\Dop}\Dop$.

\subsection{The change of potentials}
\label{b:sub:potentials}

The local inverse of Part~I integrates twice.  It first produces
$\pot=B/y$ from a closed form built out of $g$, and then produces $A$ from a
second form that already contains the function $\pot$ just constructed.  That
is the source of the apparent asymmetry between the two integrations, and it
is what the following change of unknowns removes.

\begin{definition}[The two potentials]
\label{b:def:potentials}
For $f=\I(F)\in\SR(\Om_U)$ with $F=A+\iota B$, set
\begin{equation}
 \pot=\frac{B}{y}\ \ \text{(the slope potential)},\qquad
 \potV=A-\frac{xB}{y}=A-x\pot\ \ \text{(the constant potential)} .
 \label{b:CV}
\end{equation}
Both are $\HH$-valued.  Then $F(z)=z\pot(z)+\potV(z)$ and
\begin{equation}
 f(x+Iy)=(x+Iy)\pot(x,y)+\potV(x,y)=q\pot+\potV .
 \label{b:reconstruct}
\end{equation}
\end{definition}

The passage from $(A,B)$ to $(\pot,\potV)$ is a real-linear isomorphism at
each point of $U$, because $y>0$; nothing is lost.  Its point is that
\eqref{b:reconstruct} exhibits $f$ in the shape of the affine kernel with
variable coefficients, so that the two coefficient functions are precisely the
two objects whose global existence can fail.

\subsection{The two closed one-forms}
\label{b:sub:twoforms}

\begin{definition}[The slope form and the constant form]
\label{b:def:forms}
For $g=P+IQ\in\AM(\Om_U)$ put
\begin{align}
 \om_g&=\tfrac12\bigl(-P\dd x+Q\dd y\bigr),
 \label{b:omega}\\
 \thet_g&=\tfrac12\bigl((xP+yQ)\dd x+(yP-xQ)\dd y\bigr).
 \label{b:theta}
\end{align}
These are $\HH$-valued one-forms on $U$, integrated componentwise over real
curves; since the coefficients multiply real differentials, no ordering
ambiguity arises.  Both are defined directly from the target $g$: no primitive,
no branch, and no derivative of $g$ occurs in them.
\end{definition}

\begin{lemma}[Closedness and the potential identities]
\label{b:lem:forms}
Let $g=P+IQ\in\AM(\Om_U)$.
\begin{enumerate}[label=\textup{(\roman*)}]
\item $\om_g$ and $\thet_g$ are closed on $U$.
\item If $f=\I(F)\in\SR(\Om_U)$ satisfies $\lap f=g$, then the potentials
\eqref{b:CV} satisfy
\begin{equation}
 \dd\pot=\om_g,\qquad \dd\potV=\thet_g,\qquad f(q)=q\pot+\potV .
 \label{b:potential-identities}
\end{equation}
\item Conversely, if $\pot,\potV:U\to\HH$ are any $C^1$ functions with
$\dd\pot=\om_g$ and $\dd\potV=\thet_g$, then $A=x\pot+\potV$, $B=y\pot$ define a
holomorphic stem $F=A+\iota B$ on $U$, and $\lap\I(F)=g$.
\end{enumerate}
\end{lemma}

\begin{proof}
(i) The coefficient of $\dd x\wedge\dd y$ in $\dd\om_g$ is $\tfrac12(Q_x+P_y)$,
which vanishes by the second Vekua equation in \eqref{b:vekua}.  For the second
form, twice the corresponding coefficient is
\[
 \partial_y(xP+yQ)-\partial_x(yP-xQ)
 =x(P_y+Q_x)+2Q+y(Q_y-P_x)
 =-\bigl\{y(P_x-Q_y)-2Q-x(Q_x+P_y)\bigr\},
\]
which vanishes by both equations of \eqref{b:vekua}.  All computations are
componentwise in the real vector space $\HH$.

(ii) From $B=y\pot$ and the Fueter formula \eqref{b:fueter}, $Q=2B_y/y-2B/y^2$
gives $\pot_y=B_y/y-B/y^2=Q/2$, while the Cauchy--Riemann equation $B_x=-A_y$
together with $P=2A_y/y$ gives $\pot_x=B_x/y=-P/2$.  Hence $\dd\pot=\om_g$.
For the second potential,
\[
 \potV_x=A_x-\pot-x\pot_x=B_y-\pot+\frac{xP}{2}=\frac{xP+yQ}{2},
 \qquad
 \potV_y=A_y-x\pot_y=\frac{yP-xQ}{2},
\]
using $A_x=B_y$ and $B_y=\pot+y\pot_y=\pot+yQ/2$.  Hence $\dd\potV=\thet_g$.

(iii) With $A=x\pot+\potV$ and $B=y\pot$,
\[
 A_x=\pot+x\pot_x+\potV_x=\pot+\frac{yQ}{2}=B_y,
 \qquad
 A_y=x\pot_y+\potV_y=\frac{yP}{2}=-B_x .
\]
So $F=A+\iota B$ is holomorphic, and substituting into \eqref{b:fueter}
returns $P+IQ$.
\end{proof}

\subsection{Elimination of the first primitive}
\label{b:sub:elimination}

The second form of the classical local construction (Part~I) is the
primitive-\emph{dependent} one-form
\begin{equation}
 \etaC=\Bigl(\pot+\frac{yQ}{2}\Bigr)\dd x+\frac{yP}{2}\dd y ,
 \label{b:etaC}
\end{equation}
whose primitive is $A$ and whose very definition presupposes that a first
primitive $\pot$ of $\om_g$ has already been chosen.  The two second forms are
different one-forms, and they are related by an exact term.

\begin{proposition}[The elimination identity]
\label{b:prop:elimination}
Wherever a primitive $\pot$ of $\om_g$ exists,
\begin{equation}
 \boxed{\ \etaC-\dd(x\pot)=\thet_g\ }
 \label{b:elimination}
\end{equation}
In particular $\thet_g$ is the primitive-independent representative of the
cohomology class of $\etaC$ on any subdomain where $\om_g$ is exact.
\end{proposition}

\begin{proof}
Using $\dd\pot=\om_g$,
\[
 \dd(x\pot)=\pot\dd x+x\,\om_g
 =\Bigl(\pot-\frac{xP}{2}\Bigr)\dd x+\frac{xQ}{2}\dd y .
\]
Subtracting this from \eqref{b:etaC} gives
$\tfrac12(xP+yQ)\dd x+\tfrac12(yP-xQ)\dd y=\thet_g$.
\end{proof}

\begin{remark}[The two obstructions are simultaneous, not sequential]
\label{b:rem:simultaneous}
Identity \eqref{b:elimination} is the pivot of the whole theory, and it is
worth saying exactly what it buys and what it does not.  The form $\etaC$ is
defined only after $\om_g$ has been integrated, so on a domain where $\om_g$
is not exact it is not even globally defined; the construction then stalls at
the first step and the second obstruction cannot be formulated at all.  By
contrast $\thet_g$ in \eqref{b:theta} is written down from $P$ and $Q$ alone
and is closed on all of $U$ by Lemma~\ref{b:lem:forms}(i), whatever the
periods of $\om_g$ may be.  Consequently the two obstructions of
Theorem~\ref{b:thm:period} are simultaneous conditions on $g$, not a condition
followed by a conditional second condition, and no choice of first primitive
can remove a nonzero period of $\thet_g$: by \eqref{b:elimination} a different
choice of $\pot$ changes $\etaC$ by an exact form and leaves $\thet_g$
untouched.  This is also the reason that \emph{two} periods, rather than one,
are the correct invariant: they are the two coefficient functions of
\eqref{b:reconstruct}, and Example~(i) of Part~X separates them.
\end{remark}

\subsection{The global period criterion}
\label{b:sub:criterion}

\begin{definition}[The period pair]
\label{b:def:periods}
For a closed piecewise $C^1$ curve $\gamma\subset U$ and $g\in\AM(\Om_U)$ put
\begin{equation}
 a_\gamma(g)=\oint_\gamma\om_g,\qquad
 b_\gamma(g)=\oint_\gamma\thet_g,\qquad
 \Per_g(\gamma)=(a_\gamma(g),b_\gamma(g))\in\HH^2 .
 \label{b:periods}
\end{equation}
We call $a_\gamma$ the \emph{slope period} and $b_\gamma$ the \emph{constant
period}.  By Lemma~\ref{b:lem:forms}(i) both depend only on the integral
homology class of $\gamma$.
\end{definition}

\begin{theorem}[The complete same-domain criterion]
\label{b:thm:period}
Let $U\subset\Cplus$ be \emph{any} connected open set and let
$g\in\AM(\Om_U)$.  The following are equivalent.
\begin{enumerate}[label=\textup{(\roman*)}]
\item There is a single-valued $f\in\SR(\Om_U)$ with $\lap f=g$.
\item Both $\om_g$ and $\thet_g$ are exact on $U$.
\item $a_\gamma(g)=b_\gamma(g)=0$ for every closed piecewise $C^1$ curve
$\gamma\subset U$.
\end{enumerate}
When these hold, fix $z_*\in U$ and $c_0,e_0\in\HH$; every solution is
obtained exactly once from
\begin{equation}
 \pot(z)=c_0+\int_{z_*}^{z}\om_g,\qquad
 \potV(z)=e_0+\int_{z_*}^{z}\thet_g,\qquad
 f(x+Iy)=(x+Iy)\pot(z)+\potV(z),
 \label{b:global-inverse}
\end{equation}
the integrals being taken along any path in $U$.  The solution is unique
modulo $\Aff$.
\end{theorem}

\begin{proof}
(i)$\Leftrightarrow$(ii) is Lemma~\ref{b:lem:forms}(ii)--(iii): a primitive
supplies primitives of the two forms, and primitives of the two forms supply a
holomorphic stem on $U$, hence, after reflecting by
$F(\bar z)=\kappa F(z)$, a symmetric stem on $U\cup\overline U$ and a
single-valued $f\in\SR(\Om_U)$ with $\lap f=g$.

(ii)$\Rightarrow$(iii) is immediate.  For (iii)$\Rightarrow$(ii), note that a
connected open subset of the plane is polygonally path connected; define the
two path integrals in \eqref{b:global-inverse}.  Two paths with the same
endpoints differ by a closed curve, so (iii) makes both values independent of
the path, for each of the four real components separately.  On a small disk
the usual differentiation of a path integral shows that the resulting
functions are primitives of the corresponding forms.

For uniqueness, two pairs of primitives differ by constants $a,b\in\HH$, and
\eqref{b:reconstruct} turns that difference into $qa+b$; conversely every such
affine addition is annihilated by $\lap$ (Corollary~\ref{b:cor:kernel}).
\end{proof}

\begin{remark}[No topological and no boundary hypothesis]
\label{b:rem:nohypothesis}
Theorem~\ref{b:thm:period} assumes nothing about $U$ beyond connectedness and
openness.  In particular it does not assume finite connectivity, a rectifiable
or accessible boundary, boundedness, or simple connectivity of the
complementary components; it is not restricted to a rectangle or a disk, and
it is an equivalence rather than a sufficient condition.  The price is that
condition~(iii) quantifies over all loops.  Converting it into a
\emph{finite} computation requires a topological hypothesis, and exactly that
hypothesis --- a finite winding basis --- is isolated in Part~IV; the
cokernel computation lives there, not here.
\end{remark}

\begin{corollary}[The affine kernel]
\label{b:cor:kernel}
On connected $U$,
\begin{equation}
 \ker\bigl(\lap|_{\SR(\Om_U)}\bigr)=\Aff=\{q\mapsto qa+b:\ a,b\in\HH\}.
 \label{b:kernel}
\end{equation}
The coefficients are \emph{quaternions}, not elements of $\HC$.
\end{corollary}

\begin{proof}
If $\lap f=0$ then $g=0$, so $\dd\pot=\om_0=0$ and $\dd\potV=\thet_0=0$ by
Lemma~\ref{b:lem:forms}(ii); connectedness makes $\pot\equiv a$ and
$\potV\equiv b$ constant quaternions, and \eqref{b:reconstruct} gives
$f(q)=qa+b$.  The converse is \eqref{b:fueter}.
\end{proof}

The restriction to \emph{real} quaternionic coefficients in \eqref{b:kernel}
is not a matter of taste: enlarging $\Aff$ to the complex-affine stems would
double the kernel and erase the obstruction.  The relevant counterexample is
the stem $F=\iota$, which induces $f(q)=I$ and has
$\lap f=-2I/y^2\neq0$; it reappears in Part~III as
Theorem~\ref{b:thm:range}, where the whole eight-real-dimensional family
$\lap\I(\iota(za+b))$ is identified as the kernel of the holomorphic
invariant.

\begin{corollary}[Affine monodromy]
\label{b:cor:monodromy}
Without any vanishing hypothesis, a local Fueter primitive of
$g\in\AM(\Om_U)$ continues analytically along every path in $U$.  Continuation
around a based loop $\gamma$ changes it by
\begin{equation}
 F\longmapsto F+z\,a_\gamma(g)+b_\gamma(g),
 \qquad
 f\longmapsto f+q\,a_\gamma(g)+b_\gamma(g).
 \label{b:monodromy}
\end{equation}
Consequently the full monodromy of the local inverse is the additive
homomorphism
\begin{equation}
 \Per_g:H_1(U;\ZZ)\longrightarrow\HH\oplus\HH,
 \qquad [\gamma]\longmapsto(a_\gamma,b_\gamma),
 \label{b:monodromy-map}
\end{equation}
and there are no further local compatibility obstructions.
\end{corollary}

\begin{proof}
Pull the two closed forms back to the universal cover of $U$, where they are
exact; equivalently, integrate them along paths from a fixed base point while
retaining the homotopy class.  Lemma~\ref{b:lem:forms}(iii) makes the
resulting $F=z\pot+\potV$ holomorphic on each chart, and it agrees with the
prescribed primitive after a correct choice of the two initial constants.
Appending $\gamma$ to a path adds $a_\gamma$ to $\pot$ and $b_\gamma$ to
$\potV$, which is \eqref{b:monodromy} by \eqref{b:reconstruct}.
Concatenation of loops adds periods and reversal negates them, and closed
forms integrate homologically, so \eqref{b:monodromy-map} is a well-defined
homomorphism.  The last sentence is the content of
Theorem~\ref{b:thm:period}: the two periods are the only obstructions, and
the construction supplies an actual continued primitive along every path.
\end{proof}

\begin{remark}[What is, and is not, noncommutative here]
\label{b:rem:abelian}
The monodromy \eqref{b:monodromy} is a \emph{translation} inside the fixed
eight-real-dimensional space $\Aff$.  It is not a quaternionic
fractional-linear action, not a path-ordered exponential, and not a
multiplicative holonomy; noncommutativity of $\HH$ plays no role in the
composition law, because the coefficients remain on the right and the group
operation is addition.  In particular the obstruction group is abelian, and
every commutator loop in $\pi_1(U)$ acts trivially.  What is genuinely
quaternionic is the \emph{size} of the obstruction --- two quaternionic
coordinates rather than one complex residue --- and the geometry of the
conjugacy spheres, which re-enters essentially in the singularity norms and
the energy law of Parts~VII and~VIII.
\end{remark}

Theorem~\ref{b:thm:period} and Corollary~\ref{b:cor:monodromy} may be
assembled into an exact sequence of right quaternionic vector spaces,
\begin{equation}
 0\longrightarrow\HH^2\xrightarrow{\ (a,b)\mapsto qa+b\ }
 \SR(\Om_U)\xrightarrow{\ \lap\ }\AM(\Om_U)
 \xrightarrow{\ \Per\ }\operatorname{Hom}\bigl(H_1(U;\ZZ),\HH^2\bigr),
 \label{b:exact-general}
\end{equation}
exact at $\SR(\Om_U)$ and at $\AM(\Om_U)$.  We do \emph{not} assert here that
the last arrow is surjective for an arbitrary, possibly infinitely connected,
$U$; surjectivity is proved in Part~IV under an explicit finite-topology
hypothesis, and the modes that realize the periods are constructed there.

\subsection{Normalization: eight real conditions, not four}
\label{b:sub:normalization}

\begin{corollary}[The normalized primitive]
\label{b:cor:normalization}
Fix $z_*=x_*+\iota y_*\in U$ and suppose $g\in\AM(\Om_U)$ has a global
primitive.  The following three descriptions single out one and the same
primitive $f$.
\begin{enumerate}[label=\textup{(\alph*)}]
\item $\pot(z_*)=\potV(z_*)=0$, i.e.\ the choice $c_0=e_0=0$ in
\eqref{b:global-inverse}.
\item $F(z_*)=0$ in $\HC$, which is eight real conditions on the stem.
\item $f$ vanishes on the entire base sphere $x_*+y_*\SSS\subset\Om_U$.
\end{enumerate}
It is the unique primitive with any one of these properties.
\end{corollary}

\begin{proof}
(a)$\Leftrightarrow$(b): the map $\HH^2\to\HC$, $(a,b)\mapsto z_*a+b$, is a
real isomorphism because $y_*>0$ --- it is the change of potentials
\eqref{b:CV} evaluated at $z_*$ --- so $F(z_*)=z_*\pot(z_*)+\potV(z_*)$
vanishes precisely when both potentials do.
(a)$\Rightarrow$(c): $f(x_*+Iy_*)=(x_*+Iy_*)\pot(z_*)+\potV(z_*)=0$ for every
$I\in\SSS$.
(c)$\Rightarrow$(a): if $A(z_*)+IB(z_*)=0$ for all $I\in\SSS$, subtract the
values at $I$ and $-I$ to get $2IB(z_*)=0$, hence $B(z_*)=0$ and then
$A(z_*)=0$.  Uniqueness: two primitives differ by $qa+b$ by
Corollary~\ref{b:cor:kernel}, and an affine function vanishing at all of
$x_*+y_*\SSS$ has $a=b=0$ by the same subtraction.
\end{proof}

\begin{remark}[A normalization that is easy to get wrong]
\label{b:rem:eightnotfour}
Descriptions (b) and (c) look like conditions of different sizes and are
often stated as if they were: (b) is manifestly eight real conditions, (c)
looks like the vanishing of a function on a two-sphere.  They coincide, and
both are strictly stronger than the four-real-dimensional condition
$f(q_*)=0$ at a \emph{single} quaternionic point $q_*=x_*+I_*y_*$.
That weaker condition does not determine $f$: the elements of $\Aff$ vanishing
at one fixed $q_*$ are exactly $q\mapsto qa-q_*a$ with $a\in\HH$ arbitrary
(note $q_*a\in\HH$), a four-real-dimensional space, so a pointwise
normalization leaves a four-parameter ambiguity and one must pass to the whole
conjugacy sphere $x_*+y_*\SSS$ to pin the primitive down.  Installing the
pointwise normalization in place of
\textup{(b)} or \textup{(c)} produces a normalized inverse operator that is
not well defined.  The quantitative consequences of the normalization ---
continuity of the resulting inverse on the zero-period subspace and its
domain-dependent constants --- belong to Part~VI.
\end{remark}

\subsection{Change of the real origin}
\label{b:sub:translation}

\begin{proposition}[Translation covariance of the period pair]
\label{b:prop:translation}
Let $c\in\RR$ and let $x'=x-c$, $y'=y$ be the translated planar coordinates,
so that $q'=q-c$.  For the same field, $\om'_g=\om_g$ and
\begin{equation}
 \thet'_g=\thet_g+c\,\om_g,
 \qquad\text{hence}\qquad
 a'_\gamma=a_\gamma,\qquad b'_\gamma=b_\gamma+c\,a_\gamma .
 \label{b:translation}
\end{equation}
Consequently, for data attached to a sphere with centre $p=u+\iota v$, whose
centre becomes $u'=u-c$, the pair
\[
 (a,\;ua+b)\quad\text{is invariant, since } u'a+b'=(u-c)a+(b+ca)=ua+b .
\]
\end{proposition}

\begin{proof}
$\om_g$ contains no $x$.  Substituting $x=x'+c$ in \eqref{b:theta},
\[
 \thet_g=\tfrac12\bigl((x'P+yQ)\dd x'+(yP-x'Q)\dd y\bigr)
 +\tfrac{c}{2}\bigl(P\dd x'-Q\dd y\bigr)
 =\thet'_g-c\,\om_g ,
\]
which is \eqref{b:translation}.  Integrate over $\gamma$.
\end{proof}

Thus $a$ and $b$ are affine-monodromy coordinates, not two independently
translation-invariant scalars; it is the combination $ua+b$, and not $b$
itself, that is intrinsic to the sphere.  This is exactly why the residue
norm $\Rp(a,b)^2=\abs{ua+b}^2+v^2\abs{a}^2$ of Part~VIII is
translation-invariant, and why the energy coefficient it controls does not
depend on a choice of real origin.

% ===========================================================================

\section{The holomorphic invariant, its exact image and kernel}
\label{b:sec:invariant}

The period criterion of Part~II uses only first-order real data and no
derivatives of $g$.  For meromorphic and rational examples, and for every
residue computation in Parts~VII and~IX, a second description is more
efficient: a single holomorphic $\HC$-valued function attached to $g$, whose
ordinary complex residues are the two quaternionic periods and whose higher
Laurent coefficients are the higher multipoles.

\subsection{Definition, and why five formulas define one object}
\label{b:sub:Hdef}

\begin{definition}[The holomorphic invariant]
\label{b:def:H}
For $g=P+IQ\in\AM(\Om_U)$ put
\begin{equation}
 \Hg(z)=-\tfrac12\bigl(P(x,y)+yP_y(x,y)\bigr)-\frac{\iota y}{2}P_x(x,y).
 \label{b:H}
\end{equation}
\end{definition}

\begin{lemma}[One object, two formulas]
\label{b:lem:Hforms}
On $\AM(\Om_U)$,
\begin{equation}
 \Hg=\tfrac12\bigl(-P+yQ_x\bigr)-\frac{\iota y}{2}P_x .
 \label{b:H-alt}
\end{equation}
\end{lemma}

\begin{proof}
The second Vekua equation $Q_x+P_y=0$ gives $yQ_x=-yP_y$.
\end{proof}

This is worth stating explicitly, because the two displays \eqref{b:H} and
\eqref{b:H-alt} occur in the literature as if they were definitions of two
different invariants; they are the same function of $g$, and every identity
below may be verified from either.  Note also that $\Hg$ depends on $Q$ only
through a derivative that the Vekua system re-expresses in $P$: the invariant
is determined by the single axial component $P$ together with axial
monogenicity.

\subsection{Holomorphy, the second derivative, and the complex periods}
\label{b:sub:Hthm}

\begin{theorem}[The target-derived second derivative]
\label{b:thm:invariant}
Let $g\in\AM(\Om_U)$.  Then $\Hg$ is a single-valued holomorphic
$\HC$-valued function on $U$; every local holomorphic stem $F$ of a Fueter
primitive of $g$ satisfies
\begin{equation}
 F''=\Hg ,
 \label{b:H-second}
\end{equation}
and for every closed piecewise $C^1$ curve $\gamma\subset U$,
\begin{equation}
 a_\gamma(g)=\oint_\gamma\Hg(z)\dd z,
 \qquad
 b_\gamma(g)=-\oint_\gamma z\,\Hg(z)\dd z .
 \label{b:complex-periods}
\end{equation}
In particular both of these $\HC$-valued integrals lie in the real subspace
$\HH\subset\HC$.
\end{theorem}

\begin{proof}
\emph{First proof of holomorphy (structural).}  By
Theorem~\ref{b:thm:period} a primitive exists on every disk $\subset U$.
With $F=A+\iota B$, $A=x\pot+\potV$, $B=y\pot$ and the derivative identities
in the proof of Lemma~\ref{b:lem:forms},
\begin{equation}
 F'=A_x+\iota B_x=\pot+\mathcal M_g,
 \qquad
 \mathcal M_g:=\frac y2\,Q-\frac{\iota y}{2}\,P .
 \label{b:M}
\end{equation}
Differentiating once more in $x$ and using $\pot_x=-P/2$ and $Q_x=-P_y$,
\[
 F''=-\frac P2+\frac y2 Q_x-\frac{\iota y}{2}P_x
 =-\frac12(P+yP_y)-\frac{\iota y}{2}P_x=\Hg .
\]
The right-hand side is defined directly from $g$ on all of $U$, and local
primitives differ by $za+b$, whose second derivative vanishes; so the local
identifications agree on overlaps and $\Hg$ is globally holomorphic and
independent of every choice.

\emph{Second proof of holomorphy (direct, without any primitive).}  Write
$\Hg=J_0+\iota J_1$ with
$J_0=\tfrac12(-P+yQ_x)$ and $J_1=-\tfrac12 yP_x$, and verify the
Cauchy--Riemann equations from \eqref{b:vekua} alone.  First,
\[
 (J_0)_x-(J_1)_y
 =\tfrac12(-P_x+yQ_{xx})+\tfrac12(P_x+yP_{xy})
 =\tfrac y2\bigl(Q_{xx}+P_{xy}\bigr)=0,
\]
by differentiating $Q_x+P_y=0$ in $x$.  Second,
\[
 (J_0)_y+(J_1)_x
 =\tfrac12(-P_y+Q_x+yQ_{xy})-\tfrac y2P_{xx}
 =Q_x+\tfrac y2\bigl(Q_{xy}-P_{xx}\bigr)=0,
\]
using $-P_y=Q_x$ in the first bracket and, for the second, the $x$-derivative
of $P_x-Q_y=2Q/y$, namely $P_{xx}-Q_{xy}=2Q_x/y$.  This proof uses no
primitive at all, so the global holomorphy of $\Hg$ is logically independent
of the period theorem; it is the reason $\Hg$ may be used to \emph{prove}
statements about primitives rather than only to record them.

\emph{The complex period formulas.}  As identities of $\HC$-valued one-forms
on $U$,
\begin{equation}
 \Hg\dd z=\om_g+\dd\mathcal M_g,
 \qquad
 z\,\Hg\dd z=\dd\bigl(z\mathcal M_g\bigr)-\thet_g .
 \label{b:corrections}
\end{equation}
The first is \eqref{b:M} differentiated, since $\dd F'=F''\dd z=\Hg\dd z$ and
$\dd\pot=\om_g$.  For the second, $F=z\pot+\potV$ gives
$zF'-F=z\mathcal M_g-\potV$, and differentiating,
$z\Hg\dd z=\dd(zF'-F)=\dd(z\mathcal M_g)-\thet_g$.
The functions $\mathcal M_g$ and $z\mathcal M_g$ are globally single-valued
on $U$, being written directly in terms of $P,Q$; hence their differentials
integrate to zero around any closed curve, and \eqref{b:corrections} yields
\eqref{b:complex-periods}.  The final reality assertion is immediate from the
original definition \eqref{b:periods}, in which the integrands are
$\HH$-valued.
\end{proof}

\begin{remark}[The mechanism, and what it costs]
\label{b:rem:mechanism}
The pair \eqref{b:corrections} is the exact mechanism converting the two real
period integrals into complex contour integrals: $\Hg\dd z$ and $\om_g$
differ by an \emph{exact} $\HC$-valued form, as do $-z\Hg\dd z$ and
$\thet_g$, and the correction $\mathcal M_g$ is built from $g$ with no
integration.  The cost of the holomorphic description is two derivatives:
$\om_g,\thet_g$ are algebraic in $(P,Q)$, whereas $\Hg$ involves $P_x$ and
$P_y$.  For a numerically sampled field the real forms are therefore the
stable objects, and $\Hg$ is the efficient one when a closed formula is
available.
\end{remark}

\subsection{The exact image: the moment-real class}
\label{b:sub:image}

Theorem~\ref{b:thm:invariant} says that the two contour integrals of $\Hg$
land in $\HH$ although $\Hg$ itself is $\HC$-valued.  This is a genuine
restriction, and it turns out to be the only one.

\begin{definition}[The moment-real class]
\label{b:def:omom}
Let
\begin{equation}
 \Omom(U)=\Bigl\{h:U\to\HC\ \text{holomorphic}\ :\
 \oint_\gamma h\dd z\in\HH\ \text{and}\
 -\oint_\gamma z\,h\dd z\in\HH\ \text{for every closed }\gamma\subset U\Bigr\}.
 \label{b:omom}
\end{equation}
\end{definition}

The name records that the two lowest complex \emph{moments} of $h$ along every
cycle are real-quaternionic.  It does \emph{not} say that $h$ is
$\HH$-valued, nor that $h$ is $\kappa$-symmetric; on a simply connected $U$
the condition is vacuous and $\Omom(U)$ is the full space
$\mathcal O(U;\HC)$.

\begin{theorem}[Exact image and exact kernel of the invariant]
\label{b:thm:range}
The map $g\mapsto\Hg$ sends $\AM(\Om_U)$ \emph{onto} $\Omom(U)$, and its
kernel is exactly the eight-real-parameter family
\begin{equation}
 \lap\I\bigl(\iota(za+b)\bigr)(x+Iy)
 =-\frac{2a}{y}-I\,\frac{2(xa+b)}{y^2},
 \qquad a,b\in\HH .
 \label{b:Hkernel}
\end{equation}
Thus $\ker(g\mapsto\Hg)$ has real dimension exactly eight on a product
domain, and each of its members is itself a global Fueter image, with global
stem $\iota(za+b)$.
\end{theorem}

\begin{proof}
\emph{Necessity of the moment condition} is \eqref{b:complex-periods}.

\emph{Surjectivity.}  Let $h\in\Omom(U)$.  Integrate twice on the universal
cover $\widetilde U$ to obtain a holomorphic $F$ with $F''=h$.  Continuation
around a lifted loop $\gamma$ leaves $F''$ unchanged, so it changes $F$ by
$za+b$ with $a,b\in\HC$ a priori.  Integrating $\dd F'=h\dd z$ and
$\dd(zF'-F)=z\,h\dd z$ identifies these coefficients as
\[
 a=\oint_\gamma h\dd z,\qquad b=-\oint_\gamma z\,h\dd z,
\]
which lie in $\HH$ by hypothesis.  By Corollary~\ref{b:cor:kernel} the local
fields $\lap\I(F)$ are therefore unchanged by every continuation; they glue to
a single-valued $g\in\AM(\Om_U)$, and $\Hg=F''=h$ locally, hence globally.

\emph{Kernel.}  If $\Hg=0$ then every local primitive stem has $F''=0$, so
$F=zc+d$ with $c,d\in\HC$.  Write $c=a_0+\iota a$ and $d=b_0+\iota b$ with
$a_0,a,b_0,b\in\HH$.  The term $za_0+b_0$ is killed by the Fueter map
(Corollary~\ref{b:cor:kernel}); for the rest, $\iota(za+b)$ has components
$A=-ya$, $B=xa+b$, and \eqref{b:fueter} gives
\[
 P=\frac{2A_y}{y}=-\frac{2a}{y},
 \qquad
 Q=\frac{2B_y}{y}-\frac{2B}{y^2}=-\frac{2(xa+b)}{y^2},
\]
which is \eqref{b:Hkernel}.  Conversely each such field satisfies
\eqref{b:vekua} and has $P_x=0$ and $P+yP_y=0$, so $\Hg=0$.  The parameters
are determined by the field --- $a$ from $P$, then $b$ from $Q$ --- so local
choices agree on overlaps and are global on connected $U$; this also shows
the family is exactly eight-real-dimensional.
\end{proof}

\begin{remark}[The invariant is not injective, and where the kernel is
prohibited]
\label{b:rem:notinjective}
Theorem~\ref{b:thm:range} identifies $\AM(\Om_U)$ with a constrained
holomorphic function \emph{plus eight real parameters}; $\Hg=0$ does not
imply $g=0$.  Equivalently, solving \eqref{b:H} directly gives
$P=c/y$, $Q=(cx+d)/y^2$ with $c=-2a$, $d=-2b$ constant on connected $U$.
Three points deserve emphasis.
\begin{enumerate}[label=\textup{(\alph*)}]
\item There is no conflict with Theorem~\ref{b:thm:period}: the fields
\eqref{b:Hkernel} have both periods zero, being global Fueter images.  The
invariant loses exactly the part of $g$ that carries no obstruction.
\item The special case $a=0$, $b=1$ is the stem $F=\iota$, whose induced
function is $f(q)=I$ with $\lap f=-2I/y^2\neq0$.  This is the standing
counterexample to enlarging $\Aff$ to complex-affine stems.
\item On a slice domain meeting $\RR$ the kernel is \emph{vacuous}: every
member of \eqref{b:Hkernel} blows up like $y^{-1}$ or $y^{-2}$ as $y\to0$, so
smooth continuation through a real interval forces $a=b=0$.  On such domains
$g\mapsto\Hg$ is injective.  The reflected statement is taken up in Part~V.
\end{enumerate}
\end{remark}

\begin{remark}[A constructive test for proposed holomorphic encodings]
\label{b:rem:test}
Theorem~\ref{b:thm:range} is the precise sense in which axially monogenic
fields are ``complexified''.  It is also a usable test: an arbitrary
holomorphic $\HC$-valued function is \emph{not} the invariant of an axially
monogenic field.  Arbitrary complex residues need not satisfy \eqref{b:omom},
and Theorem~\ref{b:thm:reality} below computes exactly which ones do at an
isolated nonreal point.
\end{remark}

\subsection{Residues, multipoles, and the reality constraint}
\label{b:sub:residues}

Fix $p=u+\iota v$ with $v>0$, let $0<R<v$, and let $g$ be axially monogenic on
the circularization of the punctured disk $0<\abs{z-p}<R$ --- that is, on a
punctured tubular neighbourhood of the sphere $\Sp=u+v\SSS$.  Recall from
Part~VII that the \emph{affine spherical residue} of $g$ at $\Sp$ is the pair
\[
 \Resaff_{\Sp}(g)=(a,b)
 =\Bigl(\oint_{\abs{z-p}=s}\om_g,\ \oint_{\abs{z-p}=s}\thet_g\Bigr),
 \qquad 0<s<R,
\]
independent of $s$ by closedness, and that Part~VII's hypothesis-free
normal form writes $g$ uniquely as a regular field, plus the logarithmic mode
$G_{p;a,b}$ of Part~IV, plus $\lap\I\bigl(\sum_{k\ge1}(z-p)^{-k}c_k\bigr)$
with $c_k\in\HC$.  We now recover \emph{all} of these data from $\Hg$ by
ordinary contour integration.

\begin{theorem}[Contour extraction of every multipole, and the reality
constraint]
\label{b:thm:reality}
Let $\Hg(z)=\sum_{n\in\ZZ}h_n(z-p)^n$, $h_n\in\HC$, be the Laurent expansion
on $0<\abs{z-p}<R$.  Then, with the circle positively oriented,
\begin{align}
 a&=2\pi\iota\,h_{-1},
 &
 b&=-2\pi\iota\bigl(p\,h_{-1}+h_{-2}\bigr),
 \label{b:residue-pair}\\
 c_k&=\frac{h_{-k-2}}{k(k+1)}
 =\frac{1}{2\pi\iota\,k(k+1)}\oint_{\abs{z-p}=s}(z-p)^{k+1}\Hg(z)\dd z,
 &&k\ge1 .
 \label{b:multipoles}
\end{align}
Equivalently, $a$ and $-b$ are $2\pi\iota$ times the ordinary complex
residues at $p$ of $\Hg$ and of $z\Hg$, computed componentwise in $\HC$;
both results lie in $\HH$ although the individual coefficients $h_n$ do not.
Moreover the two lowest negative coefficients are \emph{not} independent:
\begin{equation}
 \boxed{\ h_{-1}=\frac{\iota}{v}\,\ReC\bigl(h_{-2}\bigr)\ }
 \label{b:reality}
\end{equation}
Explicitly, writing $h_{-2}=c+\iota d$ with $c,d\in\HH$,
\begin{equation}
 a=-\frac{2\pi}{v}\,c,
 \qquad
 b=2\pi d+\frac{2\pi u}{v}\,c .
 \label{b:one-coefficient}
\end{equation}
Conversely, a holomorphic $h$ on the punctured disk is the invariant of some
axially monogenic field there if and only if its Laurent coefficients satisfy
\eqref{b:reality}; the remaining coefficients are unrestricted in $\HC$.
\end{theorem}

\begin{proof}
Formulas \eqref{b:residue-pair} are the componentwise complex residue theorem
applied to \eqref{b:complex-periods}, the punctured disk having a single
homology generator.  Only the Laurent expansion is used, so they hold also
when $\Hg$ has an essential singularity at $p$.

For \eqref{b:multipoles}, apply $g\mapsto\Hg$ to the normal form of Part~VII.
The regular summand has a holomorphic invariant.  The mode $G_{p;a,b}$ has
the rational invariant
$\frac{1}{2\pi\iota}\bigl(\frac{a}{z-p}-\frac{pa+b}{(z-p)^2}\bigr)$ of
Part~IV, which contributes nothing below index $-2$.  Finally, by
\eqref{b:H-second} applied to the stem $(z-p)^{-k}c_k$,
\[
 \Hg\bigl[\lap\I\bigl((z-p)^{-k}c_k\bigr)\bigr]
 =\frac{\dd^2}{\dd z^2}\bigl((z-p)^{-k}c_k\bigr)
 =k(k+1)(z-p)^{-k-2}c_k ,
\]
so the coefficient of $(z-p)^{-k-2}$ in $\Hg$ is $k(k+1)c_k$ for $k\ge1$;
extracting it by the Cauchy coefficient formula gives the displayed contour
integral.

For \eqref{b:reality}: by \eqref{b:residue-pair} and $a\in\HH$, the first
formula forces $h_{-1}=\iota\alpha$ with $\alpha=-a/(2\pi)\in\HH$.  Write
$h_{-2}=c+\iota d$ with $c,d\in\HH$.  Since $p=u+\iota v$,
\[
 p\,h_{-1}+h_{-2}=\iota u\alpha-v\alpha+c+\iota d
 =(c-v\alpha)+\iota(d+u\alpha).
\]
The requirement $b=-2\pi\iota(p h_{-1}+h_{-2})\in\HH$ kills the
$\HH$-component of $p h_{-1}+h_{-2}$, that is, $c=v\alpha$.  Hence
$\alpha=c/v$, so $h_{-1}=\iota c/v=(\iota/v)\ReC(h_{-2})$, and
$a=-2\pi\alpha=-(2\pi/v)c$, while
$b=-2\pi\iota\cdot\iota(d+u\alpha)=2\pi d+(2\pi u/v)c$.  This is
\eqref{b:one-coefficient}.

Conversely, on a punctured disk every loop is homologous to a multiple of one
circle, so the two conditions defining $\Omom$ reduce to
$a\in\HH$ and $b\in\HH$ for that circle, which by the computation just made is
exactly \eqref{b:reality}.  Theorem~\ref{b:thm:range} then produces an
axially monogenic $g$ with $\Hg=h$.
\end{proof}

\begin{corollary}[The sharp vanishing criterion at a punctured tube]
\label{b:cor:h2}
For $g$ axially monogenic on a punctured tube around $\Sp$, the following are
equivalent:
\begin{enumerate}[label=\textup{(\roman*)}]
\item $g$ has a single-valued slice-regular Fueter primitive on the punctured
tube;
\item $(a,b)=(0,0)$;
\item $h_{-2}=0$;
\item $h_{-1}=h_{-2}=0$.
\end{enumerate}
In particular $\Hg$ can never have an isolated pole of order exactly one at a
nonreal $p$, and a single coefficient $h_{-2}\in\HC$ carries all eight real
obstruction coordinates.
\end{corollary}

\begin{proof}
(i)$\Leftrightarrow$(ii) is Theorem~\ref{b:thm:period} on the punctured disk,
whose first homology is generated by one circle.  (ii)$\Leftrightarrow$(iv)
is \eqref{b:residue-pair}, since $2\pi\iota$ is invertible in $\HC$.
(iii)$\Rightarrow$(iv) is \eqref{b:reality}, and (iv)$\Rightarrow$(iii) is
trivial.  For the last two assertions: a simple pole would mean $h_{-2}=0$ and
$h_{-1}\ne0$, which \eqref{b:reality} forbids; and \eqref{b:one-coefficient}
is a real-linear bijection $h_{-2}\mapsto(a,b)$ from $\HC\cong\RR^8$ onto
$\HH^2\cong\RR^8$, with inverse $c=-va/(2\pi)$,
$d=(b+ua)/(2\pi)$ --- note that it is $ua+b$, in accordance with
Proposition~\ref{b:prop:translation}, that appears.
\end{proof}

\begin{remark}[A true but redundant form of the criterion]
\label{b:rem:redundant}
Criterion (iv) --- the conjunction $h_{-1}=h_{-2}=0$, with the two polar
coefficients presented as independent data --- is the form in which the
vanishing criterion is most often stated, and it is correct.  The reality
constraint \eqref{b:reality} shows it is redundant: $h_{-2}=0$ alone already
forces $h_{-1}=0$ and hence both periods to vanish.  The sharp form is (iii).
The same constraint shows that the double-pole coefficient $h_{-2}$ is free in
$\HC$ while the simple-pole coefficient is not free at all, so a
classification that treats $(h_{-1},h_{-2})$ as sixteen free real parameters
overcounts the local obstruction by a factor of two.
\end{remark}

\begin{remark}[What the constraint does \emph{not} make redundant]
\label{b:rem:notredundant}
Constraint \eqref{b:reality} is specific to an \emph{isolated nonreal
parameter point}, where the local first homology is cyclic and the two
periods are attached to the same circle.  It does not make the two one-forms
of Part~II redundant.  On a general $U$ an arbitrary loop still carries two
independent quaternionic periods --- Part~IV realizes every pair
$(a,b)\in\HH^2$ at a single hole --- and testing the ordinary residue
$h_{-1}$ alone still fails to detect the constant obstruction, since
$h_{-1}=0$ with $h_{-2}=\iota d\neq0$ gives $a=0$ and $b=2\pi d\ne0$.  What
\eqref{b:reality} says is only that the eight real coordinates of the pair
$(a,b)$ are repackaged, at one isolated sphere, into the eight real
coordinates of one element of $\HC$.
\end{remark}

\subsection{Faithful detection of singularities}
\label{b:sub:detect}

\begin{corollary}[Removability is a statement about $\Hg$, with no growth
hypothesis]
\label{b:cor:removable-iff}
Let $g$ be axially monogenic on a punctured tube around $\Sp$.  Then $g$
extends to an axially monogenic field across $\Sp$ if and only if $\Hg$
extends holomorphically across $p$.  No growth, integrability, or
meromorphy hypothesis is needed in either direction.
\end{corollary}

\begin{proof}
If $g$ extends, \eqref{b:H} defines $\Hg$ on the full disk and
Theorem~\ref{b:thm:invariant} makes it holomorphic there.  Conversely,
suppose $\Hg$ extends holomorphically to the disk $\abs{z-p}<R$.  Choose a
holomorphic double primitive $F_0$ of the extended function on that
(simply connected) disk and put $g_0=\lap\I(F_0)$, which is axially monogenic
on the whole circularized disk, the disk staying in $\{y>0\}$.  Then
$g-g_0$ has vanishing invariant on the punctured disk, so by
Theorem~\ref{b:thm:range} it equals $\lap\I(\iota(za+b))$ for some
$a,b\in\HH$; that field is smooth across the centre.  Hence $g=g_0+
\lap\I(\iota(za+b))$ extends.
\end{proof}

\begin{remark}[Scope]
\label{b:rem:scope}
Corollary~\ref{b:cor:removable-iff} is a clean equivalence, but it transfers
the problem rather than solving it: deciding whether $\Hg$ extends still
requires information about $g$.  What the classification of Part~VII adds is
the dictionary between hypotheses on $g$ --- local integrability, a pointwise
growth bound, finite energy --- and the pole order of $\Hg$.  Two boundary
cases should be kept in view here.  First, $\Hg$ may have an essential
singularity at $p$: \eqref{b:residue-pair} and hence
Corollary~\ref{b:cor:h2} remain valid, so a field with an essential
singularity can have both periods zero and a single-valued primitive on the
punctured tube while being wildly nonremovable.  Second, the statements of
this part are proved in the \emph{axial} class and at a \emph{nonreal}
$p$; the constants in \eqref{b:one-coefficient} contain $1/v$ and degenerate
as $v\to0$, and no claim is made here about real points, about nonaxial
Cauchy--Fueter fields, or about uniformity as the sphere collapses onto the
real axis.
\end{remark}

%% ==================== part4_5.tex ====================

% =====================================================================
% PART IV.  Logarithmic modes, the split exact sequence, and the cokernel
% PART V.   Reflection: domains meeting the real axis
% Writer key: c
% =====================================================================

\section{Logarithmic modes and the split exact sequence}
\label{c:sec:modes}

Part~II produced a criterion: a single-valued slice-regular Fueter
primitive exists on $\Om_U$ exactly when the two closed one-forms
$\om_g$ and $\thet_g$ have vanishing periods on $U$.  A criterion by
itself does not say that every formally admissible period vector is
attained, nor does it produce a normal form for the failure.  This part
supplies both.  We exhibit, for each pair $(a,b)\in\HH^2$ and each
puncture $p$, a single explicit globally single-valued axially monogenic
field whose period pair is $(a,b)$; we assemble these fields into a right
inverse of the period map; and we conclude that the cokernel of the
Fueter map on a base with $m$ independent holes has real dimension
exactly $8m$, with the direct-sum decomposition realized by concrete
functions rather than by an abstract quotient.

Throughout, $U\subset\Cplus$ is a connected open set,
$\Om_U=\{x+Iy:\ x+\iota y\in U,\ I\in\SSS\}$ is its circularization,
$\SR(\Om_U)$ is the space of slice-regular functions induced by
holomorphic stems on $U$, and $\AM(\Om_U)$ is the space of axially
monogenic (axially Fueter-regular) fields $g=P+IQ$ on $\Om_U$.  All
quaternionic coefficients multiply on the right, and $\iota$ is the
central complex unit of $\HC$, never an element of $\SSS$.

\subsection{The two-parameter logarithmic mode}

Fix $p=u+\iota v\in\CC\setminus U$ and $a,b\in\HH$.  On any simply
connected patch of $\CC\setminus\{p\}$ choose a branch of $\log(z-p)$
and set
\begin{equation}
  F_{p;a,b}(z)\;=\;\frac{za+b}{2\pi\iota}\,\log(z-p),
  \qquad a,b\in\HH .
  \label{c:eq:logstem}
\end{equation}
This is a holomorphic $\HC$-valued stem on the patch.  Two branches
differ by $2\pi\iota n$ with $n\in\ZZ$, so the two stems differ by
$n(za+b)$, which is a complex-affine stem with \emph{quaternionic}
coefficients; by the affine-kernel corollary of Part~II (Corollary~\ref{b:cor:kernel}) it is annihilated by $\lap\I(\cdot)$.  Hence the local fields $\lap\I(F_{p;a,b})$ agree on
overlaps and define one single-valued field.

\begin{proposition}[Branch-free formula, potentials, and period
normalization]
\label{c:prop:modes}
Let $p=u+\iota v$, let $\xi=x-u$ and $\eta=y-v$ be the offsets from the
centre, and let $\rho=\abs{z-p}=(\xi^2+\eta^2)^{1/2}$.  Write
$X=xa+b$.  Then the local fields $\lap\I(F_{p;a,b})$ patch to a single
smooth axially monogenic field
\[
  G_{p;a,b}\;=\;\lap\I\!\left(\frac{za+b}{2\pi\iota}\log(z-p)\right)
  \ \in\ \AM\bigl(\Om_{U}\bigr)
  \qquad\text{for every }U\subset\Cplus\setminus\{p\},
\]
whose axial components are given by the branch-free expressions
\begin{align}
  P_{p;a,b}&=\frac1{\pi y}\left[\frac{X\xi+y\,a\,\eta}{\rho^{2}}
       +a\log\rho\right],
  \label{c:eq:mode-P}\\[2pt]
  Q_{p;a,b}&=\frac{y\,a\,\xi-X\eta}{\pi y\rho^{2}}
       +\frac{X\log\rho}{\pi y^{2}} .
  \label{c:eq:mode-Q}
\end{align}
Its two potentials, computed from any branch $\log(z-p)=\log\rho+\iota\vartheta$,
are
\begin{equation}
  \pot=\frac{1}{2\pi}\left(a\vartheta-\frac{X\log\rho}{y}\right),
  \qquad
  \potV=\frac{1}{2\pi}\left(b\vartheta
        +\Bigl(ya+\frac{xX}{y}\Bigr)\log\rho\right).
  \label{c:eq:mode-potentials}
\end{equation}
For every closed piecewise-$C^1$ curve $\gamma\subset U$,
\begin{equation}
  \left(\oint_\gamma\om_{G_{p;a,b}},\ \oint_\gamma\thet_{G_{p;a,b}}\right)
  \;=\;\Ind(\gamma,p)\,(a,b).
  \label{c:eq:mode-periods}
\end{equation}
The map $(a,b)\mapsto G_{p;a,b}$ is right $\HH$-linear.  If
$\Ima_{\CC}p<0$, a branch of $\log(z-p)$ exists on all of $\Cplus$, so
$G_{p;a,b}$ is defined on all of $\Om_{\Cplus}$ and both of its periods
vanish there.
\end{proposition}

\begin{proof}
Write $\log(z-p)=\log\rho+\iota\vartheta$ on a branch.  Separating
\eqref{c:eq:logstem} into $F=A+\iota B$ and using $1/\iota=-\iota$,
\[
  A=\frac{X\vartheta+y\,a\log\rho}{2\pi},
  \qquad
  B=\frac{y\,a\,\vartheta-X\log\rho}{2\pi},
\]
where the coefficient $za+b$ has been split as $(xa+b)+\iota ya=X+\iota ya$.
The change of potentials $\pot=B/y$, $\potV=A-xB/y$ of Part~II, \eqref{b:CV}, gives
\eqref{c:eq:mode-potentials} at once; in $\potV$ the multivalued terms
combine as $X\vartheta-xa\vartheta=b\vartheta$.

For the components, use $P=-2\pot_x$ and $Q=2\pot_y$, which is the content of $\dd\pot=\om_g$ in \eqref{b:potential-identities}.  With
$\vartheta_x=-\eta/\rho^2$, $\vartheta_y=\xi/\rho^2$,
$(\log\rho)_x=\xi/\rho^2$, $(\log\rho)_y=\eta/\rho^2$ and $X_x=a$,
\[
  \pot_x=\frac{1}{2\pi}\left[-\frac{a\eta}{\rho^2}-\frac{a\log\rho}{y}
      -\frac{X\xi}{y\rho^2}\right],
  \qquad
  \pot_y=\frac{1}{2\pi}\left[\frac{a\xi}{\rho^2}
      -\frac{X\eta}{y\rho^2}+\frac{X\log\rho}{y^2}\right],
\]
which are exactly \eqref{c:eq:mode-P}--\eqref{c:eq:mode-Q}.  Every term
involves only the real logarithm of a positive number: the argument
$\vartheta$ has cancelled, so \eqref{c:eq:mode-P}--\eqref{c:eq:mode-Q}
are unambiguous functions on $U$, and they are the promised single-valued
representative of the branchwise construction.

The periods are read off from \eqref{c:eq:mode-potentials}.  Continuing
along $\gamma$ returns $\log\rho$ to its initial value and increases
$\vartheta$ by $2\pi\Ind(\gamma,p)$.  Hence $\pot$ gains
$\Ind(\gamma,p)\,a$ and $\potV$ gains $\Ind(\gamma,p)\,b$, and since
$\oint_\gamma\om_g$ and $\oint_\gamma\thet_g$ are precisely these
increments, \eqref{c:eq:mode-periods} follows.  (Equivalently, continuation
of \eqref{c:eq:logstem} changes the stem by $\Ind(\gamma,p)(za+b)$, and
the affine-monodromy corollary of Part~II, Corollary~\ref{b:cor:monodromy}, identifies the period pair.)
Right linearity is visible in \eqref{c:eq:logstem}.  The final assertion
holds because a branch of the logarithm defined on the whole simply
connected half-plane makes $F_{p;a,b}$ a genuine single-valued stem
there, so $G_{p;a,b}$ is itself a global Fueter image.
\end{proof}

\begin{remark}[One formula, nine sources]
Formulas \eqref{c:eq:mode-P}--\eqref{c:eq:mode-Q} are the single
branch-free expression that the source manuscripts wrote in nine
different variable conventions.  Each of them either splits the mode
into two one-parameter kernels carrying respectively the slope and the
constant period, or writes the two-parameter family directly; the
two-parameter form is adopted here because it makes
\eqref{c:eq:mode-periods} a single line and removes an inversion hazard
present in the split notations, where the letter $A$ carries the slope
charge in one convention and the constant charge in another.  The
specializations are recovered by
$G_{p;a,b}=G_{p;1,0}\,a+G_{p;0,1}\,b$.
The potentials \eqref{c:eq:mode-potentials} give a second, independent
check of \emph{both} period normalizations, since $\pot$ gains $a$ and
$\potV$ gains $b$ on one circuit; they are also the form in which the
mode should be evaluated numerically near the real axis, where
\eqref{c:eq:mode-Q} suffers cancellation.
\end{remark}

\begin{corollary}[Holomorphic encoding of a mode]
\label{c:cor:mode-encoding}
With $\Hg$ the holomorphic invariant of Part~III,
\begin{equation}
  \mathcal H_{G_{p;a,b}}(z)
  \;=\;\left(\frac{za+b}{2\pi\iota}\log(z-p)\right)''
  \;=\;\frac{1}{2\pi\iota}\left(\frac{a}{z-p}
        -\frac{pa+b}{(z-p)^{2}}\right).
  \label{c:eq:mode-encoding}
\end{equation}
\end{corollary}

\begin{proof}
Theorem~\ref{b:thm:invariant} of Part~III proves $\Hg=F''$ for every local Fueter primitive stem $F$ of $g$, and \eqref{c:eq:logstem} is such a stem for $G_{p;a,b}$ on every
patch; the right-hand side of \eqref{c:eq:mode-encoding} is single-valued,
so the identity is global.  Differentiating \eqref{c:eq:logstem} twice,
\[
  \Bigl(\tfrac{za+b}{2\pi\iota}\log(z-p)\Bigr)'
   =\frac{a\log(z-p)}{2\pi\iota}+\frac{za+b}{2\pi\iota(z-p)},
\]
and one more differentiation gives
$\frac{2a}{2\pi\iota(z-p)}-\frac{za+b}{2\pi\iota(z-p)^2}$, which equals
the stated expression because $za+b=(z-p)a+(pa+b)$.
\end{proof}

Formula \eqref{c:eq:mode-encoding} is the point of contact between this
part and Part~III: the Laurent coefficients of $\mathcal H_{G_{p;a,b}}$
at $p$ are $h_{-1}=a/(2\pi\iota)$ and $h_{-2}=-(pa+b)/(2\pi\iota)$, and
one checks directly that they satisfy the reality constraint
$h_{-1}=(\iota/v)\ReC(h_{-2})$ of Theorem~\ref{b:thm:reality}.  In particular the eight
real parameters of a mode are already carried by the single $\HC$
coefficient $h_{-2}$.

\begin{remark}[The mode functions themselves are not claimed to be new]
\label{c:rem:known-modes}
For $p=\iota$ the reflected logarithm of Part~V has derivative
$\pi^{-1}(1+z^2)^{-1}$, so a local primitive is
$\pi^{-1}\arctan z$ up to a real constant, and the functions with local
primitives $(2\pi)^{-1}\arctan z$ and $(2\pi)^{-1}z\arctan z$ are the
spherical Cauchy kernels $N^{+}$ and $N^{-}$ of the inverse-Fueter
literature.  In that normalization
$G^{\mathrm{sym}}_{\iota;0,1}=2N^{+}$ and
$G^{\mathrm{sym}}_{\iota;1,0}=2N^{-}$.  We do not claim the invention of
these kernels, and we record that the source manuscripts of this merge
disagree about which published article is the primary locus for them:
some cite a 2011 paper on sphere-integrated Cauchy kernels, others a
2014 example on inverse Fueter kernels.  That attribution is not settled
here, and the identification with $N^{\pm}$ is asserted only with the
normalization just spelled out; see the literature paragraph of
Part~XI.  What is proposed here is their normalization by \emph{two}
global periods, the resulting completeness statement
\eqref{c:eq:exact}, and the classification and energy laws of the later
parts.
\end{remark}

\subsection{Finite winding bases}

The splitting theorem needs a topological hypothesis.  We state it as a
hypothesis about cycles and winding numbers, not about the regularity of
the complement, so that no boundary smoothness is imposed on an arbitrary
closed set.

\begin{definition}[Finite winding basis]
\label{c:def:winding}
A connected open $U\subset\Cplus$ \emph{has a finite winding basis of
size $m$} if there are points $p_1,\dots,p_m\in\Cplus\setminus U$ and
piecewise-$C^1$ cycles $\gamma_1,\dots,\gamma_m\subset U$ whose homology
classes form a $\ZZ$-basis of $H_1(U;\ZZ)$ and satisfy
\begin{equation}
  \Ind(\gamma_j,p_k)=\delta_{jk},\qquad 1\le j,k\le m .
  \label{c:eq:winding-basis}
\end{equation}
The case $m=0$ means $H_1(U;\ZZ)=0$.
\end{definition}

Every upper domain obtained by deleting $m$ points from a simply
connected upper domain has such a basis, as does one obtained by deleting
$m$ disjoint closed Jordan regions, with one point chosen inside each.
The hypothesis is used only through \eqref{c:eq:winding-basis}; nothing
below requires the complementary set to have rectifiable or even
locally connected boundary.

\subsection{The split range theorem}

\begin{theorem}[Split exact sequence and the $8m$ cokernel]
\label{c:thm:split}
Let $U\subset\Cplus$ be connected with a finite winding basis
$(\gamma_j,p_j)_{j=1}^m$, and define the period map
\[
  \Per(g)=\bigl(a_{\gamma_j}(g),\,b_{\gamma_j}(g)\bigr)_{j=1}^{m}
   \in\HH^{2m},
  \qquad
  a_\gamma(g)=\oint_\gamma\om_g,\quad b_\gamma(g)=\oint_\gamma\thet_g .
\]
Then the sequence of right quaternionic vector spaces
\begin{equation}
  0\longrightarrow\Aff\longrightarrow\SR(\Om_U)
   \xrightarrow{\ \lap\I(\cdot)\ }\AM(\Om_U)
   \xrightarrow{\ \Per\ }\HH^{2m}\longrightarrow0
  \label{c:eq:exact}
\end{equation}
is exact, where the first map is $(a,b)\mapsto(q\mapsto qa+b)$.  The
final surjection has the explicit right $\HH$-linear section
\begin{equation}
  S\bigl((a_j,b_j)_{j=1}^m\bigr)=\sum_{j=1}^{m}G_{p_j;a_j,b_j},
  \qquad \Per\circ S=\operatorname{id}_{\HH^{2m}} .
  \label{c:eq:section}
\end{equation}
Consequently
\begin{equation}
  \AM(\Om_U)=\lap\I\bigl(\SR(\Om_U)\bigr)\ \oplus\ S(\HH^{2m}),
  \qquad
  \dim_{\RR}\frac{\AM(\Om_U)}{\lap\I(\SR(\Om_U))}=8m,
  \label{c:eq:directsum}
\end{equation}
every $g\in\AM(\Om_U)$ has the unique decomposition
\begin{equation}
  g=\lap\I(F)+\sum_{j=1}^{m}G_{p_j;\,a_{\gamma_j}(g),\,b_{\gamma_j}(g)},
  \label{c:eq:decomposition}
\end{equation}
with $F$ unique modulo a stem $za+b$, and
$\Pi=\operatorname{id}-S\circ\Per$ is a right-linear projection of
$\AM(\Om_U)$ onto the space of targets possessing a single-valued global
Fueter primitive.
\end{theorem}

\begin{proof}
Exactness at $\SR(\Om_U)$ is Corollary~\ref{b:cor:kernel} of Part~II: a
slice-regular $f$ induced by a stem $F$ satisfies $\lap f=0$ on a
connected base exactly when $\dd\pot=\dd\potV=0$, i.e.\ when $F=za+b$
with $a,b\in\HH$; the coefficients are quaternionic and not merely
elements of $\HC$, which is what makes the sequence's left-hand term
eight-real-dimensional rather than sixteen.

Exactness at $\AM(\Om_U)$ is the global period criterion Theorem~\ref{b:thm:period} of Part~II, combined with the basis hypothesis.  The period functionals $a_\gamma$
and $b_\gamma$ depend only on the homology class of $\gamma$, because
$\om_g$ and $\thet_g$ are closed; so $\Per(g)=0$ forces both periods to
vanish on every cycle of $U$, and the criterion produces a single-valued
primitive.  Conversely the periods of a Fueter image vanish, since
$\om_g=\dd\pot$ and $\thet_g=\dd\potV$ globally in that case.

Surjectivity of $\Per$ and the section identity follow from
Proposition~\ref{c:prop:modes} together with \eqref{c:eq:winding-basis}:
the $k$th period pair of $S((a_j,b_j)_j)$ is
$\sum_j\Ind(\gamma_k,p_j)(a_j,b_j)=(a_k,b_k)$.  Because $\gamma_k$ lies
in $U$ and $p_j\notin U$, each mode in the sum is regular on $\Om_U$, so
the sum is a legitimate element of $\AM(\Om_U)$.

For the direct sum, write $g=(g-S\Per g)+S\Per g$; the first summand
lies in $\ker\Per=\lap\I(\SR)$ by the previous paragraph.  If a vector
$\mathbf w\in\HH^{2m}$ has $S\mathbf w\in\ker\Per$ then
$\mathbf w=\Per S\mathbf w=0$, so the sum is direct.  Applying $\Per$ to
any decomposition of the shape \eqref{c:eq:decomposition} recovers the
mode coefficients, which proves their uniqueness; the remaining ambiguity
in $F$ is the affine kernel.  Finally $\HH^{2m}$ has real dimension $8m$,
and $\Pi^2=\Pi$ because $\Per S=\operatorname{id}$.
\end{proof}

\begin{remark}[What is canonical and what is a choice]
\label{c:rem:canonical}
The quotient $\AM(\Om_U)/\lap\I(\SR(\Om_U))$ and the period coordinates
on it are canonical: they depend only on $U$ and on the chosen homology
basis, not on the auxiliary points.  The \emph{complement}
$S(\HH^{2m})$ is not canonical.  Moving $p_j$ to another point $p_j'$
of the same complementary component changes the representative
$G_{p_j;a,b}$ to $G_{p_j';a,b}$, and the difference has zero periods on
every cycle of $U$ — the two points have the same winding numbers against
$\gamma_1,\dots,\gamma_m$ — hence is itself a global Fueter image.  So
the splitting moves but the decomposition class does not.  Continuity of
$\Pi$ and $S$ for compact-open convergence, and the corresponding closed
range statement, are proved in Part~VI; nothing in
Theorem~\ref{c:thm:split} asserts boundedness in any weighted or Hilbert
norm.
\end{remark}

\begin{example}[Neither obstruction can be dropped]
\label{c:ex:independent}
Let $U=\Cplus\setminus\{p\}$ and let $\gamma$ be a small positively
oriented circle about $p$.  By \eqref{c:eq:mode-periods},
\[
  \Per\bigl(G_{p;0,1}\bigr)=(0,1),\qquad
  \Per\bigl(G_{p;1,0}\bigr)=(1,0).
\]
Thus $G_{p;0,1}$ has an exact first form, so the slope potential $\pot$
exists globally on $U$ and the first of the two classical integrations
succeeds; nevertheless no global slice-regular Fueter primitive exists,
because the second period is $1$.  The field $G_{p;1,0}$ fails already at
the first integration.  Neither test implies the other, and neither may
be omitted.  Right multiplication by $1,i,j,k$ turns each of the two
families into four real directions, giving the eight independent real
directions counted by \eqref{c:eq:directsum} at a single hole.  This also
shows that the second obstruction is not a residue of a badly chosen
first primitive: by the elimination identity $\etaC-\dd(x\pot)=\thet_g$ of Part~II (Proposition~\ref{b:prop:elimination}), $\thet_g$ is computed from the
target alone, and it is nonexact here for every choice of $\pot$.
\end{example}


\section{Reflection: domains meeting the real axis}
\label{c:sec:reflection}

Everything so far took place over a base $U\subset\Cplus$, so the
circularized domain $\Om_U$ avoided the real axis of $\HH$.  One might
suspect that the obstruction of Theorem~\ref{c:thm:split} is an artifact
of that avoidance — that a domain which contains real points, and whose
slice functions are therefore rigid there, would have no room for a
period.  It does not.  The obstruction survives verbatim, with a count
that is \emph{not} the first Betti number of the base and is completely
insensitive to holes that conjugation fixes.

\subsection{Symmetric bases, parity, and invariance of the two forms}

Let $D\subset\CC$ be open, connected, invariant under
$\tau(z)=\bar z$, and meeting $\RR$.  Its circularization
\[
  \Om_D=\{x+Iy:\ x+\iota y\in D,\ I\in\SSS\}\subset\HH
\]
contains the real interval $D\cap\RR$.  A holomorphic stem on $D$ is
required to satisfy the reflection-reality condition
$F(\bar z)=\kappa(F(z))$, where $\kappa$ is the conjugation of $\HC$
fixing $\HH$; equivalently $A$ is even and $B$ is odd in $y$.  This is
the usual notion of a slice-regular function on an axially symmetric
slice domain.

For $g\in\AM(\Om_D)$ the axial components are recovered from a single
slice by
\begin{equation}
  P(x,y)=\frac{g(x+Iy)+g(x-Iy)}2,\qquad
  Q(x,y)=-\frac I2\bigl(g(x+Iy)-g(x-Iy)\bigr)
  \label{c:eq:parity}
\end{equation}
for any fixed $I\in\SSS$, which exhibits $P$ as even and $Q$ as odd in
$y$ and smooth across $y=0$.

\begin{lemma}[The two forms extend and are $\tau$-invariant]
\label{c:lem:invariance}
For $g\in\AM(\Om_D)$ the one-forms
$\om_g=\tfrac12(-P\,\dd x+Q\,\dd y)$ and
$\thet_g=\tfrac12\bigl((xP+yQ)\dd x+(yP-xQ)\dd y\bigr)$
are smooth on all of $D$, closed on all of $D$, and satisfy
\begin{equation}
  \tau^{*}\om_g=\om_g,\qquad \tau^{*}\thet_g=\thet_g .
  \label{c:eq:tauinvariance}
\end{equation}
\end{lemma}

\begin{proof}
Neither form contains a division by $y$, so smoothness across the axis
follows from \eqref{c:eq:parity}.  Closedness is Lemma~\ref{b:lem:forms}(i) of Part~II, valid where the Vekua system holds, i.e.\ for $y\ne0$; the
$\dd x\wedge\dd y$ coefficients $(Q_x+P_y)/2$ and
$\tfrac12\{y(P_x-Q_y)-2Q-x(Q_x+P_y)\}$ are continuous on $D$ and vanish
off the axis, hence vanish on it by continuity.  For
\eqref{c:eq:tauinvariance}: under $\tau$ one has $\dd x\mapsto\dd x$ and
$\dd y\mapsto-\dd y$, while $P$ is even and $Q$ is odd.  The $\dd x$
coefficients $-P/2$ and $(xP+yQ)/2$ are therefore even, and the $\dd y$
coefficients $Q/2$ and $(yP-xQ)/2$ are odd; each form is thus preserved.
\end{proof}

\begin{lemma}[The upper meridian is connected]
\label{c:lem:upper}
If $D$ is connected, open and $\tau$-invariant, then
$U=D\cap\Cplus$ is connected.
\end{lemma}

\begin{proof}
Let $z_0,z_1\in U$ and join them by a path in $D$.  Fold the path upward
by replacing $x+\iota y$ with $x+\iota\abs y$; the result is a continuous
path in $D\cap\{y\ge0\}$, because $D$ is $\tau$-invariant.  Its image is
compact and therefore has positive distance $\delta$ from the closed
complement of $D$.  Translating the whole folded path by
$\iota\delta/2$ keeps it in $D$ and places it strictly above $\RR$.
Two short vertical segments, from $z_0$ and $z_1$ up to the translated
endpoints, stay in $D\cap\{y>0\}$ for $\delta$ small enough.  The
concatenation is a path in $U$.
\end{proof}

\subsection{The criterion across the real axis}

\begin{theorem}[Real-axis period criterion]
\label{c:thm:realaxis}
Let $D$ be connected, open, $\tau$-invariant and meeting $\RR$, and let
$g\in\AM(\Om_D)$.  There exists $f\in\SR(\Om_D)$ with $\lap f=g$ if and
only if $\om_g$ and $\thet_g$ have vanishing periods on $D$, equivalently
if and only if they have vanishing periods on the upper meridian
$U=D\cap\Cplus$.  Such an $f$ is unique modulo $q\mapsto qa+b$,
$a,b\in\HH$.  No finite-topology hypothesis and no boundary regularity
hypothesis is used.
\end{theorem}

\begin{proof}
Necessity is the restriction of the criterion of Part~II to $\Om_U$,
together with the observation that for a global primitive the potentials
$\pot=B/y$ and $\potV=A-xB/y$ extend smoothly and evenly across $y=0$:
locally $B(x,y)/y=\int_0^1B_y(x,ty)\,\dd t$ because $B$ is odd, and the identities \eqref{b:potential-identities} then hold on all of $D$ by continuity.  This also gives uniqueness, since a difference of two
primitives has both potentials locally constant, hence constant on the
connected $D$.

For sufficiency assume only the weaker hypothesis, that both periods
vanish on $U$; the construction will produce single-valued potentials on
all of $D$, so that the periods on $D$ vanish as well and the two
conditions are indeed equivalent.  Lemma~\ref{c:lem:upper} makes $U$
connected, and Part~II supplies potentials $\pot,\potV$ on $U$ with
$\dd\pot=\om_g$, $\dd\potV=\thet_g$.

The extension across the axis is the symmetrization argument.  Fix a real
point $x_*\in D\cap\RR$ and a small disk $\Delta\subset D$ centred there
and invariant under $\tau$.  By Lemma~\ref{c:lem:invariance} both forms
are smooth and closed on $\Delta$, which is simply connected, so each has
a primitive there; call them $\pot_0,\potV_0$, normalized to agree with
$\pot,\potV$ on the (connected) upper half-disk.  Replacing $\pot_0$ by
$\tfrac12(\pot_0+\pot_0\circ\tau)$ and $\potV_0$ likewise does not change
their differentials, by \eqref{c:eq:tauinvariance}, and makes both even
in $y$; the normalization is unaffected because the averaged function
differs from the original by a constant that vanishes on the real
diameter.  Hence $\pot,\potV$ extend to $\Delta$ as even functions, and
any two such local extensions agree on overlaps because they agree above
the axis and are even.  Finally extend $\pot,\potV$ to the remaining
lower part of $D$ by $\pot(\tau z)=\pot(z)$, $\potV(\tau z)=\potV(z)$;
this is consistent with the extensions already made near the axis and, by
\eqref{c:eq:tauinvariance}, still satisfies $\dd\pot=\om_g$ and
$\dd\potV=\thet_g$.  The potentials are now single-valued on all of $D$,
so both forms are exact there and every period on $D$ vanishes.

Set $A=x\pot+\potV$ and $B=y\pot$ on the extended domain.  Then $A$ is
even, $B$ is odd, and the Cauchy--Riemann system $A_x=B_y$, $A_y=-B_x$
holds off the axis by Lemma~\ref{b:lem:forms}(iii); both sides are
continuous, so it holds on the axis too.  Therefore $F=A+\iota B$ is a
holomorphic stem on $D$ satisfying $F(\bar z)=\kappa(F(z))$, its induced
slice function $f$ is slice-regular on $\Om_D$, and $\lap f=g$ off the
axis, hence everywhere by continuity.
\end{proof}

\subsection{Which holes actually obstruct}

We now fix a finite-type class in which the count can be made.  The
complement of $D$ in the Riemann sphere has finitely many components;
each bounded component (a \emph{hole}) is a point or a closed Jordan
region, and the unbounded component contains $\infty$ and may reduce to
$\{\infty\}$.  No assertion about wild boundaries is needed, and the
same conclusions hold for a finite winding basis in the sense of
Definition~\ref{c:def:winding}.  Conjugation permutes the holes: it
either fixes a hole or exchanges it with a distinct one.  Write $s$ for
the number of fixed holes and $m$ for the number of exchanged pairs, so
that the first Betti number of $D$ is $s+2m$.

\begin{lemma}[Reflection on planar homology]
\label{c:lem:reflection}
The classes of positively oriented loops around the holes freely generate
$H_1(D;\ZZ)$, and for a hole $K$ with class $e_K$,
\[
  \tau_*e_K=-e_{\tau K}.
\]
Consequently a $\tau$-invariant closed form has period zero around a
conjugation-fixed hole, and around a conjugate pair its counterclockwise
periods are negatives of one another.  In particular
\begin{equation}
  \dim_{\RR}H^1_{\mathrm{dR}}(D;\RR)^{+}=m,
  \label{c:eq:invariantcohom}
\end{equation}
where the superscript denotes the $+1$ eigenspace of $\tau^*$.
\end{lemma}

\begin{proof}
Cutting $D$ along disjoint arcs joining each hole to the unbounded
complementary component leaves a simply connected region; restoring the
arcs one at a time produces exactly one independent loop per hole, and
these classes are detected by the integer winding numbers about one
selected point in each hole.  Conjugation reverses planar orientation and
permutes the selected points, which gives $\tau_*e_K=-e_{\tau K}$.  For
an invariant closed form $\alpha$ and a cycle $\gamma$,
$\int_{\tau\gamma}\alpha=\int_\gamma\tau^*\alpha=\int_\gamma\alpha$;
applied to $e_K$ with $\tau K=K$ this gives
$\int_{e_K}\alpha=-\int_{e_K}\alpha$, hence $0$, there being no
$2$-torsion in $(\HH,+)$ or $(\RR,+)$.  Applied to an exchanged pair it
gives the opposite-period relation.  Those are the only relations imposed
on the free basis, so one real coordinate survives per exchanged pair,
which is \eqref{c:eq:invariantcohom}.
\end{proof}

The two families of source manuscripts state the resulting count in two
different ways: some count $m$ as the rank of a winding basis of the
upper meridian $U$, others as the number of conjugation-exchanged pairs
of holes of $D$.  Neither family proves that the two indices agree.  The
missing link is the following elementary observation, which we isolate
because without it the reflection theorem and the product-domain theorem
read as two different theorems about two different integers.

\begin{lemma}[Invariant holes meet the real axis]
\label{c:lem:invariant-holes}
Let $K\subset\CC$ be nonempty, connected and $\tau$-invariant.  Then
$K\cap\RR\ne\emptyset$.  Consequently every conjugation-fixed hole of $D$
meets $\RR$, while the two members of an exchanged pair are disjoint from
$\RR$.
\end{lemma}

\begin{proof}
Suppose $K\cap\RR=\emptyset$.  Then
$K=(K\cap\Cplus)\sqcup(K\cap\{y<0\})$ is a partition of $K$ into two
relatively open subsets.  Both are nonempty: if $K$ were contained in one
open half-plane, $\tau K=K$ would be contained in the other, forcing
$K=\emptyset$.  This contradicts connectedness.  For the last claim, if
$K$ is one member of an exchanged pair and $x\in K\cap\RR$, then
$x=\tau(x)\in\tau K$, so $K\cap\tau K\ne\emptyset$; but distinct
complementary components are disjoint.
\end{proof}

\begin{proposition}[The two hole counts agree]
\label{c:prop:bridge}
Let $D$ be as above, with $s$ fixed holes $K_1,\dots,K_s$ and $m$
exchanged pairs, and choose the upper member $L_j\subset\Cplus$ of the
$j$th pair together with a point $p_j\in L_j$.  Then
$U=D\cap\Cplus$ is connected and has a finite winding basis of size $m$,
realized by $p_1,\dots,p_m$ and small positively oriented loops
$\gamma_1,\dots,\gamma_m$ in $U$ encircling $L_1,\dots,L_m$.  In
particular the rank of a winding basis of the upper meridian equals the
number of conjugation-exchanged pairs of holes of $D$.
\end{proposition}

\begin{proof}
Connectedness is Lemma~\ref{c:lem:upper}.  Compute the complement of $U$
in the Riemann sphere:
\[
  \widehat\CC\setminus U
   =\bigl(\widehat\CC\setminus D\bigr)\cup\bigl(\{y\le0\}\cup\{\infty\}\bigr).
\]
The set $\{y\le0\}\cup\{\infty\}$ is connected and contains $\infty$, so
it lies in one component together with the unbounded complementary
component of $D$.  By Lemma~\ref{c:lem:invariant-holes} each fixed hole
$K_i$ meets $\RR\subset\{y\le0\}$ and is connected, so it lies in that
same component.  Again by Lemma~\ref{c:lem:invariant-holes} each $L_j$ is
a connected compact subset of the open half-plane $\Cplus$, disjoint from
$\{y\le0\}$, from the other holes, and from the unbounded component;
hence each $L_j$ is a component of $\widehat\CC\setminus U$ on its own,
as is the lower member $\tau L_j$ — which, however, lies in $\{y<0\}$ and
so has already been absorbed into the big component.  Therefore
$\widehat\CC\setminus U$ has exactly $m+1$ components, and $H_1(U;\ZZ)$
is free of rank $m$, generated by loops separating $L_1,\dots,L_m$ from
$\infty$.  Small loops $\gamma_j$ around $L_j$ inside $U$ represent those
generators and satisfy $\Ind(\gamma_j,p_k)=\delta_{jk}$.
\end{proof}

\subsection{Reflected modes and the range theorem across the axis}

The modes of Part~IV are defined on $\Cplus\setminus\{p\}$ and need not
extend across $\RR$.  Their reflected differences do.

\begin{proposition}[Reflected modes]
\label{c:prop:reflected}
For $p=u+\iota v$ with $v>0$ and $a,b\in\HH$, put
\begin{equation}
  G^{\mathrm{sym}}_{p;a,b}=G_{p;a,b}-G_{\bar p;a,b}
  \quad\text{on }\ y>0,
  \qquad
  F^{\mathrm{sym}}_{p;a,b}(z)=\frac{za+b}{2\pi\iota}
      \log\frac{z-p}{z-\bar p}.
  \label{c:eq:symlog}
\end{equation}
Then $G^{\mathrm{sym}}_{p;a,b}$ extends to a smooth axially monogenic
field on $\HH\setminus\Sp$, where $\Sp=u+v\SSS$; it is locally
$\lap\I(F^{\mathrm{sym}}_{p;a,b})$.  Its period pair on an upper
counterclockwise loop $\gamma$ is $\Ind(\gamma,p)(a,b)$, and on a
counterclockwise loop about $\bar p$ it is $(-a,-b)$.  Its restriction to
the real axis is
\begin{equation}
  G^{\mathrm{sym}}_{p;a,b}(x)
   =-\frac{4v}{\pi d}\,a+\frac{4v(x-u)}{\pi d^{2}}\,(xa+b),
  \qquad d=(x-u)^2+v^2 .
  \label{c:eq:realtrace-general}
\end{equation}
\end{proposition}

\begin{proof}
Away from $\RR$ the field is the difference of two modes, hence axially
monogenic; by the last clause of Proposition~\ref{c:prop:modes} the
subtracted term $G_{\bar p;a,b}$ is regular on all of $\Om_{\Cplus}$ and
has zero periods there, so the period statement on upper loops follows
from \eqref{c:eq:mode-periods}.  Near a real point $x_*\in\RR$ the
quotient $R(z)=(z-p)/(z-\bar p)$ is holomorphic and nonzero, and
$\abs{R}=1$ on $\RR$; choose the local branch $L$ of $\log R$ with purely
imaginary value at $x_*$.  Uniqueness of the branch then gives
$L(\bar z)=-\overline{L(z)}$, so $L/(2\pi\iota)$ is real on the real
diameter, and \eqref{c:eq:symlog} is a reflection-real holomorphic stem
on that disk.  Its Fueter image is smooth across $\RR$.  Two admissible
branches differ by an integer multiple of $za+b$, whose image is zero, so
the local images agree with each other and, above the axis, with
$G_{p;a,b}-G_{\bar p;a,b}$.  For the period around $\bar p$, the winding
numbers of numerator and denominator of $R$ are exchanged.

For \eqref{c:eq:realtrace-general}, use the real-trace identity \eqref{e:eq:realtrace} of Part~IX, $(\lap\I F)(x)=-2F''(x)$, together with
Corollary~\ref{c:cor:mode-encoding} applied to both $p$ and $\bar p$:
\[
  \bigl(F^{\mathrm{sym}}_{p;a,b}\bigr)''(x)
  =\frac1{2\pi\iota}\left[\frac{a}{x-p}-\frac{pa+b}{(x-p)^2}
    -\frac{a}{x-\bar p}+\frac{\bar pa+b}{(x-\bar p)^2}\right].
\]
Put $\xi=x-u$, $w=\xi-\iota v$, $w'=\xi+\iota v$, so $ww'=d$; and split
the quaternionic data as $\alpha=a$, $\beta=ua+b$, whence
$pa+b=\beta+\iota v\alpha$ and $\bar pa+b=\beta-\iota v\alpha$.  All the
complex factors below are central in $\HC$, so the right-hand
coefficients may be kept on the right throughout.  From
$w'-w=2\iota v$ and $w'+w=2\xi$,
\[
  \frac1w-\frac1{w'}=\frac{2\iota v}{d},\qquad
  \frac1{w^{2}}-\frac1{w'^{2}}=\frac{4\iota v\xi}{d^{2}},\qquad
  \frac1{w^{2}}+\frac1{w'^{2}}=\frac{2(\xi^{2}-v^{2})}{d^{2}} .
\]
Therefore the bracket equals
\[
  \frac{2\iota v}{d}\alpha
  -\frac{4\iota v\xi}{d^{2}}\beta
  -\iota v\,\frac{2(\xi^{2}-v^{2})}{d^{2}}\alpha
  =\frac{2\iota v}{d^{2}}
   \bigl[(d-\xi^{2}+v^{2})\alpha-2\xi\beta\bigr]
  =\frac{4\iota v}{d^{2}}\bigl[v^{2}\alpha-\xi\beta\bigr],
\]
using $d-\xi^2+v^2=2v^2$.  Dividing by $2\pi\iota$ and multiplying by
$-2$,
\[
  G^{\mathrm{sym}}_{p;a,b}(x)
   =-\frac{4v}{\pi d^{2}}\bigl(v^{2}a-\xi(ua+b)\bigr),
\]
which is \eqref{c:eq:realtrace-general}: indeed
$-\frac{4v}{\pi d}a+\frac{4v\xi}{\pi d^{2}}(xa+b)
 =\frac{4v}{\pi d^{2}}\bigl[(\xi^{2}-d)a+\xi(ua+b)\bigr]
 =\frac{4v}{\pi d^{2}}\bigl[-v^{2}a+\xi(ua+b)\bigr]$,
since $xa+b=\xi a+(ua+b)$.  The same
formula follows, without the trace identity, by expanding the
single-valued imaginary part of \eqref{c:eq:symlog},
$-\tfrac1{4\pi}\log\frac{(x-u)^2+(y-v)^2}{(x-u)^2+(y+v)^2}
 = \tfrac{vy}{\pi d}+O(y^3)$, in \eqref{c:eq:mode-P}.
\end{proof}

\begin{theorem}[Reflection-sensitive range theorem]
\label{c:thm:reflection-range}
Let $D$ be as in Proposition~\ref{c:prop:bridge}, with $s$ conjugation
fixed holes and $m$ exchanged pairs, and let $\gamma_j,p_j$ be the upper
detecting loops and points.  Then
\begin{equation}
  0\longrightarrow\Aff\longrightarrow\SR(\Om_D)
   \xrightarrow{\ \lap\I(\cdot)\ }\AM(\Om_D)
   \xrightarrow{\ \Per\ }\HH^{2m}\longrightarrow0
  \label{c:eq:exact-sym}
\end{equation}
is exact, with $\Per(g)=(a_{\gamma_j}(g),b_{\gamma_j}(g))_{j=1}^m$ and
with the explicit right inverse
$S((a_j,b_j)_j)=\sum_j G^{\mathrm{sym}}_{p_j;a_j,b_j}$.  Hence every
$g\in\AM(\Om_D)$ is uniquely
\begin{equation}
  g=\lap\I(F)+\sum_{j=1}^{m}G^{\mathrm{sym}}_{p_j;a_{\gamma_j}(g),\,
     b_{\gamma_j}(g)},
  \qquad
  \dim_{\RR}\operatorname{coker}\lap\I(\cdot)=8m,
  \label{c:eq:8m}
\end{equation}
and the cokernel is canonically
$H^1_{\mathrm{dR}}(D;\RR)^{+}\otimes_{\RR}\HH^{2}$
via $g\mapsto([\om_g],[\thet_g])$.  The count $8m$ is independent of
$s$: it is \emph{not} $8(s+2m)$, and it is \emph{not} eight times the
first Betti number of $D$.
\end{theorem}

\begin{proof}
Exactness at $\SR(\Om_D)$ is the affine kernel again, now on the
connected base $D$.  Exactness at $\AM(\Om_D)$: if the $m$ selected upper
period pairs vanish, then by Lemma~\ref{c:lem:reflection} all periods of
the invariant forms $\om_g,\thet_g$ vanish on every generator of
$H_1(D;\ZZ)$ — automatically around each fixed hole, and by the
opposite-period relation around each lower hole — so
Theorem~\ref{c:thm:realaxis} produces a global primitive.  Conversely the
periods of a Fueter image vanish.

For surjectivity, $G^{\mathrm{sym}}_{p_j;a,b}$ is regular on $\Om_D$
because its singular sphere $S_{p_j}$ lies over the missing hole $L_j$,
and by Proposition~\ref{c:prop:reflected} its selected period pair is
$(a,b)$ at the $j$th pair and $(0,0)$ at all the others.  Thus
$\Per S=\operatorname{id}$ and the argument of
Theorem~\ref{c:thm:split} applies verbatim, giving \eqref{c:eq:8m}.

The cohomological description is a restatement of what has just been
proved, not an appeal to an unproved quaternionic sheaf theory: by
Lemma~\ref{c:lem:invariance} the classes $[\om_g],[\thet_g]$ lie in the
$\tau$-invariant part of $H^1_{\mathrm{dR}}(D;\HH)$, and
\eqref{c:eq:invariantcohom} gives
$\dim_\RR\bigl(H^1_{\mathrm{dR}}(D;\RR)^{+}\otimes\HH^2\bigr)=8m$.
Explicit representatives are furnished by $S$.
\end{proof}

\begin{corollary}[Exact topological criterion for surjectivity]
\label{c:cor:onto}
In this finite-type class, $\lap\I(\cdot):\SR(\Om_D)\to\AM(\Om_D)$ is
surjective if and only if every complementary component of $D$ is fixed
by complex conjugation, i.e.\ iff $m=0$.  Simple connectivity of the
components of $D$ is sufficient but not necessary.
\end{corollary}

\begin{proof}
Immediate from \eqref{c:eq:8m}; for necessity, a nonzero reflected mode
placed in a missing pair is an explicit non-image.
\end{proof}

\subsection{Worked cases}

\begin{example}[The four-dimensional shell: a hole that does not obstruct]
\label{c:ex:shell}
Let $D=\{z:R_1<\abs z<R_2\}$ with $0<R_1<R_2$, so
$\Om_D=\{q\in\HH:R_1<\abs q<R_2\}$.  Its one bounded hole
$\{\abs z\le R_1\}$ is conjugation-fixed, so $s=1$, $m=0$, and every
axially monogenic field on the shell has a global slice-regular Fueter
primitive — although $b_1(D)=1$.  The upper-meridian picture agrees:
$U$ is parametrized by $(\log\abs z,\arg z)$ with $\arg z\in(0,\pi)$ and
is simply connected, so it has a winding basis of size $0$, as
Proposition~\ref{c:prop:bridge} predicts.  A blanket rule
``$8\dim_\RR H^1(D;\RR)$'' would give $8$ here, and would be wrong.
\end{example}

\begin{corollary}[Deleted spheres and deleted real points]
\label{c:cor:deleted}
Let $p_1,\dots,p_m\in\Cplus$ be distinct and let $c_1,\dots,c_\ell\in\RR$
be distinct.  On
\[
  \Om=\HH\setminus\Bigl(\bigcup_{j=1}^{m}S_{p_j}\cup
     \{c_1,\dots,c_\ell\}\Bigr)
\]
the Fueter image in $\AM(\Om)$ has real codimension $8m$, independent of
$\ell$.  In particular, for $m=0$ every axially monogenic field on
$\HH$ minus finitely many real points has a global slice-regular Fueter
primitive there.
\end{corollary}

\begin{proof}
The symmetric planar base is
$D=\CC\setminus(\{p_j,\bar p_j\}_j\cup\{c_k\}_k)$.  Its holes are the
$\ell$ real punctures, each conjugation-fixed, and the $m$ exchanged
pairs $\{p_j,\bar p_j\}$.  Apply Theorem~\ref{c:thm:reflection-range};
equivalently, $U=\Cplus\setminus\{p_1,\dots,p_m\}$ by
Proposition~\ref{c:prop:bridge}, with small circles as a winding basis of
size $m$.  A real puncture lies on the boundary of the upper parameter
domain, not inside it, and so creates no affine-monodromy cycle.
\end{proof}

\begin{remark}[What Corollary~\ref{c:cor:deleted} does \emph{not} say]
It does not assert that real-point singularities are removable, and it
does not assert that a field on $\HH$ minus a real point has no principal
part there.  It asserts only that such a field lies in the image of the
Fueter map on that same punctured domain.  The distinction is exhibited
by the identity, valid for real $a$,
\begin{equation}
  \lap\bigl((q-a)^{-1}\bigr)
   =-\frac{4\,\overline{q-a}}{\abs{q-a}^{4}},
  \label{c:eq:realkernel}
\end{equation}
whose right-hand side is, up to the constant $-8\pi^2$, the
Cauchy--Fueter kernel at $a$: it is a nonremovable singular axially
monogenic field on $\HH\setminus\{a\}$ that nevertheless has the global
slice-regular primitive $(q-a)^{-1}$.  Global invertibility on a
punctured domain and removability through the puncture are different
statements, and only the first is being claimed.
\end{remark}

\begin{example}[One sphere removed: an eight-dimensional cokernel]
\label{c:ex:sphere}
Let $D=\CC\setminus\{\iota,-\iota\}$, so $\Om_D=\HH\setminus\SSS$.
There is exactly one exchanged pair, hence $m=1$ and
$\dim_\RR\operatorname{coker}=8$ by
Theorem~\ref{c:thm:reflection-range}.  The two normalized generators are
\[
  g_b=G^{\mathrm{sym}}_{\iota;0,1},\qquad
  g_a=G^{\mathrm{sym}}_{\iota;1,0},
\]
with upper period pairs $(0,1)$ and $(1,0)$ respectively.  Both are
smooth axially monogenic on all of $\HH\setminus\SSS$, including the real
axis, yet neither is $\lap f$ for any single-valued $f\in\SR(\Om_D)$.
Specializing \eqref{c:eq:realtrace-general} at $u=0$, $v=1$, $d=1+x^2$,
\begin{equation}
  g_b(x)=\frac{4x}{\pi(1+x^{2})^{2}},
  \qquad
  g_a(x)=-\frac{4}{\pi(1+x^{2})^{2}} .
  \label{c:eq:realtraces}
\end{equation}
The obstruction is therefore not caused by a concealed singularity on
$\RR$.
\end{example}

\subsection{The obstructed field on \texorpdfstring{$\HH\setminus\SSS$}{H minus S} is an averaged Cauchy--Fueter kernel}

It is worth insisting that $g_b$ is a concrete, branch-free function and
not a multivalued object disguised by notation.  Its first potential is
globally single-valued.  Writing $q=x+Iy$ and
$d_\mp^2=x^2+(y\mp1)^2$,
\begin{equation}
  \pot_b(x,y)=-\frac1{2\pi y}\log\frac{d_-}{d_+}
   =-\frac{1}{4\pi y}\log\frac{x^{2}+(y-1)^{2}}{x^{2}+(y+1)^{2}},
  \qquad g_b=-2(\pot_b)_x+2I(\pot_b)_y .
  \label{c:eq:Cb}
\end{equation}
The apparent singularity at $y=0$ is removable, with
$\pot_b(x,0)=1/(\pi(1+x^2))$; expanding the logarithm to first order in
$y$ gives this at once, and differentiating recovers
\eqref{c:eq:realtraces}.  Thus even the \emph{first} potential can be
globally single-valued while the second is not — the concrete form of
Example~\ref{c:ex:independent} on a domain meeting $\RR$.

\begin{proposition}[Sphere-average representation]
\label{c:prop:sphereaverage}
Let $\dd\sigma$ be surface measure on $\SSS$, with $\sigma(\SSS)=4\pi$,
and let $E(q)=\bar q/(2\pi^{2}\abs q^{4})$ be the Cauchy--Fueter kernel.
Then, for $q\notin\SSS$,
\begin{equation}
  \pot_b(q)=\frac1{4\pi^{2}}\int_{\SSS}\frac{\dd\sigma(J)}{\abs{q-J}^{2}},
  \qquad
  g_b(q)=2\int_{\SSS}E(q-J)\,\dd\sigma(J)
       =-2\,\overline{\Dop}\,\pot_b(q).
  \label{c:eq:sphereaverage}
\end{equation}
\end{proposition}

\begin{proof}
For $q=x+Iy$ integrate over the angle $t=\langle I,J\rangle$:
\[
  \int_{\SSS}\frac{\dd\sigma(J)}{\abs{q-J}^{2}}
  =2\pi\int_{-1}^{1}\frac{\dd t}{x^{2}+y^{2}+1-2yt}
  =\frac{\pi}{y}\log\frac{x^{2}+(y+1)^{2}}{x^{2}+(y-1)^{2}},
\]
which after division by $4\pi^2$ is exactly \eqref{c:eq:Cb}.  For the
second formula, $-2\overline{\Dop}\bigl(4\pi^{2}\abs{q-J}^{2}\bigr)^{-1}
=2E(q-J)$, and differentiation under the integral sign is legitimate on
every compact set disjoint from $\SSS$.  The relation
$g_b=-2\overline{\Dop}\pot_b$ is \eqref{c:eq:Cb} rewritten in
quaternionic form.
\end{proof}

So the obstructed target is a Cauchy--Fueter kernel averaged over the
singular sphere; smoothness and monogenicity off $\SSS$ are transparent
from \eqref{c:eq:sphereaverage}, since the kernel is monogenic away from
its pole.  This is the special case, computed by hand, of the general layer representation, Theorem~\ref{e:thm:layer} of Part~VIII.

\begin{proposition}[An independent complex-analytic verification]
\label{c:prop:complexcheck}
There is no reflection-real holomorphic stem $F$ on
$D=\CC\setminus\{\iota,-\iota\}$ with $\lap\I(F)=g_b$.
\end{proposition}

\begin{proof}
The real-trace identity \eqref{e:eq:realtrace} and
\eqref{c:eq:realtraces} would force
$F''(x)=-2x\bigl(\pi(1+x^{2})^{2}\bigr)^{-1}$ for all real $x\in D$.
Applying the identity theorem to the four complex components of the
holomorphic function $F''$ on the connected set $D$ extends this to
\[
  F''(z)=-\frac{2z}{\pi(1+z^{2})^{2}}
   =\left(\frac1{\pi(1+z^{2})}\right)',
  \qquad\text{hence}\qquad
  F'(z)=\frac1{\pi(1+z^{2})}+A
\]
for some constant $A\in\HC$, the difference having vanishing derivative
on a connected set.  But $F'$ is the derivative of a single-valued
function, so $\oint_\gamma F'(z)\,\dd z=0$ for every closed
$\gamma\subset D$, whereas around a small positively oriented circle
about $\iota$ the displayed right-hand side integrates to
$2\pi\iota\cdot\frac1{2\pi\iota}=1$.  Contradiction.
\end{proof}

This argument uses no period theorem, no reflection lemma and no
classification: it is an independent proof that same-domain surjectivity
fails on $\HH\setminus\SSS$.

\begin{remark}[The editorial position taken by this document]
\label{c:rem:editorial}
Corollary~\ref{c:cor:onto} gives a necessary and sufficient topological
criterion, and Example~\ref{c:ex:sphere} together with
Proposition~\ref{c:prop:complexcheck} exhibits explicit axially monogenic
fields on $\HH\setminus\SSS$ that are not Fueter images on that same
domain.  It follows that an \emph{unqualified} assertion of same-domain
surjectivity is false as stated.  The sufficient condition ``every
component of the symmetric planar set is simply connected'' is correct as
hypothesized, and Corollary~\ref{c:cor:onto} shows it is not necessary.
We deliberately do not adjudicate whether any particular published
theorem is incorrect: that would require auditing its own domain,
locality and branch hypotheses, which we have not done.  The source
manuscripts of this merge are split on exactly this point — one of them
names a published sentence as requiring a topological qualification,
while four others explicitly decline to assert that any published
theorem is wrong — and the position adopted here is the second.  See
Part~XI.
\end{remark}

%% ==================== part6_7.tex ====================

%% =====================================================================
%% Part VI and Part VII of the merged document.
%% Writer key: d.  Every label defined below begins with "d:".
%% Lead sources: 05, 10, 01; consulted 02, 04, 06, 07, 08, 09.
%% =====================================================================

\section{Quantitative theory: stability, closed range, and period-corrected
approximation}
\label{d:sec:quantitative}

The criterion of Part~II (Theorem~\ref{b:thm:period}) is an exact
dichotomy: a single-valued
slice-regular Fueter primitive exists on the circularization $\Om_U$ if and
only if both period functionals annihilate $g$.  A dichotomy alone is
fragile.  It does not say whether a computed period of size $10^{-6}$ is a
genuine obstruction or an artifact of quadrature, and it does not say
whether a target with a nonzero period might nevertheless be approximated,
arbitrarily well and uniformly, by Fueter images of globally single-valued
primitives.  This part answers both questions.  The period functionals are
continuous in the crudest reasonable topology, the range is closed, and a
nonzero period is a quantitative \emph{error floor} on every detecting
loop.  The part closes with the approximation theorem that this floor
forces: Runge~\cite{Runge} approximation is possible for every axially monogenic target
only after a fixed correction space of real dimension exactly $8m$ is
adjoined.

Throughout, $U\subset\Cplus$ is a connected open set, $g=P+IQ\in\AM(\Om_U)$,
and $\om_g$, $\thet_g$ are the two closed one-forms of Part~II
(Definition~\ref{b:def:forms}, Lemma~\ref{b:lem:forms}),
\begin{equation}
 \om_g=\tfrac12\bigl(-P\,\dd x+Q\,\dd y\bigr),
 \qquad
 \thet_g=\tfrac12\bigl((xP+yQ)\,\dd x+(yP-xQ)\,\dd y\bigr).
 \label{d:eq:forms}
\end{equation}
Recall from Part~II that $\thet_g$ is computed from the target alone, and
is related to the primitive-dependent second form $\etaC$ of the classical
two-integration construction by the elimination identity
\eqref{b:elimination}, $\etaC-\dd(x\pot)=\thet_g$; nothing below uses
$\etaC$.

\subsection{The two norms on an axial field}
\label{d:sec:norms}

The sources measure an axial field in two inequivalent ways, and the
difference is the single most dangerous point in this theory.  We name both
and never conflate them.

\begin{definition}[Axial pair norm and spherical supremum]
\label{d:def:norms}
For $g=P+IQ$ axial on $\Om_U$ put, at a base point $z=x+\iota y\in U$,
\begin{equation}
 \axn{g}(z)=\bigl(\abs{P(x,y)}^2+\abs{Q(x,y)}^2\bigr)^{1/2},
 \qquad
 \abs{g}_{\mathrm{sph}}(z)=\sup_{I\in\SSS}\abs{P(x,y)+IQ(x,y)} .
 \label{d:eq:twonorms}
\end{equation}
For a compact $K\subset U$ we abbreviate $\sup_K\axn{g}=\sup_{z\in K}\axn{g}(z)$.  Near
an isolated singular sphere (Part~VII) the two associated radial maxima are
\begin{equation}
 \Max_g(\rho)=\max_{\abs{z-p}=\rho}\axn{g}(z),
 \qquad
 \Msph_g(\rho)=\max_{\abs{z-p}=\rho}\ \max_{I\in\SSS}\abs{P+IQ} .
 \label{d:eq:radialnorms}
\end{equation}
\end{definition}

\begin{lemma}[Norms over the sphere of imaginary units]
\label{d:lem:sphere-norm}
Let $P,Q\in\HH$ and let $\dd\sigma$ be ordinary surface measure on
$\SSS$, of total mass $4\pi$.  Then
\begin{align}
 \frac1{4\pi}\int_{\SSS}\abs{P+IQ}^2\,\dd\sigma(I)
   &=\abs P^2+\abs Q^2,
   \label{d:eq:sphere-mean}\\
 \max_{I\in\SSS}\abs{P+IQ}^2
   &=\abs P^2+\abs Q^2+2\,\bigl\lvert\ImH(P\bar Q)\bigr\rvert .
   \label{d:eq:sphere-max}
\end{align}
Consequently
\begin{equation}
 \axn{g}(z)\ \le\ \abs{g}_{\mathrm{sph}}(z)\ \le\ \sqrt2\,\axn{g}(z),
 \label{d:eq:norm-equivalence}
\end{equation}
and both inequalities are attained.  Moreover, for $X,Y\in\HH$ and real
$\alpha,\beta$,
\begin{equation}
 \ImH\bigl((\alpha X-\beta Y)\overline{(\beta X+\alpha Y)}\bigr)
 =(\alpha^2+\beta^2)\,\ImH(X\bar Y).
 \label{d:eq:rotation-cross}
\end{equation}
Finally, for all $x,y\in\RR$ and $P,Q\in\HH$,
\begin{equation}
 \abs{xP+yQ}^2+\abs{yP-xQ}^2=(x^2+y^2)\bigl(\abs P^2+\abs Q^2\bigr)
 =\abs z^2\axn{g}(z)^2 .
 \label{d:eq:rotation-real}
\end{equation}
\end{lemma}

\begin{proof}
Expanding the quaternionic square,
$\abs{P+IQ}^2=\abs P^2+\abs Q^2-2\ReH(P\bar QI)$, and the last term is
twice the Euclidean inner product of the imaginary vector $\ImH(P\bar Q)$
with $I\in\SSS$.  Its mean over $\SSS$ vanishes and its maximum is twice
the length of that vector, which gives \eqref{d:eq:sphere-mean} and
\eqref{d:eq:sphere-max}.  The left inequality in
\eqref{d:eq:norm-equivalence} follows from \eqref{d:eq:sphere-mean} (a
supremum dominates a mean), the right from \eqref{d:eq:sphere-max} and
$2\abs{\ImH(P\bar Q)}\le2\abs P\abs Q\le\abs P^2+\abs Q^2$.  Equality on
the left holds exactly when $\ImH(P\bar Q)=0$, on the right exactly when
$\abs P=\abs Q$ and $P\bar Q$ is purely imaginary.  For
\eqref{d:eq:rotation-cross}, expand the product: the terms $X\bar X$ and
$Y\bar Y$ are real and drop out of $\ImH$, while
$\ImH(Y\bar X)=-\ImH(X\bar Y)$; collecting gives the stated factor.
Identity \eqref{d:eq:rotation-real} is the same computation with the real
orthogonal matrix $\left(\begin{smallmatrix}x&y\\ y&-x\end{smallmatrix}
\right)$, whose two rows are orthogonal of common length $\abs z$, so the
cross terms $2xy\ReH(P\bar Q)$ cancel.
\end{proof}

\begin{remark}[The equivalence of norms is not an equivalence of constants]
\label{d:rem:norm-warning}
Inequality \eqref{d:eq:norm-equivalence} makes $\axn\cdot$ and
$\abs\cdot_{\mathrm{sph}}$ equivalent, so they define the same topology,
the same notion of closed range, and the same class of continuous
functionals.  They do \emph{not} give the same asymptotic constants.  The
two exact leading-size laws at a critical singular sphere
(Theorem~\ref{d:thm:leading} below) have genuinely different limits in the
two norms, differing by the term $2v\abs{\ImH(b\bar a)}$ under the square
root.  Every estimate in this part is stated in the axial pair norm, which
is the sharper of the two on the period side; every asymptotic law in
Part~VII is stated twice, once in each norm.
\end{remark}

\subsection{Period bounds and the distance to the range}

\begin{proposition}[Quantitative stability of the obstruction]
\label{d:prop:period-bounds}
Let $\gamma\subset U$ be a closed piecewise $C^1$ curve of length
$L_\gamma$, and put $R_\gamma=\max_{z\in\gamma}\abs z>0$.  Write
$(a_\gamma,b_\gamma)=\bigl(\int_\gamma\om_g,\int_\gamma\thet_g\bigr)$ for
the period pair of Definition~\ref{b:def:periods}.  Then
\begin{equation}
 \abs{a_\gamma}\le\frac{L_\gamma}{2}\,\sup_\gamma\axn g,
 \qquad
 \abs{b_\gamma}\le\frac{L_\gamma R_\gamma}{2}\,\sup_\gamma\axn g .
 \label{d:eq:period-bounds}
\end{equation}
Consequently, for \emph{every} globally single-valued $f\in\SR(\Om_U)$,
\begin{equation}
 \sup_\gamma\axn{g-\lap f}\ \ge\
 \max\left\{\frac{2\abs{a_\gamma}}{L_\gamma},
             \frac{2\abs{b_\gamma}}{L_\gamma R_\gamma}\right\} .
 \label{d:eq:distance-lower}
\end{equation}
\end{proposition}

\begin{proof}
Parametrize $\gamma$ by $(x(t),y(t))$ with unit speed.  The value of
$2\om_g$ on the tangent is $-P\dot x+Q\dot y$, a right
$\HH$-linear combination of $P,Q$ with real coefficient vector
$(-\dot x,\dot y)$ of Euclidean length $1$; Cauchy--Schwarz in
$\HH\oplus\HH\cong\RR^8$ gives
$\abs{\om_g(\dot\gamma)}\le\tfrac12\axn g$.  For $\thet_g$ the coefficient
vector of $(P,Q)$ in $2\thet_g(\dot\gamma)$ is
$(x\dot x+y\dot y,\ y\dot x-x\dot y)$, whose Euclidean length is $\abs z$
by \eqref{d:eq:rotation-real}; hence
$\abs{\thet_g(\dot\gamma)}\le\tfrac12\abs z\axn g$.  Integrating over
$\gamma$ proves \eqref{d:eq:period-bounds}.  By Theorem~\ref{b:thm:period}
every $\lap f$ with $f\in\SR(\Om_U)$ has vanishing periods, so
$g-\lap f$ has the periods of $g$; applying
\eqref{d:eq:period-bounds} to $g-\lap f$ and solving for
$\sup_\gamma\axn{g-\lap f}$ gives \eqref{d:eq:distance-lower}.
\end{proof}

Inequality \eqref{d:eq:distance-lower} is the quantitative content of the
whole theory.  A target with a nonzero period around $\gamma$ cannot be
uniformly approximated on the circularized loop by Fueter images of global
primitives: there is a positive error floor, computable from the two period
integrals and the elementary geometry of $\gamma$ alone, below which no
such approximation can go.  This is strictly stronger than the statement
that no exact global solution exists.

\begin{remark}[Two weaker forms in circulation]
\label{d:rem:weaker-distance}
Manuscript~10 states the lower bound in the combined form
$2(\abs{a_\gamma}^2+\abs{b_\gamma}^2)^{1/2}
 \bigl(L_\gamma\sqrt{1+R_\gamma^2}\bigr)^{-1}$.
That is implied by, and in general strictly weaker than, the max-form
\eqref{d:eq:distance-lower}, since
$\max\{s,t\}\ge(s^2+t^2)^{1/2}/\sqrt2$ and the denominators differ.  The
bounds \eqref{d:eq:period-bounds} are also frequently written with the
spherical supremum in place of $\axn\cdot$; by
\eqref{d:eq:norm-equivalence} that version is weaker by a factor at most
$\sqrt2$.  The pair-norm form is adopted here because it is the sharper
one and because it is what the singular asymptotics of Part~VII are
calibrated against.
\end{remark}

\begin{corollary}[Closed range, with no topological hypothesis]
\label{d:cor:closed-range}
On every connected open $U\subset\Cplus$ --- including one whose
complement has infinitely many components --- the range $\lap\SR(\Om_U)$ is
closed in $\AM(\Om_U)$ for locally uniform convergence, and a fortiori for
the $C^\infty$ compact-open topology.
\end{corollary}

\begin{proof}
By Theorem~\ref{b:thm:period},
$\lap\SR(\Om_U)=\bigcap_\gamma\ker\bigl(g\mapsto(a_\gamma,b_\gamma)\bigr)$,
the intersection taken over all closed piecewise $C^1$ curves in $U$.  Each
such functional is right $\HH$-linear and, by
\eqref{d:eq:period-bounds}, continuous for uniform convergence on the
compact set $\gamma$.  An intersection of kernels of continuous functionals
is closed.  No finiteness of the winding basis is used, and no boundary
regularity of $U$ is used.
\end{proof}

\begin{corollary}[Continuity of the splitting, in the finite case]
\label{d:cor:projection}
Suppose $U$ carries a finite winding basis of size $m$ in the sense of
Definition~\ref{c:def:winding}, with omitted points $p_1,\dots,p_m$ and
dual loops $\gamma_1,\dots,\gamma_m$, and let $S$ and
$\Pi=\mathrm{id}-S\,\Per$ be the section \eqref{c:eq:section} and the
projection of the split range Theorem~\ref{c:thm:split}.  Then $\Per$,
$S$ and $\Pi$ are continuous for the compact-open topology and for the
$C^\infty$ compact-open topology, $\Pi^2=\Pi$, $\Pi$ has range
$\lap\SR(\Om_U)$ and kernel $S(\HH^{2m})$, and for every compact
$K\subset U$
\begin{equation}
 \sup_K\axn{S\Per g}\ \le\
 \sum_{j=1}^m\Bigl(
   \abs{a_{\gamma_j}(g)}\,\sup_K\axn{G_{p_j;1,0}}
  +\abs{b_{\gamma_j}(g)}\,\sup_K\axn{G_{p_j;0,1}}\Bigr).
 \label{d:eq:correction-bound}
\end{equation}
\end{corollary}

\begin{proof}
Each of the $2m$ period coordinates is continuous by
\eqref{d:eq:period-bounds}.  $S$ sends a coefficient vector to a finite
right $\HH$-combination of the \emph{fixed} smooth modes
$G_{p_j;1,0},G_{p_j;0,1}$ of Proposition~\ref{c:prop:modes}, so it is continuous from $\HH^{2m}$
with any norm; \eqref{d:eq:correction-bound} is the triangle inequality in
$\HH\oplus\HH$ together with $\abs{uw}=\abs u\abs w$ for quaternions.  The
algebraic identities follow from $\Per S=\mathrm{id}$.
\end{proof}

\subsection{The normalized inverse as a continuous operator}

Closedness of the range says nothing about how large a primitive must be.
The next theorem packages the inverse as a named operator with a stated
domain, a stated normalization, and an explicit estimate --- together with
an explicit statement of what that estimate does \emph{not} give.

\begin{theorem}[The normalized inverse and its pathwise estimate]
\label{d:thm:normalized-inverse}
Fix $z_*=x_*+\iota y_*\in U$.  On the zero-period subspace
$\AM_0(\Om_U)=\ker\Per\subset\AM(\Om_U)$ there is a unique right
$\HH$-linear inverse
\[
 \mathcal N:\AM_0(\Om_U)\longrightarrow\SR(\Om_U),
 \qquad \lap\,\mathcal N g=g,
\]
normalized by $\pot(z_*)=\potV(z_*)=0$ for the two potentials of
Definition~\ref{b:def:potentials};
equivalently, $\mathcal N g$ is the unique primitive vanishing on the
entire base sphere $x_*+y_*\SSS$, equivalently the unique primitive whose
stem satisfies $F(z_*)=0$ in $\HC$ --- eight real conditions, not the four
real conditions of vanishing at a single quaternionic point
(Corollary~\ref{b:cor:normalization} and Remark~\ref{b:rem:eightnotfour}).

Let $z\in U$ be joined to $z_*$ by a path $\sigma\subset U$ of length at
most $L$, all of whose points satisfy $\abs w\le R_*$, and suppose
$\axn g\le M$ on $\sigma$.  Then
\begin{equation}
 \abs{\pot(z)}\le\frac{LM}{2},
 \qquad
 \abs{\potV(z)}\le\frac{LR_*M}{2},
 \qquad
 \abs{(\mathcal N g)(q)}\le\frac{(\abs q+R_*)LM}{2}
 \quad (q=x+Iy),
 \label{d:eq:inverse-bound}
\end{equation}
and the stem derivative obeys
$\norm{F'(z)}_{\HC}\le\tfrac12(L+y)M$.
For every compact $K\subset U$ there exist a compact $K_*\subset U$
containing $K\cup\{z_*\}$ and a finite $L_K$ for which the path hypothesis
holds simultaneously at every $z\in K$; with
$R_*=\max_{K_*}\abs w$ one obtains the operator estimate
\begin{equation}
 \sup_{z\in K,\ I\in\SSS}\abs{(\mathcal N g)(x+Iy)}
 \ \le\ L_K R_*\,\sup_{K_*}\axn g .
 \label{d:eq:operator-estimate}
\end{equation}
Finally $\mathcal N$ is continuous for the $C^\infty$ compact-open
topologies.
\end{theorem}

\begin{proof}
Existence and the formula $\mathcal Ng(q)=q\pot(x,y)+\potV(x,y)$ are the
Theorem~\ref{b:thm:period}, with $\pot$ and $\potV$ obtained by
path integration of the exact forms $\om_g$ and $\thet_g$ from $z_*$.  For
uniqueness of the normalization: an element $q\mapsto qa+b$ of $\Aff$
vanishing on the whole sphere $x_*+y_*\SSS$ has, on subtracting its values
at $I$ and $-I$, $2Iy_*a=0$, and $y_*>0$ forces $a=0$ and then $b=0$;
equivalently the real-linear map $(a,b)\mapsto z_*a+b$ from $\HH^2$ to
$\HC$ is an isomorphism precisely because $y_*>0$, which is the
reconciliation of the "vanishes on a sphere" and "$F(z_*)=0$" descriptions
of the same eight conditions.

For the estimates, integrate the pointwise bounds established in the proof
of Proposition~\ref{d:prop:period-bounds} along $\sigma$:
$\abs{\om_g(\dot\sigma)}\le\tfrac12\axn g$ and
$\abs{\thet_g(\dot\sigma)}\le\tfrac12\abs w\axn g\le\tfrac12R_*\axn g$.
This gives the first two inequalities in \eqref{d:eq:inverse-bound}, and
$\mathcal Ng=q\pot+\potV$ with $\abs q=\abs z$ gives the third.  The stem
bound follows from the identity \eqref{b:M} of Part~III,
$F'=\pot+\tfrac y2(Q-\iota P)$, together with
$\norm{Q-\iota P}_{\HC}=\axn g$.

For the uniform path family, cover $K$ by finitely many open disks with
compact closures in $U$, join $z_*$ to each centre by a fixed polygonal
path in $U$, and append a final segment inside the relevant disk.  The
union of these paths with the closed disks is a compact $K_*\subset U$, and
their lengths have a common bound $L_K$.  Specializing
\eqref{d:eq:inverse-bound} with $\abs q\le R_*$ gives
\eqref{d:eq:operator-estimate}.  Smooth continuity follows because
$\dd\pot=\om_g$ and $\dd\potV=\thet_g$ express every first derivative of
the potentials as an explicit coefficient of $P,Q$, and every higher
derivative as a derivative of those coefficients; on compact subsets of $U$
the factor $1/y$ stays bounded, so finitely many $C^{k-1}$ seminorms of $g$
control the $C^k$ seminorms of $\mathcal Ng$, the zeroth order being
\eqref{d:eq:operator-estimate}.  Applying all of this to differences of
targets gives continuity.
\end{proof}

\begin{remark}[What the estimate does not assert]
\label{d:rem:no-uniform}
The constants $L_K$ and $R_*$ in \eqref{d:eq:operator-estimate} depend on
the domain, on the chosen base point, and on the path family; nothing above
gives a bound uniform over a family of domains, and in particular nothing
is asserted about domains that degenerate, about a singular sphere
collapsing onto the real axis, or about bases with a shrinking
inradius.  Likewise, Corollary~\ref{d:cor:closed-range} is a statement about
the compact-open topology only: the splitting $S$ is \emph{not} asserted to
be bounded in any weighted $L^2$, Bergman or Hilbert-space norm on
$\AM(\Om_U)$.  Sharper forms of both statements are recorded as open
problems in Part~XI, \S\ref{f:sec:scope}.  The refinement $(\abs q+R_*)/2$ in
\eqref{d:eq:inverse-bound} reduces to the cruder $LR_*M$ of several sources
exactly when $\abs q=R_*$; it is genuinely smaller for $q$ nearer the
origin.
\end{remark}

\subsection{Period-corrected Runge approximation}

The error floor \eqref{d:eq:distance-lower} says that a target with a
nonzero period is inaccessible from the range.  It does not say that the
range plus a finite correction is inaccessible, and in fact adjoining the
$2m$ modes of Part~IV restores a full Runge theorem.  Two distinct
rational classes occur in the sources; they are \emph{not} interchangeable,
and the distinction is entirely a matter of the pole set.

\begin{definition}[Two rational stem classes]
\label{d:def:rational}
(i) \emph{One-sided (product-domain) class.}  Let $U\subset\Cplus$ satisfy
the hypothesis of Definition~\ref{c:def:winding} with omitted points
$p_1,\dots,p_m$, one in each bounded component of the complement of $U$ in
the Riemann sphere.  Let $\mathrm{Rat}^{\mathrm{one}}$ be the $\HC$-valued
rational functions with poles only in $\{p_1,\dots,p_m,\infty\}$.  These
are \emph{upper-stem} rational functions.

(ii) \emph{Symmetric (slice-domain) class.}  Let $D\subset\CC$ be
conjugation-invariant with the analogous finite topology, let
$E=\{p_1,\dots,p_m,\bar p_1,\dots,\bar p_m\}$, and let $\mathrm{Rat}_E$ be
the $\HC$-valued rational $R$ with poles only in $E\cup\{\infty\}$
satisfying the stem symmetry $R(\bar z)=\kappa(R(z))$, where
$\kappa(A+\iota B)=A-\iota B$.
\end{definition}

\begin{theorem}[Period-corrected rational approximation]
\label{d:thm:runge}
Let $\Per(g)=\bigl((a_j,b_j)\bigr)_{j=1}^m$ and let $S$ be the section of
Part~IV.
\begin{enumerate}[label=\textup{(\roman*)}]
\item \textup{(Product domain.)}  For every $g\in\AM(\Om_U)$ there are
$R_n\in\mathrm{Rat}^{\mathrm{one}}$ with
\begin{equation}
 \lap\I(R_n)+S\Per(g)\longrightarrow g
 \label{d:eq:runge}
\end{equation}
locally uniformly on $\Om_U$, together with every fixed order of real
derivatives.
\item \textup{(Symmetric slice domain.)}  The same conclusion holds on
$\Om_D$ with $R_n\in\mathrm{Rat}_E$ and with the reflected modes
$G^{\mathrm{sym}}$ of Proposition~\ref{c:prop:reflected} used in $S$.  Here a common denominator may be
chosen to be a \emph{real} polynomial with zeros in $E$, and the numerator
then has right quaternionic coefficients; the $R_n$ are genuine rational
slice functions.
\end{enumerate}
In both cases every approximant has \emph{exactly} the periods of $g$, and
\[
 \overline{\lap\I(\mathcal R)}^{\,\mathrm{loc}}=\lap\SR
 \qquad\text{inside }\AM .
\]
In particular, uncorrected approximation --- a locally uniform limit of
Fueter images of global slice-regular functions --- is possible if and only
if $\Per(g)=0$.
\end{theorem}

\begin{proof}
By Theorem~\ref{c:thm:split}, $g-S\Per(g)$ has vanishing periods,
hence equals $\lap\I(F)$ for a single-valued holomorphic stem $F$ (on $U$
in case~(i), on $D$ in case~(ii)).  Classical scalar Runge approximation,
applied to the four complex components of $F$ with the prescribed pole set
--- which meets every complementary component of the relevant planar set,
including the one containing $\infty$ --- produces rational $R_n^0$
converging to $F$ uniformly on compact subsets.

In case~(ii) replace $R_n^0$ by the symmetrization
\begin{equation}
 R_n(z)=\tfrac12\bigl(R_n^0(z)+\kappa\bigl(R_n^0(\bar z)\bigr)\bigr),
 \label{d:eq:symmetrization}
\end{equation}
which is still rational, moves no pole into $D$, satisfies
$R_n(\bar z)=\kappa(R_n(z))$, and still converges locally uniformly on the
symmetric compact sets.  Taking the common denominator
$d(z)=\prod_j(z-p_j)^{m_j}(z-\bar p_j)^{m_j}$, which has real coefficients
and is therefore fixed by $\kappa$-symmetrization, the numerator
$N=R_nd$ inherits $N(\bar z)=\kappa(N(z))$, and a polynomial with that
property has all coefficients in $\HH$.

Cauchy estimates on slightly larger compact sets upgrade uniform
convergence of the stems to uniform convergence of all complex
derivatives.  Since compact subsets of $\Om_U$ (respectively of $\Om_D$
away from the axis, and, across the axis, via the convergent symmetric
Taylor expansions of the stems at real points) keep $y$ in a bounded range,
the axial formula \eqref{a:eq:Fu-first} of Part~I converts this into
convergence of $\lap\I(R_n)$ with all fixed real derivatives.  Each
$\lap\I(R_n)$ is a global Fueter image and so has zero periods; adding the
fixed field $S\Per(g)$ therefore reproduces exactly the periods of $g$.
The closure statement follows by applying the construction to a zero-period
target for one inclusion and Corollary~\ref{d:cor:closed-range} for the
other; the final equivalence combines that closure with
\eqref{d:eq:distance-lower}.
\end{proof}

\begin{remark}[Why the two pole-set hypotheses must both be displayed]
\label{d:rem:pole-sets}
Read carelessly, the two halves of Theorem~\ref{d:thm:runge} contradict one
another: one says the approximating slice functions need \emph{not} be
rational with coefficients in $\HH$, the other that they are.  Both are
correct, and the difference is exactly the pole set.  In case~(i) the poles
are \emph{one-sided} --- only the chosen upper points --- so after
reflection to the lower half-plane the induced slice function is defined
piecewise and is not in general a single rational function with
quaternionic coefficients on a connected symmetric planar domain.  In
case~(ii) the pole set $E$ is conjugation-symmetric, and it is precisely
that symmetry that makes the argument of \eqref{d:eq:symmetrization} yield
a real denominator and an $\HH$-coefficient numerator.  A merge that
recorded only one of the two statements, or recorded both without their
pole-set hypotheses, would be either incomplete or apparently
self-contradictory.
\end{remark}

\begin{proposition}[Minimality of the correction space]
\label{d:prop:minimality}
Under the finite-winding-basis hypothesis, let $W\subset\AM(\Om_U)$ be any
fixed finite-dimensional \emph{real} subspace such that
$\lap\SR(\Om_U)+W$ is dense in $\AM(\Om_U)$ for locally uniform
convergence.  Then $\dim_\RR W\ge8m$.  The $2m$ quaternionic modes of the
section $S$ therefore form a correction space of the smallest possible
dimension.
\end{proposition}

\begin{proof}
The period map $\Per:\AM(\Om_U)\to\HH^{2m}$ is continuous
(Proposition~\ref{d:prop:period-bounds}), surjective (Theorem~\ref{c:thm:split}), and
annihilates $\lap\SR(\Om_U)$.  Hence
$\Per\bigl(\lap\SR(\Om_U)+W\bigr)=\Per(W)$, a finite-dimensional and
therefore closed real subspace of $\HH^{2m}\cong\RR^{8m}$.  Density of
$\lap\SR+W$ and continuity of $\Per$ force
$\Per(W)\supseteq\overline{\Per(W)}\supseteq\Per(\AM)=\HH^{2m}$, so
$\Per(W)=\HH^{2m}$ and $\dim_\RR W\ge\dim_\RR\HH^{2m}=8m$.
\end{proof}

\begin{remark}[Scope of the approximation statement]
\label{d:rem:runge-scope}
The only approximation theorem imported above is the classical scalar Runge
theorem; nothing here is a new approximation result for monogenic
functions.  What is specifically quaternionic is the exact finite
augmentation that the Fueter map requires, its exact preservation of every
period at every stage of the approximation, and the sharp count $8m$ of
Proposition~\ref{d:prop:minimality}.  No density statement is made for
bases with infinitely many holes, where $\Per$ takes values in an infinite
product and Proposition~\ref{d:prop:minimality} has no finite-dimensional
content, although Corollary~\ref{d:cor:closed-range} continues to apply.
\end{remark}

\medskip
\noindent\textbf{The reconstruction procedure implied by this part.}
Integrate the two explicit coefficient pairs of \eqref{d:eq:forms}
numerically around one loop per hole; subtract
$\sum_jG_{p_j;a_j,b_j}$ using the branch-free formulas
\eqref{c:eq:mode-P}--\eqref{c:eq:mode-Q};
then integrate the now-exact forms along any convenient paths and assemble
$f=q\pot+\potV$.  The two bounds \eqref{d:eq:period-bounds} serve as the
numerical diagnostic in both directions: they certify that a computed
period is too large to be quadrature error, and
\eqref{d:eq:distance-lower} certifies that the residual cannot be reduced
below an explicit floor by any better choice of primitive.

% =======================================================================

\section{Singularities along a nonreal sphere: normal forms and sharp
thresholds}
\label{d:sec:singularities}

\subsection{Geometry, and the affine spherical residue}

Fix $p=u+\iota v\in\Cplus$ with $v>0$, let $\Sp=u+v\SSS$ be the
corresponding conjugacy two-sphere in $\HH$, and choose $0<R<v$.  Write
\begin{equation}
 \Tube_R=\Om_{\{\abs{z-p}<R\}}
 =\bigl\{x+Iy:\ (x-u)^2+(y-v)^2<R^2,\ I\in\SSS\bigr\},
 \qquad
 \xi=x-u,\quad \eta=y-v,
 \label{d:eq:tube}
\end{equation}
and $\rho=\abs{z-p}=(\xi^2+\eta^2)^{1/2}$.  The letters $s$ and $t$ are not
used in this part: three of the source manuscripts gave the pair $(s,t)$
three mutually exclusive meanings.

\begin{lemma}[$\rho$ is the Euclidean distance to the sphere]
\label{d:lem:distance}
For $q=x+Iy\in\Tube_R$ one has
$\operatorname{dist}(q,\Sp)=\sqrt{(x-u)^2+(y-v)^2}=\rho$, the nearest point
of $\Sp$ being $p(q)=u+Iv$, with the \emph{same} imaginary unit as $q$.
\end{lemma}

\begin{proof}
For $J\in\SSS$, $\abs{q-(u+vJ)}^2=(x-u)^2+y^2+v^2-2vy\langle I,J\rangle$,
where $\langle\cdot,\cdot\rangle$ is the Euclidean inner product on
imaginary quaternions.  Since $y>0$ and $v>0$ this is minimized exactly at
$J=I$, where it equals $(x-u)^2+(y-v)^2$.
\end{proof}

Thus a punctured planar disk in the base describes a genuine tubular
puncture along a two-sphere in four real dimensions.  The singular set has
real \emph{codimension two}, not four.  Every threshold below is calibrated
to that geometry and must not be compared with the isolated-point
Cauchy--Fueter thresholds of Part~I: the sharp isolated-point removability
exponent there is $o(\abs q^{-3})$ (Theorem~\ref{a:thm:fueter-removable}),
and the Cauchy--Fueter kernel
$\bar q/(2\pi^2\abs q^4)$ has size $\abs q^{-3}$, whereas the sharp
exponent here is $1$.

\begin{definition}[Affine spherical residue]
\label{d:def:resaff}
For $g\in\AM(\Tube_R\setminus\Sp)$ put
\begin{equation}
 \Resaff_{\Sp}(g)=(a,b)
 =\left(\int_{\abs{z-p}=\rho}\om_g,\ \int_{\abs{z-p}=\rho}\thet_g\right),
 \qquad 0<\rho<R,
 \label{d:eq:resaff}
\end{equation}
the circle positively oriented.  Closedness of both forms and Stokes'
theorem on an annulus make the pair independent of $\rho$.
\end{definition}

The pair $(a,b)\in\HH^2$ is by construction the affine monodromy coefficient
pair of Part~II, evaluated on the meridian circle.  It is \emph{not} the
slice Laurent residue $\Res_c$ at a real isolated singularity, and it is
\emph{not} the three-dimensional Cauchy--Fueter flux residue $\Res^F$ of
an isolated point (Theorem~\ref{a:thm:residue-slice} and
Definition~\ref{a:def:flux-residue}); the three notions
are kept apart throughout.  Under a real translation $q\mapsto q-c$ one has
$a\mapsto a$ and $b\mapsto b+ca$, so the pair $(a,ua+b)$ is
translation-invariant, and every invariant below is a function of it.

\subsection{(A) The normal form, with no hypothesis whatever}

\begin{theorem}[Full spherical normal form]
\label{d:thm:normal}
Every $g\in\AM(\Tube_R\setminus\Sp)$ has a unique representation
\begin{equation}
 g=g_{\mathrm{reg}}+G_{p;a,b}
   +\lap\I\left(\sum_{k\ge1}(z-p)^{-k}c_k\right),
 \qquad (a,b)\in\HH^2,\ \ c_k\in\HC,
 \label{d:eq:normal-form}
\end{equation}
where $g_{\mathrm{reg}}$ extends axially monogenically through the whole of
$\Sp$, $(a,b)=\Resaff_{\Sp}(g)$, and the stem series
$\sum_{k\ge1}(z-p)^{-k}c_k$ converges as a holomorphic $\HC$-valued
function, with all complex derivatives, uniformly on compact subannuli; the
transformed series converges with all real derivatives on compact
subannuli of the punctured tube.  In particular $g$ extends regularly
through $\Sp$ if and only if $a=b=0$ and every $c_k=0$.
\end{theorem}

\begin{proof}
Let $(a,b)$ be the residue pair and put $g_0=g-G_{p;a,b}$.  By the period
normalization \eqref{c:eq:mode-periods} of the logarithmic modes, $g_0$ has
vanishing periods on the punctured disk, whose first homology is generated by one
circle; by Theorem~\ref{b:thm:period} there is a single-valued
holomorphic stem $F_0$ on $0<\abs{z-p}<R$ with $\lap\I(F_0)=g_0$.  Apply
the scalar Laurent theorem to the four complex components of $F_0$:
$F_0=F_++\sum_{k\ge1}(z-p)^{-k}c_k$ with $F_+$ holomorphic on the full
disk.  Set $g_{\mathrm{reg}}=\lap\I(F_+)$.  Laurent series and all their
derivatives converge normally on compact subannuli, and the axial formula
for $\lap\I(\cdot)$ involves only $F'$, $\Ima_{\CC}F$ and the factors
$1/y,1/y^2$, which are bounded there; so the transform may be taken
termwise.

For uniqueness, the periods already determine $(a,b)$, since a field
regular across $\Sp$ has closed forms on the full disk and hence zero
circle periods.  Suppose a negative stem series $H$ has Fueter image
regular across $\Sp$.  That extended field has a holomorphic primitive $K$
on the full disk by Part~II, and on the punctured disk
$\lap\I(H-K)=0$, so $H-K=za'+b'$ with $a',b'\in\HH$ by
Corollary~\ref{b:cor:kernel}.  The right-hand side and $K$ are holomorphic through
$p$, so every negative Laurent coefficient of $H$ vanishes.  This gives
uniqueness of all $c_k$ and then of $g_{\mathrm{reg}}$.
\end{proof}

\begin{remark}[The logarithmic sector is not an optional rewriting]
\label{d:rem:log-sector}
Every term $\lap\I((z-p)^{-k}c_k)$ has a single-valued holomorphic stem on
the punctured disk and therefore vanishing affine residues.  No sum, finite
or convergent, of such terms can represent a $G_{p;a,b}$ with
$(a,b)\neq(0,0)$ on an annulus carrying the period cycle.  Equivalently:
the logarithmic sector of \eqref{d:eq:normal-form} cannot be produced by
any single-valued holomorphic source on the same punctured tube, and it
must be removed \emph{before} an ordinary Laurent expansion of the
primitive is available.  This is a classification of \emph{axially
monogenic} singularities along a nonreal sphere, expressed through their
slice-regular primitives; it is not a claim to introduce quaternionic
Laurent expansions in general.
\end{remark}

\subsection{(B) The $L^1$ theorem: the weakest hypothesis making the two
residues complete}

The classification under the pointwise critical bound
$\Msph_g(\rho)=O(\rho^{-1})$ is a corollary of the theorem proved in this
subsection under the strictly weaker hypothesis of local integrability
across the sphere.  We record the logical relation explicitly, because
seven of the ten source manuscripts state only the pointwise version.

\begin{lemma}[Reduction to the planar base]
\label{d:lem:reduction}
For an axial $g=P+IQ$ on $\Tube_R\setminus\Sp$ and $1\le s<\infty$, the
four-dimensional condition $g\in L^s(\Tube_{R_0})$ is equivalent to the
planar condition $P,Q\in L^s(\{\abs{z-p}<R_0\})$, for every $R_0<R$.
\end{lemma}

\begin{proof}
The volume element is $\dd V=y^2\,\dd x\,\dd y\,\dd\sigma(I)$ with
$\int_{\SSS}\dd\sigma=4\pi$, and on a small tube $y$ lies between $v-R_0$
and $v+R_0$, both positive.  The averages
$P=(4\pi)^{-1}\int_{\SSS}g\,\dd\sigma$ and
$Q=-(4\pi)^{-1}\int_{\SSS}Ig\,\dd\sigma$, together with Jensen's
inequality, bound the planar norms by the four-dimensional one; the
pointwise bound $\abs{P+IQ}\le\abs P+\abs Q$ bounds it in the other
direction.
\end{proof}

\begin{theorem}[$L^1$ spherical normal form]
\label{d:thm:L1}
Let $g\in\AM(\Tube_R\setminus\Sp)$ be locally integrable across $\Sp$, that
is $\int_{\Tube_{R_0}}\abs{g}\,\dd V<\infty$ for some $R_0<R$.  Then every
$c_k$ in \eqref{d:eq:normal-form} vanishes, and
\begin{equation}
 \boxed{\ g=G_{p;a,b}+g_0\ }
 \label{d:eq:L1-normal}
\end{equation}
with $(a,b)=\Resaff_{\Sp}(g)$ and $g_0$ extending as a smooth axially
monogenic field through the whole sphere.  The coefficients are unique.  In
particular
\begin{equation}
 g\ \text{is removable across }\Sp
 \iff a=b=0,
 \label{d:eq:L1-removal}
\end{equation}
and the space of locally $L^1$ axially monogenic germs at $\Sp$, modulo
removable germs, is a right quaternionic vector space of dimension exactly
two, with coordinates $(a,b)$ and real dimension eight.
\end{theorem}

Two genuinely different proofs are recorded.  Neither is a variant of the
other: the first removes the singularity of the \emph{primitive} at the
endpoint case of Sobolev regularity, the second computes the
\emph{distributional defect} of the field itself and shows it is a point
mass.  The second proof produces the surface source of Part~VIII directly,
rather than deducing it afterwards.

The first proof uses the two endpoint lemmas of
Appendix~\ref{f:app:distributions}: Lemma~\ref{f:lem:W11}, that a puncture
does not obstruct an integrable gradient in two dimensions, so a $C^1$
function off the origin of a planar disk with $L^1$ gradient extends to
$W^{1,1}_{\mathrm{loc}}$ with the same weak gradient; and
Lemma~\ref{f:lem:weakCR}, weak Cauchy--Riemann removability for
$W^{1,1}_{\mathrm{loc}}$ data.  Neither is restated here.

\begin{proof}[First proof of Theorem~\textup{\ref{d:thm:L1}} (the $W^{1,1}$
route)]
By Lemma~\ref{d:lem:reduction} the planar coefficients $P,Q$ are $L^1$ near
$p$.  The branch-free formulas
\eqref{c:eq:mode-P}--\eqref{c:eq:mode-Q} give
$\abs{P_{G}}+\abs{Q_{G}}=O(\rho^{-1}+1+\abs{\log\rho})$ for the components
of $G_{p;a,b}$, so those modes are themselves $L^1$ on a small tube.
Subtract the residue-normalized mode and call the remainder $g_1$; it is
$L^1$ and has both periods zero, so by Theorem~\ref{b:thm:period} it has a
single-valued
holomorphic stem $F=A+\iota B$ on the punctured disk with
$\lap\I(F)=g_1$, and the potentials $\pot=B/y$, $\potV=A-xB/y$ satisfy
$\dd\pot=\om_{g_1}$, $\dd\potV=\thet_{g_1}$.  The coefficients of those two
forms are, by \eqref{d:eq:forms}, real-linear combinations of $P,Q$ with
coefficients $1$, $x$ and $y$, all bounded on a smaller disk; hence
$\nabla\pot$ and $\nabla\potV$ are in $L^1$.  Apply
Lemma~\ref{f:lem:W11} to all eight real components.  Therefore $A=\potV+x\pot$ and $B=y\pot$
extend to $W^{1,1}_{\mathrm{loc}}$ functions across $p$, with weak
derivatives equal to the original ones, so their weak Cauchy--Riemann
equations hold on the full disk.  Lemma~\ref{f:lem:weakCR} extends $F$
holomorphically through $p$; its induced slice function and its Fueter
image extend smoothly through the entire sphere.  A smooth extension has
zero periods by Stokes' theorem on a meridian disk, which forces the
coefficients of any such decomposition to be exactly \eqref{d:eq:resaff}
and gives uniqueness and \eqref{d:eq:L1-removal}.  Every pair $(a,b)$ is
realized by the corresponding mode, which is $L^1$, so the germ quotient
has quaternionic dimension exactly two.
\end{proof}

The second proof runs on the planar Vekua operator
\begin{equation}
 \mathcal L(P,Q)=\Bigl(P_x-Q_y-\frac{2Q}{y},\ \ Q_x+P_y\Bigr),
 \label{d:eq:Lop}
\end{equation}
so that $g=P+IQ$ is axially monogenic exactly when $\mathcal L(P,Q)=0$; its
coefficients, including $1/y$, are smooth and bounded on
$\{\abs{z-p}<R\}$ because $y\ge v-R>0$ there.  Three facts about
$\mathcal L$ are proved in Appendix~\ref{f:app:distributions} and are used
here without restatement.  First, by Lemma~\ref{f:lem:point-support} a
distribution supported at one point is a finite sum of derivatives of the
point mass.  Second, Lemma~\ref{f:lem:defect-scaling}: if $P,Q\in L^1$ near
$p$ and $\mathcal L(P,Q)=0$ away from $p$, then
\begin{equation}
 \mathcal L(P,Q)=(\sigma_0,\sigma_1)\,\delta_p
 \label{d:eq:defect}
\end{equation}
for unique $\sigma_0,\sigma_1\in\HH$, and \emph{no} derivative of
$\delta_p$ occurs --- the scaling argument there is precisely where the
$L^1$ hypothesis and the first order of $\mathcal L$ are both consumed, and
it is what makes the residue classification exhaustive rather than a list
of examples.  Third, the mode defect \eqref{f:eq:mode-defect},
\begin{equation}
 \mathcal L(P_G,Q_G)=\Bigl(\frac{2(ua+b)}{v},\ 2a\Bigr)\delta_p
 \qquad\text{for }G=G_{p;a,b},
 \label{d:eq:mode-defect}
\end{equation}
computed there by a single circle integral from the leading term
\eqref{d:eq:leading} below together with the normalization
$\bar\partial\bigl((z-p)^{-1}\bigr)=\pi\delta_p$.  Finally the $L^1$ Weyl
Lemma~\ref{f:lem:weyl}, applied to each of the four real components
together with the factorization $\overline{\Dop}\Dop=\lap$, shows there
that weak axial monogenicity is removable: an $L^1$ pair with
$\mathcal L(P,Q)=0$ distributionally on the \emph{full} disk has a smooth
axially monogenic representative on the full tube $\Tube_R$.

\begin{proof}[Second proof of Theorem~\textup{\ref{d:thm:L1}} (the
distributional route)]
By Lemma~\ref{d:lem:reduction} the planar coefficients are $L^1$, and
Lemma~\ref{f:lem:defect-scaling} gives a defect
$(\sigma_0,\sigma_1)\delta_p$ as in \eqref{d:eq:defect}.  Put
\begin{equation}
 a=\frac{\sigma_1}{2},
 \qquad
 b=\frac{v\sigma_0-u\sigma_1}{2},
 \label{d:eq:source-to-residue}
\end{equation}
so that $2(ua+b)/v=\sigma_0$.  By \eqref{d:eq:mode-defect} the field
$g-G_{p;a,b}$ has zero planar defect; it is still $L^1$, so the $L^1$ Weyl
Lemma~\ref{f:lem:weyl}, in the axial form recalled above, provides the
smooth monogenic extension $g_0$.  Read backwards,
\eqref{d:eq:source-to-residue} identifies the defect coefficients with the
affine surface source of Part~VIII, \S\ref{e:sec:surface-source}.
Since a field regular on the full disk has closed forms there and hence
zero circle periods, applying the period normalization of the modes shows
that the coefficients in \eqref{d:eq:L1-normal} are exactly
$\Resaff_{\Sp}(g)$, whence \eqref{d:eq:source-to-residue} also identifies
the defect coefficients in terms of the residues.  Uniqueness and the
dimension count follow as before.
\end{proof}

\begin{remark}[$L^1$ versus $O(\rho^{-1})$: which hypothesis is used where]
\label{d:rem:L1-vs-O}
Local integrability across $\Sp$ is a formally \emph{weaker} hypothesis
than the pointwise critical bound $\Msph_g(\rho)=O(\rho^{-1})$: since the
transverse area element is $\rho\,\dd\rho\,\dd\theta$ and $y$ is bounded,
$O(\rho^{-1})$ implies local integrability, while the converse is not
available a priori.  Theorem~\ref{d:thm:L1} therefore subsumes the
critical-growth classification, which is recovered as
Corollary~\ref{d:cor:critical} below.  A posteriori the two classes
coincide --- Theorem~\ref{d:thm:L1} shows that every $L^1$ germ is a mode
plus a regular field, and the modes are $O(\rho^{-1})$ --- but that is a
\emph{consequence} of the theorem, not a hypothesis of it, and the
implication must not be used to justify the weaker statement.  The boundary
is exactly where Proposition~\ref{d:prop:L1sharp} places it: the field
$\lap\I((z-p)^{-1}c)$ has zero periods, exact order $\rho^{-2}$, fails to
be locally $L^1$ (its planar integral diverges like
$\int_0\rho^{-1}\dd\rho$), and is not removable.
\end{remark}

\begin{remark}[What the contour formulas do and do not require]
\label{d:rem:contour-scope}
Definition~\ref{d:def:resaff} makes sense for every
$g\in\AM(\Tube_R\setminus\Sp)$, with no integrability or growth hypothesis
at all; the pair $(a,b)$ is always defined and always radius-independent.
What Theorem~\ref{d:thm:L1} adds is \emph{completeness}: in the $L^1$ class
no further singular invariant survives modulo regular germs.  Outside that
class the residues remain defined but stop being a complete invariant, in
both directions --- upward through the filtration of
Theorem~\ref{d:thm:growth}, and downward through
Proposition~\ref{d:prop:L1sharp}.
\end{remark}

\subsection{(C) The growth filtration}

\begin{lemma}[Growth of a single-valued primitive]
\label{d:lem:primitive-growth}
Let $g_1\in\AM(\Tube_R\setminus\Sp)$ have vanishing periods and satisfy
$\Max_{g_1}(\rho)=O(\rho^{-N})$ for an integer $N\ge1$.  Then every
single-valued holomorphic primitive stem $F_1$ on $0<\abs{z-p}<R$ satisfies,
uniformly in the polar angle,
\begin{equation}
 \norm{F_1(z)}_{\HC}=
 \begin{cases}
  O\bigl(1+\abs{\log\rho}\bigr), & N=1,\\[2pt]
  O\bigl(\rho^{1-N}\bigr), & N\ge2 .
 \end{cases}
 \label{d:eq:primitive-growth}
\end{equation}
\end{lemma}

\begin{proof}
Use the global potentials $\pot,\potV$ of
Definition~\ref{b:def:potentials}.  On a small disk
$\abs z$ is bounded, so the pointwise bounds of
Proposition~\ref{d:prop:period-bounds} give
$\abs{\nabla\pot}\le c_1\rho^{-N}$ and $\abs{\nabla\potV}\le c_2\rho^{-N}$.
Fix an outer circle $\abs{z-p}=\rho_0<R$.  To reach
$p+\rho e^{\iota\theta}$, follow the outer circle from a base point to
angle $\theta$ --- a uniformly bounded contribution --- and then the radial
segment inward, contributing at most a constant times
$\int_\rho^{\rho_0}\varsigma^{-N}\,\dd\varsigma$, which is
$O(1+\abs{\log\rho})$ for $N=1$ and $O(\rho^{1-N})$ for $N\ge2$.  Since
$F_1=z\pot+\potV$ with $z$ bounded, \eqref{d:eq:primitive-growth} follows;
adding an element of $\Aff$ does not change it.
\end{proof}

\begin{lemma}[Growth of the holomorphic invariant]
\label{d:lem:detector-growth}
If $\Msph_g(\rho)=O(\rho^{-N})$ for an integer $N\ge1$, then the invariant
$\Hg$ of Definition~\ref{b:def:H} satisfies $\Hg(z)=O(\rho^{-N-1})$, so $\Hg$ is
meromorphic at $p$ with pole order at most $N+1$.
\end{lemma}

\begin{proof}
Every real component of $g$ is harmonic, since
$\overline{\Dop}\Dop g=\lap g=0$.  A ball of radius $\rho/2$ centred at a
point at distance $\rho$ from $\Sp$ stays in the punctured tube, with
distance to $\Sp$ between $\rho/2$ and $3\rho/2$ throughout.  The interior
gradient estimate for harmonic functions --- which follows from the volume
mean-value property and the divergence theorem in the form
$\abs{\partial_\mu h(w)}\le(4/r)\sup_{B(w,r)}\abs h$ --- therefore gives
$\abs{\nabla g}=O(\rho^{-N-1})$ uniformly in $I$.  Evaluating on a fixed
pair of opposite slices recovers $P,Q$ and their planar derivatives from
$g$ and its derivatives with bounded constants, and $y$ stays away from
zero, so the defining formula
$\Hg=-\tfrac12(P+yP_y)-\tfrac{\iota y}2P_x$ of Definition~\ref{b:def:H}
inherits the bound.  The Laurent coefficient estimate then kills every coefficient below
degree $-N-1$.
\end{proof}

\begin{theorem}[Finite-order principal parts and their exact dimension]
\label{d:thm:growth}
Let $N\ge1$ be an integer and $g\in\AM(\Tube_R\setminus\Sp)$.  The
following are equivalent.
\begin{enumerate}[label=\textup{(\roman*)}]
\item $\Msph_g(\rho)=O(\rho^{-N})$ uniformly as $\rho\downarrow0$.
\item $\Max_g(\rho)=O(\rho^{-N})$ uniformly as $\rho\downarrow0$.
\item The normal form \eqref{d:eq:normal-form} of $g$ terminates:
\begin{equation}
 g=g_{\mathrm{reg}}+G_{p;a,b}
   +\lap\I\left(\sum_{k=1}^{N-1}(z-p)^{-k}c_k\right),
 \label{d:eq:finite-normal}
\end{equation}
the sum being empty for $N=1$, with $a,b\in\HH$ and $c_k\in\HC$ arbitrary
and mutually independent.
\item $\Hg$ is meromorphic at $p$ with pole order at most $N+1$.
\end{enumerate}
Writing $\mathcal M_N$ for the space of germs at $\Sp$ satisfying these
conditions and $\mathcal M_0$ for the regular germs,
\begin{equation}
 \mathcal M_N/\mathcal M_0\ \cong\ \HH^2\oplus\HC^{\,N-1},
 \qquad
 \dim_{\HH}\bigl(\mathcal M_N/\mathcal M_0\bigr)=2N,
 \qquad
 \dim_{\RR}\bigl(\mathcal M_N/\mathcal M_0\bigr)=8N .
 \label{d:eq:local-dim}
\end{equation}
\end{theorem}

\begin{proof}
(i)$\Leftrightarrow$(ii) is the norm equivalence
\eqref{d:eq:norm-equivalence}, which changes the implied constant by at
most $\sqrt2$ and no exponent.

(ii)$\Rightarrow$(iii).  The branch-free formulas of Part~IV give
$\abs{G_{p;a,b}}=O(\rho^{-1})$: the numerators $\xi,\eta$ of the
$\rho^{-2}$ terms are $O(\rho)$, $y$ and $xa+b$ remain bounded, and
$\abs{\log\rho}=O(\rho^{-1})$.  Subtracting the mode with the actual
residue pair preserves the bound and produces a zero-period remainder
$g_1$; choose a single-valued primitive stem $F_1$.  By
Lemma~\ref{d:lem:primitive-growth} and the Cauchy coefficient estimate,
$\norm{c_k}_{\HC}\le\rho^{k}\max_{\abs{z-p}=\rho}\norm{F_1}_{\HC}$, which
tends to zero for every $k\ge1$ when $N=1$, and for every $k\ge N$ when
$N\ge2$.  Only the coefficients displayed in \eqref{d:eq:finite-normal}
survive.

(iii)$\Rightarrow$(i).  For $F_k=(z-p)^{-k}c_k$ both stem components are
$O(\rho^{-k})$ and their first derivatives $O(\rho^{-k-1})$; the axial
formula \eqref{a:eq:Fu-first}, with $y$ bounded away from zero, gives
$\lap\I(F_k)=O(\rho^{-k-1})=O(\rho^{-N})$ for $k\le N-1$.  With the
$O(\rho^{-1})$ mode and the bounded regular term this proves (i).

(i)$\Rightarrow$(iv) is Lemma~\ref{d:lem:detector-growth}.
(iv)$\Rightarrow$(iii): subtract the mode determined by the two periods;
by the residue formulas \eqref{b:residue-pair} this cancels the entire
polar part of
$\Hg$ of order $\le2$, so the remainder $g_1$ has zero periods and $\Hg$
with a pole of order at most $N+1$.  Its single-valued primitive stem
$F_1=\sum_nd_n(z-p)^n$ has $F_1''=\Hg$ (Theorem~\ref{b:thm:invariant}),
and the coefficient of
$(z-p)^{n-2}$ in $F_1''$ is $n(n-1)d_n$, nonzero for $n<0$; the pole-order
bound therefore forces $d_n=0$ for $n<-(N-1)$, which is
\eqref{d:eq:finite-normal}.

Finally, the data $\bigl((a,b),(c_k)_{k<N}\bigr)$ map right
$\HH$-linearly onto $\mathcal M_N/\mathcal M_0$; injectivity is the
uniqueness in Theorem~\ref{d:thm:normal} and surjectivity is
\eqref{d:eq:finite-normal}.  Since $\dim_\RR\HH^2=8$ and $\dim_\RR\HC=8$,
the real dimension is $8+8(N-1)=8N$.
\end{proof}

\begin{remark}[Two proofs, and what the count is not]
\label{d:rem:two-proofs}
The implications (ii)$\Rightarrow$(iii) and (iv)$\Rightarrow$(iii) are two
independent proofs of the vanishing of the high Laurent coefficients.  The
first bounds the \emph{primitive} and applies Cauchy's coefficient
estimate; the second bypasses the primitive entirely, bounding the
target-derived invariant $\Hg$ by the harmonic interior gradient estimate.
The second route is available because each real component of an axially
monogenic field is harmonic, so it uses no period theory at all.

The count in \eqref{d:eq:local-dim} is $8N$: eight real dimensions from the
affine residue pair $(a,b)\in\HH^2$, plus $8(N-1)$ from the complexified
Laurent coefficients $c_1,\dots,c_{N-1}\in\HC$.  It is emphatically not
$4N$ --- the Laurent coefficients live in $\HC$, not in $\HH$ --- and not
$8N+8$, because the residue pair is not an additional datum beyond the
$N=1$ level but is precisely that level.
\end{remark}

\begin{corollary}[An explicit real basis and exact orders]
\label{d:cor:basis}
With $e_0=1$, $e_1=\mathbf i$, $e_2=\mathbf j$, $e_3=\mathbf k$, a real
basis of $\mathcal M_N/\mathcal M_0$ is
\[
 G_{p;e_\ell,0},\quad G_{p;0,e_\ell}\quad(0\le\ell\le3),
 \qquad
 \lap\I\bigl((z-p)^{-k}e_\ell\bigr),\quad
 \lap\I\bigl((z-p)^{-k}\iota e_\ell\bigr)
 \quad(1\le k\le N-1,\ 0\le\ell\le3).
\]
Every nonzero critical residue class has exact uniform order $\rho^{-1}$,
and every class whose highest Laurent index is $k$ has exact uniform order
$\rho^{-k-1}$: the Fueter map loses exactly one power of the distance at a
nonreal sphere.
\end{corollary}

\begin{proof}
The basis is \eqref{d:eq:local-dim} written out.  The exact orders are the
positive limits of Theorem~\ref{d:thm:leading}(ii) and
Corollary~\ref{d:cor:higher-order} below.
\end{proof}

\subsection{(D) Sharp transverse removability}

\begin{corollary}[Critical growth and the sharp threshold]
\label{d:cor:critical}
Let $g\in\AM(\Tube_R\setminus\Sp)$.
\begin{enumerate}[label=\textup{(\roman*)}]
\item If $\Msph_g(\rho)=O(\rho^{-1})$ then $g=G_{p;a,b}+g_{\mathrm{reg}}$,
and $g$ is removable if and only if $\Resaff_{\Sp}(g)=(0,0)$.
\item If $\Msph_g(\rho)=o(\rho^{-1})$ uniformly, then $g$ extends axially
monogenically through $\Sp$, and the extension is unique.
\item The little-$o$ in \textup{(ii)} cannot be weakened to big-$O$: any
$G_{p;a,b}$ with $(a,b)\ne(0,0)$ satisfies the big-$O$ bound and is not
removable.
\end{enumerate}
Quantitatively, on the circle $\abs{z-p}=\rho$,
\begin{equation}
 \abs a\le\pi\rho\,\Msph_g(\rho),
 \qquad
 \abs b\le\pi\rho\bigl(\abs p+\rho\bigr)\Msph_g(\rho).
 \label{d:eq:circle-bounds}
\end{equation}
\end{corollary}

\begin{proof}
Part~(i) is the case $N=1$ of Theorem~\ref{d:thm:growth} together with the
uniqueness of the normal form.  For (ii), the meridian circle has length
$2\pi\rho$ and $\max_{\abs{z-p}=\rho}\abs z\le\abs p+\rho$; substituting
into \eqref{d:eq:period-bounds}, and using
$\axn g\le\abs g_{\mathrm{sph}}$, gives \eqref{d:eq:circle-bounds}.  The
two periods are independent of $\rho$, so the little-$o$ hypothesis forces
them to vanish; part~(i) then supplies the extension, and uniqueness is
continuity.  Part~(iii) is the period normalization of the modes.
\end{proof}

Note the codimension.  The exponent $1$ here reflects a singular set of
real codimension two together with the axial structure.  It is not the
exponent $3$ of the isolated-point Cauchy--Fueter removability theorem of
Part~I, where the distance is to a single point of $\HH$ and the kernel has
size $\abs q^{-3}$.  Mixing the two geometries produces an incorrect
comparison of thresholds and of residues.

\subsection{(E) $L^2$ removability}

\begin{theorem}[Sharp $L^2$ removability across a nonreal sphere]
\label{d:thm:L2}
Let $g$ be Cauchy--Fueter regular on a punctured neighbourhood of $\Sp$ and
locally square integrable across $\Sp$.  Then $g$ extends uniquely to a
smooth Cauchy--Fueter-regular field through $\Sp$.  Axiality is not needed
for this statement; if $g$ is axial, so is the extension.  The exponent $2$
is sharp: for every $1\le s<2$ the nonzero modes $G_{p;a,b}$ are locally
$L^s$ and not removable.
\end{theorem}

\begin{proof}
Assign arbitrary values on the measure-zero set $\Sp$, so that $g\in
L^2_{\mathrm{loc}}$ on the full tube.  Choose a logarithmic cutoff
$\chi_\varepsilon$ equal to $0$ for $\rho\le\varepsilon^2$, to $1$ for
$\rho\ge\varepsilon$, and affine in $\log\rho$ in between, smoothed without
changing the bounds
\[
 \abs{\nabla\chi_\varepsilon}\le\frac{C}{\rho\log(1/\varepsilon)},
 \qquad
 \norm{\nabla\chi_\varepsilon}_{L^2}^2\le\frac{C'}{\log(1/\varepsilon)}
 \longrightarrow0 ,
\]
the second using the codimension-two volume element
$O(\rho\,\dd\rho\,\dd\theta\,\dd S)$, which makes
$\int\rho^{-2}\log(1/\varepsilon)^{-2}\rho\,\dd\rho$ a multiple of
$1/\log(1/\varepsilon)$.  Test $\Dop g=0$ on the punctured tube against
$\chi_\varepsilon\phi$ with $\phi$ smooth and compactly supported.  The
term carrying $\chi_\varepsilon\nabla\phi$ converges to the desired
distributional pairing, and the cutoff error is at most
$C_\phi\norm g_{L^2}\norm{\nabla\chi_\varepsilon}_{L^2}\to0$.  Hence
$\Dop g=0$ distributionally on the full tube; applying
$\overline{\Dop}$ and Weyl's lemma gives a smooth harmonic representative,
which then satisfies the first-order equation pointwise.  Two continuous
extensions agree on the dense complement of $\Sp$, giving uniqueness, and
axiality passes to the limit.  For sharpness, $G_{p;a,b}=O(\rho^{-1})$ and
the transverse polar element $\rho\,\dd\rho$ give
$\int_0^{R}\rho^{1-s}\,\dd\rho<\infty$ for $s<2$, while a nonzero residue
pair obstructs removability by \eqref{d:eq:L1-removal}.
\end{proof}

\begin{remark}[Three routes to the same theorem]
\label{d:rem:L2-routes}
The cutoff proof above is adopted because it is the shortest, needs no
classification machinery, no growth bound and --- as several sources
independently observe --- no axiality; it is a standard first-order
elliptic capacity argument, and the new content here is not the mechanism
but its exact relation to the two periods and to the energy coefficient in
\eqref{e:eq:energylaw}.  Two alternatives are recorded.  (a)~A \emph{linear} shell
cutoff, vanishing within distance $\varepsilon$ and equal to one beyond
$2\varepsilon$ with $\abs{\nabla\chi_\varepsilon}\le C/\varepsilon$: a
compact portion of a codimension-two submanifold has shell volume
$O(\varepsilon^2)$, so $\norm{\nabla\chi_\varepsilon}_{L^2}\le C$ is merely
bounded, and the error still tends to zero, this time by absolute
continuity of $\int\abs g^2$ over the shrinking shells.  (b)~A route
through the classification: local $L^2$ implies, by the harmonic mean-value
property and Cauchy--Schwarz, the bound $\abs g=O(\rho^{-2})$, hence the
case $N=2$ of Theorem~\ref{d:thm:growth}; the surviving Laurent
coefficient is then killed because a nonzero Laurent mode cannot have
finite local $L^2$ energy (compare Proposition~\ref{e:prop:poleenergy}),
and the residues are killed by the energy law~\eqref{e:eq:energylaw}.  Route~(b) is correct but needs the whole $8N$ classification
and is restricted to axial fields, so it is not taken as primary.  Within
the $L^1$ class the chain of equivalences
\[
 \text{removability}
 \iff (a,b)=(0,0)
 \iff \text{finite local }L^2\text{ mass}
\]
holds, the last equivalence being the energy law~\eqref{e:eq:energylaw}.
\end{remark}

\subsection{(F) The sharp $L^p$ trichotomy}

\begin{proposition}[Sharpness of the $L^1$ hypothesis]
\label{d:prop:L1sharp}
For every $0<s<1$ there is an axially monogenic field near $\Sp$, locally
in $L^s$, with $\Resaff_{\Sp}=(0,0)$, that is not removable.  One may take
$g=\lap\I(F)$ with the single-valued stem $F(z)=(z-p)^{-1}c$,
$0\ne c\in\HC$.
\end{proposition}

\begin{proof}
The stem is single valued on the punctured disk, so $g$ is a global Fueter
image there, both periods vanish, and $g$ has a global single-valued
slice-regular primitive on the punctured tube.  The explicit components and
the exact size are computed in Example~\ref{f:ex:rational-mode} of Part~X:
for $c=1$ one finds $\abs P^2+\abs Q^2=4/(v^2\rho^4)+O(\rho^{-3})$, so
$\Max_g(\rho)\sim2/(v\rho^2)$ uniformly in the meridian angle and in
$I\in\SSS$, with no exceptional quaternionic direction.  Hence $g$ is
locally $L^s$ exactly when $\int_0\rho^{1-2s}\,\dd\rho<\infty$, that is for
$s<1$, and $g$ is \emph{not} locally $L^1$.  It is not removable:
otherwise the local inverse theorem on the full meridian disk would give a
holomorphic stem $F_\sharp$ with $\lap\I(F_\sharp)=g$ there, and
Corollary~\ref{b:cor:kernel} would force $F-F_\sharp=za'+b'$ on the
punctured disk, impossible since $F$ has a pole.  This last step needs no
pointwise lower bound in a possibly exceptional quaternionic direction.
\end{proof}

\begin{corollary}[Sharp $L^p$ trichotomy at a nonreal sphere]
\label{d:cor:Lp}
For locally axially monogenic germs along $\Sp$:
\begin{enumerate}[label=\textup{(\roman*)}]
\item if $1\le s<2$, the locally $L^s$ germs modulo removable germs form a
right quaternionic vector space of dimension two, with coordinates
$(a,b)$ --- real dimension eight;
\item if $2\le s\le\infty$, every locally $L^s$ germ is removable;
\item for every $0<s<1$ the residues are \emph{not} a complete invariant
modulo removable germs.
\end{enumerate}
\end{corollary}

\begin{proof}
The modes are locally $L^s$ for $s<2$ by the computation in
Theorem~\ref{d:thm:L2}.  For $s\ge1$, local $L^s$ implies local $L^1$ on a
finite-volume tube by H\"older, so Theorem~\ref{d:thm:L1} applies and gives
(i).  For $s\ge2$, local $L^s$ implies local $L^2$ and
Theorem~\ref{d:thm:L2} gives (ii).  Part~(iii) is
Proposition~\ref{d:prop:L1sharp}.
\end{proof}

\begin{remark}[Two different thresholds]
\label{d:rem:two-thresholds}
The exponents $1$ and $2$ are two distinct sharp thresholds describing two
distinct phenomena, and neither is a consequence of the other.  At $L^1$
the pair of periods \emph{becomes a complete invariant} of the singularity;
at $L^2$ the periods are \emph{forced to vanish}.  Below $L^1$ the periods
survive as invariants but no longer classify.  The example separating
$L^1$ from $L^2$ is the mode itself, which is $L^s$ for all $s<2$ and never
removable; the example separating $L^s$, $s<1$, from $L^1$ is
Proposition~\ref{d:prop:L1sharp}, which shows in addition that a global
single-valued primitive on a punctured domain is not the same thing as
removability through the missing sphere.
\end{remark}

\subsection{(G) Exact leading constants, in both norms}

The classification decides whether a singularity is present.  We now
determine its exact size, which is where the sphere $\SSS$ of imaginary
units must be handled explicitly rather than replaced by one complex
section.  Recall the residue norm and pair norm
\begin{equation}
 \Rp(a,b)^2=\abs{ua+b}^2+v^2\abs a^2=\norm{pa+b}_{\HC}^2,
 \qquad
 \nn(\alpha,\beta)=\max_{I\in\SSS}\abs{\beta+I\alpha}
 =\bigl(\abs\alpha^2+\abs\beta^2
        +2\abs{\ImH(\beta\bar\alpha)}\bigr)^{1/2},
 \label{d:eq:residue-norm}
\end{equation}
the second evaluation being \eqref{d:eq:sphere-max}.  By
\eqref{d:eq:norm-equivalence},
$(\abs\alpha^2+\abs\beta^2)^{1/2}\le\nn(\alpha,\beta)
 \le\sqrt2(\abs\alpha^2+\abs\beta^2)^{1/2}$, so $\nn$ is a genuine norm on
the eight-real-dimensional pair space.  $\Rp$ is positive definite because
$v>0$, and it is unchanged by a real translation of the origin, since
$ua+b$ is.

\begin{lemma}[The uniform leading term]
\label{d:lem:leading}
Let $g$ have residue pair $(a,b)$ and satisfy the critical bound of
Corollary~\ref{d:cor:critical}\textup{(i)}, or merely be locally $L^1$
across $\Sp$.  Put $c_0=ua+b$ and $c_1=va$, so that $pa+b=c_0+\iota c_1$,
and write $z=p+\rho e^{\iota\theta}$.  Then, uniformly in $\theta$ and in
$I\in\SSS$,
\begin{align}
 P&=\frac{\cos\theta\,c_0+\sin\theta\,c_1}{\pi v\rho}
     +O\bigl(1+\abs{\log\rho}\bigr),
 &
 Q&=\frac{-\sin\theta\,c_0+\cos\theta\,c_1}{\pi v\rho}
     +O\bigl(1+\abs{\log\rho}\bigr),
 \label{d:eq:leading-PQ}
\end{align}
equivalently
\begin{equation}
 P+\iota Q=\frac{pa+b}{\pi v(z-p)}+O\bigl(1+\abs{\log\rho}\bigr),
 \qquad
 g(x+Iy)=\frac{e^{-I\theta}\bigl((ua+b)+Iva\bigr)}{\pi v\rho}
         +O\bigl(1+\abs{\log\rho}\bigr).
 \label{d:eq:leading}
\end{equation}
The quaternionic coefficients remain on the right throughout; no
commutation with $I$ is implicit.
\end{lemma}

\begin{proof}
By Theorem~\ref{d:thm:L1} (or Corollary~\ref{d:cor:critical}(i)) the
remainder $g_0$ is bounded on a smaller tube, so only $G_{p;a,b}$
contributes.  Take the local stem $F=(za+b)\ell_p$ with
$\ell_p=(2\pi\iota)^{-1}\log(z-p)$ on a branch whose argument is bounded at
the point being estimated; a change of branch adds an element of $\Aff$ and
does not change $\lap\I(F)$.  Then
\[
 F'=a\ell_p+\frac{za+b}{2\pi\iota(z-p)}
   =\frac{c_0+\iota c_1}{2\pi\iota\rho e^{\iota\theta}}
    +O\bigl(1+\abs{\log\rho}\bigr),
 \qquad
 \norm F_{\HC}=O\bigl(1+\abs{\log\rho}\bigr).
\]
The axial Fueter formula \eqref{a:eq:Fu-first} of Part~I, written on stems
as
$P=-\tfrac2y\Ima_{\CC}F'$ and
$Q=\tfrac2y\Rea_{\CC}F'-\tfrac2{y^2}\Ima_{\CC}F$, together with
$y=v+O(\rho)$ and $(\iota e^{\iota\theta})^{-1}
 =-\sin\theta-\iota\cos\theta$, gives \eqref{d:eq:leading-PQ}.  Assembling
$P+\iota Q$ produces $(c_0+\iota c_1)/(\pi v\rho e^{\iota\theta})$, which
is the first formula in \eqref{d:eq:leading}, and
$P+IQ=e^{-I\theta}(c_0+Ic_1)/(\pi v\rho)$ by expanding
$e^{-I\theta}=\cos\theta-I\sin\theta$.  One may also read
\eqref{d:eq:leading-PQ} off the branch-free formulas
\eqref{c:eq:mode-P}--\eqref{c:eq:mode-Q} by
substituting $\xi=\rho\cos\theta$, $\eta=\rho\sin\theta$, $y=v+O(\rho)$ and
$xa+b=c_0+O(\rho)$, the logarithmic terms being $O(1+\abs{\log\rho})$.
\end{proof}

\begin{theorem}[Exact leading size at a critical singular sphere]
\label{d:thm:leading}
Let $g$ be as in Lemma~\textup{\ref{d:lem:leading}}, with residue pair
$(a,b)$.
\begin{enumerate}[label=\textup{(\roman*)}]
\item \textup{(Axial pair norm.)}  Uniformly in $\theta$,
\begin{equation}
 \rho\,\axn g(p+\rho e^{\iota\theta})\longrightarrow
 \frac{\Rp(a,b)}{\pi v},
 \qquad\text{hence}\qquad
 \rho\,\Max_g(\rho)\longrightarrow\frac{\Rp(a,b)}{\pi v}.
 \label{d:eq:axial-leading}
\end{equation}
\item \textup{(Spherical supremum.)}
\begin{equation}
 \rho\,\Msph_g(\rho)\longrightarrow
 \frac1{\pi v}
 \sqrt{\Rp(a,b)^2+2v\,\bigl\lvert\ImH(b\bar a)\bigr\rvert}
 \ =\ \frac1\pi\,\nn\!\left(a,\frac{ua+b}{v}\right).
 \label{d:eq:exact-sup}
\end{equation}
\end{enumerate}
Both limits are strictly positive for every nonzero residue pair.  The two
constants are different: they agree if and only if $\ImH(b\bar a)=0$, and
their ratio never exceeds $\sqrt2$.
\end{theorem}

\begin{proof}
Write $c_0=ua+b$, $c_1=va$.  The pair in \eqref{d:eq:leading-PQ} is
obtained from $(c_0,c_1)$ by the real rotation of angle $\theta$, scaled by
$(\pi v\rho)^{-1}$; rotations preserve the sum of squared quaternionic
norms, so
\[
 \abs{P_*}^2+\abs{Q_*}^2
 =\frac{\abs{c_0}^2+\abs{c_1}^2}{\pi^2v^2\rho^2}
 =\frac{\Rp(a,b)^2}{\pi^2v^2\rho^2},
\]
independently of $\theta$.  Since the error in \eqref{d:eq:leading-PQ} is
$O(1+\abs{\log\rho})=o(\rho^{-1})$, multiplying by $\rho$ and letting
$\rho\downarrow0$ gives \eqref{d:eq:axial-leading}, uniformly in $\theta$
and hence also for the maximum over the circle.

For (ii) apply \eqref{d:eq:sphere-max} to the leading pair.  Using
\eqref{d:eq:rotation-cross} with $X=c_0$, $Y=c_1$, $\alpha=\cos\theta$ and
$\beta=-\sin\theta$, the cross term of the rotated pair is
\[
 \bigl\lvert\ImH(P_*\overline{Q_*})\bigr\rvert
 =\frac{\bigl\lvert\ImH(c_0\bar c_1)\bigr\rvert}{\pi^2v^2\rho^2}
 =\frac{v\bigl\lvert\ImH\bigl((ua+b)\bar a\bigr)\bigr\rvert}
        {\pi^2v^2\rho^2}
 =\frac{v\bigl\lvert\ImH(b\bar a)\bigr\rvert}{\pi^2v^2\rho^2},
\]
because $ua\bar a=u\abs a^2$ is real and drops out of $\ImH$.  Again the
result is independent of $\theta$.  Substituting into
\eqref{d:eq:sphere-max} gives
\[
 \max_{I\in\SSS}\abs{P_*+IQ_*}^2
 =\frac{\Rp(a,b)^2+2v\abs{\ImH(b\bar a)}}{\pi^2v^2\rho^2},
\]
and the uniform $O(1+\abs{\log\rho})$ error gives
\eqref{d:eq:exact-sup}.  The second expression in \eqref{d:eq:exact-sup}
is the same quantity rewritten with $\beta_0=(ua+b)/v$, since
$\nn(a,\beta_0)^2=\abs a^2+\abs{\beta_0}^2
 +2\abs{\ImH(\beta_0\bar a)}
 =v^{-2}\bigl(\Rp(a,b)^2+2v\abs{\ImH(b\bar a)}\bigr)$.
Positivity holds because $\Rp$ is positive definite for $v>0$, and the
comparison of the two constants is \eqref{d:eq:norm-equivalence}.
\end{proof}

\begin{remark}[The correction is governed by $\ImH(b\bar a)$, not by
commutativity]
\label{d:rem:correction}
The extra term $2v\abs{\ImH(b\bar a)}$ in \eqref{d:eq:exact-sup} measures
the mismatch between the average of $\abs g^2$ over a conjugacy sphere and
its largest value on that sphere.  It vanishes exactly when
$\ImH(b\bar a)=0$, which is \emph{not} the same as $ab=ba$.  It is
therefore never legitimate to replace the four-dimensional supremum by
evaluation on one preselected complex section: this is the single point at
which the two laws quoted in the literature under the same symbol $M_g$
genuinely diverge, and the two constants must be reported with their
respective norms.
\end{remark}

The worked instance is Example~\ref{f:ex:im-correction} of Part~X:
$p=\iota$, $a=1$, $b=\mathbf i$ give $\Rp(a,b)^2=2$ and
$\abs{\ImH(b\bar a)}=1$, hence the two different constants $2/\pi$ and
$\sqrt2/\pi$ in \eqref{d:eq:exact-sup} and \eqref{d:eq:axial-leading} ---
for a pair $a,b$ that \emph{commutes}.  There
$g(\rho+I)=(\mathbf i+I)/(\pi\rho)+O(1+\abs{\log\rho})$ vanishes at
$I=-\mathbf i$ and is maximal at $I=\mathbf i$, so pointwise cancellation
along a selected approach to the sphere does not make the spherical
singularity removable, and a numerical probe restricted to one complex
slice can report the wrong constant or no singularity at all.

\begin{corollary}[Exact leading size above the critical order]
\label{d:cor:higher-order}
Let $g\in\mathcal M_N\setminus\mathcal M_{N-1}$ with $N\ge2$, and let
$c_{N-1}=\alpha+\iota\beta\ne0$ be the top Laurent coefficient in
\eqref{d:eq:finite-normal}, $\alpha,\beta\in\HH$.  Then, with $k=N-1$,
\begin{equation}
 \rho^{k+1}\Max_g(\rho)\longrightarrow
 \frac{2k}{v}\norm{c_k}_{\HC}
 =\frac{2k}{v}\bigl(\abs\alpha^2+\abs\beta^2\bigr)^{1/2},
 \qquad
 \rho^{k+1}\Msph_g(\rho)\longrightarrow
 \frac{2k}{v}\,\nn(\alpha,\beta) .
 \label{d:eq:pole-asymptotic}
\end{equation}
Equivalently, in terms of the Laurent coefficients $h_n$ of the invariant
$\Hg$ of Part~III, where $c_k=h_{-k-2}/\bigl(k(k+1)\bigr)$, the invariant
$\Hg$ has a pole of order exactly $N+1$ and
\begin{equation}
 \Max_g(\rho)\sim\frac{2\norm{h_{-N-1}}_{\HC}}{vN}\,\rho^{-N}.
 \label{d:eq:J-asymptotic}
\end{equation}
\end{corollary}

\begin{proof}
Put $w=z-p$ and $F_k=w^{-k}c_k$.  With $\phi=(k+1)\theta$, the components
of $F_k'=-kw^{-k-1}c_k$ are
$A_x=-k\rho^{-k-1}(\cos\phi\,\alpha+\sin\phi\,\beta)$ and
$B_x=-k\rho^{-k-1}(\cos\phi\,\beta-\sin\phi\,\alpha)$; using
$A_y=-B_x$, $B_y=A_x$ and $y=v+O(\rho)$ in the axial Fueter formula gives
\begin{equation}
 \lap\I(F_k)(x+Iy)
 =\frac{2k}{v\rho^{k+1}}e^{-I(k+1)\theta}\bigl(\beta-I\alpha\bigr)
 +O(\rho^{-k}),
 \label{d:eq:highest-term}
\end{equation}
uniformly in $\theta$ and $I$.  Equivalently, in complexified form,
$Q-\iota P=-\tfrac{2k}{v}w^{-k-1}c_k+O(\rho^{-k})$, where replacing $y$ by
$v$ and the term $\Ima_{\CC}F_k/y^2$ each cost only $O(\rho^{-k})$.  All
lower polar terms, the logarithmic mode (which is $O(\rho^{-1})$) and the
regular term contribute $O(\rho^{-k})$.  Since
$\norm{Q-\iota P}_{\HC}=\axn g$ and central complex multiplication scales
$\norm\cdot_{\HC}$, the first limit in \eqref{d:eq:pole-asymptotic}
follows; the second follows from \eqref{d:eq:highest-term} and
\eqref{d:eq:sphere-max}, since $e^{-I(k+1)\theta}$ is a unit quaternion and
$\max_I\abs{\beta-I\alpha}=\nn(\alpha,\beta)$, the maximum being unchanged
by the sign of $I\alpha$.  Both limits are positive by
\eqref{d:eq:norm-equivalence} and $c_k\ne0$.  The statement about $\Hg$
follows from the multipole extraction formula
$c_k=h_{-k-2}/(k(k+1))$ of \eqref{b:multipoles} with $k=N-1$, which gives
$h_{-N-1}=N(N-1)c_{N-1}$ and hence
$2\norm{h_{-N-1}}_{\HC}/(vN)=2(N-1)\norm{c_{N-1}}_{\HC}/v$.
\end{proof}

\begin{corollary}[No fractional singularity orders]
\label{d:cor:fractional}
Let $\alpha\ge0$ be a \emph{real} exponent.  A bound
$\Max_g(\rho)=O(\rho^{-\alpha})$ --- equivalently, by
\eqref{d:eq:norm-equivalence}, $\Msph_g(\rho)=O(\rho^{-\alpha})$ ---
allows precisely the principal parts of \eqref{d:eq:finite-normal} with
$N=\lfloor\alpha\rfloor$, where $N=0$ means no principal part at all.  The
corresponding singular quotient
$\mathcal M_{\lfloor\alpha\rfloor}/\mathcal M_0$ has real dimension
$8\lfloor\alpha\rfloor$.  In particular every nonremovable singularity of
finite power growth has a positive \emph{integer} growth order, and the
possible uniform orders are exactly $\rho^{-1},\rho^{-2},\rho^{-3},\dots$
\end{corollary}

\begin{proof}
Apply Theorem~\ref{d:thm:growth} with the integer
$\max(1,\lceil\alpha\rceil)$ to obtain a finite normal form.  If the
surviving principal part had least integer order $N>\alpha$, then
Theorem~\ref{d:thm:leading} (for $N=1$) or
Corollary~\ref{d:cor:higher-order} (for $N\ge2$) would give a strictly
positive limit for $\rho^{N}\Max_g(\rho)$, contradicting
$\Max_g(\rho)=O(\rho^{-\alpha})$ with $\alpha<N$.  Conversely each such
normal form does satisfy the bound, by the same theorems, and regular germs
are bounded.  The dimension is \eqref{d:eq:local-dim}.
\end{proof}

\begin{remark}[What zero residues do and do not say, and the limits of
this classification]
\label{d:rem:scope-VII}
Three notions are separated by the results above and must not be conflated:
existence of a global single-valued slice-regular primitive on the
punctured tube; vanishing of the affine spherical residues; and
removability through $\Sp$.  The first two are equivalent on a punctured
disk, by Theorem~\ref{b:thm:period}.  Neither implies the third without a hypothesis: for
$0\ne c\in\HC$ the field $\lap\I((z-p)^{-1}c)$ has zero periods, a
single-valued primitive, exact uniform order $\rho^{-2}$, and is not
removable (Proposition~\ref{d:prop:L1sharp} and
Example~\ref{f:ex:rational-mode}); and
$\lap\I\bigl(\exp((z-p)^{-1})\bigr)$ has zero periods, an essential
singularity of its invariant $\Hg$, and admits no finite power bound at
all, so it lies outside every $\mathcal M_N$
(Example~\ref{f:ex:essential}).  What supplies the missing
hypothesis is local integrability (Theorem~\ref{d:thm:L1}) or a finite
growth order (Theorem~\ref{d:thm:growth}), and in the presence of either
the residue pair is a complete or a partial invariant respectively.

Finally, the scope.  Every classification statement in this part is within
the \emph{axial} class; general Cauchy--Fueter-regular fields near a
two-sphere are not classified here, and only the $L^2$ removability
Theorem~\ref{d:thm:L2} is free of the axiality hypothesis.  All constants
above contain a factor $1/v$ or $1/v^2$ and therefore degenerate as the
sphere collapses onto the real axis; no uniform statement as $v\downarrow0$
is claimed, and a real point is a different local problem with different
thresholds.  Nothing in this part asserts a classification at a singular
set that is not a single conjugacy sphere.
\end{remark}

%% ==================== part8_9.tex ====================

% =====================================================================
% Part VIII.  The surface source, the conserved fluxes, and the exact
%             energy law.
% Lead sources: 07, 10, 01, with 05, 02 and 08 where the contract marks a
% result as shared.  Notation follows the merge contract without exception:
% p = u + \iota v, \xi = x-u, \eta = y-v, \rho = |z-p|; \thet is the
% primitive-independent second form; \Rp is the residue norm; \Msph and
% \Max are kept rigorously apart.
% =====================================================================

\section{The surface source, the conserved fluxes, and the logarithmic
energy law}
\label{e:sec:source-energy}

Parts~II--VII produced the affine spherical residue
$\Resaff_{\Sp}(g)=(a,b)\in\HH^2$ in three guises: as a pair of periods of
the two closed forms $\om_g$ and $\thet_g$, as a pair of complex moments
of the holomorphic invariant $\Hg$, and as the coefficients of the
logarithmic mode in the local normal form. All three are contour
constructions in the planar base. The present part identifies the same
eight real parameters by three four-dimensional measurements made on the
field itself: a distributional charge carried by the singular sphere, two
conserved three-dimensional boundary fluxes, and the exact coefficient of
a logarithmic energy divergence. None of the three requires a primitive,
a branch, or a choice of loop.

Throughout, $p=u+\iota v$ with $v>0$, $\Sp=u+v\SSS$, and
\[
 \Tube_\rho(\Sp)=\{q\in\HH:\operatorname{dist}(q,\Sp)<\rho\},\qquad
 \Tube_{\varepsilon,R}
   =\{q:\varepsilon<\operatorname{dist}(q,\Sp)<R\},\qquad 0<\varepsilon<R<v.
\]
For $q=x+Iy$ the transverse distance is $\operatorname{dist}(q,\Sp)=\abs{z-p}=\rho$
with $z=x+\iota y$. We use repeatedly the two exact measure identities
for the tube,
\begin{equation}
 \dd S_{\partial\Tube_\rho}=y^2\rho\,\dd\theta\,\dd\sigma(I),
 \qquad
 \dd V=y^2\,\dd x\,\dd y\,\dd\sigma(I)
      =y^{2}\rho\,\dd\rho\,\dd\theta\,\dd\sigma(I),
 \label{e:eq:tube-measures}
\end{equation}
where $x=u+\rho\cos\theta$, $y=v+\rho\sin\theta$, and $\sigma$ is the
ordinary surface measure on $\SSS$ with $\sigma(\SSS)=4\pi$. Part~VII's
uniform leading term is quoted so often below that we display it once, in
the form it will be used:
\begin{equation}
 g(x+Iy)=\frac{e^{-I\theta}\bigl(c_0+Ic_1\bigr)}{\pi v\rho}
             +O\bigl(1+\abs{\log\rho}\bigr),
 \qquad c_0=ua+b,\quad c_1=va,
 \label{e:eq:leading}
\end{equation}
uniformly in $\theta$ and in $I\in\SSS$; equivalently
$P+\iota Q=(pa+b)/(\pi v(z-p))+O(1+\abs{\log\rho})$ in $\HC$. The
quaternionic coefficients stand on the right and no commutation with $I$
is implicit. The residue norm of the contract is
\begin{equation}
 \Rp(a,b)^2=\abs{ua+b}^2+v^2\abs a^2
           =\norm{pa+b}_{\HC}^2=\abs{c_0}^2+\abs{c_1}^2,
 \label{e:eq:residue-norm}
\end{equation}
positive definite in $(a,b)$ precisely because $v>0$.

\subsection{The affine surface source}
\label{e:sec:surface-source}

For a compact smooth surface $S\subset\HH$ and a quaternion-valued
density $j$ on it, $j\,\delta_S$ denotes the $\HH$-valued distribution
$\langle j\delta_S,\varphi\rangle=\int_Sj(q)\varphi(q)\dd S(q)$ acting on
real scalar test functions; the operator $\Dop$ carries its quaternionic
units on the left, in the distributional pairing as in the classical one.

\begin{theorem}[The singular sphere as an affine Cauchy--Fueter source]
\label{e:thm:source}
Let $g$ be axially monogenic on a punctured tube about $\Sp$ and
\emph{locally integrable across $\Sp$}, with affine spherical residue
$\Resaff_{\Sp}(g)=(a,b)$. Extend $g$ arbitrarily over the measure-zero
set $\Sp$. Then, on a sufficiently small full tube,
\begin{equation}
 \Dop g=\sigma_p\,\delta_{\Sp},
 \qquad
 \sigma_p(s)=\frac2v\,(sa+b)\quad (s\in\Sp),
 \label{e:eq:source}
\end{equation}
in the sense of distributions. Written out at $s=u+Iv$,
\begin{equation}
 \sigma_p(u+Iv)=\frac2v\bigl(c_0+Ic_1\bigr)
   =2\left(\frac{ua+b}{v}+I\,a\right).
 \label{e:eq:source-expanded}
\end{equation}
Conversely, the map $(a,b)\mapsto\sigma_p$ is injective, and its image is
exactly the space of densities that are \emph{affine} in the point of the
sphere: one constant angular mode and one mode linear in $I$.
\end{theorem}

\begin{proof}
Let $\varphi$ be a real test function supported in the tube. Because $g$
is locally integrable, integration by parts on the complement of the
shrinking tube $\Tube_\rho(\Sp)$ is legitimate and gives
\[
 \langle\Dop g,\varphi\rangle
   =\lim_{\rho\downarrow0}\int_{\partial\Tube_\rho}
        n_\rho\,g\,\varphi\,\dd S,
 \qquad n_\rho=\cos\theta+I\sin\theta=e^{I\theta}.
\]
The sign is positive because the outward normal of the punctured region
along its inner boundary is $-n_\rho$. Multiplying \eqref{e:eq:leading}
on the left by $n_\rho$ cancels $e^{-I\theta}$ exactly, so
\[
 n_\rho g=\frac{c_0+Ic_1}{\pi v\rho}+O\bigl(1+\abs{\log\rho}\bigr).
\]
By \eqref{e:eq:tube-measures} the surface element carries a factor
$y^2\rho=v^2\rho+O(\rho^2)$, so the remainder contributes
$O(\rho(1+\abs{\log\rho}))\to0$, while $\varphi(x+Iy)\to\varphi(u+Iv)$
uniformly. The limit is
\[
 \frac{v}{\pi}\int_{\SSS}\int_0^{2\pi}
   (c_0+Ic_1)\,\varphi(u+Iv)\,\dd\theta\,\dd\sigma(I)
 =2v\int_{\SSS}(c_0+Ic_1)\varphi(u+Iv)\dd\sigma(I),
\]
and since $\dd S_{\Sp}=v^2\dd\sigma(I)$ this equals
$\int_{\Sp}(2/v)(sa+b)\varphi(s)\dd S(s)$, which is
\eqref{e:eq:source}. Injectivity: subtracting the values of $\sigma_p$ at
$u+Iv$ and $u-Iv$ gives $4a$, so $a=0$, and then $b=0$. Surjectivity onto
the affine densities is the identity \eqref{e:eq:source-expanded} read
backwards.
\end{proof}

\begin{remark}[Where axiality has content]
\label{e:rem:affine-density}
The hypothesis used in Theorem~\ref{e:thm:source} is the weakest of the
family: local integrability across $\Sp$, not $\Max_g(\rho)=O(\rho^{-1})$.
The conclusion is nevertheless a finite-dimensional restriction on what a
sphere can carry. An arbitrary surface density on $\Sp$ produces a
locally integrable field with $\Dop$-source supported on $\Sp$; the
axially monogenic ones use only the eight-real-dimensional affine
subspace $\{s\mapsto(2/v)(sa+b)\}$, which is exactly the trace on $\Sp$ of
the affine kernel $\Aff$. Fields with general densities exist and are not
classified here.
\end{remark}

\subsection{The Cauchy--Fueter layer representation}
\label{e:sec:layer}

Let $E(q)=\bar q/(2\pi^2\abs q^4)$ be the normalized Cauchy--Fueter
kernel of Part~I, so that $\Dop E=\delta_0$ in distributions; the
normalization is fixed by $nE=1/(2\pi^2r^3)$ on the sphere $\abs q=r$,
whose area is $2\pi^2r^3$, so the total flux is $1$.

\begin{theorem}[Canonical layer representation, with uniqueness]
\label{e:thm:layer}
For all $a,b\in\HH$ and $q\notin\Sp$,
\begin{equation}
 G_{p;a,b}(q)=\frac2v\int_{\Sp}E(q-s)\,(sa+b)\,\dd S(s).
 \label{e:eq:layer}
\end{equation}
The right-hand side is the \emph{unique} locally integrable field on
$\HH$ that tends to zero at infinity and whose distributional
Cauchy--Fueter source is $\sigma_p\delta_{\Sp}$. The uniqueness step does
not assume that a competing field is axial.
\end{theorem}

\begin{proof}
Call the right-hand side $K$. Differentiation under the integral sign
makes $K$ smooth and monogenic off $\Sp$. It is locally integrable in
four variables: by Fubini, $\int_{\mathrm{cpt}}\abs K\dd V$ is bounded by
a constant times $\int_{\Sp}\int_{\mathrm{cpt}}\abs{q-s}^{-3}\dd V_q\dd
S(s)$, finite and uniform in $s$ because $\abs q^{-3}$ is integrable in
$\RR^4$. Convolving the fundamental solution with the compactly supported
surface density gives $\Dop K=\sigma_p\delta_{\Sp}$, and the bound for
$E$ gives $K(q)=O(\abs q^{-3})$ at infinity.

By Part~IV the left-hand side $G_{p;a,b}$ is smooth on $\HH\setminus\Sp$
and, by Part~VII, locally integrable across $\Sp$; Theorem~\ref{e:thm:source}
gives it the same source, and Proposition~\ref{e:prop:exterior} below gives
it the same decay. Hence $G_{p;a,b}-K$ is globally $L^1_{\mathrm{loc}}$,
distributionally monogenic across $\Sp$, and tends to zero at infinity.
Applying $\overline{\Dop}$ makes every real component weakly harmonic, so
Weyl's lemma (Appendix~A) produces a smooth entire harmonic
representative. It is bounded, hence constant by harmonic Liouville, and
the constant is zero by the decay. The same three steps applied to the
difference of \emph{any} two locally integrable fields with the given
source and decay prove uniqueness; nowhere was axiality used.
\end{proof}

\begin{remark}
Theorem~\ref{e:thm:layer} is what turns the abstract period obstruction
into a concrete object: the mode $G_{p;a,b}$ that Part~IV manufactured
from a multivalued logarithm is, without any correction term, a
single-layer Cauchy--Fueter potential with an affine density. It also
explains the count: the layer has two free quaternionic parameters
because the admissible density has one constant angular mode and one mode
linear in $I$. The relation to the sphere-supported Cauchy kernels of the
literature is recorded in Part~XI; see also the sphere-average
computation of the $\HH\setminus\SSS$ counterexample in Part~V, which is
this theorem's simplest instance computed by hand.
\end{remark}

\subsection{Two conserved fluxes that recover both residues}
\label{e:sec:fluxes}

The ordinary Cauchy--Fueter flux of $g$ through a tube boundary is
conserved but blind to half of the residue.

\begin{proposition}[Total flux, and what it misses]
\label{e:prop:flux0}
Let $g$ be as in Theorem~\ref{e:thm:source}. For every sufficiently small
$\rho$,
\begin{equation}
 \Phi_0:=\int_{\partial\Tube_\rho(\Sp)}n_\rho\,g\,\dd S
       =8\pi v\,(ua+b),
 \label{e:eq:flux0}
\end{equation}
independently of $\rho$. Consequently a nonremovable singularity can have
zero total flux: take $a\ne0$ and $b=-ua$.
\end{proposition}

\begin{proof}
The componentwise divergence theorem gives
$\int_{\partial V}n\,g\,\dd S=\int_V\Dop g\,\dd V$, which vanishes on an
annular tube where $g$ is regular; hence $\Phi_0$ is
$\rho$-independent. Its small-radius limit is the total mass of the
source \eqref{e:eq:source}, and $\int_{\SSS}I\dd\sigma(I)=0$ leaves
$2v\cdot4\pi c_0=8\pi vc_0$. For the last claim, $a\ne0$ and $b=-ua$ give
$c_0=0$ but $c_1=va\ne0$, so $\Rp(a,b)\ne0$ and the germ is not
removable by Part~VII.
\end{proof}

The missing component is recovered by a \emph{weighted} flux. The weight
must be chosen with care, and this is not a cosmetic point.

\begin{theorem}[Two conserved fluxes]
\label{e:thm:twoflux}
Put
\begin{equation}
 H_u(q)=3(\Rea q-u)+\Ima q,
 \label{e:eq:Hu}
\end{equation}
and define, for small $\rho$,
\begin{equation}
 \Phi_0(\rho)=\int_{\partial\Tube_\rho}n_\rho\,g\,\dd S,
 \qquad
 \Phi_1(\rho)=\int_{\partial\Tube_\rho}H_u(q)\,n_\rho\,g\,\dd S.
 \label{e:eq:twoflux}
\end{equation}
Both are independent of $\rho$, and
\begin{equation}
 \Phi_0=8\pi v\,(ua+b),
 \qquad
 \Phi_1=-8\pi v^3a,
 \qquad\text{so}\qquad
 a=-\frac{\Phi_1}{8\pi v^3},
 \quad
 b=\frac{\Phi_0}{8\pi v}-ua.
 \label{e:eq:recover-flux}
\end{equation}
Removability of an $L^1$ axial germ along $\Sp$ is therefore equivalent
to the simultaneous vanishing of these two conserved three-dimensional
boundary integrals.
\end{theorem}

\begin{proof}
The mechanism is the Green identity for a \emph{right} Fueter-regular
weight. If $H$ is smooth with $\sum_{\mu=0}^3(\partial_{x_\mu}H)e_\mu=0$,
where $e_0=1$, $e_1=\mathbf i$, $e_2=\mathbf j$, $e_3=\mathbf k$, then on
any region $V$ where $\Dop g=0$,
\begin{equation}
 \int_{\partial V}H\,n\,g\,\dd S
 =\int_V\Bigl(\sum_\mu(\partial_{x_\mu}H)e_\mu\Bigr)g\,\dd V
   +\int_VH\,\Dop g\,\dd V=0.
 \label{e:eq:greenweight}
\end{equation}
The weight \eqref{e:eq:Hu} qualifies: writing $q-u=x_0-u+x_1\mathbf
i+x_2\mathbf j+x_3\mathbf k$ one has $H_u=3(x_0-u)+x_1\mathbf
i+x_2\mathbf j+x_3\mathbf k$, whence
\[
 \sum_\mu(\partial_{x_\mu}H_u)e_\mu
   =3\cdot1+\mathbf i\cdot\mathbf i+\mathbf j\cdot\mathbf j
     +\mathbf k\cdot\mathbf k=3-3=0 .
\]
Applying \eqref{e:eq:greenweight} to an annular tube, with $H=1$ and then
with $H=H_u$, makes both integrals in \eqref{e:eq:twoflux} independent of
$\rho$.

The values are computed by the small-radius limit of the proof of
Theorem~\ref{e:thm:source}. For $H=1$ this is
Proposition~\ref{e:prop:flux0}. For $H=H_u$ note that on the source
sphere itself $\Rea(u+Iv)=u$ and $\Ima(u+Iv)=Iv$, so
$H_u(u+Iv)=Iv$, and therefore
\[
 \Phi_1=2v\int_{\SSS}Iv\,(c_0+Ic_1)\dd\sigma(I)
      =2v^2\int_{\SSS}\bigl(Ic_0-c_1\bigr)\dd\sigma(I)
      =-8\pi v^2c_1=-8\pi v^3a,
\]
using $I^2=-1$ and $\int_{\SSS}I\dd\sigma=0$. Solving the two linear
relations gives \eqref{e:eq:recover-flux}; vanishing of both fluxes
forces $a=b=0$, which by Part~VII is removability in the $L^1$ class.
\end{proof}

\begin{remark}[The weight cannot be simplified]
\label{e:rem:naiveweight}
On $\Sp$ one has $H_u(q)=q-u$, so \eqref{e:eq:recover-flux} looks like
the first angular moment
$M_1=\int_{\Sp}(q-u)\sigma_p(q)\dd S_q=-8\pi v^3a$ of the surface
density. The two numbers agree, but the two formulas are not
interchangeable. The moment presupposes that one already knows the
density --- that is, the residue --- whereas $\Phi_1$ is a boundary
integral of $g$ itself, computable at any finite radius. Moreover the
naive weight $q-u$ does \emph{not} give a conserved flux, because it is
not right Fueter regular:
\[
 \sum_{\mu=0}^3\partial_{x_\mu}(q-u)\,e_\mu
 =1+\mathbf i\,\mathbf i+\mathbf j\,\mathbf j+\mathbf k\,\mathbf k=-2\ne0 .
\]
Only the coefficient $3$ in \eqref{e:eq:Hu} repairs this. A presentation
of the angular moment as a computable flux would therefore be incorrect.
\end{remark}

\begin{remark}[The weight is the first entire axially monogenic
polynomial]
\label{e:rem:Huispoly}
With $\mathcal P_k$ as in Lemma~\ref{e:lem:polynomial} below,
\begin{equation}
 H_u(q)=-\tfrac14\lap\bigl((q-u)^3\bigr)=3\,\mathcal P_1(q-u),
 \label{e:eq:Hu-poly}
\end{equation}
because $\lap(q^3)=-4(3\Rea q+\Ima q)$ and
$\mathcal P_1(q)=-\lap(q^3)/12$. The constant weight $1$ is
$\mathcal P_0$. Thus the two residues are recovered by pairing $g$
against the two lowest entire axially monogenic polynomials, $\mathcal
P_0$ and $\mathcal P_1$ --- precisely the two that also span the kernel
of the period map at infinity in Part~IX. This identification links the
flux calculus of this part to the polynomial calculus of the next; it is
not present in any single source manuscript.
\end{remark}

\subsection{The exact logarithmic energy law}
\label{e:sec:energylaw}

Here and throughout, \emph{energy} means the squared $L^2$ norm
$\int\abs g^2\dd V$ of the field, never a Dirichlet integral
$\int\abs{\nabla g}^2$. That distinction is what makes the critical
exponent $2$ and the divergence logarithmic. Set
\begin{equation}
 \Eg(\varepsilon,R)=\int_{\Tube_{\varepsilon,R}}\abs{g(q)}^2\dd V,
 \qquad
 \Jg(\rho)=\int_{\partial\Tube_\rho}\abs{g(q)}^2\dd S .
 \label{e:eq:energydef}
\end{equation}
The spherical mean identity of Part~VI,
$\tfrac1{4\pi}\int_{\SSS}\abs{P+IQ}^2\dd\sigma(I)=\abs P^2+\abs Q^2$,
converts both into planar integrals exactly:
\begin{equation}
 \Eg(\varepsilon,R)
 =4\pi\int_\varepsilon^R\!\!\rho\int_0^{2\pi}
    y^2\bigl(\abs P^2+\abs Q^2\bigr)\dd\theta\,\dd\rho,
 \qquad
 \Jg(\rho)=4\pi\rho\int_0^{2\pi}y^2\bigl(\abs P^2+\abs Q^2\bigr)\dd\theta .
 \label{e:eq:energyreduction}
\end{equation}

\begin{theorem}[Exact energy law with a finite renormalized constant]
\label{e:thm:energy}
Let $g$ be axially monogenic on a punctured tube about $\Sp$ and locally
integrable across $\Sp$, with residue $(a,b)$. Then the shell density has
the exact trace
\begin{equation}
 \lim_{\rho\downarrow0}\rho\,\Jg(\rho)=8\,\Rp(a,b)^2,
 \label{e:eq:shelltrace}
\end{equation}
and for every fixed sufficiently small $R$,
\begin{equation}
 \Eg(\varepsilon,R)
  =8\,\Rp(a,b)^2\log\frac R\varepsilon+\Eren_g(R)+o(1)
  \qquad(\varepsilon\downarrow0),
 \label{e:eq:energylaw}
\end{equation}
with a \emph{finite} real constant given in closed form by the convergent
integral
\begin{equation}
 \Eren_g(R)=\int_0^{R}\left(\Jg(s)-\frac{8\Rp(a,b)^2}{s}\right)\dd s .
 \label{e:eq:Eren}
\end{equation}
Subtracting the logarithm therefore leaves an actual limit, not merely a
bounded remainder. The coefficient admits three further descriptions:
\begin{align}
 8\Rp(a,b)^2
 &=32\pi^2\,\norm{h_{-2}}_{\HC}^2
 &&\text{(Laurent form, with $h_{-2}$ the coefficient of Part~III),}
 \label{e:eq:coef-laurent}\\
 8\Rp(a,b)^2
 &=\frac1{2\pi}\int_{\Sp}\abs{\sigma_p(s)}^2\dd S(s)
 &&\text{(geometric form),}
 \label{e:eq:coef-source}\\
 8\Rp(a,b)^2
 &=\frac1{8\pi^2}\left(\frac{\abs{\Phi_0}^2}{v^2}
                       +\frac{\abs{\Phi_1}^2}{v^4}\right)
 &&\text{(flux form).}
 \label{e:eq:coef-flux}
\end{align}
In particular the coefficient is strictly positive exactly when the
singularity is nonremovable.
\end{theorem}

\begin{proof}
By the $L^1$ normal form of Part~VII the hypothesis puts
$g=G_{p;a,b}+g_0$ with $g_0$ axially monogenic across $\Sp$, hence
bounded on a smaller tube; in particular the leading estimate
\eqref{e:eq:leading} is available. Writing $g_*$ for the first term of
\eqref{e:eq:leading} and using the real rotation
\[
 P_*=\frac{\cos\theta\,c_0+\sin\theta\,c_1}{\pi v\rho},
 \qquad
 Q_*=\frac{-\sin\theta\,c_0+\cos\theta\,c_1}{\pi v\rho},
\]
one gets, with no commutation hypothesis on $c_0,c_1$ because the
cancellation uses only the real inner product,
\begin{equation}
 \abs{P_*}^2+\abs{Q_*}^2
 =\frac{\abs{c_0}^2+\abs{c_1}^2}{\pi^2v^2\rho^2}
 =\frac{\Rp(a,b)^2}{\pi^2v^2\rho^2},
 \label{e:eq:leadingnorm}
\end{equation}
uniformly in $\theta$. Since $y^2=v^2+O(\rho)$, substitution into
\eqref{e:eq:energyreduction} gives
\begin{equation}
 \Jg(\rho)=\frac{8\Rp(a,b)^2}{\rho}
   +O\bigl(1+\abs{\log\rho}+\rho(1+\abs{\log\rho})^2\bigr),
 \label{e:eq:Jerror}
\end{equation}
where the three error terms come, in order, from $y^2-v^2$, from the
cross term $2\abs{g_*}\abs{g-g_*}=O(\rho^{-1}(1+\abs{\log\rho}))$, and
from $\abs{g-g_*}^2=O((1+\abs{\log\rho})^2)$, each multiplied by the
normal area factor $\rho$. This proves \eqref{e:eq:shelltrace}. The error
in \eqref{e:eq:Jerror} is absolutely integrable at $\rho=0$; integrating
in $\rho$ therefore yields \eqref{e:eq:energylaw} with
\eqref{e:eq:Eren}, the integral converging precisely because the
subtracted singular term matches the trace.

For \eqref{e:eq:coef-laurent}, Part~III gives $pa+b=-2\pi\iota h_{-2}$,
so $\Rp(a,b)^2=\norm{pa+b}_{\HC}^2=4\pi^2\norm{h_{-2}}_{\HC}^2$. For
\eqref{e:eq:coef-source}, $\sigma_p(u+Iv)=2(c_0+Ic_1)/v$ and
$\dd S=v^2\dd\sigma(I)$, so
$\int_{\Sp}\abs{\sigma_p}^2\dd S=4\int_{\SSS}\abs{c_0+Ic_1}^2\dd\sigma
=16\pi(\abs{c_0}^2+\abs{c_1}^2)=16\pi\Rp^2$, using the spherical mean
identity again. For \eqref{e:eq:coef-flux}, substitute
$\Phi_0=8\pi vc_0$ and $\Phi_1=-8\pi v^2c_1$ from
Theorem~\ref{e:thm:twoflux}:
$\abs{\Phi_0}^2/v^2+\abs{\Phi_1}^2/v^4=64\pi^2(\abs{c_0}^2+\abs{c_1}^2)$.
Positivity is $v>0$ in \eqref{e:eq:residue-norm}.
\end{proof}

\begin{remark}[Hypothesis, and what is weaker than what]
\label{e:rem:energy-hyp}
Theorem~\ref{e:thm:energy} is stated under local integrability across
$\Sp$, which is strictly weaker as a hypothesis than the pointwise
condition $\Max_g(\rho)=O(\rho^{-1})$ under which several of the source
manuscripts stated the law: $O(\rho^{-1})$ implies $L^1$ because
$\dd V\sim\rho\,\dd\rho$, but not conversely \emph{a priori}. That the
two classes coincide \emph{a posteriori} is a consequence of the $L^1$
normal form of Part~VII, not an available hypothesis. Separately, the
assertion that the remainder converges --- rather than being merely
$O(1)$ --- is genuinely stronger than a bounded-remainder statement, and
is what makes \eqref{e:eq:Eren} meaningful as a number attached to $g$.
We do not assert that $\Eren_g(R)$ has a closed form in $(a,b,p,R)$ for a
general $g$; \eqref{e:eq:Eren} is an explicit convergent integral, not an
evaluation.
\end{remark}

\begin{remark}[Numerical corroboration]
Two of the source manuscripts evaluated the coefficient
$8\Rp(a,b)^2$ against an independent quadrature of
$\int\abs g^2\dd V$ rather than against its own formula. Those
certificates, and three further ones, are reconciled in the single table
of Part~X; all five reproduce from \eqref{e:eq:residue-norm}.
\end{remark}

\subsection{A growth-free lower bound with the sharp leading constant}
\label{e:sec:lowerbound}

The energy law above assumes enough regularity to have a normal form. The
following bound assumes nothing at all beyond smoothness off the sphere,
and still sees the exact constant.

\begin{theorem}[Sharp lower bound, no growth hypothesis]
\label{e:thm:lowerenergy}
Let $g$ be any smooth axially monogenic field on the punctured tube, with
period pair $(a,b)$; no growth bound and no integrability hypothesis is
assumed. For $0<\varepsilon<R<v$,
\begin{equation}
 \Eg(\varepsilon,R)\ \ge\
 8\int_\varepsilon^R\frac{(v-\rho)^2}{\rho}
 \left(\abs a^2+\frac{\abs{ua+b}^2}{v^2+\rho^2}\right)\dd\rho .
 \label{e:eq:lowerenergy}
\end{equation}
Consequently
\begin{equation}
 \liminf_{\varepsilon\downarrow0}
 \frac{\Eg(\varepsilon,R)}{\log(R/\varepsilon)}\ \ge\ 8\,\Rp(a,b)^2,
 \label{e:eq:liminf}
\end{equation}
and the constant $8\Rp(a,b)^2$ is sharp for every residue pair. In
particular, if $\Eg(0,R)<\infty$ for one $R$, then $a=b=0$, with no
pointwise hypothesis whatever.
\end{theorem}

\begin{proof}
Parameterize the transverse circle of radius $\rho$ by
$x=u+\rho\cos\varphi$, $y=v+\rho\sin\varphi$ and substitute
$\dd x=-\rho\sin\varphi\dd\varphi$, $\dd y=\rho\cos\varphi\dd\varphi$
into the two period integrals of Part~II. The first form gives
directly
\begin{equation}
 a=\frac\rho2\int_0^{2\pi}\bigl(P\sin\varphi+Q\cos\varphi\bigr)\dd\varphi .
 \label{e:eq:circlea}
\end{equation}
For the second, form $\oint(\thet_g+u\,\om_g)$, whose value is $ua+b$;
expanding the integrand the terms linear in $\rho P$ cancel and the
coefficient of $P$ becomes $v\cos\varphi$, that of $Q$ becomes
$-v\sin\varphi-\rho$:
\begin{equation}
 ua+b=\frac\rho2\int_0^{2\pi}
   \bigl(vP\cos\varphi-(v\sin\varphi+\rho)Q\bigr)\dd\varphi .
 \label{e:eq:circlec}
\end{equation}
Regard $(P,Q)$ as an element of $L^2([0,2\pi];\RR^2)$ in each of the four
real quaternionic components. The two coefficient vectors
\[
 f_1(\varphi)=(\sin\varphi,\cos\varphi),
 \qquad
 f_2(\varphi)=(v\cos\varphi,\,-v\sin\varphi-\rho)
\]
are \emph{exactly orthogonal}:
$\int_0^{2\pi}(v\sin\varphi\cos\varphi-v\sin\varphi\cos\varphi
-\rho\cos\varphi)\dd\varphi=0$. Their squared norms are $2\pi$ and
$2\pi(v^2+\rho^2)$, so the functionals in
\eqref{e:eq:circlea}--\eqref{e:eq:circlec} have squared norms
$\pi\rho^2/2$ and $\pi\rho^2(v^2+\rho^2)/2$. Bessel's inequality applied
componentwise and summed over the four real components gives
\begin{equation}
 \int_0^{2\pi}\bigl(\abs P^2+\abs Q^2\bigr)\dd\varphi
 \ \ge\ \frac2{\pi\rho^2}
   \left(\abs a^2+\frac{\abs{ua+b}^2}{v^2+\rho^2}\right).
 \label{e:eq:bessel}
\end{equation}
Insert this into \eqref{e:eq:energyreduction} with $y^2\ge(v-\rho)^2$ and
integrate in $\rho$; the factor $4\pi\rho\cdot2/(\pi\rho^2)=8/\rho$
produces \eqref{e:eq:lowerenergy}. As $\rho\downarrow0$ the bracket in
\eqref{e:eq:lowerenergy} tends to $\abs a^2+\abs{ua+b}^2/v^2$ and the
prefactor to $v^2$, so the radial integrand differs from
$8\Rp(a,b)^2/\rho$ by $O(1)$, which is integrable; this gives
\eqref{e:eq:liminf}. Sharpness is Theorem~\ref{e:thm:energy} applied to
$G_{p;a,b}$ itself, whose proof uses only the explicit mode formulas and
not the present bound, so there is no circularity. Finally, if the energy
were finite then \eqref{e:eq:liminf} would force $\Rp(a,b)=0$, hence
$a=b=0$ because $v>0$.
\end{proof}

\begin{corollary}[An explicit finite-radius form]
\label{e:cor:explicitlower}
With $W=\Rp(a,b)^2$ and $0<\varepsilon<R<v$,
\begin{equation}
 \Eg(\varepsilon,R)\ \ge\
 8W\int_\varepsilon^R\left(\frac{v-\rho}{v+\rho}\right)^2\frac{\dd\rho}\rho
 =8W\left[\log\frac R\varepsilon+\frac{4v}{v+R}-\frac{4v}{v+\varepsilon}
  \right].
 \label{e:eq:explicitlower}
\end{equation}
No optimality is claimed for the finite-radius correction in
\eqref{e:eq:explicitlower}; only its leading logarithmic coefficient is
sharp.
\end{corollary}

\begin{proof}
Since $(v+\rho)^2>v^2+\rho^2$ for $v,\rho>0$, one has
$v^2/(v+\rho)^2<1$ and $1/(v+\rho)^2<1/(v^2+\rho^2)$, so the integrand of
\eqref{e:eq:explicitlower} is pointwise dominated by that of
\eqref{e:eq:lowerenergy}; the corollary is therefore weaker than the
theorem at every radius. The closed form follows from the elementary
identity
\[
 \frac1\rho\left(\frac{v-\rho}{v+\rho}\right)^2
 =\frac1\rho-\frac{4v}{(v+\rho)^2},
\]
integrated between $\varepsilon$ and $R$.
\end{proof}

\begin{remark}[Two routes to the same liminf, and one that is not sharp]
Inequality \eqref{e:eq:lowerenergy} is obtained from an exact orthogonal
decomposition and Bessel; \eqref{e:eq:explicitlower} is obtained by
bounding the operator norm of the $2\times2$ matrix relating
$(va,ua+b)$ to $(yP,yQ)$ by $(v+\rho)/(v-\rho)$ and applying
Cauchy--Schwarz, which discards the orthogonality. Both give the sharp
liminf; the first dominates the second at every radius. A third route
--- bounding the two periods separately by Cauchy--Schwarz on the circle
and taking the maximum --- yields a shell bound with leading constant
$8v^2\max\{\abs a^2,\abs b^2/\abs p^2\}$, which is strictly smaller than
$8\Rp(a,b)^2$ whenever $a$ and $b$ are both nonzero (with $u=0$, $v=1$,
$a=b=1$ it gives $8$ against the sharp $16$). That route is not adopted
here.
\end{remark}

\subsection{The energy scale of a higher principal part}
\label{e:sec:poleenergy}

Above the critical order the divergence is a power, and its exponent and
coefficient are completely determined.

\begin{proposition}[Leading $L^2$ scale of a rational principal part]
\label{e:prop:poleenergy}
Let $g$ be axially monogenic on the punctured tube with the normal form
of Part~VII, and suppose $k\ge1$ is the largest index with $c_k\ne0$ in
its principal part $\lap\I\bigl(\sum_{k\ge1}(z-p)^{-k}c_k\bigr)$. Then
\begin{equation}
 \Eg(\varepsilon,R)
 =16\pi^2k\,\norm{c_k}_{\HC}^2\,\varepsilon^{-2k}
      +o\bigl(\varepsilon^{-2k}\bigr).
 \label{e:eq:poleenergy}
\end{equation}
Thus logarithmic defects and rational principal parts have completely
specified and mutually distinct leading $L^2$ scales: $\log(1/\varepsilon)$
for the former and $\varepsilon^{-2k}$ for the latter.
\end{proposition}

\begin{proof}
For a holomorphic stem $F=A+\iota B$ the complexified axial pair of its
Fueter image is
\begin{equation}
 P+\iota Q=\frac{2\iota}{y}F'(z)-\frac{2\iota}{y^2}B(z),
 \label{e:eq:complexified-fueter}
\end{equation}
which is the identity of Part~II. The top pole of the stem is
$(z-p)^{-k}c_k$, whose derivative is $-k(z-p)^{-k-1}c_k$; the term
$2\iota B/y^2$ and every lower pole are of strictly lower order, and the
logarithmic sector contributes $O(\rho^{-1}\log\rho)$. Since
$\norm{P+\iota Q}_{\HC}^2=\abs P^2+\abs Q^2$ and $y=v+O(\rho)$,
\[
 \abs P^2+\abs Q^2
 =\frac{4k^2}{v^2}\norm{c_k}_{\HC}^2\,\rho^{-2k-2}\bigl(1+O(\rho)\bigr)
\]
uniformly in the angle. Substituting into \eqref{e:eq:energyreduction}
with $y^2=v^2+O(\rho)$ gives
\[
 32\pi^2k^2\norm{c_k}^2\int_\varepsilon^R\rho^{-2k-1}\dd\rho
 =16\pi^2k\norm{c_k}^2\varepsilon^{-2k}+O(1),
\]
and the lower-order contributions are $o(\varepsilon^{-2k})$ by splitting
the radial integral at a fixed small radius and then letting that radius
tend to zero.
\end{proof}

\begin{remark}[Why the meromorphic sector alone already excludes finite
energy]
\label{e:rem:poleenergy2}
Proposition~\ref{e:prop:poleenergy} is quantitative; the following
qualitative statement isolates the same mechanism for a single mode and
is worth recording because it separates the meromorphic sector from the
logarithmic one. If $F=(z-p)^{-m}c$ with $c\in\HC$ and $m\ge1$, then by
\eqref{e:eq:complexified-fueter} and the spherical average
$\frac1{4\pi}\int_{\SSS}\bigl|\I(e^{\iota t}c)(I)\bigr|^2\dd\sigma(I)
=\norm c_{\HC}^2$, the shell integral of $\abs{\lap\I(F)}^2$ is a
positive multiple of $\rho^{-2m-1}$, so the local energy is a positive
multiple of $\int_0\rho^{-2m-1}\dd\rho=\infty$. Thus no nonzero Laurent
mode has finite local $L^2$ energy. This is an alternative to the cutoff
proof of $L^2$ removability in Part~VII; it is not needed for that
theorem, but it is the only argument in the family that quantifies the
exclusion of the meromorphic part separately from the logarithmic part.
\end{remark}

\subsection{Invariance, and an audit of the constant}
\label{e:sec:audit}

\begin{remark}[Coordinate invariance of the coefficient]
\label{e:rem:translation}
Under a real translation $q'=q-t$, $t\in\RR$, the affine increment
$qa+b$ becomes $q'a+(b+ta)$ and the centre becomes $u'=u-t$. Hence the
slope residue is unchanged and $u'a+b'=ua+b$, so both
$\Rp(a,b)^2$ and the two fluxes \eqref{e:eq:recover-flux} are unchanged.
The invariant pair is $(a,\,ua+b)$; neither $a$ nor $b$ separately is a
translation-invariant scalar, and $\Rp$ is not an artifact of the choice
of real origin. Equation \eqref{e:eq:coef-source} gives the coefficient
an intrinsic geometric meaning without choosing period coordinates at
all.
\end{remark}

\begin{remark}[Where the constant $8$ comes from]
\label{e:rem:audit8}
It is worth isolating the arithmetic, because the same letter $8$ appears
in the cokernel count $8m$ of Part~IV for an entirely unrelated reason.
In \eqref{e:eq:energyreduction} the factor $4\pi$ is the total mass of
$\sigma$ on $\SSS$, the angular integration over $\theta$ contributes
$2\pi$, and the leading squared norm \eqref{e:eq:leadingnorm} carries
$1/\pi^2$ after the factors $v^2$ from $y^2$ and $1/v^2$ from
\eqref{e:eq:leading} cancel. Thus
\[
 8=\frac{(4\pi)\times(2\pi)}{\pi^2},
\]
a constant of the codimension-two normal-circle geometry. The $2\pi^2$
of the point Cauchy--Fueter kernel $E$ belongs to the three-sphere
geometry of an isolated point in $\RR^4$ and does not appear here. The
same caution applies to the critical exponents: the threshold exponent
$1$ at a nonreal sphere reflects a singular set of real codimension two
and must not be confused with the isolated-point exponent $3$ of Part~I.
Further sign and normalization conventions --- $\Dop$ with no factor
$\tfrac12$, $\lap q^2=-4$, coefficients on the right, normal circle
length $2\pi\rho$, sphere area $4\pi v^2$, tube boundary element
$y^2\rho\,\dd\theta\,\dd\sigma$, and the derived constants $2/v$,
$8\pi v$, $8$, $16\pi^2k$ --- are collected in Appendix~C.
\end{remark}


% =====================================================================
% Part IX.  Global classification: prescribed data, finite principal
%           parts, and the far field.
% Lead sources: 01, 04, 05, 06.
% =====================================================================

\section{Global classification: prescribed data, realization, and decay}
\label{e:sec:global}

The local data of Part~VII --- an affine residue pair and a finite or
infinite Laurent principal part at each singular sphere --- are now
assembled into global statements. Three questions are answered: which
global configurations of local data actually occur (a Mittag--Leffler
realization theorem), what the resulting global function space looks like
when the transverse orders and the growth at infinity are prescribed (a
normal form with an exact dimension), and how the local residues are
visible from infinity (a moment hierarchy). The entire axially monogenic
polynomials appear at both ends of the story, as the growth sector of the
normal form and as the flux weights of Part~VIII.

\subsection{Entire axially monogenic polynomials}
\label{e:sec:polynomials}

\begin{lemma}[Polynomial growth and the normalized basis]
\label{e:lem:polynomial}
Define
\begin{equation}
 \mathcal P_k(q)=-\frac{\lap(q^{k+2})}{2(k+1)(k+2)},\qquad k\ge0 .
 \label{e:eq:Pk}
\end{equation}
Then $\mathcal P_k$ is entire axially monogenic with real-axis trace
$\mathcal P_k(x)=x^k$. If $h$ is entire axially monogenic and
$\abs{h(q)}\le C(1+\abs q)^\alpha$ for some $\alpha\ge0$, then $h$ is a
polynomial in the four real coordinates of degree at most
$N=\lfloor\alpha\rfloor$, and there are unique $c_0,\dots,c_N\in\HH$ with
\begin{equation}
 h(q)=\sum_{k=0}^N\mathcal P_k(q)\,c_k .
 \label{e:eq:polybasis}
\end{equation}
That space has real dimension $4(N+1)$.
\end{lemma}

\begin{proof}
Since $\overline{\Dop}\Dop=\lap$, every real component of $h$ is entire
harmonic. The interior derivative estimate for a harmonic $u$ on a ball,
$\abs{\partial^\nu u(q_0)}\le C_\nu R^{-\abs\nu}\sup_{B(q_0,R)}\abs u$
--- obtained by writing $u$ as its convolution with a smooth radial
averaging kernel supported in the ball and transferring the derivatives
to the kernel, so that $C_\nu$ depends on neither $u$ nor $R$ --- gives,
on letting $R\to\infty$, that every derivative of order $\abs\nu>\alpha$
vanishes. Real analyticity makes $h$ a polynomial of degree at most $N$.

For the basis, let $F$ be an entire holomorphic stem with
$\lap\I(F)=h$; one exists because the upper half-plane is simply
connected and the period criterion of Part~II is vacuous there, and it
extends across the real axis by the reflection theorem of Part~V. The
real-trace identity
\begin{equation}
 \bigl(\lap\I(F)\bigr)(x)=-2F''(x),\qquad x\in\RR,
 \label{e:eq:realtrace}
\end{equation}
follows from Taylor expansion at a real point,
$A(x,y)=F(x)-y^2F''(x)/2+O(y^4)$ and
$B(x,y)=yF'(x)-y^3F'''(x)/6+O(y^5)$, inserted into the Fueter formula
for $\lap$ on stems. Writing the real-axis polynomial as
$h(x)=\sum_{k\le N}x^kc_k$ gives $F''(z)=-\tfrac12\sum_kz^kc_k$ by the
holomorphic identity theorem applied to the four central-complex
coordinates; integrating twice and applying $\lap\I$ proves
\eqref{e:eq:polybasis}. The same computation gives
$\lap(q^{k+2})\big|_{\RR}=-2(k+2)(k+1)x^k$, hence
$\mathcal P_k(x)=x^k$, and the linear independence of the real powers
over $\HH$ gives uniqueness and the dimension.
\end{proof}

The Liouville-type conclusion is standard harmonic analysis; the
normalized family \eqref{e:eq:Pk} is recorded to make the classification
below explicit, and is not offered as a new family of polynomials. Its
two lowest members are the two flux weights of Part~VIII:
$\mathcal P_0=1$ and $3\mathcal P_1(q-u)=H_u(q)$ by
\eqref{e:eq:Hu-poly}.

\subsection{Exterior expansions of the modes}
\label{e:sec:exterior}

On $\HH$ itself the relevant modes are the reflected ones
$G^{\mathrm{sym}}_{p;a,b}$ of Part~V, built from the symmetric logarithm
$\ell^{\mathrm s}_p(z)=\frac1{2\pi\iota}\log\frac{z-p}{z-\bar p}$, which
is smooth across the real axis; near $\Sp$ it differs from
$G_{p;a,b}$ by a field regular there, so it carries the same affine
residue $(a,b)$. Likewise the reflection-compatible rational stems
\begin{equation}
 R_{p,k,c}(z)=\frac{c}{(z-p)^k}+\frac{\bar c}{(z-\bar p)^k},
 \qquad c\in\HC,\ k\ge1,
 \label{e:eq:pairedrational}
\end{equation}
(the bar being complex-stem conjugation only) satisfy
$R_{p,k,c}(\bar z)=\overline{R_{p,k,c}(z)}$, so they are genuine stems on
$\CC\setminus\{p,\bar p\}$, with upper principal part exactly
$(z-p)^{-k}c$; the reflected summand is holomorphic near $p$.

\begin{proposition}[Exterior expansion and uniform decay]
\label{e:prop:exterior}
For $\abs q>\abs p$,
\begin{equation}
 G^{\mathrm{sym}}_{p;a,b}(q)=\sum_{n\ge1}\lap(q^{-n})\,d_n(p;a,b),
 \qquad
 d_n(p;a,b)=-\frac1\pi\left(
   \frac{\Ima_{\CC}(p^{n+1})}{n+1}\,a
  +\frac{\Ima_{\CC}(p^{n})}{n}\,b\right),
 \label{e:eq:exterior}
\end{equation}
the coefficients $\Ima_{\CC}(p^n)$ being ordinary real scalars. The
series and all its derivatives converge uniformly on compact subsets of
$\{\abs q>\abs p\}$. Each of $G^{\mathrm{sym}}_{p;a,b}$ and
$\lap\I(R_{p,k,c})$ is $O(\abs q^{-3})$ at infinity, uniformly in the
quaternionic direction, including directions arbitrarily close to the
real axis where the axial formulas have apparent denominators.
\end{proposition}

\begin{proof}
Near infinity choose the branch of the symmetric logarithm vanishing
there:
\[
 \ell^{\mathrm s}_p(z)
 =\frac1{2\pi\iota}\bigl(\log(1-p/z)-\log(1-\bar p/z)\bigr)
 =-\sum_{n\ge1}\frac{p^n-\bar p^{\,n}}{2\pi\iota n}z^{-n}
 =-\sum_{n\ge1}\frac{\Ima_{\CC}(p^n)}{\pi n}z^{-n},
\]
whose coefficients are real. Multiplying by $za+b$ produces the constant
term $-\Ima_{\CC}(p)a/\pi=-va/\pi$, which is real-affine and is killed by
$\lap$, followed by the coefficients \eqref{e:eq:exterior}. The
reflection-compatible rational stem likewise has a convergent exterior
Laurent series $\sum_{n\ge1}z^{-n}a_n$ with $a_n\in\HH$.

For the uniform decay, apply $\lap\I$ term by term. Each
$\lap(q^{-n})$ is homogeneous of degree $-n-2$ and smooth on the unit
three-sphere. Concretely, for $q\ne0$ the first and second real
derivatives of $q^{-1}$ are bounded by $C\abs q^{-2}$ and
$C\abs q^{-3}$, so the product rule bounds second derivatives of
$q^{-n}$ by $Cn^2\abs q^{-n-2}$ uniformly on spheres; the Laurent
coefficients have geometric bounds, so the differentiated series
converges uniformly outside any larger ball. The constant term has zero
Laplacian, so the first surviving term is $O(\abs q^{-3})$. Homogeneity,
not the axial formula, is what makes the estimate uniform near the real
axis.
\end{proof}

We record the identification that anchors the leading term:
$\lap(q^{-1})=-4\bar q/\abs q^4=-8\pi^2E(q)$, the kernel link of
Part~I.

\subsection{A Mittag--Leffler realization theorem}
\label{e:sec:ML}

Parts~IV--VII compute what can occur and what the obstruction is. The
following converse --- the only one in the family --- says that every
formally admissible global singular configuration is actually realized,
with no compatibility conditions between spheres.

\begin{theorem}[Realization of locally finite spherical data]
\label{e:thm:ML}
Let $V\subset\Cplus$ be a \emph{simply connected} open set and let
$Z=\{p_j\}_{j\in A}\subset V$ be discrete and locally finite in $V$. At
each $p_j$ prescribe, independently and with no relation to the data at
any other point,
\[
 a_j,b_j\in\HH,
 \qquad
 R_j(z)=\sum_{m=1}^{M_j}(z-p_j)^{-m}c_{jm},\quad c_{jm}\in\HC,
\]
with $M_j\ge0$ finite and the empty sum allowed. Then there exists
$g\in\AM(\Om_{V\setminus Z})$ whose local normal form at every
$S_{p_j}$ has exactly the affine residue $(a_j,b_j)$ and exactly the
Laurent principal part $R_j$, and which has no other singularity in
$\Om_V$. Two such functions differ by a field in $\AM(\Om_V)$, and $g$
admits a single-valued slice-regular Fueter primitive on
$\Om_{V\setminus Z}$ if and only if $a_j=b_j=0$ for every $j$.
\end{theorem}

\begin{proof}
The construction prescribes the \emph{holomorphic invariant}, not the
field. At $p_j$ prescribe the meromorphic principal part
\begin{equation}
 \Hg^{(j)}(z)=
 \sum_{m=1}^{M_j}m(m+1)(z-p_j)^{-m-2}c_{jm}
 +\frac{a_j}{2\pi\iota(z-p_j)}
 -\frac{p_ja_j+b_j}{2\pi\iota(z-p_j)^2}.
 \label{e:eq:MLdata}
\end{equation}
Applied to the four central-complex coordinates, the classical scalar
Mittag--Leffler theorem produces a meromorphic $H:V\to\HC$ with precisely
these principal parts and no other poles; no absolute convergence of the
uncorrected sum is required, the theorem supplying the holomorphic
corrections.

For any closed curve $\gamma$ in $V\setminus Z$ the residue theorem ---
legitimate because $V$ is simply connected and $Z$ has no accumulation
point in $V$ --- gives
\begin{equation}
 \int_\gamma H(z)\dd z=\sum_j\Ind(\gamma,p_j)\,a_j,
 \qquad
 -\int_\gamma zH(z)\dd z=\sum_j\Ind(\gamma,p_j)\,b_j,
 \label{e:eq:MLperiods}
\end{equation}
only finitely many terms being nonzero for a fixed $\gamma$. Indeed
$R_j''$ and $zR_j''$ have no simple pole, the simple-pole coefficient of
$\Hg^{(j)}$ is $a_j/(2\pi\iota)$, and that of $z\Hg^{(j)}$ is
$(p_ja_j-(p_ja_j+b_j))/(2\pi\iota)=-b_j/(2\pi\iota)$. Both integrals in
\eqref{e:eq:MLperiods} therefore lie in $\HH$, so $H$ belongs to the
moment-real class $\Omom(V\setminus Z)$ of Part~III. By the surjectivity
half of Part~III's range theorem, $H=\Hg$ for some
$g\in\AM(\Om_{V\setminus Z})$: integrate $H\dd z$ and $-zH\dd z$ on the
universal cover to obtain $\pot$ and $\potV$, set $F=z\pot+\potV$ and
descend $\lap\I(F)$, which is single-valued because the monodromy of $F$
lies in the affine kernel $\Aff$. Its affine periods are exactly
\eqref{e:eq:MLperiods}.

It remains to check the \emph{complete} local data, not merely the
periods. Near $p_j$ consider the multivalued holomorphic expression
$F_j=R_j+\frac{za_j+b_j}{2\pi\iota}\log(z-p_j)$. A direct
differentiation shows $F_j''=\Hg^{(j)}$, so $H-F_j''$ extends
holomorphically to a disk about $p_j$; choose a holomorphic double
primitive $K_j$ of that difference on the disk. On a branch over the
punctured disk $(F-F_j-K_j)''=0$, so the difference is complex-affine,
and its image under $\lap\I$ is regular on the full disk because the
disk stays away from the real axis. Hence
$g-\lap\I(R_j)-G_{p_j;a_j,b_j}$ extends across $S_{p_j}$, which is
exactly the assertion about the normal form. Away from $Z$ the
construction is locally $\lap\I$ of a holomorphic stem, so there are no
extra singularities.

If two fields have the same data, the uniqueness in Part~VII's normal
form makes their difference removable at every $p_j$; local finiteness
lets these local extensions be combined into a member of $\AM(\Om_V)$.
Finally, if every pair vanishes then both integrals in
\eqref{e:eq:MLperiods} vanish on every loop, so Part~II gives a global
primitive; conversely a small loop about $p_j$ detects its pair, which
must then vanish.
\end{proof}

\begin{remark}[Why the second derivative is the right thing to prescribe]
\label{e:rem:whysecond}
The formal sum
$\sum_j\bigl(R_j+\frac{za_j+b_j}{2\pi\iota}\log(z-p_j)\bigr)$ need not
converge, and it cannot be repaired by silently choosing global branches
of the logarithms. Prescribing \eqref{e:eq:MLdata} avoids both problems
at once: the polar data are single-valued meromorphic data, classical
Mittag--Leffler supplies the convergence-forcing holomorphic
corrections, and the two residue moments guarantee that the surviving
monodromy of the primitive lies in exactly the real-affine kernel
$\Aff$. This is why arbitrary residue pairs at arbitrarily many spheres
can coexist without producing a multivalued target. Simple connectivity
of $V$ is used twice, for the residue theorem and to exclude periods of
$H$ coming from the topology of $V$ itself; over a base with holes the
realization statement would have to be combined with the split range
theorem of Part~IV.
\end{remark}

\begin{remark}[The prescribed data automatically satisfy the reality
constraint]
\label{e:rem:MLreality}
Part~III proves that the two polar coefficients of $\Hg$ at a nonreal
$p$ are not independent: $h_{-1}=(\iota/v)\ReC(h_{-2})$. The data
\eqref{e:eq:MLdata} comply automatically, and it is worth seeing why,
since that is what makes the prescription consistent. Here
$h_{-1}=a/(2\pi\iota)=-\iota a/(2\pi)$ and
$h_{-2}=-(pa+b)/(2\pi\iota)=\iota(pa+b)/(2\pi)
=\bigl(-va+\iota(ua+b)\bigr)/(2\pi)$, so
$\ReC(h_{-2})=-va/(2\pi)$ and $(\iota/v)\ReC(h_{-2})=-\iota a/(2\pi)
=h_{-1}$. Equivalently, with $h_{-2}=c+\iota d$ one recovers
$a=-(2\pi/v)c$ and $b=2\pi d+(2\pi u/v)c$. The free datum at each sphere
is therefore the single $\HC$ coefficient $h_{-2}$, which carries all
eight real obstruction coordinates, together with the higher coefficients
$c_{jm}$; there is no independent choice of $h_{-1}$ to be made, and no
pole of order exactly one can be prescribed.
\end{remark}

\begin{example}[Independent defects and poles at three spheres]
\label{e:ex:several}
Over a simply connected upper base containing distinct $p_1,p_2,p_3$,
prescribe the first-period defect $(\mathbf i,0)$ at $p_1$, the
second-period defect $(0,\mathbf j)$ at $p_2$, and the zero-period
principal part $(z-p_3)^{-2}(1+\iota\mathbf k)$ at $p_3$.
Theorem~\ref{e:thm:ML} produces a single axially monogenic target with
exactly these classes. By Part~VII the first two spheres carry exact
transverse order $\rho^{-1}$ and the third exact order $\rho^{-3}$;
after subtracting $G_{p_1;\mathbf i,0}+G_{p_2;0,\mathbf j}$ the remainder
has zero periods and hence a global primitive on the punctured base,
although it is still singular at $S_{p_3}$. This is the sharpest
available illustration that global invertibility on a punctured domain is
not removability through the missing sphere.
\end{example}

\subsection{The global normal form with mixed transverse orders}
\label{e:sec:globalnormal}

We now fix finitely many spheres in $\HH$ and prescribe both the
transverse orders and the growth at infinity.

\begin{theorem}[Rational--logarithmic global classification]
\label{e:thm:globalnormal}
Let $p_1,\dots,p_n\in\Cplus$ be distinct, $S_j=u_j+v_j\SSS$, and choose
integers $m_j\ge1$ and $d\ge0$. Let $\mathcal M(\boldsymbol m,d)$ be the
real vector space of smooth axially monogenic fields on
$\HH\setminus\bigcup_jS_j$ with transverse growth
$\Max_g(\rho_j)=O(\rho_j^{-m_j})$ near each $S_j$ and
$\abs{g(q)}=O(\abs q^{\,d})$ at infinity. Every
$g\in\mathcal M(\boldsymbol m,d)$ has a unique representation
\begin{equation}
 g(q)=\sum_{k=0}^{d}\mathcal P_k(q)\,c_k
   +\sum_{j=1}^{n}G^{\mathrm{sym}}_{p_j;a_j,b_j}(q)
   +\sum_{j=1}^{n}\sum_{i=1}^{m_j-1}
        \lap\I\bigl(R_{p_j,i,c_{j,i}}\bigr)(q),
 \label{e:eq:globalnormal}
\end{equation}
with $c_k,a_j,b_j\in\HH$ and $c_{j,i}\in\HC$, and all choices of these
coefficients occur. Consequently
\begin{equation}
 \dim_{\RR}\mathcal M(\boldsymbol m,d)=4(d+1)+8\sum_{j=1}^{n}m_j .
 \label{e:eq:globaldim}
\end{equation}
A field in this class has a global slice-regular Fueter primitive on
$\HH\setminus\bigcup_jS_j$ exactly when every $a_j$ and every $b_j$
vanishes.
\end{theorem}

\begin{proof}
\emph{Remove the local singular classes.} Apply Part~VII's normal form
near each $p_j$; it yields unique $(a_j,b_j)$ and unique
$c_{j,1},\dots,c_{j,m_j-1}$. Near $p_j$ the symmetric logarithm differs
from $\ell_{p_j}$ by a function holomorphic up to a harmless real branch
constant, so $G^{\mathrm{sym}}_{p_j;a_j,b_j}$ has the same local singular
class as $G_{p_j;a_j,b_j}$; and $R_{p_j,i,c}$ has exactly the required
upper principal part by \eqref{e:eq:pairedrational}. The terms attached
to the other spheres are regular near $S_j$. Subtracting every singular
term in \eqref{e:eq:globalnormal} leaves a field $g_0$ that extends
axially and monogenically across every sphere, hence to all of $\HH$.

\emph{Use the growth at infinity.} By
Proposition~\ref{e:prop:exterior} the subtracted terms are
$O(\abs q^{-3})$, so $g_0=O(\abs q^{\,d})$. Lemma~\ref{e:lem:polynomial}
makes $g_0=\sum_{k\le d}\mathcal P_k c_k$ with unique $c_k$.

\emph{Uniqueness, converse, and count.} The local theorem determines
$(a_j,b_j)$ and all $c_{j,i}$ uniquely; the polynomial coefficients are
unique because $\lap(q^{k+2})\big|_{\RR}=-2(k+2)(k+1)x^k$ and the real
powers are independent over $\HH$. Conversely every term in
\eqref{e:eq:globalnormal} has the stated regularity and local growth, and
Proposition~\ref{e:prop:exterior} supplies the exterior bound. Counting:
$4(d+1)$ from the $c_k\in\HH$, $8$ from each pair
$(a_j,b_j)\in\HH^2$, and $8(m_j-1)$ from the $c_{j,i}\in\HC$, giving
\eqref{e:eq:globaldim}. Finally the only terms with nonzero periods are
the logarithmic ones, so the reflection form of the period criterion
(Part~V) gives the last assertion.
\end{proof}

\begin{remark}[Which hypothesis each statement needs]
\label{e:rem:globalhyp}
Theorem~\ref{e:thm:globalnormal} is stated under the pointwise
transverse hypothesis $\Max_g(\rho_j)=O(\rho_j^{-m_j})$, and that is
genuinely what it needs: for $m_j\ge2$ the higher Laurent coefficients
are extracted by a growth argument and no integrability hypothesis
replaces it. At the bottom of the filtration, $m_j=1$, the hypothesis can
be weakened to local integrability across $S_j$, which is strictly
weaker; that is the content of the next corollary, and it is the only
place in this part where the $L^1$ theorem of Part~VII is used in
preference to a growth bound. We do not assert that the mixed-order
theorem holds under an $L^1$-type hypothesis at the higher-order spheres.
Note also that the two norms $\Max$ and $\Msph$ define the same $O(\cdot)$
classes here, since $\axn g\le\Msph_g\le\sqrt2\,\axn g$; the constants in
the exact asymptotics of Part~VII do differ between them, and are not
used in this section.
\end{remark}

\subsection{The decaying class, the layer inverse, and additivity}
\label{e:sec:decaying}

\begin{corollary}[Decaying fields with critical spherical growth]
\label{e:cor:decaying}
Let $\Sigma=\bigcup_{j=1}^{n}S_{p_j}$ with $p_j=u_j+\iota v_j$ distinct,
and let $\mathcal M_{1,0}(\Sigma)$ be the right quaternionic vector space
of smooth axially monogenic fields on $\HH\setminus\Sigma$ that are
\emph{locally integrable across every sphere} and satisfy
$\abs{g(q)}\to0$ as $\abs q\to\infty$. Then the residue map is an
isomorphism
\begin{equation}
 \Resaff:\ \mathcal M_{1,0}(\Sigma)\ \xrightarrow{\ \sim\ }\ \HH^{2n},
 \qquad
 g\longmapsto\bigl((a_j,b_j)\bigr)_{j\le n},
 \label{e:eq:globaliso}
\end{equation}
of real dimension $8n$, with explicit inverse
\begin{equation}
 g(q)=\sum_{j=1}^{n}G^{\mathrm{sym}}_{p_j;a_j,b_j}(q)
     =\sum_{j=1}^{n}\frac2{v_j}\int_{S_{p_j}}
          E(q-s)\,(sa_j+b_j)\,\dd S(s).
 \label{e:eq:globalfield}
\end{equation}
Every such field decays at least as $O(\abs q^{-3})$, and more precisely
\begin{equation}
 g(q)=8\pi E(q)\sum_{j=1}^{n}v_j\,(u_ja_j+b_j)+O(\abs q^{-4})
     =\frac4\pi\,\frac{\bar q}{\abs q^4}
       \sum_{j=1}^{n}v_j\,(u_ja_j+b_j)+O(\abs q^{-4}).
 \label{e:eq:farfieldfirst}
\end{equation}
The only field in this class with a global slice-regular Fueter primitive
is $0$. Dropping the decay hypothesis, every smooth axially monogenic
field on $\HH\setminus\Sigma$ locally integrable across the spheres is
uniquely $g_{\mathrm{ent}}+\sum_jG^{\mathrm{sym}}_{p_j;a_j,b_j}$ with
$g_{\mathrm{ent}}$ entire axially monogenic.
\end{corollary}

\begin{proof}
Extract unique residues in disjoint small tubes by the $L^1$ normal form
of Part~VII, and subtract the corresponding finite sum of reflected
modes. Near each sphere the other modes are smooth and the singular part
at that sphere has been removed, so the remainder extends to an entire
axially monogenic field; by Proposition~\ref{e:prop:exterior} it still
tends to zero at infinity, so harmonic Liouville makes it zero. Without
the decay hypothesis the same subtraction leaves an entire field, which
is the last assertion. Conversely every finite mode sum lies in the class
and has the prescribed residue vector, which gives the isomorphism; the
layer form of \eqref{e:eq:globalfield} is Theorem~\ref{e:thm:layer}, and
the asymptotic \eqref{e:eq:farfieldfirst} is the $n=1$ term of
\eqref{e:eq:exterior} summed over $j$, using
$D_1=-\pi^{-1}\sum_jv_j(u_ja_j+b_j)$ and
$\lap(q^{-1})=-4\bar q/\abs q^4=-8\pi^2E(q)$. A global primitive forces
all periods to vanish, hence $g=0$; and if all pairs vanish in the
non-decaying version the field is entire, its upper base is the simply
connected half-plane, and Parts~II and~V give a primitive on all of
$\HH$.
\end{proof}

\begin{remark}[Local integrability at finite points, not global $L^1$]
The hypothesis in Corollary~\ref{e:cor:decaying} is local integrability
across each sphere together with decay at infinity. It is not an
assertion that $g\in L^1(\HH)$, and none is used.
\end{remark}

\begin{remark}[Two finite-dimensional objects, two different questions]
\label{e:rem:twoobjects}
The split exact sequence of Part~IV computes a \emph{cokernel}: the
obstruction to inverting $\lap$ on a fixed punctured domain, of real
dimension $8m$. Corollary~\ref{e:cor:decaying} computes a \emph{space of
actual fields}, of real dimension $8n$. They are linked by the same
explicit representatives $G^{\mathrm{sym}}_{p_j;a_j,b_j}$, and the
numerical coincidence of the two counts is not accidental, but they
answer different questions and neither implies the other.
\end{remark}

\begin{remark}[Real punctures are excluded on purpose]
\label{e:rem:realpunctures}
Theorem~\ref{e:thm:globalnormal} and Corollary~\ref{e:cor:decaying}
remove only nonreal spheres. Real punctures carry no additional
\emph{inversion} obstruction --- that is the content of Part~V's
corollary on $\HH$ minus finitely many real points, where the codimension
stays $8m$ independently of the number of deleted reals --- but they do
carry independent singular principal parts, which would have to be
listed in any classification theorem that also deleted real points. The
distinction is exactly the one Part~V insists on: $q^{-1}$ is a global
primitive of a nonremovable singularity.
\end{remark}

\subsection{Far-field moments: what a charge at infinity can and cannot
see}
\label{e:sec:farfield}

\begin{theorem}[Exact moment conditions for faster decay]
\label{e:thm:moments}
In the decaying class of Corollary~\ref{e:cor:decaying} set
\begin{equation}
 D_N=-\frac1\pi\sum_{j=1}^{n}\left(
    \frac{\Ima_{\CC}(p_j^{\,N+1})}{N+1}\,a_j
   +\frac{\Ima_{\CC}(p_j^{\,N})}{N}\,b_j\right),\qquad N\ge1 .
 \label{e:eq:globalmoments}
\end{equation}
Then for $\abs q>\max_j\abs{p_j}$,
\begin{equation}
 g(q)=\sum_{N\ge1}\lap(q^{-N})\,D_N ,
 \label{e:eq:globalexterior}
\end{equation}
and for every integer $K\ge1$,
\begin{equation}
 g(q)=O\bigl(\abs q^{-K-3}\bigr)\ \text{uniformly at infinity}
 \quad\Longleftrightarrow\quad
 D_1=\dots=D_K=0 .
 \label{e:eq:momentequivalence}
\end{equation}
\end{theorem}

\begin{proof}
Summing the finitely many expansions \eqref{e:eq:exterior} gives
\eqref{e:eq:globalexterior}. If the first $K$ coefficients vanish, then
homogeneity of $\lap(q^{-N})$ (degree $-N-2$) and the uniform
convergence of Proposition~\ref{e:prop:exterior} give the stated uniform
decay. Conversely evaluate \eqref{e:eq:globalexterior} on the positive
real axis: by the real-trace identity \eqref{e:eq:realtrace} applied to
$z^{-N}$,
\[
 \lap(q^{-N})\big|_{q=x}=-2N(N+1)\,x^{-N-2}.
\]
A first nonvanishing $D_N$ with $N\le K$ would therefore produce a
nonzero term of order $x^{-N-2}$, larger than $x^{-K-3}$, contradicting
the assumed decay; uniqueness of Laurent coefficients (or successive
multiplication by powers of $x$ and passage to the limit) makes the
argument componentwise rigorous.
\end{proof}

\begin{corollary}[A single charge at infinity cannot separate the two
residues]
\label{e:cor:chargeinfinity}
There exist fields in the decaying class with nonzero affine spherical
residue and vanishing first far-field moment. For $p=\iota$, the mode
$G^{\mathrm{sym}}_{\iota;1,0}$ has residue $(1,0)$ and
$D_1=-\pi^{-1}v(ua+b)=0$ since $u=0$ and $b=0$, so it decays as
$O(\abs q^{-4})$, whereas $G^{\mathrm{sym}}_{\iota;0,1}$ has residue
$(0,1)$, $D_1=-1/\pi\ne0$, and a genuine $\abs q^{-3}$ term.
\end{corollary}

This is the far-field counterpart of the flux statement of Part~VIII:
$\Phi_0$ alone misses the slope residue $a$, and $D_1$, being
$-\pi^{-1}\sum_jv_j(u_ja_j+b_j)$, is $-\Phi_0^{\mathrm{tot}}/(8\pi^2)$ up
to the obvious normalization and misses exactly the same information.
Both are cured in the same way: by pairing against $\mathcal P_1$ as well
as $\mathcal P_0$.

\subsection{Additivity of the singular energy}
\label{e:sec:additivity}

\begin{corollary}[Energy additivity over disjoint tubes]
\label{e:cor:additivity}
Let $g$ be axially monogenic on $\HH\setminus\Sigma$ and locally
integrable across every $S_{p_j}$, with residues $(a_j,b_j)$. For
pairwise disjoint sufficiently small tubes, with fixed outer radii
$R_j$,
\begin{equation}
 \sum_{j=1}^{n}\int_{\varepsilon<\operatorname{dist}(q,S_{p_j})<R_j}
        \abs{g(q)}^2\dd V
 =8\sum_{j=1}^{n}\mathcal R_{p_j}(a_j,b_j)^2\,\log\frac1\varepsilon
   +C+o(1)
 \qquad(\varepsilon\downarrow0),
 \label{e:eq:additivity}
\end{equation}
with $C$ a finite real constant. In particular the total logarithmic
energy coefficient of a configuration is the sum of the individual
residue norms, with no interaction term.
\end{corollary}

\begin{proof}
Near each sphere the other modes, and any entire remainder, are smooth
and therefore contribute a bounded amount to the shell density.
Theorem~\ref{e:thm:energy} applies inside each tube separately; adding
the $n$ laws and absorbing the fixed terms $8\mathcal
R_{p_j}^2\log R_j$ and the $n$ renormalized constants into $C$ gives
\eqref{e:eq:additivity}. There is no cross term because the tubes are
disjoint.
\end{proof}

\begin{remark}[What the global statements do not assert]
\label{e:rem:globalnotasserted}
Three limitations are worth stating explicitly, since the merged
statements above are deliberately sharp. First, every classification here
is within the \emph{axial} class; the uniqueness step in
Theorem~\ref{e:thm:layer} is the only statement in Parts~VIII--IX that is
proved without axiality, and nonaxial fields with the same surface source
are not classified. Second, all constants degenerate as $v\to0$: they
contain $1/v$, $1/v^3$ and $1/v^4$, so nothing here survives a sphere
collapsing onto the real axis, and no uniform statement over such a
degeneration is claimed. Third, Theorem~\ref{e:thm:ML} needs simple
connectivity of the base, and we do not assert a realization theorem over
a base with holes; over such a base the prescribed data must in addition
be compatible with the period lattice computed by the split range theorem
of Part~IV.
\end{remark}

%% ==================== part10_11.tex ====================

% ===========================================================================
%  Part X, Part XI, and the Appendices of the merged document.
%  Writer key: f.  All labels defined here begin with "f:".
%
%  Notation is the merged contract's: p = u + \iota v with v>0, offsets
%  \xi = x-u and \eta = y-v, \rho = |z-p|; \om_g and \thet_g for the two
%  closed forms (\etaC only for the primitive-dependent form of the local
%  two-integration construction); \Msph and \Max for the two DIFFERENT
%  maxima; \Rp for the residue norm; \Hg for the holomorphic invariant.
%  The meridian angle is written \varphi, not \theta, because \thet_g is
%  the second one-form.  E(q) denotes the Cauchy-Fueter kernel; the energy
%  is \Eg and never E.
% ===========================================================================

\section{Examples and separations}
\label{f:sec:examples}

The nine research manuscripts merged here carried, between them, some forty
worked examples, most of them duplicates of one another. The eight
collected below are exactly the ones that \emph{separate} concepts: each is
the unique reason that some hypothesis in Parts~I--IX cannot be dropped, or
that two notions which look interchangeable are not. Each is proved where
it is first needed; this section does not reprove any of them. What it adds
is the comparison --- which hypothesis each one kills, and which pairs of
notions each one pulls apart --- together with a closing table that reads
the whole catalogue against the theorems.

Throughout, $p=u+\iota v$ with $v>0$, $\Sp=u+v\SSS$, $\xi=x-u$,
$\eta=y-v$, $\rho=\abs{z-p}=\operatorname{dist}(q,\Sp)$, and
$G_{p;a,b}=\lap\I\bigl((za+b)(2\pi\iota)^{-1}\log(z-p)\bigr)$ is the
logarithmic mode of \eqref{c:eq:logstem}, with period pair $(a,b)$ by
\eqref{c:eq:mode-periods}.

\subsection{$G_{p;0,1}$ against $G_{p;1,0}$: the two obstructions are
independent}
\label{f:ex:two-obstructions}

Example~\ref{c:ex:independent} exhibits the pair
$\Per(G_{p;0,1})=(0,1)$ and $\Per(G_{p;1,0})=(1,0)$ on
$U=\Cplus\setminus\{p\}$.

\emph{What it separates.} It separates the \emph{first} integration from
the \emph{second}. For $G_{p;0,1}$ the form $\om_g$ is exact, so the slope
potential $\pot=B/y$ does have a global single-valued primitive and the
first of the two classical integrations succeeds --- yet there is no global
slice-regular primitive, because $\thet_g$ is not exact. A construction
that verified only the exactness of the first closed form would therefore
report success and be wrong. For $G_{p;1,0}$ the failure is already at the
first step. Neither test implies the other, and by
Proposition~\ref{b:prop:elimination} the second is not an artifact of a
badly chosen first primitive: $\thet_g$ is computed from the target alone.
Right multiplication by $1,\mathbf i,\mathbf j,\mathbf k$ turns each family
into four real directions, which is how the count $8m$ of
\eqref{c:eq:directsum} becomes visible one coordinate at a time.

\emph{A second separation carried by the same pair.} The obstruction is
global and leaves no local trace on $\RR$: by \eqref{c:eq:realtraces} the
real-axis traces of $G_{\iota;0,1}$ and $G_{\iota;1,0}$ are the perfectly
regular functions $4x/(\pi(1+x^2)^2)$ and $-4/(\pi(1+x^2)^2)$. Both fields
have $\Rp(a,b)^2=1$, so by \eqref{d:eq:exact-sup} and
\eqref{e:eq:energylaw} they obey the \emph{same} leading law
$\rho\Msph_g(\rho)\to1/\pi$ and the same energy law with coefficient $8$.
The two obstructions are independent but isometric, and no size
measurement can tell them apart --- only the pair of periods can.

\subsection{The stem $F=\iota$: the kernel is $\Aff$, not the
complex-affine stems}
\label{f:ex:affine-kernel}

The stem $F\equiv\iota$ has $A=0$, $B=1$, hence $f(x+Iy)=I$ and, by the
Fueter formula \eqref{a:eq:Fu-first}, $\lap\I(\iota)=-2I/y^2\ne0$; more
generally $\lap\I(\iota(za+b))=-2a/y-2I(xa+b)/y^2$, which is
\eqref{a:eq:complex-affine}, and which reappears as the kernel of the
holomorphic invariant in Remark~\ref{b:rem:notinjective}.

\emph{What it separates.} Three things at once.
\begin{enumerate}[label=\textup{(\roman*)}]
\item It separates $\Aff=\{qa+b:a,b\in\HH\}$ from the stems affine over
$\HC$. Replacing $\HH$ by $\HC$ would double the kernel and erase the
obstruction entirely, since the period pair takes values in $\HH^2$
precisely because the kernel does. This is the single most frequently
repeated warning in the source manuscripts, and this one-line computation
is the whole of its content.
\item It separates the period criterion from the holomorphic invariant.
The field \eqref{a:eq:complex-affine} has zero periods on every loop --- it
is a global Fueter image --- and it spans $\ker\Hg$ by
Remark~\ref{b:rem:notinjective}. So $\Hg$ is not injective on a product
domain, and ``$\Hg$ extends'' is a statement about $g$ only modulo this
eight-dimensional space.
\item It separates product domains from slice domains meeting $\RR$. The
factor $y^{-2}$ blows up at the axis, so smooth continuation across a real
interval forces $a=b=0$; $\ker\Hg$ is eight-real-dimensional on a product
domain and vacuous on a symmetric slice domain. Both statements are needed,
because Part~III states the range theorem in both settings.
\end{enumerate}

\subsection{$\lap\I((z-p)^{-1}c)$: zero periods without removability}
\label{f:ex:zero-period-pole}

\begin{example}[Exact order and exact integrability of the rational mode]
\label{f:ex:rational-mode}
Let $0\ne c\in\HC$ and let $g=\lap\I(F)$ with $F(z)=(z-p)^{-1}c$, on the
punctured tube $\Tube_R\setminus\Sp$. Since $F$ is single valued there, $g$
is a global Fueter image, so both of its periods vanish and it has a global
single-valued slice-regular primitive; its holomorphic invariant is
$\Hg=F''=2(z-p)^{-3}c$, a pole of order exactly three with
$h_{-1}=h_{-2}=0$. For $c=1$ the components are explicit: from
$(\xi+\iota\eta)^{-1}=(\xi-\iota\eta)/\rho^2$ one has
$A=\xi/\rho^2$, $B=-\eta/\rho^2$, and the Fueter formula gives
\begin{equation}
 P=-\frac{4\xi\eta}{y\rho^{4}},\qquad
 Q=\frac{2(\eta^{2}-\xi^{2})}{y\rho^{4}}+\frac{2\eta}{y^{2}\rho^{2}} .
 \label{f:eq:rational-PQ}
\end{equation}
Because $4\xi^2\eta^2+(\eta^2-\xi^2)^2=(\xi^2+\eta^2)^2=\rho^4$, the
leading squared pair norm is
\begin{equation}
 \abs P^{2}+\abs Q^{2}
 =\frac{16\xi^{2}\eta^{2}+4(\eta^{2}-\xi^{2})^{2}}{v^{2}\rho^{8}}
 +O(\rho^{-3})
 =\frac{4}{v^{2}\rho^{4}}+O(\rho^{-3}),
 \label{f:eq:rational-size}
\end{equation}
so $\Max_g(\rho)\sim2/(v\rho^2)$: the order $\rho^{-2}$ is attained
\emph{uniformly}, with no dependence on the meridian angle and no
exceptional direction $I\in\SSS$ along which the field is smaller.
Consequently $g$ is locally $L^s$ exactly for $0<s<1$, since
$\int_0\rho^{1-2s}\dd\rho<\infty$ only then, and in particular $g$ is
\emph{not} locally $L^1$: its planar integral diverges like
$\int_0\rho^{-1}\dd\rho$. It is nonetheless not removable, and this needs
no quantitative input: a removable $g$ would have a stem primitive
$F_\sharp$ on the full disk, and the affine-kernel theorem would force
$F-F_\sharp=za'+b'$ on the punctured disk, impossible because $F$ has a
pole.
\end{example}

\emph{What it separates.} It pulls apart three notions that the $L^1$
theorem identifies and that nothing else does:
\[
 \text{a global primitive on }\Tube_R\setminus\Sp,
 \qquad
 \Resaff_{\Sp}(g)=(0,0),
 \qquad
 \text{removability through }\Sp .
\]
The first two are equivalent on a punctured tube for \emph{every} axially
monogenic $g$, with no hypothesis --- that is the content of the period
criterion. Neither implies the third. This is exactly why the classification
theorem of Part~VII is stated with the local integrability hypothesis
(Theorem~\ref{d:thm:L1}) and not with the vanishing of the periods, and it
is the reason that theorem's two sharp thresholds are two \emph{different}
phenomena (Remark~\ref{d:rem:two-thresholds}): at $L^1$ the periods become
a complete invariant, at $L^2$ they must vanish, and the first threshold is
not a consequence of the second. The field above sits immediately below the
first threshold, its planar integral diverging like
$\int_0\rho^{-1}\dd\rho$.

\emph{Why the order matters.} The leading squared pair norm $4/(v^2\rho^4)$
is \emph{independent of the meridian angle and of $I\in\SSS$}. So the order
$\rho^{-2}$ is attained uniformly, with no exceptional quaternionic
direction along which the field is smaller --- which is what rules out any
attempt to rescue removability by evaluating on a well-chosen slice.

\subsection{$\lap\I(\exp((z-p)^{-1}))$: an essential spherical singularity}
\label{f:ex:essential}

The same construction with $F(z)=\exp((z-p)^{-1})$ gives zero periods, an
essential holomorphic encoding $\Hg=F''$, and therefore no bound
$\Msph_g(\rho)=O(\rho^{-N})$ for any $N$.

\emph{What it separates.} It separates the hypothesis-free normal form
$g=g_{\mathrm{reg}}+G_{p;a,b}+\lap\I\bigl(\sum_{k\ge1}(z-p)^{-k}c_k\bigr)$
of Theorem~\ref{d:thm:normal} from the finite-growth normal form of
Theorem~\ref{d:thm:growth}: the series really can have infinitely many
nonzero coefficients, so the filtration by $N$ does not exhaust the axially
monogenic germs at a sphere. Together with
Example~\ref{f:ex:zero-period-pole} it shows that the growth filtration is
a genuine filtration in both directions --- there are germs above every
finite level, and germs below $L^1$ that the residues do not see.

\subsection{The shell and the deleted real points: surjectivity on
multiply connected domains}
\label{f:ex:shell}

Example~\ref{c:ex:shell} treats $\Om_D=\{1<\abs q<2\}$, where $D$ is the
planar annulus: its single hole is conjugation invariant, the upper
meridian is simply connected, and the Fueter map is \emph{onto}, although
$b_1(D)=1$. Corollary~\ref{c:cor:deleted} treats
$\Om=\HH\setminus(\bigcup_jS_{p_j}\cup\{c_1,\dots,c_\ell\})$ with real
$c_k$, where the cokernel has real dimension $8m$ independent of $\ell$.

\emph{What it separates.} It separates the obstruction from the first Betti
number. With $s$ conjugation-fixed holes and $m$ exchanged pairs one has
$b_1(D)=s+2m$ and $\dim_\RR\operatorname{coker}\lap=8m$: not $8b_1(D)$ and
not $8(s+2m)$. The shell is the case $s=1$, $m=0$, where the naive count
predicts an eight-dimensional obstruction and the truth is zero; the
deleted real points are the same phenomenon with $\ell$ degenerate fixed
holes. The bridge that makes the two indices of Part~V agree ---
exchanged pairs on the one hand, winding basis of the upper meridian on the
other --- is Lemma~\ref{c:lem:invariant-holes}, and these two examples are
where it is felt.

\emph{What it does not say.} Corollary~\ref{c:cor:deleted} is about
\emph{inversion}, not removability. Real-point singularities keep their own
principal parts: by \eqref{c:eq:realkernel}, for real $c$ one has
$\lap((q-c)^{-1})=-4\overline{(q-c)}/\abs{q-c}^4=-8\pi^2E(q-c)$, so
$q\mapsto(q-c)^{-1}$ is, up to a constant, a global slice-regular primitive
of a field with a genuinely nonremovable isolated singularity at $c$. A
primitive exists; the singularity does not go away. Dropping this caveat
turns a true inversion theorem into a false removability theorem.

\subsection{$p=\iota$, $a=1$, $b=\mathbf i$: the correction is
$\ImH(b\bar a)$, not commutativity}
\label{f:ex:cancellation}

\begin{example}[A commuting residue pair with a nonzero correction]
\label{f:ex:im-correction}
Take $p=\iota$, so $u=0$ and $v=1$, and let $a=1$, $b=\mathbf i$. Then
\[
 \Rp(a,b)^2=\abs{ua+b}^2+v^2\abs a^2=1+1=2,
 \qquad
 \ImH(b\bar a)=\ImH(\mathbf i)=\mathbf i,
 \quad\abs{\ImH(b\bar a)}=1 ,
\]
so the two leading-size laws of Part~VII give two \emph{different}
constants,
\begin{equation}
 \lim_{\rho\downarrow0}\rho\Max_g(\rho)=\frac{\Rp(a,b)}{\pi v}
 =\frac{\sqrt2}{\pi},
 \qquad
 \lim_{\rho\downarrow0}\rho\Msph_g(\rho)
 =\frac{\sqrt{\Rp(a,b)^2+2v\abs{\ImH(b\bar a)}}}{\pi v}=\frac2\pi .
 \label{f:eq:two-constants}
\end{equation}
The mechanism is visible pointwise. Approaching the sphere along the
meridian direction $\varphi=0$, the uniform leading term of Part~VII reads
\begin{equation}
 g(\rho+I)=\frac{\mathbf i+I}{\pi\rho}+O(1+\abs{\log\rho}),
 \label{f:eq:pointwise-cancellation}
\end{equation}
which \emph{vanishes} at $I=-\mathbf i$ and attains the modulus
$2/(\pi\rho)$ at $I=\mathbf i$. Here $a=1$ and $b=\mathbf i$ commute, so
the correction term is not governed by $ab-ba$.
\end{example}

\emph{What it separates.} Two things.
\begin{enumerate}[label=\textup{(\roman*)}]
\item It separates the axial pair norm from the spherical supremum. These
are different functions of $\rho$ with different limits, and the source
manuscripts of this merge wrote both with the same symbol $M_g$, which
makes their statements contradictory as printed. Here $a=1$ and
$b=\mathbf i$ \emph{commute}, so the vanishing condition for the correction
term is not $ab=ba$; it is $\ImH(b\bar a)=0$
(Remark~\ref{d:rem:correction}). Any merged statement that quietly
identifies the two norms, or that offers commutativity as the criterion, is
wrong on this example.
\item It separates pointwise cancellation from removability. The leading
term vanishes identically along the ray $\{q=\rho-\mathbf i\}$, that is, on
an entire one-parameter family of points of the four-dimensional domain,
and the singularity is nonetheless a nonremovable critical one with a
nonzero residue pair. Evaluating an axially monogenic field on one
preselected complex slice is therefore not a substitute for the supremum
over $\SSS$, and a numerical test that sampled a single slice would report
a smaller order than the true one.
\end{enumerate}

\subsection{$a\ne0$, $b=-ua$: a nonremovable singularity of zero total
flux}
\label{f:ex:zero-flux}

By \eqref{e:eq:twoflux} this data gives
$\Phi_0=8\pi v(ua+b)=0$ at every radius, while
$\Phi_1=-8\pi v^3a\ne0$; the singularity is nonremovable, its germ carries
all eight real obstruction coordinates, and its energy diverges
logarithmically with coefficient $8v^2\abs a^2$. For $p=\iota$, $a=1$,
$b=0$ this is $G_{\iota;1,0}$, whose surface source
$\sigma_p(s)=\frac2v(sa+b)=2I$ on $\SSS$ has \emph{zero total mass} and
nonzero $L^2$ norm.

\emph{What it separates.} It separates the ordinary quaternionic flux from
the affine spherical residue. A single scalar flux is a one-quaternion
functional and cannot recover a two-quaternion invariant; the second,
conserved flux of Theorem~\ref{e:thm:twoflux} is what completes it. The
choice of weight there is not cosmetic: the naive $q-u$ satisfies
$\sum_\mu\partial_{x_\mu}(q-u)e_\mu=-2\ne0$, so it is not
right-Fueter-regular and its flux is not radius-independent, whereas
$H_u(q)=3(\Rea q-u)+\Ima q$ satisfies
$\sum_\mu\partial_{x_\mu}H_u\,e_\mu=3-3=0$
(Remark~\ref{e:rem:naiveweight}). The two expressions agree numerically
because $H_u$ restricted to $\Sp$ equals $q-u$, but the moment
presupposes knowledge of the density while the flux does not.

\emph{A unification worth repeating.} By Remark~\ref{e:rem:Huispoly} that
weight is exactly $-\tfrac14\lap((q-u)^3)=3\mathcal P_1(q-u)$, the
degree-one entire axially monogenic polynomial of \eqref{e:eq:Pk}. So the
two residues are recovered by pairing the field against the two
\emph{lowest entire axially monogenic polynomials}, $\mathcal P_0=1$ and
$\mathcal P_1$ --- which is the reason the flux pair and the period pair
carry the same information.

\subsection{The classical catalogue, and why it is still binding}
\label{f:ex:classical}

Four failures from the expository trunk of Part~I are invoked repeatedly
above and are collected here with the statements they delimit.

\begin{enumerate}[label=\textup{(\alph*)}]
\item \emph{Regular factors do not list roots}
(Example~\ref{a:ex:factors}): $(q-\mathbf i)*(q-\mathbf j)$ vanishes at
$\mathbf i$ but not at $\mathbf j$, while $(q-\mathbf i)*(q+\mathbf i)
=q^2+1$ vanishes on the whole sphere $\SSS$; the two share the
symmetrization $(q^2+1)^2$. The symmetrization sees a conjugacy class, not
a point. \emph{Binding on:} the fact that the singular set of Parts~VII--IX
is an entire two-sphere, never an isolated nonreal point, and that the
transverse geometry is codimension two.
\item \emph{The unqualified open mapping theorem fails}
(Example~\ref{a:ex:q2}): $q^2$ is not open at $\mathbf i$, because
$\Ima(q^2)=2x\Ima q$ forces $\abs{q-\mathbf i}\ge1$ whenever
$q^2=-1+t\mathbf j$ with $t\ne0$. \emph{Binding on:} the need to remove the
\emph{closure} of the degenerate set in the restricted open-map statement
of Part~I, and on the general principle that a quaternionic theorem must
name which regularity it uses.
\item \emph{Fueter regularity kills identity and open mapping}
(Example~\ref{a:ex:fueter-plane}): $L(q)=x_1-\mathbf i x_0$ is Fueter
regular, vanishes on the plane $\{x_0=x_1=0\}$ without being zero, and has
image inside one slice. \emph{Binding on:} every uniqueness argument in
Parts~VII--IX, which must therefore run through \emph{axial} data and the
planar Vekua system, and never through an accumulation-point identity
theorem for $\Dop$.
\item \emph{A zero flux residue without removability}
(Example~\ref{a:ex:dxE}): $\partial_{x_0}E$ is regular off the origin,
homogeneous of degree $-4$, and even, while the outward normal is odd, so
$\Res^F_0(\partial_{x_0}E)=0$ although the singularity is not removable.
\emph{Binding on:} Example~\ref{f:ex:zero-flux}, of which it is the exact
isolated-point ancestor, and on the standing rule that the three residue
notions of this document --- the slice Laurent residue $\Res_c$ at a real
isolated singularity (Part~I), the Cauchy--Fueter flux residue $\Res^F$
(Part~I), and the affine spherical residue $\Resaff_{\Sp}$
(Definition~\ref{d:def:resaff}) --- are three different objects and are
never to be identified. Their removability exponents belong to different
geometries: $o(\rho^{-1})$ at a nonreal sphere of codimension two, against
$o(r^{-3})$ at an isolated point in four dimensions.
\end{enumerate}

\subsection{The catalogue read against the theorems}

\begin{center}
\footnotesize
\renewcommand{\arraystretch}{1.3}
\begin{tabular}{@{}l l l@{}}
\toprule
Example & Separates & Would make true\\
\midrule
$G_{p;0,1}$ vs.\ $G_{p;1,0}$ &
first integration / second &
``$\om_g$ exact'' as a criterion\\
$F=\iota$ &
$\HH$- / $\HC$-affine kernel &
``$\ker\Hg=0$''\\
$\lap\I((z-p)^{-1}c)$ &
zero periods / removability &
Thm.~\ref{d:thm:L1} without $L^1$\\
$\lap\I(\exp((z-p)^{-1}))$ &
finite / infinite principal part &
Thm.~\ref{d:thm:growth} for all germs\\
shell, deleted real points &
obstruction / first Betti number &
``$\dim\operatorname{coker}=8b_1$''\\
$p=\iota,\,a=1,\,b=\mathbf i$ &
$\Max$ / $\Msph$; $\ImH(b\bar a)$ / $ab-ba$ &
one constant for both norms\\
$a\ne0,\ b=-ua$ &
total flux / affine residue &
``$\Phi_0$ classifies''\\
$L(q)=x_1-\mathbf i x_0$ &
Fueter / slice regularity &
an identity theorem for $\Dop$\\
$\partial_{x_0}E$ &
zero flux residue / removability &
``$\Res^F=0$ implies removable''\\
\bottomrule
\end{tabular}
\end{center}

The third column names the statement each example refutes: every one of
those nine sentences is a plausible simplification of a theorem in
Parts~II--IX, and each is false for exactly one reason, recorded above.

\section{The consolidated verification account}
\label{f:sec:verification}

The ten source manuscripts each carried a computational supplement. They
overlap almost completely, and they are consolidated here into a single
account, together with a single reconciliation of the numerical
certificates and a single disclaimer. Nothing below is offered as a proof;
the proofs are in Parts~II--IX.

\subsection{What is checked symbolically}
\label{f:sub:symbolic}

The following identities are finite algebraic identities in the real
components and can be, and have been, verified by exact symbolic
differentiation with symbolic real $x$, $y$, $u$, $v$ and symbolic
quaternionic parameters. Because every displayed identity is real-linear in
each quaternionic component separately, scalar parameters suffice for those
checks in which no product of two quaternionic unknowns occurs; where a
product does occur --- in the layer formula, in the spherical supremum and
in the energy identities --- genuinely noncommuting data are used.

\begin{enumerate}[label=\textup{(S\arabic*)}]
\item The two axial differential identities for
$\Dop(A+IB)$ and $\lap(A+IB)$, hence the Vekua system
$P_x-Q_y=2Q/y$, $Q_x+P_y=0$, and the Fueter formula
$\lap\I(F)=2A_y/y+I(2B_y/y-2B/y^2)$.
\item Both Vekua equations for the branch-free mode components
\eqref{c:eq:mode-P}--\eqref{c:eq:mode-Q}, for both
specializations $(1,0)$ and $(0,1)$ and for the reflected modes.
\item Closedness of $\om_g$ and $\thet_g$, i.e.\ vanishing of the
$\dd x\wedge\dd y$ coefficients $(Q_x+P_y)/2$ and
$\tfrac12\{y(P_x-Q_y)-2Q-x(Q_x+P_y)\}$.
\item The four potential derivative identities
$\pot_x=-P/2$, $\pot_y=Q/2$, $\potV_x=(xP+yQ)/2$, $\potV_y=(yP-xQ)/2$,
for a holomorphic test stem with genuinely complex $\HC$ coefficients.
\item The elimination identity $\etaC-\dd(x\pot)=\thet_g$.
\item The second-derivative identity $\Hg=F''$ and the two residue
moments $a_\gamma=\int_\gamma\Hg\dd z$,
$b_\gamma=-\int_\gamma z\Hg\dd z$, together with the equality of the two
printed forms of $\Hg$ under $Q_x=-P_y$.
\item The vanishing of $\Hg$ on the complex-affine kernel
\eqref{a:eq:complex-affine}.
\item The leading angular energy identity and the two diagonal Gram
entries $\pi\rho^2/2$ and $\pi\rho^2(v^2+\rho^2)/2$, together with the
vanishing off-diagonal entry, which is what makes the growth-free energy
lower bound of Part~VIII sharp.
\item The real-axis limits \eqref{c:eq:realtraces} and the identity
$\lap(q^3)=-4(3x_0+\Ima q)$, hence
$H_u=-\tfrac14\lap((q-u)^3)=3\mathcal P_1(q-u)$.
\end{enumerate}

\subsection{What is checked numerically}
\label{f:sub:numeric}

\begin{enumerate}[label=\textup{(N\arabic*)}]
\item \emph{Periods by contour integration.} Both forms are integrated
around circles in the planar base by periodic quadrature. The independent
runs behind this merge used a periodic trapezoidal rule with
$4\,096$, $16\,384$ and $32\,768$ nodes and, in one case, $45$-decimal
working precision. Circles both enclosing and not enclosing the singular
point were used, as were circles of several radii about the same point
(to test radius-independence), reflected modes (to test smoothness across
$\RR$), and a two-hole superposition with different noncommuting
quaternionic data at each hole (to test the full $16$-entry period matrix
of a two-hole basis). Observed maximum componentwise discrepancies from
the predicted pairs $(1,0)$, $(0,1)$, $(0,0)$ ranged from
$\sim10^{-16}$ to $\sim10^{-13}$ in double precision and below $10^{-46}$
in the high-precision run.
\item \emph{The layer formula.} The four-dimensional sphere integral
$\frac2v\int_{\Sp}E(q-s)(sa+b)\dd S(s)$ is evaluated by independent
spherical quadrature implementing the Hamilton product --- not by replacing
quaternionic products with scalar multiplication --- and compared against
the explicit mode at real and nonreal test points with noncommuting
$(a,b)$; observed component errors were of order $10^{-15}$ to $10^{-12}$.
\item \emph{The two conserved fluxes.} $\Phi_0$ and $\Phi_1$ are evaluated
on tubes of several radii with the spherical average taken analytically and
the meridian integral by quadrature; observed radius-independence and
agreement with $8\pi v(ua+b)$ and $-8\pi v^3a$ to about $10^{-12}$.
\item \emph{Shell energy.} The shell density
$\Jg(\rho)=\int_{\partial \Tube_\rho}\abs g^2\dd S$ and the truncated
energy $\Eg(\varepsilon,R)$ are computed by Gauss--Legendre or midpoint
rules in the meridian angle and in $\log\rho$, using the exact spherical
average $\frac1{4\pi}\int_{\SSS}\abs{P+IQ}^2\dd\sigma=\abs P^2+\abs Q^2$
rather than sampling $\SSS$. Both $\rho\Jg(\rho)\to8\Rp(a,b)^2$ and the
convergence of $\Eg(\varepsilon,R)-8\Rp^2\log(R/\varepsilon)$ to a finite
limit were observed.
\item \emph{Higher-order energy scale.} The leading shell energy for top
Laurent coefficients of orders $k=1,2,3$ with genuinely complexified
$\HC$ coefficients was compared with $16\pi^2k\norm{c_k}^2_{\HC}
\varepsilon^{-2k}$; observed relative errors below $5\times10^{-10}$ at
$\rho=10^{-4}$.
\item \emph{Supremum constant.} $\rho\Msph_g(\rho)$ is computed using the
exact algebraic formula
$\max_{I}\abs{P+IQ}^2=\abs P^2+\abs Q^2+2\abs{\ImH(P\bar Q)}$, so that only
the meridian angle is sampled and no sampling error over $\SSS$ is
incurred.
\end{enumerate}

\subsection{Reconciliation of the numerical certificates}
\label{f:sub:certificates}

Six independent parameter sets were used across the source manuscripts,
with no two the same. They are reconciled here: every one of them
reproduces from the two formulas
\[
 8\Rp(a,b)^2=8\bigl(\abs{ua+b}^2+v^2\abs a^2\bigr),
 \qquad
 \lim_{\rho\downarrow0}\rho\Msph_g(\rho)
 =\frac{1}{\pi v}\sqrt{\Rp(a,b)^2+2v\abs{\ImH(b\bar a)}} .
\]

\begin{center}
\small
\renewcommand{\arraystretch}{1.3}
\begin{tabular}{@{}l l r l@{}}
\toprule
$(u,v)$ & $(a,b)$ & $8\Rp(a,b)^2$ &
$\lim_{\rho\downarrow0}\rho\Msph_g(\rho)$\\
\midrule
$(0.3,\,1.2)$ &
$\begin{aligned}[t]
 a&=1-2\mathbf i+\tfrac12\mathbf j+\mathbf k\\[-2pt]
 b&=\tfrac12+\mathbf i-\mathbf j+2\mathbf k
\end{aligned}$ &
$253/2=126.5$ & $1.4724221755937\ldots$\\
$(0.3,\,1.7)$ &
$\begin{aligned}[t]
 a&=1-2\mathbf i+\tfrac12\mathbf j+3\mathbf k\\[-2pt]
 b&=-0.4+\mathbf i+2\mathbf j-\mathbf k
\end{aligned}$ &
$367.88$ & $1.611890081993\ldots$\\
$(0.3,\,1.7)$ &
$\begin{aligned}[t]
 a&=1+\tfrac12\mathbf i-\tfrac14\mathbf j+\tfrac34\mathbf k\\[-2pt]
 b&=0.2-0.4\mathbf i+0.3\mathbf j-0.1\mathbf k
\end{aligned}$ &
$46.38$ & ---\\
$(0.5,\,2)$ &
$\begin{aligned}[t]
 a&=1+2\mathbf i-\mathbf j\\[-2pt]
 b&=-\mathbf i+3\mathbf j+\mathbf k
\end{aligned}$ &
$252$ & ---\\
$(0.4,\,1.3)$ &
$\begin{aligned}[t]
 a&=1+2\mathbf i-\mathbf j+\tfrac12\mathbf k\\[-2pt]
 b&=\tfrac14-\mathbf i+3\mathbf j+2\mathbf k
\end{aligned}$ &
$181$ & ---\\
$(0.4,\,1.7)$ &
$\begin{aligned}[t]
 a&=1+0.2\mathbf i-0.3\mathbf j+0.7\mathbf k\\[-2pt]
 b&=-0.4+0.5\mathbf i+0.9\mathbf j-1.1\mathbf k
\end{aligned}$ &
$50.392$ & ---\\
\bottomrule
\end{tabular}
\end{center}

Two of these deserve comment. In the first row,
$\Rp^2=253/16$ and $\ImH(b\bar a)=4\mathbf i+3.75\mathbf j+3\mathbf k$ has
modulus $6.25$, so the spherical constant is
$\sqrt{253/16+15}/(1.2\pi)=1.47242\ldots$, whereas the \emph{pair-norm}
constant for the same data is $\Rp/(\pi v)=1.05480\ldots$; the two differ
by nearly forty percent and are limits of different quantities. In the
second row, $\Rp^2=45.985$ and
$\abs{\ImH(b\bar a)}=\sqrt{68.42}=8.27164\ldots$, giving
$\sqrt{45.985+28.12357}/(1.7\pi)=1.61189\ldots$. These two certificates
were produced independently, by manuscripts that did not share a formula;
that they reproduce from the \emph{same} spherical-supremum expression is
the strongest available evidence that the exact leading-size law of
Part~VII is correctly normalized.

The observed convergence is of the expected order. For the fifth row
(with $R=0.3$) the recorded shell density and renormalized remainder were
\begin{center}
\small
\renewcommand{\arraystretch}{1.2}
\begin{tabular}{@{}rrrr@{}}
\toprule
$\varepsilon$ & $\Eg(\varepsilon,R)$ &
$\Eg(\varepsilon,R)-181\log(R/\varepsilon)$ &
$\rho\Jg(\rho)$ at $\rho=\varepsilon$\\
\midrule
$10^{-2}$ & $626.9031484$ & $11.28642233$ & $181.2151049$\\
$10^{-3}$ & $1043.8028211$ & $11.41819314$ & $181.0049681$\\
$10^{-4}$ & $1460.5735515$ & $11.42102173$ & $181.0000892$\\
$10^{-5}$ & $1877.3415023$ & $11.42107075$ & $181.0000014$\\
\bottomrule
\end{tabular}
\end{center}
The middle column is the numerical appearance of $\Eren$: it stabilizes,
rather than merely staying bounded, which is exactly the distinction
between the weaker $O(1)$ form of the energy law and the finite-limit form
adopted in Part~VIII. For the first row the corresponding table of
$\rho\Msph_g(\rho)$ and $\rho\Jg(\rho)$ reads
$1.49118,\,1.47469,\,1.47269,\,1.47245$ and
$126.669,\,126.504,\,126.5001,\,126.50000$ at
$\rho=10^{-2},\dots,10^{-5}$, approaching the proved limits
$1.4724221756$ and $126.5$ at the predicted rate $O(\rho\log(1/\rho))$.

\begin{remark}[What the computations do and do not establish]
\label{f:rem:disclaimer}
These are checks of formulas, signs, normalizations and multiplication
order. They are not proofs. In particular they do not establish the
topological statements (connectedness of the upper meridian, the
identification of the two hole counts, the exactness of the range
sequence), the classification theorems (the hypothesis-free normal form,
the $L^1$ theorem, the growth filtration), the realization theorem, or any
priority claim; all of those rest on the written arguments of
Parts~II--IX. Quadrature discrepancies reported above are \emph{observed}
floating-point errors, not certified interval bounds. Symbolic checks
verify finite algebraic identities only. No proof assistant was used on any
statement in this document.

Two specific failure modes are worth recording as excluded by the checks
rather than by the prose. First, the central unit $\iota$ of $\HC$ appearing
in a holomorphic logarithm is never identified with a variable quaternionic
$I\in\SSS$; the period tests use quaternionic data whose imaginary parts are
not collinear, so a computation that silently restricted all coefficients to
one complex slice would fail them. Second, all quaternionic coefficients
multiply on the right, and the layer quadrature implements the Hamilton
product; a computation that replaced quaternionic multiplication by scalar
multiplication would agree on the commuting test data and fail on these.
\end{remark}

\section{Scope, literature boundary, and open problems}
\label{f:sec:scope}

This section states once, in one editorial voice, what this document claims,
what it declines to claim, and what it leaves open. The ten source
manuscripts each carried a version of it; they did not all take the same
position, and the reconciliation below is deliberate rather than a
composite.

\subsection{The contribution ledger}
\label{f:sub:ledger}

\emph{Established background, for which no priority is claimed.}
The stem formalism for slice functions and the representation formula;
Fueter's mapping theorem $\Dop(\lap f)=0$ and the biharmonicity
$\lap^2f=0$; the affine kernel $\ker(\lap|_{\SR})=\Aff$; the axial Vekua
system; the local inverse obtained by two successive integrations on a
simply connected base, together with its own period caveat; the spherical
Cauchy and arctangent kernels attached to $p=\iota$; classical one-variable
Runge~\cite{Runge} and Mittag--Leffler approximation and interpolation; Weyl's lemma;
and the standard cutoff/capacity method for $L^2$ removability of a
first-order elliptic system. The entire classical apparatus assembled in
Part~I --- the slice Cauchy formula, real-centred Taylor and Laurent
theory, the regular product and symmetrization, the zero geometry and the
argument principle, Schwarz--Pick and the intrinsic Riemann mapping
theorem, the Cauchy--Fueter Green identity and its consequences --- is
likewise expository synthesis of the cited sources.

\emph{Proposed contributions.} The connected package is:
\begin{enumerate}[label=\textup{(\alph*)}]
\item the two \emph{simultaneously} defined, primitive-independent closed
one-forms $\om_g$ and $\thet_g$, computed from the target alone, together
with the elimination identity $\etaC-\dd(x\pot)=\thet_g$ that shows the
second obstruction is not a residual of the first;
\item the same-domain solvability criterion with \emph{no} topological and
no boundary-regularity hypothesis, and its identification with the affine
monodromy homomorphism $H_1(U;\ZZ)\to\HH\oplus\HH$;
\item the split range theorem with its explicit section, its
$8m$ cokernel, and --- via the reflection analysis --- the correct hole
count, together with the lemma that every conjugation-invariant hole meets
the real axis, which is what reconciles the winding-basis index with the
exchanged-pair index;
\item the holomorphic invariant $\Hg$, its exact image $\Omom(U)$, its
exact kernel, the recovery of every multipole by an ordinary contour
integral, and the reality constraint
$h_{-1}=(\iota/v)\ReC(h_{-2})$ with its consequences (no isolated simple
pole of $\Hg$ at a nonreal $p$; the sharp criterion $h_{-2}=0$);
\item the local classification: the hypothesis-free normal form, the $L^1$
theorem with two independent proofs, the growth filtration with its $8N$
count, the sharp removability thresholds, the $L^p$ trichotomy, and the
exact leading constants in both norms;
\item the surface source $\Dop g=\frac2v(qa+b)\delta_{\Sp}$, the layer
representation with uniqueness, the two conserved three-dimensional fluxes,
and the energy law $\Eg(\varepsilon,R)=8\Rp(a,b)^2\log(R/\varepsilon)
+\Eren+o(1)$ with its sharp growth-free lower bound;
\item the realization (Mittag--Leffler) theorem and the global normal form
with mixed transverse orders and polynomial growth.
\end{enumerate}

\subsection{The literature boundary, and a single editorial position}
\label{f:sub:literature}

The literature audit behind this document was carried out on
21~September~2026 and was targeted rather than exhaustive. Inspected
primary sources include the slice-function framework of Ghiloni--Perotti~\cite{GP};
the inverse Fueter mapping theorem of Colombo--Sabadini--Sommen~\cite{CSS} and its
integral form using spherical monogenics; the integral-formula paper of
Colombo--Pe\~na~Pe\~na--Sabadini--Sommen, whose Theorem~3 is stated on a
\emph{rectangular} axial base and whose Example~2 exhibits the spherical
kernels; the contour surjectivity formulation of Dong--Qian; Perotti's
local Cauchy integral formula, whose Theorem~2 asserts surjectivity under
the hypothesis that every component of the symmetric planar set is simply
connected; Krau\ss har--Perotti on eigenvalue problems for slice functions;
and the 2025--2026 preprints of Colombo--De~Martino--Sabadini and
De~Martino--Pinton, which concern Laurent expansions of axially harmonic
and polyanalytic classes and a polyanalytic Cauchy--Kovalevskaya extension
rather than the ordinary four-dimensional Laplacian studied here.
Common--Sommen's axial construction and the 2016 inversion paper of
Dong--Kou--Qian--Sabadini were bibliographic leads whose complete texts
were not inspected. The absence of a matching statement from a bounded
search is not a proof of absence from the published literature, and a
specialist review remains necessary before any priority claim.

\begin{remark}[The position taken on published surjectivity statements]
\label{f:rem:position}
The source manuscripts of this merge did not agree on this point, and a
single position is adopted here.
\begin{enumerate}[label=\textup{(\roman*)}]
\item What this document replaces is an \emph{unqualified same-domain
surjectivity assertion}: the statement that every axially monogenic
function on an axially symmetric domain is the Fueter image of a
single-valued slice-regular function on that same domain. Part~V exhibits
explicit obstructed fields on $\HH\setminus\SSS$, so such an assertion is
false as stated, and Parts~II and~IV supply the necessary-and-sufficient
criterion and the exact size of the failure.
\item The simply-connected sufficient condition is correct as hypothesized.
What is removed here is the hypothesis, not the theorem.
\item No claim is made that any particular published theorem, stated with
its own domain, locality and branch hypotheses, is incorrect. Statements
called ``inverse'' or ``surjectivity'' theorems in this literature occur
with different domain and continuation conventions --- rectangular bases,
contour formulations, local integral formulas --- and neither the word
``inverse'' nor the existence of a local integral formula by itself removes
monodromy. Adjudicating any individual published sentence would require
auditing its own domain, locality and branch hypotheses, which this
document does not do.
\end{enumerate}
\end{remark}

\begin{remark}[Attribution of the logarithmic modes]
\label{f:rem:attribution}
The two-parameter family $G_{p;a,b}$ is not claimed as new. For $p=\iota$
one has $L'_{\iota}(z)=\pi^{-1}(1+z^2)^{-1}$, so a local primitive is
$\pi^{-1}\arctan z$ up to an additive real constant, and the fields with
local primitive stems $(2\pi)^{-1}\arctan z$ and
$(2\pi)^{-1}z\arctan z$ are the spherical Cauchy kernels $N^+$ and $N^-$;
in that normalization $G_{\iota;0,1}=2N^+$ and $G_{\iota;1,0}=2N^-$. This
identification is internally consistent with the normalizations used
throughout this document and is stated \emph{with} them, because it is
normalization-dependent. The attribution itself is unsettled in the source
material: some of the merged manuscripts cite
Colombo--Sabadini--Sommen~(2011), Theorem~5.1, for the arctangent
primitives and sphere-integrated Cauchy kernels, and others cite
Colombo--Pe\~na~Pe\~na--Sabadini--Sommen~(2014), Example~2. At most one of
these is the primary locus. The bibliography of this document should not
print either attribution without direct verification against both papers;
until that is done, the safe statement is that these kernels are known and
that what is developed here is their \emph{period normalization}, their use
as a complete obstruction basis, and the classification, source, flux and
energy laws attached to them.
\end{remark}

\subsection{Boundaries: what the theorems do not assert}
\label{f:sub:boundaries}

\begin{enumerate}[label=\textup{(B\arabic*)}]
\item \emph{Axiality.} Every classification statement in Parts~VII and~IX
is in the \emph{axial} class. It is not asserted that an arbitrary
Cauchy--Fueter-regular field is the Laplacian of a slice-regular function,
nor that a general monogenic field with a spherical singularity is
classified by two quaternions. A nonaxial field can carry far richer
tangential data along the sphere; the admissible source density there need
not be affine, and an infinite-dimensional analysis is to be expected. The
exceptions, explicitly noted where they occur, are the $L^2$ removability
theorem of Part~VII and the uniqueness half of the layer representation of
Part~VIII, neither of which uses axiality.
\item \emph{Positive radius.} All spherical results concern a nonreal
sphere $\Sp$ of radius $v>0$. The constants $2/v$, $8\pi v$, $8\pi v^3$ and
the residue norm $\Rp(a,b)^2=\abs{ua+b}^2+v^2\abs a^2$ all degenerate as
$v\downarrow0$, and none of the statements is uniform in that limit. The
collapse of a two-sphere to a real point is a different local problem, with
a different removability exponent; it cannot be read off by setting $v=0$
in any formula here.
\item \emph{Real points and nonaxial fields are different problems.}
Corollary~\ref{c:cor:deleted} says real punctures create no
\emph{inversion} obstruction. It does not say their singularities are
removable, and it does not classify their principal parts.
\item \emph{Finite topology.} The period criterion of Part~II requires no
finite-topology hypothesis of any kind, and the closed-range statement of
Part~VI holds on every connected base, even with infinitely many holes.
The finite-dimensional cokernel statements, by contrast, use the explicit
finite winding basis of Part~IV; no finite-dimensional cokernel is asserted
for a base with infinitely many independent holes.
\item \emph{Topology of the splitting.} The projection $\Pi=\mathrm{id}-S\Per$
and the normalized inverse are continuous for compact-open convergence, with
domain- and path-dependent constants. They are \emph{not} asserted to be
bounded in any unspecified weighted, Hardy, Bergman, Sobolev or Hilbert
norm, and the direct sum $\AM=\lap\SR\oplus S(\HH^{2m})$ is not asserted to
be orthogonal in any such norm. No uniform conditioning is claimed when
holes collide or approach the real axis.
\item \emph{Approximation.} The Runge statement of Part~VI concerns locally
uniform convergence with all derivatives on compact subsets, not
interpolation on a contractible patch, and the two settings have genuinely
different conclusions about the approximants: on a product domain with a
one-sided pole set the reflected slice functions need not be rational with
coefficients in $\HH$, whereas on a symmetric slice domain with a
conjugation-symmetric pole set they are.
\item \emph{Energy.} ``Energy'' means the squared $L^2$ norm
$\int\abs g^2\dd V$, not the Dirichlet integral $\int\abs{\nabla g}^2$.
The distinction matters for the critical exponent and for the logarithm.
\item \emph{The constant $8$ is codimension-two geometry.} It arises as the
angular length $2\pi$ of the transverse circle times the area $4\pi$ of
$\SSS$, divided by the $\pi^2$ of the logarithmic normalization
$1/(2\pi\iota)$. The $2\pi^2$ of the point Cauchy--Fueter kernel belongs to
a different geometry and does not appear in any of these constants.
\end{enumerate}

\subsection{Open problems}
\label{f:sub:open}

\begin{enumerate}[label=\textup{(O\arabic*)}]
\item \emph{Uniform estimates under degeneration.} The path-dependent
constants in the continuous inverse of Part~VI, and the $1/v$ powers in the
spherical constants, both obstruct any uniform statement for a family of
parameter domains whose holes collide, or for a sphere collapsing onto the
real axis. Find the correct rescaled limits. The exact leading asymptotics
of Part~VII, rather than the upper bounds, are the natural input.
\item \emph{Infinitely connected bases.} The criterion holds and the range
is closed for every connected base, but the splitting of Part~IV is
constructed from a finite winding basis. Is there a convergent kernel
splitting, with sharp norm control, on a base with infinitely many holes?
The naive sum of logarithmic modes diverges --- which is exactly the
difficulty the Mittag--Leffler realization theorem of Part~IX resolves in
the discrete locally finite case by prescribing the meromorphic second
derivative instead --- so the question is whether an analogous device
controls norms and not merely convergence.
\item \emph{A weighted Hilbert theory.} The energy law identifies a
canonical positive-definite quadratic form $8\Rp(a,b)^2$ on the obstruction
space, with the geometric and flux expressions
$8\Rp^2=\frac1{2\pi}\norm{\sigma_p}^2_{L^2(\Sp)}
=\frac1{8\pi^2}\bigl(\abs{\Phi_0}^2/v^2+\abs{\Phi_1}^2/v^4\bigr)$.
Is there a weighted Hilbert space of axially monogenic fields in which the
critical energy norm controls a canonical inverse after period subtraction,
so that $\Pi$ becomes an orthogonal projection and the normalized inverse a
bounded operator? Nothing proved here implies such boundedness.
\item \emph{Higher order: the Fueter--Sce map.} For Fueter--Sce operators
in a higher odd number of spatial variables the local kernel is a
polynomial space of higher degree. That suggests a corresponding higher
polynomial monodromy --- one period coordinate for each kernel coefficient
--- and a family of closed forms $\om^{(0)}_g,\dots,\om^{(N)}_g$
generalizing the pair $(\om_g,\thet_g)$, together with a higher spherical
residue and a higher energy law. The proofs here indicate the shape of the
invariant, but none of the higher-dimensional formulas is proved, and none
is needed for any statement in this document.
\item \emph{Beyond $L^1$.} Outside the locally integrable class the affine
residues stop being a complete invariant (\S\ref{f:ex:zero-period-pole}).
Classify the higher distributional derivatives of spherical sources under
weaker integrability or prescribed growth, and compare the resulting
hierarchy with the growth filtration of Part~VII.
\end{enumerate}

% ===========================================================================
\appendix
% ===========================================================================

\section{Distributional toolkit}
\label{f:app:distributions}

Part~VII proves its sharpest local theorem twice, by two genuinely
different mechanisms: a distributional route on the target, and a
$W^{1,1}$ route on the primitive. The general lemmas both routes need are
collected here, so that no part of the document repeats them. A third
subsection records the removability cutoffs, of which Part~VII uses one and
the sources supplied three.

\subsection{The distributional route}
\label{f:sub:pointmass}

\begin{lemma}[Distributions supported at one point]
\label{f:lem:point-support}
A distribution $T$ on $\RR^n$ supported at $0$ is a finite linear
combination of derivatives of $\delta_0$.
\end{lemma}
\begin{proof}
Continuity of a distribution gives, for a compact neighbourhood $K$ of $0$,
an order $r$ and a constant $C$ with
$\abs{\langle T,\phi\rangle}\le
C\max_{\abs\alpha\le r}\sup_K\abs{\partial^\alpha\phi}$
for all test functions supported in $K$. Let $\chi$ be a smooth cutoff
equal to one near $0$ and put $\chi_\varepsilon(x)=\chi(x/\varepsilon)$;
since $\operatorname{supp}T=\{0\}$ we have
$\langle T,\phi\rangle=\langle T,\chi_\varepsilon\phi\rangle$.
If all derivatives of $\phi$ through order $r$ vanish at $0$, Taylor's
theorem applied to each derivative gives
$\sup_{\abs x\le C_0\varepsilon}\abs{\partial^\beta\phi(x)}
=o(\varepsilon^{\,r-\abs\beta})$ for $\abs\beta\le r$; the product rule and
the bounds $\abs{\partial^\alpha\chi_\varepsilon}\le
C\varepsilon^{-\abs\alpha}$ then give
$\max_{\abs\alpha\le r}\sup\abs{\partial^\alpha(\chi_\varepsilon\phi)}
=o(1)$, whence $\langle T,\phi\rangle=0$. Therefore $T$ depends only on
the $r$-jet of its argument at $0$, and linear functionals on that
finite-dimensional jet space are exactly linear combinations of derivative
evaluations. The $\HH$- and $\HC$-valued statements follow componentwise
over $\RR$.
\end{proof}

The structure theorem alone is not enough. It permits derivatives of
$\delta_p$ of every order, whereas the classification of Part~VII needs
\emph{only} a point mass. The missing input is a scaling argument, and it
is exactly here that the $L^1$ hypothesis and the first order of the
operator are both consumed.

\begin{lemma}[An $L^1$ puncture creates only a point mass]
\label{f:lem:defect-scaling}
Let
\[
 \mathcal L(P,Q)=\Bigl(P_x-Q_y-\frac{2Q}{y},\ \ Q_x+P_y\Bigr)
\]
be the planar first-order Vekua operator on a disk
$B_p(R)\subset\Cplus$, so that $1/y$ is smooth and bounded there. If
$P,Q\in L^1(B_p(R);\HH)$ and $\mathcal L(P,Q)=0$ away from $p$, then there
are unique $s_0,s_1\in\HH$ with
\begin{equation}
 \mathcal L(P,Q)=(s_0,s_1)\,\delta_p
 \label{f:eq:defect}
\end{equation}
in the sense of distributions. No derivative of $\delta_p$ occurs.
\end{lemma}
\begin{proof}
The defect is supported at $p$, so by Lemma~\ref{f:lem:point-support} it is
a finite sum of derivatives of $\delta_p$. Suppose the highest nonvanishing
order were $r\ge1$. Choose a test function $\phi$ whose $r$-jet at $p$
detects a nonzero top-order coefficient and set
$\phi_\varepsilon(z)=\varepsilon^{\,r}\phi\bigl((z-p)/\varepsilon\bigr)$.
Pairing the defect against $\phi_\varepsilon$ tends to a nonzero constant,
all terms of lower order tending to $0$. On the other hand $\mathcal L$ is
\emph{first} order with bounded coefficients, so moving its derivatives
onto the test function gives
\[
 \abs{\langle\mathcal L(P,Q),\phi_\varepsilon\rangle}
 \le C\bigl(\varepsilon^{\,r-1}+\varepsilon^{\,r}\bigr)
 \int_{\abs{z-p}<C'\varepsilon}(\abs P+\abs Q)\,\dd x\,\dd y ,
\]
and the right-hand side tends to $0$ by absolute continuity of the integral
of an $L^1$ function --- \emph{including} at $r=1$, where the prefactor is
$O(1)$ and all of the decay comes from the shrinking domain of
integration. This contradiction leaves only the point mass, and uniqueness
is immediate.
\end{proof}

\begin{remark}[What is used, and where]
It is the conjunction of first order and $L^1$ that does the work: either
alone would leave the case $r=1$ open. This is the step that makes the
residue classification of Part~VII exhaustive rather than a list of
examples, and it is the exact point at which the integrability hypothesis
of that theorem is consumed. The $W^{1,1}$ route below consumes the same
hypothesis at a different place --- on the primitive rather than on the
target --- which is why the distributional route produces the surface
source of Theorem~\ref{e:thm:source} directly while the other deduces it
afterwards.
\end{remark}

The coefficient of the point mass is computed by one circle integral. Write
$W=P+\iota Q$ and $\bar\partial=\tfrac12(\partial_x+\iota\partial_y)$, so
that $\mathcal L(P,Q)=0$ becomes $2\bar\partial W=2Q/y$. For the mode
$G=G_{p;a,b}$ the uniform leading term of Part~VII gives
\[
 W=\frac{(ua+b)+\iota va}{\pi v\,(z-p)}+O(1+\abs{\log\rho}).
\]
The boundary functional for $2\bar\partial$ on the circle
$z-p=\varepsilon e^{\iota\varphi}$ is $\int(n_x+\iota n_y)W\,\dd s$; its
principal part integrates to $2\bigl((ua+b)+\iota va\bigr)/v$, the
logarithmic remainder contributes
$O(\varepsilon(1+\abs{\log\varepsilon}))$, and the lower-order term $2Q/y$
contributes nothing over a shrinking disk because it is $L^1$. Equivalently
this is the elementary normalization
$\bar\partial\bigl((z-p)^{-1}\bigr)=\pi\delta_p$, verified by the same
computation. Separating central real and imaginary parts gives the defect
of the canonical modes,
\begin{equation}
 \mathcal L(P_G,Q_G)
 =\Bigl(\frac{2(ua+b)}{v},\ \ 2a\Bigr)\delta_p .
 \label{f:eq:mode-defect}
\end{equation}
Lifted through the volume element $y^2\dd x\,\dd y\,\dd\sigma(I)$ and
$\dd S=v^2\dd\sigma(I)$ on $\Sp$, \eqref{f:eq:mode-defect} is exactly the
affine surface source $\Dop g=\frac2v(qa+b)\delta_{\Sp}$ of
Theorem~\ref{e:thm:source}.

\begin{lemma}[An $L^1$ Weyl lemma]
\label{f:lem:weyl}
If $w\in L^1_{\mathrm{loc}}(V)$ on an open $V\subset\RR^n$ and
$\Delta w=0$ distributionally, then $w$ has a smooth harmonic
representative.
\end{lemma}
\begin{proof}
Mollify: $w_\varepsilon=w*\psi_\varepsilon$ is smooth and harmonic on
slightly smaller subdomains, with local $L^1$ norms bounded uniformly by
those of $w$. For a smooth harmonic $h$, averaging the spherical mean-value
identities against a fixed smooth radial weight of total integral one
supported in a ball of radius $r$ writes $h$, on a concentric smaller ball
$B'$, as a convolution with a smooth kernel; differentiating the kernel
gives
\[
 \sup_{B'}\abs{\partial^\alpha h}\le C_{\alpha,r}\int_B\abs h\,\dd x
\]
for every multi-index $\alpha$, with $B'\Subset B$ and the kernel supports
fitting inside $B$. Applying this to $w_\varepsilon$ converts the uniform
$L^1$ bounds into uniform bounds on all derivatives on compact subsets.
Arzel\`a--Ascoli extraction gives a subsequence converging smoothly on
compact sets to a harmonic $h$, and $w_\varepsilon\to w$ in
$L^1_{\mathrm{loc}}$ identifies $h$ as a representative of $w$. Local
representatives agree almost everywhere and hence, being continuous,
everywhere.
\end{proof}

For $\HH$-valued fields the lemma is applied to each of the four real
components. Together with the factorization $\overline{\Dop}\Dop=\lap$ it
is the only elliptic regularity input used in Part~VII: it upgrades a
distributionally monogenic $L^1$ field to a smooth monogenic one, and a
distributionally monogenic $L^2$ field to a smooth one in the removability
theorem. The axial form used there --- that weak axial monogenicity across
the missing sphere is removable --- follows by mollifying $P,Q$ in the
planar variables, since pairing the lifted field with a four-dimensional
test function amounts to pairing the planar coefficients with $y^2$ times
its spherical averages, and those are again smooth planar test functions.

\subsection{The $W^{1,1}$ route}
\label{f:sub:W11}

This route never mentions distributions supported at a point. Its content
is that in two dimensions a puncture is too small to obstruct an integrable
gradient.

\begin{lemma}[A puncture does not obstruct an integrable gradient]
\label{f:lem:W11}
Let $B_R\subset\RR^2$ be a disk centred at $0$ and
$w\in C^1(B_R\setminus\{0\};\RR^N)$. If $\nabla w\in L^1(B_R)$, then $w$
extends, up to its value at the origin, to a function in
$W^{1,1}_{\mathrm{loc}}(B_R)$ whose weak gradient is its pointwise gradient
off $0$.
\end{lemma}
\begin{proof}
Fix $0<R_0<R$. In polar coordinates $(\rho,\varphi)$, radial integration
gives $\abs{w(\rho,\varphi)}\le\abs{w(R_0,\varphi)}
+\int_\rho^{R_0}\abs{w_t(t,\varphi)}\,\dd t$. Multiplying by $\rho$ and
integrating in $0<\rho<R_0$ and in $\varphi$ shows $w\in L^1(B_{R_0})$, the
contribution of $w_t$ being at most
\[
 \frac12\int_0^{2\pi}\!\!\int_0^{R_0}t^2\abs{w_t(t,\varphi)}
 \,\dd t\,\dd\varphi
 \le\frac{R_0}{2}\int_{B_{R_0}}\abs{\nabla w}\,\dd x\,\dd y .
\]
Moreover, with $H(t)=\int_0^{2\pi}\abs{w_t(t,\varphi)}\,\dd\varphi$,
\[
 \rho\int_0^{2\pi}\abs{w(\rho,\varphi)}\,\dd\varphi
 \le O(\rho)+\int_\rho^{R_0}\frac{\rho}{t}\,tH(t)\,\dd t\longrightarrow0
 \qquad(\rho\downarrow0),
\]
by dominated convergence, since $tH(t)$ is integrable. Integrating by parts
against a test function on $B_{R_0}\setminus B_\rho$, the inner boundary
term is bounded by the left-hand side above times $\sup\abs\phi$, and
therefore vanishes in the limit. This is the weak-derivative assertion on
$B_{R_0}$, and $R_0<R$ was arbitrary.
\end{proof}

\begin{lemma}[Weak Cauchy--Riemann removability]
\label{f:lem:weakCR}
If $F\in W^{1,1}_{\mathrm{loc}}(B_R;\CC^N)$ satisfies
$\partial_{\bar z}F=0$ in the sense of distributions, then $F$ agrees
almost everywhere with a holomorphic function on $B_R$.
\end{lemma}
\begin{proof}
Mollify on slightly smaller disks. The convolutions $F_\varepsilon$ are
smooth, satisfy the Cauchy--Riemann equation, and converge to $F$ in
$L^1_{\mathrm{loc}}$. For a holomorphic $H$ on a disk $K'$ and
$K\Subset K'$, the mean-value property on disks gives
$\sup_K\abs H\le C_{K,K'}\norm H_{L^1(K')}$. Applied to the differences
$F_\varepsilon-F_{\varepsilon'}$ this makes the family uniformly Cauchy on
compact subsets; the holomorphic limit agrees almost everywhere with $F$.
\end{proof}

\begin{remark}[Where the hypothesis enters]
The seam with Part~VII is the potential identities $\dd\pot=\om_g$ and
$\dd\potV=\thet_g$ of Part~II: the gradients of the eight real components
of $\pot=B/y$ and $\potV=A-xB/y$ are \emph{exactly} the coefficients of the
two one-forms, hence are $L^1$ near $p$ precisely when $g$ is. So the
hypothesis of Lemma~\ref{f:lem:W11} is the local integrability hypothesis
of the classification theorem, transported through the potentials with no
loss; Lemma~\ref{f:lem:weakCR} then extends the stem $F=A+\iota B$
holomorphically, and $\lap\I(F)$ is the regular extension of $g$. Neither
lemma uses the structure of distributions supported at a point, and neither
produces the surface source; the two routes are genuinely independent
proofs of the same theorem.
\end{remark}

\subsection{Removability cutoffs and the planar elliptic estimate}
\label{f:sub:cutoffs}

Three cutoff schemes for $L^2$ removability occur in the source material.
All are correct. Theorem~\ref{d:thm:L2} uses the first; the other two are
recorded here so that the merged document keeps them without giving them a
second proof of the theorem.

\begin{enumerate}[label=\textup{(C\arabic*)}]
\item \emph{Logarithmic cutoff} (the proof adopted in Part~VII, because it
is shortest and needs no classification machinery, no growth bound and no
axiality). With $\chi_\varepsilon=0$ for $\rho\le\varepsilon^2$, $=1$ for
$\rho\ge\varepsilon$ and affine in $\log\rho$ between,
$\norm{\nabla\chi_\varepsilon}^2_{L^2}\le C/\log(1/\varepsilon)\to0$, so
the cutoff error itself vanishes.
\item \emph{Linear-shell cutoff.} With $\chi_\varepsilon=0$ for
$\rho<\varepsilon$, $=1$ for $\rho>2\varepsilon$ and
$\abs{\nabla\chi_\varepsilon}\le C/\varepsilon$, a compact portion of a
smooth codimension-two submanifold has shell volume $O(\varepsilon^2)$, so
$\norm{\nabla\chi_\varepsilon}_{L^2}\le C$ --- bounded, but \emph{not}
tending to zero. The cutoff error
$C_\phi\norm g_{L^2(\{\varepsilon<\rho<2\varepsilon\})}
\norm{\nabla\chi_\varepsilon}_{L^2}$ still tends to zero, the decay now
coming entirely from absolute continuity of the integral of $\abs g^2$
rather than from the cutoff. This variant has the advantage of needing no
distance function adapted to $\Sp$ beyond smoothness, and it applies
verbatim to any smooth embedded two-sphere in a four-dimensional open set.
\item \emph{Planar reduction.} Writing $G=P+\iota Q$, the Vekua system is
equivalent to $\partial_{\bar z}G=Q/y$, and four-dimensional local $L^2$ of
$g$ is equivalent to two-dimensional local $L^2$ of $G$ (the spherical
average being exact, and $y$ bounded above and away from zero). The
linear-shell cutoff in the plane then gives the equation distributionally
across the puncture, and regularity follows from the elementary Fourier
identity
\begin{equation}
 \norm{\nabla w}_{L^2}=2\norm{\partial_{\bar z}w}_{L^2}
 \label{f:eq:elliptic2d}
\end{equation}
for compactly supported smooth $\CC^N$-valued $w$, which holds
distributionally whenever $w,\partial_{\bar z}w\in L^2$. Localizing $G$
makes its $\bar\partial$-derivative $L^2$, so $G\in H^1_{\mathrm{loc}}$;
since $1/y$ is smooth, differentiating, relocalizing and reapplying
\eqref{f:eq:elliptic2d} give $G\in H^k_{\mathrm{loc}}$ for every $k$, and
Sobolev embedding gives smoothness. This is the only one of the three that
produces the regularity without invoking Lemma~\ref{f:lem:weyl}, at the
cost of being restricted to the axial class.
\end{enumerate}

\begin{remark}[The fourth route, and the one thing it alone supplies]
A fourth argument appears in the sources: derive $\abs g=O(\rho^{-2})$ from
the harmonic mean value property and Cauchy--Schwarz, apply the $N=2$ case
of the growth filtration, kill the surviving Laurent coefficient by the
fact that a nonzero Laurent mode cannot have finite local $L^2$ energy, and
only then invoke the energy law. It is correct, but it needs the whole $8N$
classification and is restricted to axial fields, so it is not the primary
proof; see Remark~\ref{d:rem:L2-routes}. Its one irreplaceable ingredient
is the separate quantification of why the \emph{meromorphic} part, as
opposed to the logarithmic part, is excluded by finite energy, and that is
kept in Part~VIII next to the higher-order energy scale
$16\pi^2k\norm{c_k}_{\HC}^2\varepsilon^{-2k}$ of
Proposition~\ref{e:prop:poleenergy}.
\end{remark}

\section{Branch-free and polar formulas}
\label{f:app:formulas}

The formulas an implementation should use are collected here. Throughout,
$a,b\in\HH$, $d=ua+b$, $\xi=x-u=\rho\cos\varphi$,
$\eta=y-v=\rho\sin\varphi$, $\rho=\abs{z-p}$, $\ell=\log\rho$, and
$y=v+\eta$. The meridian angle is written $\varphi$, not $\theta$, because
$\thet_g$ denotes the second one-form.

\subsection{The Cartesian branch-free form, and what it hides}

The single-valued components of $G_{p;a,b}$ are
\eqref{c:eq:mode-P}--\eqref{c:eq:mode-Q},
\[
 P=\frac{1}{\pi y}\left[\frac{(xa+b)\xi+y\,a\,\eta}{\rho^{2}}
   +a\log\rho\right],
 \qquad
 Q=\frac{y\,a\,\xi-(xa+b)\eta}{\pi y\rho^{2}}
   +\frac{(xa+b)\log\rho}{\pi y^{2}},
\]
in which the multivalued real part of the logarithm has cancelled: only
$\log\rho=\ReC\log(z-p)$ survives, because a branch change contributes
$n(za+b)\in\ker\lap$. Every one of the nine research manuscripts merged
here derived this pair independently, in four mutually incompatible offset
conventions; after translation into $(\xi,\eta)$ they are literally the
same function, and that agreement is one of the load-bearing facts of the
merge.

What the Cartesian form hides is numerical: near $\Sp$ the bracket
subtracts nearly equal quantities, and the critical pole is not separated
from the logarithmic remainder. Both defects are repaired by passing to
polar coordinates.

\subsection{The polar form}

Substituting $\xi=\rho\cos\varphi$, $\eta=\rho\sin\varphi$, and replacing
the smooth coefficients $x,y,1/y$ by their values at the centre inside the
$O(\rho^{-1})$ terms, gives
\begin{align}
 P&=\frac{1}{\pi y}\left[
   \frac{d\cos\varphi+v\,a\sin\varphi}{\rho}+a(1+\ell)\right],
 \label{f:eq:polarP}\\
 Q&=\frac{1}{\pi y}\left[
   \frac{v\,a\cos\varphi-d\sin\varphi}{\rho}
   +\frac{(d+\rho\,a\cos\varphi)\,\ell}{y}\right],
 \label{f:eq:polarQ}
\end{align}
algebraically identical to \eqref{c:eq:mode-P}--\eqref{c:eq:mode-Q}. The
first bracketed term in each is the exact critical pole and the remainder
is $O(1+\abs{\log\rho})$, so the two scales never have to be recovered by
cancellation. Letting $\rho\downarrow0$ in
\eqref{f:eq:polarP}--\eqref{f:eq:polarQ} reproduces the uniform leading
term \eqref{d:eq:leading} of Part~VII, and taking the squared pair norm
reproduces $\Rp(a,b)^2/(\pi^2v^2\rho^2)$, which is the input to both
leading-size laws and to the energy law. In practice this form is what
makes the shell-energy and supremum tables of \S\ref{f:sub:certificates}
stable down to $\rho=10^{-5}$ and below; the Cartesian form loses several
digits there.

\subsection{The potentials as a second check of both normalizations}

The potentials of the same mode are \eqref{c:eq:mode-potentials},
\[
 \pot=\frac{1}{2\pi}\left(a\varphi-\frac{(xa+b)\ell}{y}\right),
 \qquad
 \potV=\frac{1}{2\pi}\left(b\varphi
   +\left(ya+\frac{x(xa+b)}{y}\right)\ell\right).
\]
They deserve a place in an implementation appendix for a reason
independent of their role in Part~IV. Since $\ell=\log\rho$ is single
valued and a positive circuit about $p$ increases $\varphi$ by $2\pi$,
these two formulas exhibit directly that $\pot$ gains exactly $a$ and
$\potV$ gains exactly $b$. Because $\dd\pot=\om_g$ and $\dd\potV=\thet_g$,
this is a \emph{second and independent} verification of both period
normalizations --- one that never integrates either form numerically, and
that fixes the sign conventions (slope period from $\om_g$, constant period
from $\thet_g$) without reference to any contour computation. A sign or
factor error anywhere in the chain
$\om_g\to\pot\to a$ or $\thet_g\to\potV\to b$ shows up here as a mismatch
of $2\pi$ or of sign, and is caught before any quadrature is run. The
formulas also display, in one line, the mechanism of
\S\ref{f:ex:two-obstructions}: for $a=0$ the multivalued term in $\pot$
disappears while the one in $\potV$ does not.

\section{An audit of signs, dimensions and normalizations}
\label{f:app:audit}

Every constant in this document is traceable to the following seven
conventions. A document assembled from sources using different ones would
silently change the model fields, the residues and the energy coefficient.

\begin{enumerate}[label=\textup{(A\arabic*)}]
\item \emph{The operator.} $\Dop=\partial_{x_0}+\mathbf i\partial_{x_1}
+\mathbf j\partial_{x_2}+\mathbf k\partial_{x_3}$, with \emph{no} factor
$\tfrac12$, and $\lap=\overline{\Dop}\Dop=\sum_{\mu=0}^3\partial_{x_\mu}^2$
is the ordinary \emph{positive} four-dimensional Laplacian. Diagnostic:
$\lap(q^2)=-4$, $\lap(q^3)=-4(3x_0+\Ima q)$, and
$\lap(q^{-1})=-4\bar q/\abs q^4=-8\pi^2E(q)$ with
$E(q)=\bar q/(2\pi^2\abs q^4)$. Using $-\lap$ instead reverses the signs of
the model fields and of both residues, but leaves the squared energy
coefficient unchanged.
\item \emph{Sides.} All quaternionic coefficients multiply on the
\emph{right}, in $\Aff$, in the period pair, in the mode $G_{p;a,b}$, in
the section $S$, and in every series. A right coefficient cannot be moved
to the left without changing the problem.
\item \emph{The central unit.} $\iota$ is central in
$\HC=\HH\oplus\iota\HH$, with $\kappa(A+\iota B)=A-\iota B$ and
$\abs{A+\iota B}^2_{\HC}=\abs A^2+\abs B^2$. It is never an element of
$\SSS$, and it is never identified with the variable $I\in\SSS$.
\item \emph{Which form gives which period.} The slope period is
$a_\gamma=\int_\gamma\om_g$ with
$\om_g=\tfrac12(-P\dd x+Q\dd y)$; the constant period is
$b_\gamma=\int_\gamma\thet_g$ with
$\thet_g=\tfrac12\bigl((xP+yQ)\dd x+(yP-xQ)\dd y\bigr)$. All loops are
counterclockwise-positive. The logarithmic normalization is
$1/(2\pi\iota)$.
\item \emph{Transverse geometry.} The normal circle $\abs{z-p}=\rho$ has
length $2\pi\rho$; $\SSS$ has area $4\pi$; the sphere $\Sp$ has area
$4\pi v^2$; the tube boundary and volume elements are
$\dd S=y^2\rho\dd\varphi\,\dd\sigma(I)$ and
$\dd V=y^2\rho\dd\rho\,\dd\varphi\,\dd\sigma(I)$, so
$\dd V=y^2\dd x\dd y\dd\sigma(I)$ in planar coordinates.
\item \emph{The derived constants.} Everything else follows:
$2/v$ in the source density $\sigma_p(s)=\frac2v(sa+b)$ (from
\eqref{f:eq:mode-defect} and $\dd S=v^2\dd\sigma$);
$8\pi v$ and $8\pi v^3$ in the two fluxes $\Phi_0=8\pi v(ua+b)$ and
$\Phi_1=-8\pi v^3a$ (from $\int_\SSS\dd\sigma=4\pi$ and the meridian
integral $2\pi$); the $8$ in
$\rho\Jg(\rho)\to8\Rp(a,b)^2$, which is $(2\pi)(4\pi)/\pi^2$; and the
$16\pi^2k$ in the higher-order energy scale.
\item \emph{Energy.} $\Eg(\varepsilon,R)=\int_{\varepsilon<\rho<R}
\abs g^2\dd V$ is the squared $L^2$ norm, not a Dirichlet integral;
$\Jg(\rho)=\int_{\partial \Tube_\rho}\abs g^2\dd S$ is the shell density,
so that $\Eg(\varepsilon,R)=\int_\varepsilon^R\Jg(s)\dd s$ and
$\Eren=\int_0^R\bigl(\Jg(s)-8\Rp^2/s\bigr)\dd s$.
\end{enumerate}

\begin{remark}[Two geometries, not one]
The factor $2\pi^2$ of the point Cauchy--Fueter kernel $E$ belongs to the
three-sphere $\abs q=r$ of area $2\pi^2r^3$ and to the isolated-point
geometry of Part~I. It does not appear in any spherical constant above,
whose geometry is codimension two: a normal circle times a two-sphere. The
removability exponents follow the same split --- $o(\rho^{-1})$ at a nonreal
sphere, $o(r^{-3})$ at an isolated point --- and mixing the two produces an
exponent that is correct in neither.
\end{remark}

\section{An implementation-oriented reconstruction procedure}
\label{f:app:procedure}

On a finitely connected base the theory gives a direct procedure that
solves no new boundary value problem. Given sampled or symbolic axial data:

\begin{enumerate}[label=\textup{\arabic*.}]
\item \emph{Recover $P,Q$.} From $g$ on the circularization,
$P=\frac1{4\pi}\int_\SSS g\dd\sigma$ and
$Q=-\frac1{4\pi}\int_\SSS Ig\dd\sigma$. These averages are exact and
inexpensive; do not sample $\SSS$ where an algebraic identity is available.
\item \emph{Check the equations.} Verify the Vekua residuals
$P_x-Q_y-2Q/y$ and $Q_x+P_y$. Arbitrary noisy pairs $(P,Q)$ need not even
define closed forms, and every statement below presupposes that they do.
\item \emph{Form the two one-forms directly.} Assemble the coefficients of
$\om_g$ and $\thet_g$ from $P$ and $Q$ alone. Do \emph{not} construct a
first primitive first: the second obstruction is computable from the target,
and the point of the elimination identity $\etaC-\dd(x\pot)=\thet_g$ is
that the two tests are simultaneous rather than sequential.
\item \emph{Integrate around one loop per conjugate pair.} On a symmetric
base, integrate only around the upper representative of each
conjugation-exchanged hole pair; conjugation-fixed holes contribute nothing
and need not be visited.
\item \emph{Subtract the modes.} Form $Sg=\sum_jG_{p_j;a_j,b_j}$ using the
branch-free formulas \eqref{c:eq:mode-P}--\eqref{c:eq:mode-Q},
or their polar version near a sphere, and replace $g$ by $\Pi g=g-S\Per(g)$.
Moving a chosen point $p_j$ inside its own hole changes the representative
by a zero-period field and does not change $\Pi g$'s membership in the
range.
\item \emph{Integrate the now-exact forms along any convenient paths.}
Compute $\pot$ and $\potV$ by path integration from a base point $z_*$ and
set $f(x+Iy)=(x+Iy)\pot+\potV$. Normalize by $\pot(z_*)=\potV(z_*)=0$,
equivalently $F(z_*)=0$ in $\HC$, which makes the primitive vanish on the
entire base sphere $x_*+y_*\SSS$ --- eight real conditions, not the four
conditions $f(x_*+I_*y_*)=0$ at one quaternionic point.
\item \emph{Use the quantitative bounds as diagnostics.} Before the
subtraction, a nonzero computed period is a \emph{certificate of
impossibility}, not a failure of the path integrator, provided it exceeds
the error budget: the bounds
$\abs{a_\gamma}\le\frac{L_\gamma}{2}\sup_\gamma\axn g$ and
$\abs{b_\gamma}\le\frac{L_\gamma R_\gamma}{2}\sup_\gamma\axn g$ convert data
error into period error, and the lower bound
$\sup_\gamma\axn{g-\lap f}\ge\max\{2\abs{a_\gamma}/L_\gamma,
2\abs{b_\gamma}/(L_\gamma R_\gamma)\}$, valid for \emph{every} global
slice-regular $f$, converts a period into a guaranteed error floor. After
the subtraction, the inverse estimate
$\abs{f(q)}\le LM(\abs q+R)/2$ bounds the reconstruction.
\end{enumerate}

A numerically small period is evidence only within specified integration and
data error bounds; the theorem requires exact vanishing. Conversely, a
period that is large compared with the bound in step~7 cannot be an artifact
of the quadrature.

\section{Theorem audit}
\label{f:app:theorem-audit}

{\small
\renewcommand{\arraystretch}{1.25}
\setlength{\LTpre}{8pt}\setlength{\LTpost}{8pt}
\begin{longtable}{@{}p{.24\textwidth}p{.33\textwidth}p{.33\textwidth}@{}}
\toprule
Statement & Essential hypothesis & What is \emph{not} asserted\\
\midrule
\endfirsthead
\toprule
Statement & Essential hypothesis & What is \emph{not} asserted\\
\midrule
\endhead
Global period criterion (Part~II) &
Smooth axial Vekua data on the circularization of a connected planar $U$;
no topology, no boundary regularity &
Inversion of arbitrary nonaxial Cauchy--Fueter-regular fields\\
Affine monodromy (Part~II) &
None beyond the above &
A multiplicative or path-ordered holonomy; the group is translation by the
fixed space $\Aff$\\
Range and kernel of $\Hg$ (Part~III) &
Axial monogenic $g$ on a connected $U$ &
Injectivity of $\Hg$: its kernel is eight-dimensional on a product domain\\
Reality constraint $h_{-1}=(\iota/v)\ReC h_{-2}$ (Part~III) &
An isolated \emph{nonreal} puncture $p$, $v>0$ &
Any analogue at a real point; the constraint uses $v>0$\\
Cokernel $8m$, product domain (Part~IV) &
An explicit finite winding basis of size $m$ &
An $8m$-dimensional function space; a finite cokernel for infinitely many
holes\\
Cokernel $8m$, symmetric domain (Part~V) &
$D$ connected, conjugation invariant, meeting $\RR$; $m$ exchanged hole
pairs &
A contribution from conjugation-fixed holes; the answer is not
$8b_1(D)=8(s+2m)$\\
Closed range and stability (Part~VI) &
Connected base; compact-open topology &
Boundedness in any weighted, Hardy, Bergman or Sobolev norm; orthogonality
of the splitting\\
Period-corrected Runge (Part~VI) &
Stated pole-set hypothesis: one-sided on a product domain,
conjugation-symmetric on a slice domain &
That the approximants are rational with $\HH$ coefficients in the
\emph{one-sided} case\\
Normal form, no hypothesis (Part~VII) &
Axial monogenic on a punctured tube &
Finite order; the principal part may be an infinite series
(Example~\ref{f:ex:essential})\\
$L^1$ classification (Part~VII) &
Local integrability across $\Sp$ --- formally weaker than
$\Msph_g=O(\rho^{-1})$, which implies it, though the two classes coincide
a posteriori &
Completeness of $(a,b)$ below $L^1$; see
\S\ref{f:ex:zero-period-pole}\\
Growth filtration $8N$ (Part~VII) &
$\Msph_g(\rho)=O(\rho^{-N})$ &
$4N$ or $8N+8$; and no fractional orders occur\\
Sharp removability (Part~VII) &
$\Msph_g(\rho)=o(\rho^{-1})$ &
That big-$O$ suffices; $G_{p;a,b}$ is the counterexample\\
$L^2$ removability (Part~VII) &
Finite local $L^2$ mass across the sphere; \emph{no} axiality, no growth
bound &
A new general elliptic removability principle; the method is standard\\
Exact leading size (Part~VII) &
Critical order, residue pair $(a,b)$ &
One constant: $\Max$ and $\Msph$ have \emph{different} limits, and the
correction vanishes iff $\ImH(b\bar a)=0$, not iff $ab=ba$\\
Surface source and layer (Part~VIII) &
Locally integrable axial field; uniqueness step needs no axiality &
A classification of arbitrary surface densities; only affine ones occur
here\\
Two conserved fluxes (Part~VIII) &
Right-Fueter-regular weights $1$ and
$H_u=3(\Rea q-u)+\Ima q$ &
That $\Phi_0$ alone classifies, or that the weight $q-u$ works\\
Energy law (Part~VIII) &
Critical growth, fixed outer radius $R$ &
Uniformity as $v\downarrow0$; $\Eren$ depends on $R$ and on the regular part\\
Growth-free lower bound (Part~VIII) &
Axial regularity off the sphere; \emph{no} growth hypothesis &
An optimal finite-radius correction term\\
Realization theorem (Part~IX) &
Simply connected upper base; discrete locally finite family of spheres &
Any compatibility condition between spheres --- there is none\\
Global normal form (Part~IX) &
Prescribed transverse orders $m_j$ and growth $O(\abs q^{d})$ &
A classification including real punctures, whose principal parts are
excluded\\
\bottomrule
\end{longtable}}

All statements listed are proved in the corresponding parts. The
computational supplement of \S\ref{f:sec:verification} checks identities,
normalizations and examples only; it is not a proof-assistant certificate,
and the qualified novelty assessment of \S\ref{f:sub:literature} applies to
the document as a whole, including every result labelled a contribution.
\clearpage
\section*{References}
\addcontentsline{toc}{section}{References}

This list is the union of the reference lists of the ten source
manuscripts: 81 bibliography entries naming 26 distinct works.  It is
reproduced in full so that the merged document does not lose the
literature those manuscripts consulted, and the comment on each entry
records which of the ten cited it.  Three works are cited by all ten.
Only some are cited in the text above; the others were consulted by one or
more of the sources for background, or for the priority audits reported in
Part~XI.

\begin{thebibliography}{99}

%% Merged from the bibliographies of the ten source manuscripts:
% 81 bibitem entries, 26 distinct works.  The supplied expository
%% manuscript is not listed here: it is Part I of this document.

\bibitem{CPSS}
F.~Colombo, D.~Pe\~na Pe\~na, I.~Sabadini, and F.~Sommen, \emph{A new integral formula for the inverse Fueter mapping theorem}, Journal of Mathematical Analysis and Applications \textbf{417} (2014), 112--122. \href{https://arxiv.org/abs/1302.0685}{arXiv:1302.0685}. The preprint's Theorem~3 is the rectangular-domain inverse result explicitly inspected for this audit.
% cited by source manuscripts 01, 02, 03, 04, 05, 06, 07, 08, 09, 10

\bibitem{CSS}
F. Colombo, I. Sabadini, and F. Sommen, \emph{The inverse Fueter mapping theorem}, Communications on Pure and Applied Analysis \textbf{10} (2011), 1165--1181. \href{https://doi.org/10.3934/cpaa.2011.10.1165}{doi:10.3934/cpaa.2011.10.1165}. Author manuscript: \url{https://biblio.ugent.be/publication/1938075/file/1938076.pdf}.
% cited by source manuscripts 01, 02, 03, 04, 05, 06, 07, 08, 09, 10

\bibitem{GP}
R.~Ghiloni and A.~Perotti, \emph{Slice regular functions on real alternative algebras}, Advances in Mathematics \textbf{226} (2011), 1662--1691. \href{https://doi.org/10.1016/j.aim.2010.08.015}{doi:10.1016/j.aim.2010.08.015}. \href{https://arxiv.org/abs/1008.4318}{arXiv:1008.4318}.
% cited by source manuscripts 01, 02, 03, 04, 05, 06, 07, 08, 09, 10

\bibitem{DMP}
A.~De Martino and S.~Pinton, \emph{A variation of the inverse Fueter theorem and the generalized polyanalytic Cauchy--Kovalevskaya extension of order 2}, preprint (2026), version 1, 7 August 2026. \href{https://arxiv.org/abs/2608.07711}{arXiv:2608.07711}.
% cited by source manuscripts 01, 02, 05, 07, 08, 09, 10

\bibitem{Perottib}
A. Perotti, \emph{A local Cauchy integral formula for slice-regular functions}, Computational Methods and Function Theory \textbf{24} (2024), 185--203. Published online 30 May 2023. \href{https://doi.org/10.1007/s40315-023-00485-5}{doi:10.1007/s40315-023-00485-5}. \href{https://arxiv.org/abs/2105.07041}{arXiv:2105.07041}. The simply connected surjectivity statement is Theorem~2 in the published article, Theorem~3 in the 2021 preprint.
% cited by source manuscripts 02, 04, 07, 08

\bibitem{CSS13}
F.~Colombo, I.~Sabadini, and F.~Sommen, \emph{The inverse Fueter mapping theorem in integral form using spherical monogenics}, Israel Journal of Mathematics \textbf{194} (2013), 485--505. \href{https://doi.org/10.1007/s11856-012-0090-4}{doi:10.1007/s11856-012-0090-4}.
% cited by source manuscripts 01, 07, 09

\bibitem{DQ}
B.~Dong and T.~Qian, \emph{Further Aspects of the Fueter Mapping Theorem}, preprint (2018). \href{https://arxiv.org/abs/1805.02966}{arXiv:1805.02966}. The primary preprint formulation, especially Theorem~3.6, is the version consulted here.
% cited by source manuscripts 05, 06, 10

\bibitem{GPS}
R.~Ghiloni, A.~Perotti, and C.~Stoppato, \emph{Singularities of slice regular functions over real alternative $*$-algebras}, Advances in Mathematics \textbf{305} (2017), 1085--1130. \href{https://doi.org/10.1016/j.aim.2016.10.009}{doi:10.1016/j.aim.2016.10.009}. \href{https://arxiv.org/abs/1606.03609}{arXiv:1606.03609}.
% cited by source manuscripts 03, 06, 09

\bibitem{GS}
G.~Gentili and D.~C.~Struppa, \emph{A new theory of regular functions of a quaternionic variable}, Advances in Mathematics \textbf{216} (2007), 279--301. \href{https://doi.org/10.1016/j.aim.2007.05.010}{doi:10.1016/j.aim.2007.05.010}.
% cited by source manuscripts 03, 08, 09

\bibitem{CDMS}
F.~Colombo, A.~De Martino, and I.~Sabadini, \emph{New Taylor and Laurent series of axially harmonic, Fueter regular and polyanalytic functions}, preprint, \href{https://arxiv.org/abs/2505.06555v2}{arXiv:2505.06555v2}, revised 28 August 2026. Relevant comparison: Section 8, on Fueter-regular Laurent expansions obtained from slice-hyperholomorphic sources.
% cited by source manuscripts 04, 08

\bibitem{Runge}
C. Runge, \emph{Zur Theorie der eindeutigen analytischen Functionen}, Acta Mathematica \textbf{6} (1885), 229--244. \href{https://archive.ymsc.tsinghua.edu.cn/pacm_paperurl/20170108202936251283802} {Author's article in the Acta Mathematica archive}.
% cited by source manuscripts 04, 10

\bibitem{BS}
C.~Bisi and C.~Stoppato, \emph{The Schwarz--Pick lemma for slice regular functions}, Indiana University Mathematics Journal \textbf{61} (2012), 297--317. \href{https://doi.org/10.1512/iumj.2012.61.5076}{doi:10.1512/iumj.2012.61.5076}; \href{https://arxiv.org/abs/1209.2060}{arXiv:1209.2060}.
% cited by source manuscripts 03

\bibitem{CGS}
F.~Colombo, G.~Gentili, and I.~Sabadini, \emph{A Cauchy kernel for slice regular functions}, Annals of Global Analysis and Geometry \textbf{37} (2010), 361--378. \href{https://doi.org/10.1007/s10455-009-9191-7}{doi:10.1007/s10455-009-9191-7}; \href{https://arxiv.org/abs/0902.4771}{arXiv:0902.4771}.
% cited by source manuscripts 03

\bibitem{Common}
A.~K.~Common and F.~Sommen, \emph{Axial monogenic functions from holomorphic functions}, Journal of Mathematical Analysis and Applications \textbf{179} (1993), 610--629. \href{https://doi.org/10.1006/jmaa.1993.1372}{doi:10.1006/jmaa.1993.1372}. Bibliographic lead; complete text not inspected in this audit.
% cited by source manuscripts 05

\bibitem{DDG}
A. De Martino, K. Diki, and A. Guzm\'an Ad\'an, \emph{The Fueter--Sce mapping and the Clifford--Appell polynomials}, Proceedings of the Edinburgh Mathematical Society \textbf{66} (2023), 642--688. \href{https://doi.org/10.1017/S0013091523000329}{doi:10.1017/S0013091523000329}. Preprint consulted: \href{https://arxiv.org/abs/2305.06998}{arXiv:2305.06998}. In particular, Section~4 studies weighted power-series range and kernel questions.
% cited by source manuscripts 04

\bibitem{DDGb}
A.~De Martino, K.~Diki, and A.~Guzm\'an Ad\'an, \emph{On the connection between the Fueter--Sce--Qian theorem and the generalized CK-extension}, preprint (2022). \href{https://arxiv.org/abs/2203.03490}{arXiv:2203.03490}.
% cited by source manuscripts 10

\bibitem{DKQS}
B.~Dong, K.~I.~Kou, T.~Qian, and I.~Sabadini, \emph{On the inversion of Fueter's theorem}, Journal of Geometry and Physics \textbf{108} (2016), 102--116. \href{https://www.sciencedirect.com/science/article/pii/S0393044016301401}{Publisher record}. Bibliographic lead; complete text not inspected in this audit.
% cited by source manuscripts 05

\bibitem{GGBS26}
J.~O.~Gonz\'alez-Cervantes, D.~Gonz\'alez-Campos, J.~Bory-Reyes, and B.~Schneider, \emph{A connection between quaternionic slice regular and Fueter hyperholormorphic functions}, preprint (2026). \href{https://arxiv.org/abs/2607.19496}{arXiv:2607.19496}.
% cited by source manuscripts 02

\bibitem{GSrecent}
R.~Ghiloni and C.~Stoppato, \emph{A manifold Fueter--Sce phenomenon in one hypercomplex variable}, preprint (2025). \href{https://arxiv.org/abs/2511.04771}{arXiv:2511.04771}.
% cited by source manuscripts 06

\bibitem{GStOpen}
G.~Gentili and C.~Stoppato, \emph{The open mapping theorem for regular quaternionic functions}, Annali della Scuola Normale Superiore di Pisa, Classe di Scienze (5) \textbf{8} (2009), 805--815. \href{https://doi.org/10.2422/2036-2145.2009.4.07}{doi:10.2422/2036-2145.2009.4.07}; \href{https://arxiv.org/abs/0802.3861}{arXiv:0802.3861}.
% cited by source manuscripts 03

\bibitem{GStZeros}
G.~Gentili and C.~Stoppato, \emph{Zeros of regular functions and polynomials of a quaternionic variable}, Michigan Mathematical Journal \textbf{56} (2008), 655--667. \href{https://doi.org/10.1307/mmj/1231770366}{doi:10.1307/mmj/1231770366}.
% cited by source manuscripts 03

\bibitem{Ghiloni}
R.~Ghiloni, \emph{Slice Fueter-Regular Functions}, The Journal of Geometric Analysis \textbf{31} (2021), 11988--12033. \href{https://doi.org/10.1007/s12220-021-00709-x}{doi:10.1007/s12220-021-00709-x}. \href{https://arxiv.org/abs/1911.06037}{arXiv:1911.06037}.
% cited by source manuscripts 06

\bibitem{KP}
R.~S.~Krau\ss har and A.~Perotti, \emph{Eigenvalue problems for slice functions}, Annali di Matematica Pura ed Applicata \textbf{201} (2022), 2519--2548. \href{https://doi.org/10.1007/s10231-022-01208-8}{doi:10.1007/s10231-022-01208-8}.
% cited by source manuscripts 06

\bibitem{Perotti}
A.~Perotti, \emph{A Local Cauchy Integral Formula for Slice-Regular Functions}, Computational Methods and Function Theory \textbf{24} (2024), 185--203; first published online 30 May 2023. \href{https://doi.org/10.1007/s40315-023-00485-5}{doi:10.1007/s40315-023-00485-5}.
% cited by source manuscripts 01

\bibitem{Rational}
D. Alpay, F. Colombo, A. De Martino, K. Diki, and I. Sabadini, \emph{On axially rational regular functions and Schur analysis in the Clifford--Appell setting}, Analysis and Mathematical Physics \textbf{14} (2024), article 41. \href{https://doi.org/10.1007/s13324-024-00902-5}{doi:10.1007/s13324-024-00902-5}.
% cited by source manuscripts 04

\bibitem{Sudbery}
A.~Sudbery, \emph{Quaternionic analysis}, Mathematical Proceedings of the Cambridge Philosophical Society \textbf{85} (1979), 199--225. \href{https://dougsweetser.github.io/Q/Stuff/pdfs/Quaternionic-analysis.pdf}{Author's paper, electronic reprint}.
% cited by source manuscripts 03

\end{thebibliography}
\end{document}
